\documentclass[UTF8,12pt,a4paper,fontset=fandol]{ctexbook}

% 支持从项目根目录或当前笔记目录编译。
\makeatletter
\def\input@path{{../}{../../}{../../../}{../../../../}}
\makeatother
\usepackage{researchnotes}

\title{概率论与数理统计（缪柏其）笔记}
\author{}
\date{}

\begin{document}
\maketitle

\frontmatter
\tableofcontents

\chapter{绪论}

概率论与统计学都研究随机现象和不确定性，但侧重点不同。\emph{概率论（probability theory）}从给定的假设和规则出发，研究各种结果出现的可能性及其规律；\emph{统计学（statistics）}研究如何收集和利用数据，描述已经观察到的现象、判断未知情况，并评价判断的可靠程度。\emph{前者侧重从模型推导可能的结果，后者的一个主要任务是利用数据认识模型及其未知部分。}

以重复抛掷硬币为例，假定每次正面出现的机会保持不变，各次抛掷的结果互不影响。如果还知道正反面出现的机会相同，研究“抛掷 100 次，恰好出现 60 次正面的可能性有多大”，就是典型的概率论问题。如果正面出现的机会未知，实际抛掷 100 次，记录到 60 次正面，再据此估计这一机会，并考察估计可能有多大误差，就是典型的统计学问题。

\emph{“记录中正面占六成”是在描述已观察到的事实，并不意味着硬币每次正面出现的真实机会一定是六成。}即使正反面的机会相同，抛掷 100 次也不保证正面恰好出现 50 次。因此，根据数据判断硬币的性质时，需要考虑偶然波动，不能仅凭次数偏离 50 次就断言硬币不公平。

概率论为统计方法提供数学基础，使我们能够研究数据和判断中的随机波动；统计学则利用数据认识未知情况，并检查模型是否适合实际问题。\emph{数理统计（mathematical statistics）}着重研究统计方法的数学原理及其性质，是统计学的重要组成部分。两者合在一起学习，体现了从认识随机规律到利用数据作出判断的知识联系。\emph{学习时应分清模型的假设、实际的观察记录，以及据此能够得到的结论。}

本笔记以缪柏其、张伟平主编的《概率论与数理统计》（高等教育出版社，2022 年第 1 版）为学习线索，按教材章节顺序整理学习思考与习题解答。

\mainmatter

\chapter{事件及其概率}
\label{chap:study-events}

% 本章依据指定教材第 1 章（印刷页 1--45）整理；保留原有专题笔记。
本章研究如何把随机现象转化为可以计算的数学模型。基本路线是先说明一次试验可能出现哪些结果，再用事件表示所关心的条件，最后在事件上赋予概率。计数、事件分解、条件化和独立性是本章的主要工具；它们分别解决“如何描述结果”“如何拆分所求事件”“如何利用已有信息”和“何时可以把概率相乘”的问题。

\section{概率论简史}
\label{sec:study-probability-history}

教材以赌博中的分配与计数问题说明概率论的早期动机：结果尚未发生，并不妨碍人们比较不同结果的可能性，并据此讨论合理的决策。1654 年帕斯卡（Blaise Pascal）与费马（Pierre de Fermat）关于赌本分配的通信，以及惠更斯（Christiaan Huygens）1657 年出版的《论赌博中的计算》，是这一发展过程中的代表性工作。伯努利（Jacob Bernoulli）在《猜度术》（1713 年出版）中研究频率与概率的联系，拉普拉斯（Pierre-Simon Laplace）的《分析概率论》（1812 年）进一步发展了古典概率理论。

从具体问题走向统一理论，还需要说明什么样的事件函数可以称为概率。科尔莫戈罗夫（Andrey Kolmogorov）于 1933 年提出以测度论为基础的公理化体系，使有限计数、连续模型和极限过程能够在同一框架下研究。教材随后介绍许宝騄等学者以及随机过程方向的工作。这里应把握的发展线索是：\emph{具体问题推动概念形成，频率规律提供经验依据，公理体系则统一推理规则。}下面保留赌本分配问题的完整讨论；其中条件期望与二项分布的名称会在后续章节进一步系统学习，此处只使用有限种结果的加权计算。

\subsection{中止比赛的赌本公平分配}
\label{subsec:points-fair-allocation}

甲、乙各拿出 $150$ 元，共 $300$ 元作为赌本，约定先赢三局者拿走全部赌本。每局没有平局，双方获胜的机会相同，各局结果相互独立。现在甲已赢两局、乙已赢一局，比赛却因外部原因中止。这 $300$ 元应怎样公平分配？关键在于：当前比分影响双方继续比赛后获胜的机会，不能直接按已赢局数之比 $2:1$ 分配。

\subsubsection{问题设定与分配原则}
\label{subsubsec:points-problem}

先把问题一般化。设 $A,B$ 各投入 $M/2$ 元，总赌本为 $M>0$ 元；比赛采用先赢 $s$ 局者获胜的规则，其中 $s$ 为正整数，胜者取得全部赌本，败者所得为 $0$。每局无平局，各局相互独立，$A$ 每局获胜的概率为已知常数 $p\in(0,1)$，$B$ 为 $1-p$。中止时双方分别已赢 $a,b$ 局，其中 $a,b\in\{0,1,\ldots,s-1\}$，尚未产生最终胜者。假定中止原因不提供关于未来胜负的额外信息，若比赛继续，后续各局仍遵循原有假设。

这里将“公平”规定为\emph{保持双方在当前信息下的期望所得}：提前分得的金额，应等于假想比赛继续至结束时各自所得的条件期望（conditional expectation）。因此，需要计算的是双方最终获胜的概率。记
\begin{equation}
  r=s-a,\qquad t=s-b,
  \label{eq:points-remaining}
\end{equation}
分别为 $A,B$ 还需取得的胜局数，并用 $Q(r,t)$ 表示在当前状态下 $A$ 最终获胜的概率。由于胜率已知且固定、各局独立，过去各局的具体先后顺序不影响这一概率。

$A$ 最终所得以概率 $Q(r,t)$ 为 $M$，以概率 $1-Q(r,t)$ 为 $0$，故其期望为 $M Q(r,t)$；$B$ 的期望同理。双方应分得
\begin{equation}
  x_A=M Q(r,t),\qquad x_B=M\bigl[1-Q(r,t)\bigr].
  \label{eq:points-allocation}
\end{equation}
这既保证金额非负，也保证 $x_A+x_B=M$。分配金额包含原始投入，净收益则分别为 $x_A-M/2$ 和 $x_B-M/2$。

\subsubsection{补齐比赛与一般公式}
\label{subsubsec:points-binomial}

为了计算 $Q(r,t)$，令
\begin{equation}
  N=r+t-1.
  \label{eq:points-horizon}
\end{equation}
若还没有胜者，$A$ 至多再赢 $r-1$ 局，$B$ 至多再赢 $t-1$ 局，合计至多进行 $r+t-2=N-1$ 局。因此，比赛最多再进行 $N$ 局便会结束。

现在预先生成 $N$ 个相互独立的未来比赛结果，每个结果中 $A$ 获胜的概率均为 $p$。真实比赛仍在某一方达到所需胜局数时结束，余下结果只用于把序列假想补齐，不改变已经确定的胜者。这样便可在固定的 $N$ 局中计数，而不必分别讨论比赛在不同局数结束的情形。

设随机变量 $K$ 表示这 $N$ 局中 $A$ 的胜局数。恰有 $k$ 局获胜的位置共有 $\binom Nk$ 种选法，每种结果序列的概率均为 $p^k(1-p)^{N-k}$，所以 $K$ 服从\emph{二项分布（binomial distribution）}，且
\begin{equation}
  \mathbb P(K=k)=\binom Nk p^k(1-p)^{N-k},
  \qquad k=0,1,\ldots,N.
  \label{eq:points-binomial-mass}
\end{equation}

\begin{proposition}
\label{prop:points-general-shares}
在上述假设下，令 $r=s-a$、$t=s-b$、$N=r+t-1$，则
\begin{equation}
  Q(r,t)=\sum_{k=r}^{N}\binom Nk p^k(1-p)^{N-k},
  \label{eq:points-binomial-tail}
\end{equation}
双方按当前信息下的期望所得应分得
\begin{equation}
  x_A=M\sum_{k=r}^{N}\binom Nk p^k(1-p)^{N-k},
  \qquad x_B=M-x_A.
  \label{eq:points-general-shares}
\end{equation}
\end{proposition}

\begin{proof}
关键是说明：$A$ 在真实比赛中获胜，当且仅当补齐后 $K\ge r$。若 $A$ 获胜，补齐后自然仍至少有 $r$ 个胜局；若 $B$ 获胜，则补齐后 $B$ 至少有 $t$ 个胜局，因而 $A$ 至多赢 $N-t=r-1$ 局。比赛必产生唯一胜者，故上述两个事件等价。于是，对式~\eqref{eq:points-binomial-mass} 中 $k=r,\ldots,N$ 的概率求和，再乘以赌本 $M$，便得到分配公式。
\end{proof}

\subsubsection{计算示例与适用范围}
\label{subsubsec:points-examples}

\begin{example}
\label{ex:points-original}
回到开头的问题：总赌本为 $300$ 元，双方每局胜率相等，先赢三局者获胜，中止时比分为 $2:1$。求双方应分得的金额。
\end{example}

\begin{solve}
此时 $r=1$、$t=2$，最多还需 $N=2$ 局。补齐后的两局结果为 $AA,AB,BA,BB$，每种概率均为 $1/4$；其中字母依次表示各局胜者，只有 $BB$ 对应 $B$ 最终获胜。注意，$AA,AB$ 中比赛其实在第一局后就已结束，第二个字母只是为计数而补上的结果。因此，$Q(1,2)=3/4$，应分给双方
\[
  x_A=300\times\frac34=225\text{ 元},\qquad
  x_B=300\times\frac14=75\text{ 元}.
\]
这也可直接看出：$B$ 必须连赢两局才能获胜。若每局 $A$ 的胜率改为一般的 $p$，则 $B$ 最终获胜的概率为 $(1-p)^2$，相应分配为 $x_A=300[1-(1-p)^2]$、$x_B=300(1-p)^2$。
\end{solve}

\begin{example}
\label{ex:points-biased}
双方各投入 $500$ 元，约定先赢四局者获得全部赌本。各局独立，$A$ 每局获胜的概率为 $3/5$；中止时 $A$ 已赢一局，$B$ 已赢两局。求公平分配金额。
\end{example}

\begin{solve}
此时 $M=1000$、$r=3$、$t=2$、$N=4$。由式~\eqref{eq:points-binomial-tail}，
\[
  Q(3,2)=\binom43\left(\frac35\right)^3\frac25
          +\binom44\left(\frac35\right)^4
        =\frac{297}{625}.
\]
所以 $A$ 分得 $1000\times297/625=475.2$ 元，$B$ 分得 $524.8$ 元。虽然 $A$ 每局胜率较高，但距离最终胜利还差更多局，说明分配同时取决于每局胜率和当前比分。
\end{solve}

上述“公平”以保持期望所得为准，投入相等并不意味着中止时应平分赌本。公式要求各局独立、胜率已知且固定，中止原因也不提供额外信息；若这些条件不成立，就需另建模型，不能把当前胜局比例直接当作每局胜率套用。

\section{随机试验和随机事件}
\label{sec:study-random-experiments}

\subsection{随机试验、样本空间与事件}
\label{subsec:study-sample-space}

\emph{随机试验（random experiment）}是对随机现象的一次实现或观测过程：在试验前可以描述可能结果，却不能预先确定本次将出现哪一个结果。同样条件下的重复试验有助于研究规律，但“重复”与“各次独立”是不同的要求，后者需要另外假定或论证。

\begin{definition}
\label{def:study-sample-space-event}
一次随机试验的所有可能结果组成的集合称为\emph{样本空间（sample space）}，记作 $\Omega$，其中的元素 $\omega$ 称为\emph{样本点（sample point）}。所关心的一种条件由满足该条件的样本点组成的集合 $A$ 表示，称为\emph{随机事件（random event）}；本次结果 $\omega\in A$ 就表示 $A$ 发生。单点集合 $\{\omega\}$ 在属于事件类时称为\emph{基本事件（elementary event）}。全集 $\Omega$ 是\emph{必然事件（certain event）}，空集 $\varnothing$ 是\emph{不可能事件（impossible event）}。
\end{definition}

样本点 $\omega$ 是一个结果，基本事件 $\{\omega\}$ 是一个集合，两者应在记号上区分。连续掷三次硬币时，可以取 $\Omega=\{H,T\}^{3}$，共有八个有序结果；“恰有一次正面”是 $\{HTT,THT,TTH\}$。若只记录正面次数，也可用 $\{0,1,2,3\}$ 描述结果，但即使硬币公平且各次独立，这四种次数也不等可能。记录灯管寿命时则可取 $\Omega=(0,\infty)$，事件 $\{t>500\}$ 表示寿命超过 $500$ 小时。样本空间可以有限、可数无穷或不可数，选择取决于记录的对象和精度。

在有限或可数样本空间中通常取全部子集为事件；在不可数样本空间中，需要指定允许赋予概率的集合类。其基本要求可以用\emph{$\sigma$ 代数（sigma-algebra）}表达：集合类 $\mathcal F$ 包含 $\Omega$，对取补集和可数并封闭，即 $A\in\mathcal F$ 蕴含 $A^{\mathrm c}\in\mathcal F$，且 $A_1,A_2,\ldots\in\mathcal F$ 蕴含 $\bigcup_{n=1}^{\infty}A_n\in\mathcal F$。由此也对可数交封闭。严格地说，事件是 $\mathcal F$ 的元素；本章所用的区间、有限集合及其可数运算均在相应的事件类内，无须展开不可测集合的构造。

% 教材例 1.3：三个骰子的点数和。
\begin{example}
\label{ex:study-three-dice-space}
独立掷三枚均匀骰子，为什么点数和为 $9$ 与为 $10$ 的概率不同，虽然把三个点数从小到大排列后，两种和都各有六种组合？
\end{example}

\begin{solve}
把骰子编号，取 $\Omega=\{1,\ldots,6\}^{3}$，其中 $216$ 个有序三元组等可能。无序组合对应的有序结果数并不相同：三个点数全异时有 $3!=6$ 种排列，恰有两个相同时有 $3$ 种，全相同时只有 $1$ 种。和为 $9$ 的无序组合为 $(1,2,6),(1,3,5),(1,4,4),(2,2,5),(2,3,4),(3,3,3)$，对应结果数为 $25$；和为 $10$ 的六种组合 $(1,3,6),(1,4,5),(2,2,6),(2,3,5),(2,4,4),(3,3,4)$ 对应结果数为 $27$。因此两个概率分别为 $25/216$ 和 $27/216$。\emph[caution]{压缩结果的记录方式后，必须同步确定新结果的概率，不能继续默认等可能。}
\end{solve}

\subsection{事件的关系与运算}
\label{subsec:study-event-operations}

事件包含 $A\subseteq B$ 表示 $A$ 发生必然导致 $B$ 发生；事件相等表示二者包含完全相同的样本点。\emph{并（union）} $A\cup B$ 表示至少一个发生，\emph{交（intersection）} $A\cap B$ 表示同时发生，\emph{差（difference）} $A\setminus B=A\cap B^{\mathrm c}$ 表示 $A$ 发生而 $B$ 不发生。\emph{对立事件（complementary event）} $A^{\mathrm c}=\Omega\setminus A$ 表示 $A$ 不发生。教材也把交写作 $AB$；本笔记主要保留 $\cap$，以免同概率的乘积混淆。若 $A\cap B=\varnothing$，称二者\emph{互斥（mutually exclusive）}或不相容；对立事件还要求 $A\cup B=\Omega$，因此互斥不一定对立。

并与交满足交换律、结合律和分配律，例如
\[
 A\cap(B\cup C)=(A\cap B)\cup(A\cap C),\qquad
 A\cup(B\cap C)=(A\cup B)\cap(A\cup C).
\]
证明集合恒等式时，可以固定任意样本点，逐步说明它属于左端与属于右端等价；只画出典型图形不能代替对所有样本点的论证。对于事件列，$\bigcup_n A_n$ 表示至少一个发生，$\bigcap_n A_n$ 表示每一个都发生。若需把并事件拆成互斥部分，可以令
\begin{equation}
 D_1=A_1,\qquad D_k=A_k\setminus\bigcup_{j=1}^{k-1}A_j\quad(k\ge2),
 \qquad \bigcup_k A_k=\bigcup_k D_k.
 \label{eq:study-disjointization}
\end{equation}
每个样本点只被归入它首次出现的那一项，故 $D_k$ 两两互斥。这种分解是从可加性推导一般概率不等式的常用方法。

% 教材例 1.6：用事件运算表达语言条件。
\begin{example}
\label{ex:study-three-event-language}
用三个事件 $A,B,C$ 表示“$A,B$ 发生而 $C$ 不发生”“三者不同时发生”“至多一个发生”“至少两个发生”和“恰好两个发生”。
\end{example}

\begin{solve}
前两者分别为 $A\cap B\cap C^{\mathrm c}$ 与 $(A\cap B\cap C)^{\mathrm c}$。至多一个发生包括全不发生和恰好一个发生，可以紧凑地写为
\[
 (A^{\mathrm c}\cap B^{\mathrm c})\cup(A^{\mathrm c}\cap C^{\mathrm c})
 \cup(B^{\mathrm c}\cap C^{\mathrm c}).
\]
至少两个发生为 $(A\cap B)\cup(A\cap C)\cup(B\cap C)$，其中包含三个都发生；恰好两个发生则为
\[
 (A\cap B\cap C^{\mathrm c})\cup(A\cap C\cap B^{\mathrm c})
 \cup(B\cap C\cap A^{\mathrm c}).
\]
最后三个组成部分互斥。“不同时发生”只排除三者全发生，不能与“至多一个发生”混用。
\end{solve}

\subsection{德摩根律：复合事件的对立事件}
\label{subsec:study-de-morgan}

描述随机事件时，经常需要把一个条件整体否定。例如，“两个事件至少发生一个”的否定是什么？“两个事件都发生”的否定又是什么？\emph{德摩根律（De Morgan's laws）}说明：对一个并集或交集取补集时，需要交换并与交，并对其中每个集合分别取补集。在概率论中，它把复合事件的“不发生”转化为若干较简单事件之间的关系，是寻找对立事件的基本依据。

沿用前面的样本空间与事件记号，相对于固定的 $\Omega$，事件 $A$ 的对立事件是
\[
  A^{\mathrm c}=\Omega\setminus A
  =\{\omega\in\Omega:\omega\notin A\},
\]
也就是“$A$ 不发生”，其集合称为 $A$ 的补集，有时也记作 $\overline A$。下面所有补集都相对于这个固定的 $\Omega$ 而言。

\begin{proposition}[两个事件的德摩根律]
\label{prop:study-de-morgan}
设 $A,B$ 是同一随机试验中的两个事件，则
\begin{align}
  (A\cup B)^{\mathrm c}
    &=A^{\mathrm c}\cap B^{\mathrm c},
    \label{eq:study-de-morgan-union}\\
  (A\cap B)^{\mathrm c}
    &=A^{\mathrm c}\cup B^{\mathrm c}.
    \label{eq:study-de-morgan-intersection}
\end{align}
\end{proposition}

式~\eqref{eq:study-de-morgan-union} 表示：\emph{“至少一个发生”的对立事件是“一个也不发生”}，因此必须让 $A$ 不发生且 $B$ 不发生。式~\eqref{eq:study-de-morgan-intersection} 表示：\emph{“都发生”的对立事件是“至少一个不发生”}，只要 $A,B$ 中有一个没有发生，就足以使“都发生”不成立。这是逻辑中否定与“或”“且”的关系：否定一个“或”条件得到两个否定条件的“且”，否定一个“且”条件得到两个否定条件的“或”。

\begin{example}
\label{ex:study-de-morgan-die}
掷一枚六面骰子一次，取 $\Omega=\{1,2,3,4,5,6\}$，样本点表示朝上一面的点数。设 $A$ 为“点数为偶数”，$B$ 为“点数大于 $3$”。分别写出 $A\cup B$、$A\cap B$ 及其对立事件，并核对德摩根律。
\end{example}

\begin{solve}
先按条件列出样本点：
\[
  A=\{2,4,6\},\qquad B=\{4,5,6\},\qquad
  A^{\mathrm c}=\{1,3,5\},\qquad B^{\mathrm c}=\{1,2,3\}.
\]
于是
\[
  \begin{aligned}
    A\cup B&=\{2,4,5,6\},&
    (A\cup B)^{\mathrm c}&=\{1,3\}
      =A^{\mathrm c}\cap B^{\mathrm c},\\
    A\cap B&=\{4,6\},&
    (A\cap B)^{\mathrm c}&=\{1,2,3,5\}
      =A^{\mathrm c}\cup B^{\mathrm c}.
  \end{aligned}
\]
第一行中，“既不是偶数，也不大于 $3$”只有点数 $1,3$；第二行中，“不是偶数，或者不大于 $3$”包括点数 $1,2,3,5$。特别地，点数为 $2$ 时 $A$ 发生而 $B$ 不发生，属于“$A,B$ 不都发生”，却不属于“$A,B$ 都不发生”。\emph{因此，“不都发生”不能理解为“都不发生”，也不能理解为“恰有一个不发生”，因为它还包括二者都不发生的情形。}
\end{solve}

德摩根律同样适用于有限多个事件。设 $n\ge1$ 为整数，$A_1,\ldots,A_n$ 是同一样本空间中的事件，则
\begin{equation}
  \left(\bigcup_{i=1}^{n}A_i\right)^{\mathrm c}
    =\bigcap_{i=1}^{n}A_i^{\mathrm c},
  \qquad
  \left(\bigcap_{i=1}^{n}A_i\right)^{\mathrm c}
    =\bigcup_{i=1}^{n}A_i^{\mathrm c}.
  \label{eq:study-de-morgan-finite}
\end{equation}
第一式说，“至少一个事件发生”的否定是“所有事件都不发生”；第二式说，“所有事件都发生”的否定是“至少一个事件不发生”。从样本点看，前者把“存在一个 $i$ 使 $\omega\in A_i$”的否定改写为“对每个 $i$ 都有 $\omega\notin A_i$”；后者把“对每个 $i$ 都有 $\omega\in A_i$”的否定改写为“存在一个 $i$ 使 $\omega\notin A_i$”。这既证明了有限情形，也说明同样的论证对可数无穷事件列 $A_1,A_2,\ldots$ 仍成立：
\begin{equation}
  \left(\bigcup_{i=1}^{\infty}A_i\right)^{\mathrm c}
    =\bigcap_{i=1}^{\infty}A_i^{\mathrm c},
  \qquad
  \left(\bigcap_{i=1}^{\infty}A_i\right)^{\mathrm c}
    =\bigcup_{i=1}^{\infty}A_i^{\mathrm c}.
  \label{eq:study-de-morgan-countable}
\end{equation}

德摩根律本身是\emph{事件之间的恒等式}，不要求事件互斥或相互独立，也不涉及各事件发生的概率大小。引入概率后，记 $\mathbb P(E)$ 为事件 $E$ 的概率，结合下一节将证明的对立事件关系 $\mathbb P(E^{\mathrm c})=1-\mathbb P(E)$，式~\eqref{eq:study-de-morgan-finite} 给出
\begin{equation}
  \mathbb P\!\left(\bigcup_{i=1}^{n}A_i\right)
    =1-\mathbb P\!\left(\bigcap_{i=1}^{n}A_i^{\mathrm c}\right).
  \label{eq:study-de-morgan-probability}
\end{equation}
\emph{因此，求“至少一个发生”的概率时，可以先求“一个也不发生”的概率，再用 $1$ 减去它。}\emph{不过，交事件的概率不能仅凭德摩根律写成各事件概率的乘积；能否相乘还需要相应的独立性条件。}记忆时可概括为：\emph{整体取补，并交互换，逐项取补}；理解时则应回到“至少一个”与“所有”的否定关系。

\section{概率的定义和性质}
\label{sec:study-probability-definition}

\subsection{概率的解释与公理化定义}
\label{subsec:study-probability-axioms}

概率用 $[0,1]$ 中的数刻画事件发生的可能性，但这一直观解释尚未规定如何赋值。教材从古典计数、频率和主观判断等角度讨论概率：古典概型依据有限个等可能结果计数；频率观点把概率与相同条件下独立重复试验中的长期比例联系起来；\emph{主观概率（subjective probability）}则表示依据已有信息对事件的可信程度作出的数量判断。不同解释用于说明概率的含义或来源，公理化定义则规定它们共同遵守的数学规则。

设独立重复同一试验 $n$ 次，事件 $A$ 出现 $N_n(A)$ 次，\emph{频率（relative frequency）}为 $F_n(A)=N_n(A)/n$。概率 $p=\mathbb P(A)$ 是给定模型中的常数，频率则随观测结果变化；频率接近概率并不表示误差随每次试验单调减小，也不表示有限次试验后必恰等于 $p$。在各次试验独立且事件概率相同的条件下，大数定律给出：对任意 $\varepsilon>0$，$\mathbb P(|F_n(A)-p|>\varepsilon)\to0$。其理论将在第~\ref{sec:study-limit-theorems} 节讨论，这里用它理解频率估计概率的依据。

\begin{definition}
\label{def:study-probability-space}
设 $\mathcal F$ 是样本空间 $\Omega$ 上的 $\sigma$ 代数。若实值集合函数 $\mathbb P:\mathcal F\to\mathbb R$ 满足非负性 $\mathbb P(A)\ge0$、规范性 $\mathbb P(\Omega)=1$，并且对任意两两互斥的事件列 $A_1,A_2,\ldots$ 满足\emph{可数可加性（countable additivity）}
\begin{equation}
 \mathbb P\!\left(\bigcup_{n=1}^{\infty}A_n\right)
 =\sum_{n=1}^{\infty}\mathbb P(A_n),
 \label{eq:study-probability-axioms}
\end{equation}
则称 $\mathbb P$ 为概率，三元组 $(\Omega,\mathcal F,\mathbb P)$ 称为\emph{概率空间（probability space）}。
\end{definition}

公理给出推理的起点，却不替具体问题选择模型。有限结果不等可能时仍可赋予概率：只要给各样本点指定 $p_i\ge0$ 且 $\sum_i p_i=1$，并令 $\mathbb P(A)=\sum_{\omega_i\in A}p_i$ 即可。因此，不能把“无法使用古典计数公式”理解为“无法定义概率”。主观赋值同样需要满足一致的概率规则，不能任意给互相关联的事件指定彼此矛盾的数值。

\subsection{概率的基本性质与连续性}
\label{subsec:study-probability-properties}

\begin{proposition}
\label{prop:study-probability-properties}
对同一概率空间中的事件，有 $\mathbb P(\varnothing)=0$ 和有限可加性；若 $A\subseteq B$，则
\begin{equation}
 \mathbb P(B\setminus A)=\mathbb P(B)-\mathbb P(A),\qquad
 \mathbb P(A)\le\mathbb P(B).
 \label{eq:study-probability-monotonicity}
\end{equation}
特别地，$0\le\mathbb P(A)\le1$ 且 $\mathbb P(A^{\mathrm c})=1-\mathbb P(A)$。对任意事件列还成立\emph{次可加性（subadditivity）}
\begin{equation}
 \mathbb P\!\left(\bigcup_{n=1}^{\infty}A_n\right)
 \le\sum_{n=1}^{\infty}\mathbb P(A_n).
 \label{eq:study-union-bound}
\end{equation}
\end{proposition}

\begin{proof}
对互斥事件列 $\Omega,\varnothing,\varnothing,\ldots$ 应用可数可加性，由 $\mathbb P(\Omega)=1$ 和非负性得到 $\mathbb P(\varnothing)=0$。有限个互斥事件后接空集，即得有限可加性。若 $A\subseteq B$，则 $B$ 是 $A$ 与 $B\setminus A$ 的不交并，故得到可减性，再由非负性得到单调性。取 $B=\Omega$ 即得补事件公式与上界。最后，按式~\eqref{eq:study-disjointization} 将任意并化为两两互斥的 $D_n$ 之并；因 $D_n\subseteq A_n$，可数可加性与单调性给出式~\eqref{eq:study-union-bound}。
\end{proof}

\begin{theorem}[概率的连续性]
\label{thm:study-probability-continuity}
若 $A_n\subseteq A_{n+1}$，记 $A=\bigcup_{n=1}^{\infty}A_n$，则 $\mathbb P(A_n)\to\mathbb P(A)$，称为下连续性；若 $A_n\supseteq A_{n+1}$，记 $A=\bigcap_{n=1}^{\infty}A_n$，同样有 $\mathbb P(A_n)\to\mathbb P(A)$，称为上连续性。
\end{theorem}

\begin{proof}
对递增情形，取 $D_1=A_1$、$D_n=A_n\setminus A_{n-1}$，则 $D_n$ 两两互斥，$A_n=\bigcup_{k=1}^{n}D_k$，$A=\bigcup_{k=1}^{\infty}D_k$。所以
\[
 \mathbb P(A)=\sum_{k=1}^{\infty}\mathbb P(D_k)
 =\lim_{n\to\infty}\sum_{k=1}^{n}\mathbb P(D_k)
 =\lim_{n\to\infty}\mathbb P(A_n).
\]
递减情形下 $A_n^{\mathrm c}$ 递增，利用德摩根律和补事件公式便得结论。这里的极限交换来自可数可加性和单调结构，不能对任意事件列直接照用。
\end{proof}

\emph[caution]{概率为零不等于事件为空集，概率为一也不等于事件就是全集。}例如在 $[0,1]$ 上按长度均匀取点，事件 $\{1/2\}$ 非空但概率为 $0$，其补集概率为 $1$ 却不等于 $[0,1]$。可数可加性只处理可数并；不能把区间写成不可数个单点的并后，将单点概率简单相加。

\subsection{古典概型与计数方法}
\label{subsec:study-classical-counting}

\begin{definition}
\label{def:study-classical-probability}
设随机试验的样本空间为 $\Omega=\{\omega_1,\ldots,\omega_n\}$，其中 $n\ge1$，各样本点互不相同。若每个样本点发生的概率均为 $1/n$，则称这一概率模型为\emph{古典概型（classical probability model）}。\emph{它具有两个基本特征：样本点有限，且各样本点等可能。}对任意事件 $A\subseteq\Omega$，有
\begin{equation}
  \mathbb P(A)=\frac{|A|}{|\Omega|}
  =\frac{\text{属于事件 }A\text{ 的样本点数}}{\text{全部样本点数}}.
  \label{eq:study-classical-probability}
\end{equation}
其中 $|A|$ 表示集合 $A$ 中的元素个数。
\end{definition}

上述公式来自概率的有限可加性：事件 $A$ 由 $|A|$ 个互斥的单点事件组成，每个单点事件的概率为 $1/n$，故总概率为 $|A|/n$。例如，掷一枚均匀骰子，六种点数等可能，“出现偶数”包含 $2,4,6$ 三个样本点，因此其概率为 $3/6=1/2$。应用古典概型前，应先明确样本空间并核实各样本点的等可能性，不能仅因结果个数有限就直接按数量之比计算概率。

计数时，\emph{加法原理（sum rule）}把互不重叠的各类方案数相加，\emph{乘法原理（product rule）}把分步选择数相乘，后者要求每个已完成的前序选择都有相应数目的后续选择。它们是计数结论，不是关于事件独立性的假设。对 $n$ 个不同对象取 $r$ 个，常用公式为
\begin{equation}
 \begin{array}{c|cc}
  &\text{有序}&\text{无序}\\ \hline
  \text{有放回}&n^r&\displaystyle\binom{n+r-1}{r}\\
  \text{不放回}&\displaystyle\frac{n!}{(n-r)!}&\displaystyle\binom nr
 \end{array}
 \label{eq:study-counting-formulas}
\end{equation}
其中 $n\ge1$，$r\ge0$ 为整数，不放回情形还要求 $r\le n$，约定 $0!=1$。有序不放回计数来自逐步选择 $n,n-1,\ldots,n-r+1$ 个对象；每个无序子集对应 $r!$ 个排列，故再除以 $r!$ 得组合数。无序有放回相当于求 $x_1+\cdots+x_n=r$ 的非负整数解，用 $r$ 个标记和 $n-1$ 个隔板排列，得到 $\binom{n+r-1}{r}$。这也解释了可辨球与不可辨球放入有编号盒子的四种计数模型。

若把 $n$ 个不同对象分到 $k$ 个有编号组，各组大小为 $r_1,\ldots,r_k$，其中 $r_i\ge0$ 且 $\sum_i r_i=n$，依次选组得到\emph{多项式系数（multinomial coefficient）}
\[
 \binom{n}{r_1}\binom{n-r_1}{r_2}\cdots
 =\frac{n!}{r_1!\cdots r_k!}.
\]
同一公式也计数含 $r_i$ 个第 $i$ 类符号的不同排列。\emph[caution]{计数公式并不自动保证这些方案等可能。}例如可辨球独立均匀投入盒子时，每个完整投放序列等可能，但不同占用数向量对应的序列数通常不同；不能直接用不可辨球的方案数作为概率分母。

% 教材例 1.9：有限总体不放回抽检。
\begin{example}
\label{ex:study-thesis-sampling}
$500$ 件对象中有 $20$ 件不合格，等可能地不放回抽取 $50$ 件；若其中不合格比例不超过 $5\%$ 就通过抽检，求通过的概率。
\end{example}

\begin{solve}
全部 $\binom{500}{50}$ 个子集等可能。令 $X$ 为抽中的不合格件数，通过条件为 $X\le2$，因为 $50\times5\%=2.5$ 且 $X$ 为整数。恰有 $i$ 件不合格的选法为 $\binom{20}{i}\binom{480}{50-i}$，$i=0,1,2$ 对应的事件互斥，故
\[
 \mathbb P(X\le2)=\sum_{i=0}^{2}
 \frac{\binom{20}{i}\binom{480}{50-i}}{\binom{500}{50}}
 \approx0.6775.
\]
这是一种超几何计数；抽到不同对象之间并不独立，但组合计数仍然有效。
\end{solve}

\subsection{生日问题：对立事件与对数近似}
\label{subsec:study-birthday-problem}

% 教材例 1.7：生日问题（birthday problem）。
\begin{example}
\label{ex:study-birthday}
试求 $r$ 人中至少有 $2$ 人生日相同的概率，一年按 $365$ 天计，其中 $r$ 为正整数。假设各人的生日相互独立，且每个人的生日在这 $365$ 天上均匀分布。
\end{example}

\begin{solve}
记 $A=\{\text{至少两人生日相同}\}$，则 $A^{\mathrm c}=\{\text{所有人的生日两两不同}\}$。若 $r\ge366$，必有至少两人生日相同，故 $\mathbb P(A)=1$；以下设 $1\le r\le365$。将日期编号为 $1,\ldots,365$，样本空间为 $\Omega=\{1,\ldots,365\}^{r}$，其中一个样本点 $(x_1,\ldots,x_r)$ 记录每个人的生日。由题设，$365^r$ 个样本点等可能。要使生日两两不同，各人的生日依次有 $365,364,\ldots,365-r+1$ 种选择，故
\begin{equation}
  \mathbb P(A)=1-\frac{365!}{(365-r)!\,365^r}
  =1-\prod_{k=1}^{r-1}\left(1-\frac{k}{365}\right).
  \label{eq:study-birthday-exact}
\end{equation}
其中 $r=1$ 时空乘积取 $1$，因而 $\mathbb P(A)=0$。

\emph{取自然对数可将连乘化为求和，再利用 $x$ 较小时的近似 $\ln(1-x)\approx-x$ 简化估算。}于是
\begin{equation}
  \begin{aligned}
    \ln\mathbb P(A^{\mathrm c})
      &=\sum_{k=1}^{r-1}\ln\left(1-\frac{k}{365}\right)\\
      &\approx-\frac1{365}\sum_{k=1}^{r-1}k
      =-\frac{r(r-1)}{730}.
  \end{aligned}
  \label{eq:study-birthday-log}
\end{equation}
取指数还原，得到
\begin{equation}
  \mathbb P(A)\approx1-\exp\left(-\frac{r(r-1)}{730}\right).
  \label{eq:study-birthday-approx}
\end{equation}
此近似要求各项 $k/365$ 较小且累计误差可以接受。$r=30,50$ 时，近似值分别为 $0.6963,0.9651$；按精确公式~\eqref{eq:study-birthday-exact} 计算，相应结果为 $0.7063,0.9704$（保留四位小数）。
\end{solve}

\subsection{几何概型与抽样机制}
\label{subsec:study-geometric-probability}

\emph{几何概型（geometric probability model）}用长度、面积或体积替代有限计数。设 $D\subseteq\mathbb R^d$ 是可测区域，$0<m(D)<\infty$，其中 $m$ 是相应维度的几何测度；若取点的分布相对于该测度均匀，则对可测事件 $A\subseteq D$，
\begin{equation}
 \mathbb P(A)=\frac{m(A)}{m(D)}.
 \label{eq:study-geometric-probability}
\end{equation}
这里均匀的是指定参数或区域上的取样机制，不能把“随机”直接解释为任何坐标下都均匀。不可数个样本点也不能使用有限古典概型的“有利点数除以总点数”。

% 教材例 1.4、习题 15：伯川德问题。
\begin{example}
\label{ex:study-bertrand-chord}
在单位圆中取一条弦，求弦长大于内接等边三角形边长 $\sqrt3$ 的概率，并比较三种随机取弦方式。
\end{example}

\begin{solve}
设弦中点到圆心的距离为 $d$，弦长为 $2\sqrt{1-d^2}$，故所求条件等价于 $d<1/2$。若先固定一个端点，再在圆周上按弧长均匀选另一端点，第二端点须落在与第一端点的圆心夹角介于 $2\pi/3$ 与 $4\pi/3$ 的圆弧上，概率为 $1/3$；同时随机均匀选两个端点也得到同一结果。若按面积在圆盘内均匀选弦中点，事件区域是半径 $1/2$ 的圆盘，概率为 $1/4$。若先选一个方向，再在相应直径上按长度均匀选弦中点，并作垂直于直径的弦，则中点坐标 $u$ 均匀于 $[-1,1]$，条件 $|u|<1/2$ 的概率为 $1/2$。零概率的退化位置可另作约定，不影响结果。

三个答案对应三个不同的概率模型，这就是\emph{伯川德悖论（Bertrand's paradox）}的要点：只说“随机取一条弦”不足以唯一确定概率；必须说明如何取样。
\end{solve}

\subsection{加法定理与容斥原理：重叠事件的概率}
\label{subsec:study-inclusion-exclusion}

求若干事件中“至少有一个发生”的概率，就是求这些事件的并集的概率。若事件之间存在重叠，直接相加各事件的概率，就可能把共同部分重复计算。概率的\emph{加法定理（addition theorem）}通过交替减去和补回各重交集的概率，修正这种重复计数；它是\emph{容斥原理（inclusion--exclusion principle）}在概率计算中的形式。以下各事件均属于同一概率空间，用 $\mathbb P$ 表示概率；教材中的 $A_iA_j$ 表示交集 $A_i\cap A_j$，不是概率的乘积。

先考察两个事件。记 $B\setminus A=B\cap A^{\mathrm c}$，即“$B$ 发生而 $A$ 不发生”。由于 $A\cup B$ 是 $A$ 与 $B\setminus A$ 的不交并，而 $B$ 是 $B\setminus A$ 与 $A\cap B$ 的不交并，由概率的有限可加性，
\[
  \mathbb P(A\cup B)=\mathbb P(A)+\mathbb P(B\setminus A),
  \qquad
  \mathbb P(B)=\mathbb P(B\setminus A)+\mathbb P(A\cap B).
\]
消去 $\mathbb P(B\setminus A)$，得到
\begin{equation}
  \mathbb P(A\cup B)
    =\mathbb P(A)+\mathbb P(B)-\mathbb P(A\cap B).
  \label{eq:study-inclusion-exclusion-two}
\end{equation}
当 $A,B$ 互斥，即 $A\cap B=\varnothing$ 时，交集项为零，公式退化为互斥事件的加法公式。

对于三个事件，先对 $A\cup B$ 与 $C$ 使用式~\eqref{eq:study-inclusion-exclusion-two}，再由分配律 $(A\cup B)\cap C=(A\cap C)\cup(B\cap C)$ 对交集部分使用同一公式，得
\begin{equation}
  \begin{aligned}
    \mathbb P(A\cup B\cup C)
      ={}&\mathbb P(A)+\mathbb P(B)+\mathbb P(C)\\
         &-\mathbb P(A\cap B)-\mathbb P(A\cap C)-\mathbb P(B\cap C)\\
         &+\mathbb P(A\cap B\cap C).
  \end{aligned}
  \label{eq:study-inclusion-exclusion-three}
\end{equation}
三者交集为何要加回来？属于三个事件共同部分的样本点，在单事件项中计入三次，又在两两交集项中减去三次，因此还要补回一次；只属于其中两个事件的样本点则计入两次、减去一次，已经恰好保留一次。

\begin{theorem}[有限个事件的加法定理]
\label{thm:study-inclusion-exclusion}
设 $n\ge1$ 为整数，$A_1,\ldots,A_n$ 是同一概率空间中的任意事件，则
\begin{equation}
  \mathbb P\!\left(\bigcup_{i=1}^{n}A_i\right)
    =\sum_{k=1}^{n}(-1)^{k-1}
      \sum_{1\le i_1<\cdots<i_k\le n}
      \mathbb P(A_{i_1}\cap\cdots\cap A_{i_k}).
  \label{eq:study-inclusion-exclusion-general}
\end{equation}
内层求和遍历所有由 $k$ 个不同下标组成的组合，每组只取一次，共有 $\binom nk$ 项。因此，右端依次为单事件概率之和，减去两两交集概率之和，加上三者交集概率之和，直到符号为 $(-1)^{n-1}$ 的全体事件交集项。公式不要求各事件互斥，也不要求相互独立。
\end{theorem}

\begin{proof}
用数学归纳法证明。$n=1$ 时两端均为 $\mathbb P(A_1)$；$n=2$ 时就是已经证明的式~\eqref{eq:study-inclusion-exclusion-two}。设 $n\ge3$，假设公式对任意 $n-1$ 个事件成立，令
\[
  U=\bigcup_{i=1}^{n-1}A_i,\qquad B_i=A_i\cap A_n\quad(1\le i\le n-1).
\]
由分配律，$U\cap A_n=\bigcup_{i=1}^{n-1}B_i$。再用二事件公式，得
\begin{equation}
  \mathbb P\!\left(\bigcup_{i=1}^{n}A_i\right)
    =\mathbb P(U)+\mathbb P(A_n)
      -\mathbb P\!\left(\bigcup_{i=1}^{n-1}B_i\right).
  \label{eq:study-inclusion-exclusion-induction}
\end{equation}
分别对 $A_1,\ldots,A_{n-1}$ 和 $B_1,\ldots,B_{n-1}$ 应用归纳假设。$\mathbb P(U)$ 的展开给出所有不含下标 $n$ 的交集项，$\mathbb P(A_n)$ 给出只含下标 $n$ 的一阶项。最后一项中，每个 $k$ 重交集满足
\[
  B_{i_1}\cap\cdots\cap B_{i_k}
    =A_{i_1}\cap\cdots\cap A_{i_k}\cap A_n,
  \qquad 1\le k\le n-1.
\]
它对应一个含下标 $n$ 的 $k+1$ 重交集。归纳公式中的系数为 $(-1)^{k-1}$，再乘以式~\eqref{eq:study-inclusion-exclusion-induction} 中的负号，变为 $(-1)^k$，恰好是 $k+1$ 重交集应有的符号。所有非空下标组合由此不重不漏地出现，合并便得到式~\eqref{eq:study-inclusion-exclusion-general}，归纳完成。
\end{proof}

容斥原理与对立事件法可以结合使用。由德摩根律~\eqref{eq:study-de-morgan-finite}，“所有事件都不发生”是“至少一个事件发生”的对立事件，因而
\begin{equation}
  \begin{aligned}
    \mathbb P\!\left(\bigcap_{i=1}^{n}A_i^{\mathrm c}\right)
      &=1-\mathbb P\!\left(\bigcup_{i=1}^{n}A_i\right)\\
      &=1-\sum_{k=1}^{n}(-1)^{k-1}
        \sum_{1\le i_1<\cdots<i_k\le n}
        \mathbb P(A_{i_1}\cap\cdots\cap A_{i_k}).
  \end{aligned}
  \label{eq:study-inclusion-exclusion-complement}
\end{equation}

\begin{example}
\label{ex:study-inclusion-exclusion-complement}
设 $A,B,C$ 是同一概率空间中的三个事件，已知
\[
  \mathbb P(A)=\mathbb P(B)=\frac14,\qquad \mathbb P(C)=\frac25,
  \qquad \mathbb P(A\cap B)=0,\qquad
  \mathbb P(A\cap C)=\mathbb P(B\cap C)=\frac16.
\]
求三个事件都不发生的概率 $\mathbb P(A^{\mathrm c}\cap B^{\mathrm c}\cap C^{\mathrm c})$。
\end{example}

\begin{solve}
先确定题目没有直接给出的三者交集概率。由于 $A\cap B\cap C\subseteq A\cap B$，由概率的单调性，
\[
  0\le\mathbb P(A\cap B\cap C)\le\mathbb P(A\cap B)=0,
\]
所以 $\mathbb P(A\cap B\cap C)=0$。结合对立事件关系和三事件公式~\eqref{eq:study-inclusion-exclusion-three}，得到
\begin{equation}
  \begin{aligned}
    \mathbb P(A^{\mathrm c}\cap B^{\mathrm c}\cap C^{\mathrm c})
      &=1-\mathbb P(A\cup B\cup C)\\
      &=1-\left[\frac14+\frac14+\frac25-0-\frac16-\frac16+0\right]\\
      &=1-\frac{17}{30}=\boxed{\frac{13}{30}}.
  \end{aligned}
  \label{eq:study-inclusion-exclusion-example-answer}
\end{equation}
条件 $\mathbb P(A\cap B)=0$ 不能在一般概率空间中推出 $A\cap B=\varnothing$，所以不宜直接把 $A,B$ 称为互斥事件；它已经足以推出三者交集的概率为零。此外，本题中 $\mathbb P(A)\mathbb P(B)=1/16\ne\mathbb P(A\cap B)$，故 $A,B$ 并不独立，也不能用三个对立事件概率的乘积代替所求概率。
\end{solve}

\section{条件概率}
\label{sec:study-conditional-probability}

\subsection{条件概率及其概率结构}
\label{subsec:study-conditional-definition}

当已知事件 $B$ 发生时，计算事件 $A$ 发生的概率，需要把考虑范围限制在 $B$ 内，并以 $B$ 的概率作为新的归一化基准。在这一范围内，$A$ 发生对应的部分是 $A\cap B$，因此应比较 $A\cap B$ 的概率与 $B$ 的概率。

\begin{definition}
\label{def:study-conditional-probability}
设 $A,B$ 是同一概率空间中的两个事件，且 $\mathbb P(B)>0$。称
\begin{equation}
  \mathbb P(A\mid B)
    =\frac{\mathbb P(A\cap B)}{\mathbb P(B)}
  \label{eq:study-conditional-probability}
\end{equation}
为在事件 $B$ 发生的条件下事件 $A$ 发生的\emph{条件概率（conditional probability）}。竖线右侧的 $B$ 表示已知条件，左侧的 $A$ 表示所考察的事件。
\end{definition}

定义中的分母用于将条件事件 $B$ 的概率归一化为 $1$，即 $\mathbb P(B\mid B)=1$。又因为 $A\cap B\subseteq B$，所以 $0\le\mathbb P(A\cap B)\le\mathbb P(B)$，除以正数 $\mathbb P(B)$ 得 $0\le\mathbb P(A\mid B)\le1$。特别地，在有限等可能的古典概型中，若 $B\ne\varnothing$，由式~\eqref{eq:study-classical-probability} 可得
\[
  \mathbb P(A\mid B)
    =\frac{|A\cap B|/|\Omega|}{|B|/|\Omega|}
    =\frac{|A\cap B|}{|B|},
\]
其中 $\Omega$ 为原样本空间，$|\cdot|$ 表示集合中元素的个数。这说明在古典概型中，可以只在 $B$ 包含的样本点中重新计数；一般概率模型则应使用概率之比，不能直接以样本点数之比代替。

条件 $\mathbb P(B)>0$ 不可省略。当 $\mathbb P(B)=0$ 时，由单调性也有 $\mathbb P(A\cap B)=0$，式~\eqref{eq:study-conditional-probability} 的右端成为没有定义的 $0/0$，因此本定义不能用于概率为零的条件事件。这并不意味着 $B$ 必为空集。还应注意，$\mathbb P(A\mid B)$ 与 $\mathbb P(A\cap B)$ 含义不同：前者是在已知 $B$ 发生后计算 $A$ 的概率，后者是在原概率模型中计算两者同时发生的概率。当 $\mathbb P(A)>0$ 时，$\mathbb P(B\mid A)$ 则交换了条件与所考察事件，通常也不等于 $\mathbb P(A\mid B)$；条件关系本身不表示时间先后或因果关系。

\begin{example}
\label{ex:study-conditional-probability-die}
掷一枚均匀骰子，令 $A$ 表示“点数大于 $3$”，$B$ 表示“点数为偶数”。已知出现偶数，求点数大于 $3$ 的条件概率。
\end{example}

\begin{solve}
样本空间为 $\Omega=\{1,2,3,4,5,6\}$，各样本点等可能。由于
\[
  A=\{4,5,6\},\qquad B=\{2,4,6\},\qquad A\cap B=\{4,6\},
\]
有 $\mathbb P(B)=3/6>0$、$\mathbb P(A\cap B)=2/6$，故由定义得
\[
  \mathbb P(A\mid B)=\frac{2/6}{3/6}=\frac23.
\]
也可以直接在条件事件 $B$ 的三个等可能样本点 $2,4,6$ 中计数，其中 $4,6$ 满足 $A$，所以结果为 $2/3$。原来的无条件概率为 $\mathbb P(A)=3/6=1/2$，而已知 $B$ 发生后的条件概率为 $2/3$，说明获得条件信息后，同一事件的概率可能改变。
\end{solve}

对固定的 $B$，令 $\mathbb P_B(A)=\mathbb P(A\mid B)$。这个函数仍满足概率公理：非负性与规范性由定义得到；若 $A_n$ 两两互斥，则 $A_n\cap B$ 也两两互斥，故
\[
 \mathbb P_B\!\left(\bigcup_n A_n\right)
 =\frac{\sum_n\mathbb P(A_n\cap B)}{\mathbb P(B)}
 =\sum_n\mathbb P_B(A_n).
\]
因此，补事件公式、加法定理等均可在同一个条件 $B$ 下使用，例如 $\mathbb P(A^{\mathrm c}\mid B)=1-\mathbb P(A\mid B)$。这里不能把条件事件也随意取补：$\mathbb P(A\mid B^{\mathrm c})$ 使用的是另一个概率模型。

% 教材例 1.16：三张双面卡片。
\begin{example}
\label{ex:study-double-sided-cards}
三张卡片分别为双面有圈、圈点各一面、双面有点。等可能抽一张，再等可能决定哪面朝上。已知朝上一面有圈，求另一面也有圈的概率。
\end{example}

\begin{solve}
把“卡片编号与朝上的具体一面”一起记录，共有六个等可能样本点。其中三个使上面有圈：双圈卡的两个面各对应一个样本点，圈点卡的圈面对应一个样本点。前两者的另一面也有圈，故所求概率为 $2/3$。在观察到圈后，双圈卡与圈点卡并不等可能，因为前者有两个能够产生该观察的面。条件化后必须保留观察机制带来的权重。
\end{solve}

\subsection{乘法公式与逐步条件化}
\label{subsec:study-multiplication}

\begin{theorem}[概率的乘法公式]
\label{thm:study-multiplication}
若 $\mathbb P(B)>0$，则 $\mathbb P(A\cap B)=\mathbb P(B)\mathbb P(A\mid B)$。更一般地，设 $n\ge2$，若 $\mathbb P(\bigcap_{i=1}^{n-1}A_i)>0$，则
\begin{equation}
 \mathbb P\!\left(\bigcap_{i=1}^{n}A_i\right)
 =\mathbb P(A_1)\prod_{k=2}^{n}
 \mathbb P\!\left(A_k\,\middle|\,\bigcap_{i=1}^{k-1}A_i\right).
 \label{eq:study-chain-rule}
\end{equation}
\end{theorem}

\begin{proof}
二事件公式就是条件概率定义的移项。一般情形中，假设保证每个前缀交事件均有正概率；把右侧各条件概率写成“相邻交事件概率之比”，分子分母依次约去，即得到左端。也可以逐次对二事件乘法公式归纳。
\end{proof}

公式不要求独立，但后一步须以此前所有指定事件同时发生为条件。例如 $10$ 件产品中有 $3$ 件一等品，不放回连续抽两件，则
\[
 \mathbb P(\text{两件均为一等品})=\frac3{10}\frac2{9},\qquad
 \mathbb P(\text{第二件一等}\mid\text{第一件一等})=\frac2{9}.
\]
第二次的无条件一等品概率仍为 $3/10$；条件概率改变是因为已知第一件的类别后，剩余组成改变了。

\subsection{全概率公式与划分的选择}
\label{subsec:study-total-probability}

\begin{definition}
\label{def:study-partition}
若事件 $B_1,\ldots,B_m$ 两两互斥、并为 $\Omega$ 且每个概率均为正，称它们构成\emph{完备事件群}，也称样本空间的一个\emph{划分（partition）}。一般集合划分不要求正概率；这里加上正概率，是为了使各条件概率均有定义。
\end{definition}

\begin{theorem}[全概率公式]
\label{thm:study-total-probability}
若 $B_1,\ldots,B_m$ 构成上述完备事件群，则对任意事件 $A$，
\begin{equation}
 \mathbb P(A)=\sum_{i=1}^{m}\mathbb P(B_i)\mathbb P(A\mid B_i).
 \label{eq:study-total-probability}
\end{equation}
若划分为可数无穷个正概率事件，公式同样成立，只需把有限和改为无穷级数。
\end{theorem}

\begin{proof}
$A=\bigcup_i(A\cap B_i)$ 且各部分互斥，对各部分应用乘法公式，再用可加性求和即可。
\end{proof}

全概率公式把总体概率写成各层条件概率的加权平均，权重是各层本来的概率。划分应做到不重不漏，并让条件模型比原问题容易处理；各层不必代表因果意义上的“原因”。如果某个划分事件概率为零，应直接忽略其交事件的零贡献，不能把未定义的条件概率乘以零来替代说明。

% 教材例 1.19、1.23：品牌、返修与后验概率。
\begin{example}
\label{ex:study-refrigerator-total}
商场售出三种品牌的冰箱，其占比分别为 $0.5,0.3,0.2$，相应返修率为 $0.002,0.003,0.004$。求随机选一台售出冰箱的返修概率。
\end{example}

\begin{solve}
令 $B_i$ 表示品牌为 $i$，$A$ 表示需要返修。品牌事件构成划分，所以
\[
 \mathbb P(A)=0.5\times0.002+0.3\times0.003+0.2\times0.004
 =0.0027.
\]
这三个返修率是给定品牌后的条件概率，不能把它们直接作不加权平均；销售占比与维修样本中的品牌占比也不同。
\end{solve}

% 教材例 1.20：随机化回答。
\begin{example}
\label{ex:study-randomized-response}
调查敏感属性时，每位受访者私下掷公平硬币，正面时回答一个已知总体“是”比例为 $\alpha$ 的普通问题，反面时回答是否具有敏感属性；只报告“是”或“否”。假定如实作答且硬币与个人属性独立，如何由回答估计敏感属性的总体比例 $p$？
\end{example}

\begin{solve}
设 $t$ 为随机受访者回答“是”的概率。按硬币结果划分，有 $t=\alpha/2+p/2$，所以 $p=2t-\alpha$。若有 $n$ 个独立、按同一总体抽取的回答，其中 $y$ 个为“是”，可用 $\widehat p=2y/n-\alpha$ 作频率估计。这是样本估计，并不恒等于未知总体比例；有限样本可能得到区间 $[0,1]$ 外的值，此时说明直接线性估计受抽样波动影响，不能把越界数值解释为概率。该方法用随机化保护单个回答的含义，推导则依赖题设的作答与抽样假定。
\end{solve}

% 教材例 1.22：波利亚罐子模型。
\begin{example}
\label{ex:study-polya-urn}
罐中初有 $a$ 个白球、$b$ 个黑球，其中 $a,b$ 为正整数。每次均匀抽一个球，放回并加入 $c$ 个同色球，$c\ge0$ 为整数。证明任意第 $n$ 次抽到白球的概率均为 $a/(a+b)$。
\end{example}

\begin{proof}
对次数 $n$ 作归纳，并要求归纳结论对任意正整数初始球数成立。$n=1$ 时结论由均匀抽取给出。设结论对 $n-1$ 次成立；按第一次抽到白球或黑球划分，后续过程分别相当于从初始组成 $(a+c,b)$ 或 $(a,b+c)$ 开始的同类过程。故
\[
 \mathbb P(\text{第 }n\text{ 次白球})
 =\frac{a}{a+b}\frac{a+c}{a+b+c}
  +\frac{b}{a+b}\frac{a}{a+b+c}
 =\frac{a}{a+b}.
\]
这里必须让归纳假设覆盖改变后的初始组成，才能应用于条件模型。
\end{proof}

当 $c>0$ 时，第二次在第一次为白球条件下的白球概率为 $(a+c)/(a+b+c)>a/(a+b)$。所以各次的边缘概率相同，并不意味着各次独立。若仅按“第 $n-1$ 次为白球”划分，罐内的完整组成仍不确定，不能直接把它当作一次抽球时的已知组成。

\subsection{贝叶斯公式与信息更新}
\label{subsec:study-bayes-formula}

\begin{theorem}[贝叶斯公式]
\label{thm:study-bayes}
设 $B_1,\ldots,B_m$ 构成完备事件群，且 $\mathbb P(A)>0$，则
\begin{equation}
 \mathbb P(B_j\mid A)
 =\frac{\mathbb P(B_j)\mathbb P(A\mid B_j)}
 {\sum_{i=1}^{m}\mathbb P(B_i)\mathbb P(A\mid B_i)},
 \qquad j=1,\ldots,m.
 \label{eq:study-bayes-formula}
\end{equation}
\end{theorem}

\begin{proof}
条件概率定义给出 $\mathbb P(B_j\mid A)=\mathbb P(A\cap B_j)/\mathbb P(A)$；对分子使用乘法公式，对分母使用全概率公式即得结论。
\end{proof}

\emph{先验概率（prior probability）} $\mathbb P(B_j)$ 表示观察 $A$ 前各状态的权重，$\mathbb P(A\mid B_j)$ 衡量在相应状态下出现这项证据的可能性，\emph{后验概率（posterior probability）} $\mathbb P(B_j\mid A)$ 则是观察后的权重。更新过程是“先验乘以证据在该状态下的概率，再统一归一化”，并不从概率条件关系自动推出因果关系。例~\ref{ex:study-refrigerator-total} 中，已知冰箱返修后，三品牌的概率分别为
\[
 \frac{0.5\times0.002}{0.0027}=\frac{10}{27},\qquad
 \frac{0.3\times0.003}{0.0027}=\frac{9}{27},\qquad
 \frac{0.2\times0.004}{0.0027}=\frac{8}{27}.
\]

% 教材例 1.24：仅讨论给定参数的检测模型。
\begin{example}
\label{ex:study-screening-bayes}
在一个假设模型中，受检者具有某属性的概率为 $0.005$；具有该属性时检测阳性的概率为 $0.95$，不具有时阴性的概率也为 $0.95$。求观察到阳性后具有该属性的概率。
\end{example}

\begin{solve}
令 $D$ 表示具有属性，$T$ 表示检测阳性。已知 $\mathbb P(T\mid D)=0.95$、$\mathbb P(T\mid D^{\mathrm c})=0.05$，由贝叶斯公式，
\begin{equation}
 \mathbb P(D\mid T)
 =\frac{0.95\times0.005}{0.95\times0.005+0.05\times0.995}
 =\frac{19}{218}\approx0.0872.
 \label{eq:study-screening-posterior}
\end{equation}
这里的 $8.72\%$ 是给定阳性后的属性概率；$95\%$ 是给定真实类别后的检测性能，二者条件不同。按模型权重看，$20000$ 个受检者中有属性者的理论人数为 $100$，其中阳性 $95$；无属性者理论人数为 $19900$，其中阳性 $995$，故阳性组中的比例为 $95/(95+995)$。这些是模型下的期望计数，不是有限样本必然出现的人数。

若再作一次检测，只有额外假定两次结果在给定 $D$ 后独立、在给定 $D^{\mathrm c}$ 后也独立，且各次条件性能不变，才能计算
\[
 \mathbb P(D\mid T_1\cap T_2)
 =\frac{0.005\times0.95^2}
 {0.005\times0.95^2+0.995\times0.05^2}
 =\frac{361}{560}\approx0.6446.
\]
共同的真实类别本身会使两次结果在总体上相关；不能仅说“重复检测”就默认上述条件独立性。
\end{solve}

\section{独立性}
\label{sec:study-independence}

\subsection{两个事件的独立性}
\label{subsec:study-two-event-independence}

\begin{definition}
\label{def:study-event-independence}
若事件 $A,B$ 满足
\begin{equation}
 \mathbb P(A\cap B)=\mathbb P(A)\mathbb P(B),
 \label{eq:study-event-independence}
\end{equation}
则称它们\emph{相互独立（independent）}。当 $\mathbb P(B)>0$ 时，这等价于 $\mathbb P(A\mid B)=\mathbb P(A)$；当 $\mathbb P(A)>0$ 时，也等价于 $\mathbb P(B\mid A)=\mathbb P(B)$。
\end{definition}

上述等价性由条件概率定义直接得到。独立性是对称的概率关系；“获知一个事件发生不改变另一个的概率”是它在条件概率有定义时的含义。乘积定义也适用于概率为零的事件：若 $\mathbb P(B)=0$，单调性保证 $\mathbb P(A\cap B)=0$，故 $B$ 与任意事件独立，尽管 $\mathbb P(A\mid B)$ 在本章定义下没有意义。

互斥与独立不能混淆。若 $A\cap B=\varnothing$ 且两个概率都为正，则交事件概率为零、概率乘积为正，二者不独立；若其中一个概率为零，则互斥与独立可以同时成立。判断是否独立，应以给定概率模型为准；时间先后、物理距离或相同的边缘概率均不能单独证明独立。

\begin{proposition}
\label{prop:study-independence-complements}
$A,B$ 独立当且仅当 $A,B^{\mathrm c}$ 独立，也当且仅当 $A^{\mathrm c},B$ 独立或 $A^{\mathrm c},B^{\mathrm c}$ 独立。
\end{proposition}

\begin{proof}
若 $A,B$ 独立，由可减性，
\[
 \mathbb P(A\cap B^{\mathrm c})
 =\mathbb P(A)-\mathbb P(A\cap B)
 =\mathbb P(A)[1-\mathbb P(B)]
 =\mathbb P(A)\mathbb P(B^{\mathrm c}).
\]
交换 $A,B$ 可对第一项取补，再重复此步骤可对两项取补。由于取补两次回到原事件，每种蕴含也可反向应用，故各条件等价。
\end{proof}

\subsection{相互独立、两两独立与条件独立}
\label{subsec:study-mutual-independence}

\begin{definition}
\label{def:study-mutual-independence}
设 $n\ge2$。若对任意 $k=2,\ldots,n$ 及任意下标 $1\le i_1<\cdots<i_k\le n$ 都有
\begin{equation}
 \mathbb P(A_{i_1}\cap\cdots\cap A_{i_k})
 =\prod_{j=1}^{k}\mathbb P(A_{i_j}),
 \label{eq:study-mutual-independence}
\end{equation}
则称 $A_1,\ldots,A_n$ \emph{相互独立（mutually independent）}。若只要求 $k=2$ 时的等式，称为\emph{两两独立（pairwise independent）}。事件列称为独立事件列，是指其中每个有限子族都相互独立。
\end{definition}

相互独立要求 $2^n-n-1$ 个非平凡乘积等式，既要检验全部 $n$ 个事件的交，也要检验每个较小子族。它等价于：对所有选择 $E_i\in\{A_i,A_i^{\mathrm c}\}$，均有
\begin{equation}
 \mathbb P\!\left(\bigcap_{i=1}^{n}E_i\right)
 =\prod_{i=1}^{n}\mathbb P(E_i).
 \label{eq:study-independent-sign-patterns}
\end{equation}
从相互独立出发，对某项取补时可将交事件拆成该项发生与不发生两部分，用可减性得到相应乘积，再逐项取补。反过来，将未指定下标的发生与不发生情形全部求和，就得到任意子族的乘积公式。此论证还说明：由相互独立事件的互不重叠下标组各自构造的事件仍相互独立，因为每个组内事件都是若干互斥的发生与否组合之并，对这些组合求和时概率乘积可以按组分解。

% 教材例 1.29。
\begin{example}
\label{ex:study-pairwise-not-mutual}
等可能抽取一个标记为 $\{1\},\{2\},\{3\},\{1,2,3\}$ 的球。令 $A_i$ 表示抽到的球含数字 $i$，判断这三个事件的独立性。
\end{example}

\begin{solve}
每个 $A_i$ 含两个样本点，故 $\mathbb P(A_i)=1/2$；任意两个的交只含标记 $\{1,2,3\}$ 的球，故交概率为 $1/4=(1/2)^2$，三个事件两两独立。但三者交概率为 $1/4\ne1/8=\prod_{i=1}^{3}\mathbb P(A_i)$，因此不相互独立。只检查二重交或只检查全体交，都不足以一般地判断相互独立。
\end{solve}

若 $\mathbb P(C)>0$，并且
\[
 \mathbb P(A\cap B\mid C)=\mathbb P(A\mid C)\mathbb P(B\mid C),
\]
称 $A,B$ 在给定 $C$ 下\emph{条件独立（conditionally independent）}。这是条件概率空间中的独立性，不能直接推出无条件独立。相反，无条件独立也不保证附加任意条件后仍独立。例如独立掷两次公平硬币，第一次正面与第二次正面独立；给定“两次相同”后，这两个事件的条件概率均为 $1/2$，交的条件概率也为 $1/2$，不再等于乘积 $1/4$。

\subsection{独立试验与可靠性计算}
\label{subsec:study-independent-applications}

若 $A_1,\ldots,A_n$ 相互独立，记 $p_i=\mathbb P(A_i)$，则由对立事件法和式~\eqref{eq:study-independent-sign-patterns}，
\begin{equation}
 \mathbb P\!\left(\bigcup_{i=1}^{n}A_i\right)
 =1-\prod_{i=1}^{n}(1-p_i).
 \label{eq:study-independent-at-least-one}
\end{equation}
若每次概率均为固定的 $p>0$，至少一次发生的概率为 $1-(1-p)^n\to1$；更一般地，若 $p_i\ge p_0>0$，未发生的概率不超过 $(1-p_0)^n\to0$。这才是在相应假设下“小概率事件在大量重复机会中可能出现”的数学依据；没有独立性或其他充分条件时不能照搬。一次观测中小概率事件仍可能发生，概率小也不能单独作为逻辑上否定模型的证明。

% 教材例 1.13：无穷次重复掷骰子。
\begin{example}
\label{ex:study-die-infinitely-often}
独立重复掷均匀骰子，证明出现 $1$ 点的次数以概率 $1$ 为无穷多。
\end{example}

\begin{proof}
令 $B_n$ 表示第 $n$ 次之后不再出现 $1$ 点，$C_{n,m}$ 表示第 $n+1,\ldots,n+m$ 次都不出现 $1$ 点。对固定 $n$，有 $C_{n,m}\downarrow B_n$，由独立性与上连续性，
\[
 \mathbb P(B_n)=\lim_{m\to\infty}(5/6)^m=0.
\]
仅出现有限次 $1$ 点等价于存在某个 $n$ 使 $B_n$ 发生，故该事件为 $\bigcup_{n=1}^{\infty}B_n$，由次可加性其概率为零。取补即得结论。这是概率为一的结论，而不是说结果空间中不存在永远不出现 $1$ 点的序列。
\end{proof}

% 教材例 1.25：补明全部成功事件相互独立的模型假设。
若 $10$ 人猜 $4$ 个字谜，假定“第 $i$ 人猜出第 $j$ 个谜”这 $40$ 个事件相互独立，且概率均为 $0.2$，则每个谜被至少一人猜出的概率为 $1-0.8^{10}$，四谜全被猜出的概率为 $(1-0.8^{10})^4\approx0.6349$。这里只给出每人的单谜成功率不足以确定答案；不同人或不同谜之间的相关性可能改变结果。

% 教材例 1.28：用事件表达连接关系，无须依赖电路图。
\begin{example}
\label{ex:study-reliability-four-components}
四个元件各自正常工作的概率均为 $p\in[0,1]$，工作状态相互独立。比较“两条各含两个串联元件的支路再并联”与“两个各含两个并联元件的模块再串联”的可靠性。
\end{example}

\begin{solve}
令 $A_i$ 表示第 $i$ 个元件正常。第一种系统正常事件为 $(A_1\cap A_2)\cup(A_3\cap A_4)$，第二种为 $(A_1\cup A_3)\cap(A_2\cup A_4)$。不同支路或模块使用互不重叠的独立元件，故可按组使用独立性，得到
\[
 R_1=1-(1-p^2)^2=p^2(2-p^2),\qquad
 R_2=[1-(1-p)^2]^2=p^2(2-p)^2.
\]
二者之差为 $R_2-R_1=2p^2(1-p)^2\ge0$，故在这些假设下第二种连接至少同样可靠；$0<p<1$ 时严格更可靠。若两条路径共用元件，则相关路径事件不能直接按独立事件处理。
\end{solve}

\section{扩展进阶：求概率的一些方法}
\label{sec:study-probability-methods}
% 教材将本节标注为带星号的扩展进阶内容。

\subsection{选取样本空间与利用对称性}
\label{subsec:study-probability-space-choice}

同一问题可以在不同精度的样本空间上建模，只要对应的概率权重一致，结果就应一致。应先决定记录顺序与否、对象是否可辨，再让分子分母使用同一套计数标准。若一个细模型中的所有样本点等可能，合并后的每一类也恰含同样多的样本点，则这些类仍等可能；若大小不同，就必须保留类的不同权重。

% 教材例 1.30。
\begin{example}
\label{ex:study-kth-black-ball}
袋中有 $a$ 个黑球、$b$ 个白球，其中 $a,b$ 为正整数，均匀不放回依次摸完。求第 $k$ 个球为黑球的概率，$1\le k\le a+b$。
\end{example}

\begin{solve}
给每个球编号，所有 $(a+b)!$ 个完整排列等可能。第 $k$ 位放黑球有 $a$ 种选择，其余球有 $(a+b-1)!$ 个排列，所以概率为 $a/(a+b)$。也可以只记录黑球出现的位置：每个含 $a$ 个位置的子集对应 $a!b!$ 个编号排列，故 $\binom{a+b}{a}$ 个位置子集等可能；包含第 $k$ 位的有 $\binom{a+b-1}{a-1}$ 个，比值仍为 $a/(a+b)$。两个解法说明了无条件的次序对称性，不表示在知道前面颜色后条件概率仍不变。
\end{solve}

% 教材例 1.8、1.18、1.21：合并重复的连绳问题。
\begin{example}
\label{ex:study-rope-pairings}
将 $n\ge1$ 根可辨短绳的 $2n$ 个端头等可能地两两配对连接，求恰成 $n$ 个圈与恰成一个圈的概率。
\end{example}

\begin{solve}
固定一个端头，它有 $2n-1$ 个配对选择；继续对剩余端头配对，所有配对数为
\[
 (2n-1)!!=(2n-1)(2n-3)\cdots1
 =\frac{(2n)!}{2^n n!}.
\]
最后一个公式也可由端头排列得到：每个配对结果被重复记录 $2^n n!$ 次，分别对应各对内部交换与配对组的排列，因此不能把 $(2n)!$ 直接称为不同连接结果数，除非始终把排列本身作为等可能样本点。

恰成 $n$ 个圈只有每根绳自接这一种配对，故概率为 $1/(2n-1)!!$。也可依次条件化：已知前 $k-1$ 根自接后，下一根自接的概率为 $1/[2(n-k+1)-1]$，由乘法公式相乘得到同一结果。恰成一个圈时，从第一根绳的某个固定端头出发，先沿它走到另一端，再按任意次序、任意方向串接其余 $n-1$ 根绳，最后接回起点，共有 $2^{n-1}(n-1)!$ 种配对；起点与第一根绳的行进方向已经固定，每个单圈恰被计数一次。因此
\begin{equation}
 \mathbb P(n\text{ 个圈})=\frac1{(2n-1)!!},\qquad
 \mathbb P(1\text{ 个圈})=\frac{2^{n-1}(n-1)!}{(2n-1)!!}.
 \label{eq:study-rope-probabilities}
\end{equation}
$n=1$ 时两个事件相同，公式均给出 $1$。
\end{solve}

\subsection{条件化与递推}
\label{subsec:study-probability-recursion}

递推法先选一个能够描述后续机制的状态，再按较早阶段的信息划分，用全概率公式联系相邻状态。状态必须保留后续计算需要的信息；例如波利亚模型中的当前球数比“上一次是什么颜色”包含更多信息。递推式、初值和解出的数列共同构成完整答案。

% 教材例 1.31。
\begin{example}
\label{ex:study-dice-turn-recursion}
甲先独立掷均匀骰子，每次掷出 $1$ 点后换另一人掷，否则仍由原来的人掷。求第 $n$ 次由甲掷的概率 $p_n$。
\end{example}

\begin{solve}
已知上一次由甲掷，下一次仍由甲掷的概率为 $5/6$；上一次由乙掷，下一次由甲掷的概率为 $1/6$。所以
\[
 p_1=1,\qquad p_n=\frac56p_{n-1}+\frac16(1-p_{n-1})
 =\frac16+\frac23p_{n-1}\quad(n\ge2).
\]
先求不动点 $1/2$，再减去它，得 $p_n-1/2=(2/3)(p_{n-1}-1/2)$。于是
\begin{equation}
 p_n=\frac12+\frac12\left(\frac23\right)^{n-1}\longrightarrow\frac12.
 \label{eq:study-dice-turn-recursion}
\end{equation}
甲先掷带来的优势逐渐衰减；这里求的是每一时刻的边缘概率，并未声称各时刻的掷骰者相互独立。
\end{solve}

\subsection{容斥与配对问题}
\label{subsec:study-matching-problem}

% 教材例 1.32、1.33：帽子与通知装封为同一配对模型。
\begin{example}
\label{ex:study-derangements}
$n$ 人等可能地各取一顶属于这 $n$ 人的帽子，求恰有 $k$ 人取对以及至少 $m$ 人取对的概率，其中 $n\ge1$，$0\le k,m\le n$。
\end{example}

\begin{solve}
样本空间是 $n!$ 个排列。先记 $D_s$ 为 $s$ 个对象全部配错的\emph{错排数（derangement number）}，约定 $D_0=1$。对 $s\ge1$，指定任意 $j$ 个人取对后，剩余帽子可任意安排，故这些 $j$ 个事件的交概率为 $(s-j)!/s!$。由容斥原理，无人取对的概率为
\[
 \frac{D_s}{s!}=\sum_{j=0}^{s}\frac{(-1)^j}{j!}.
\]
设 $K$ 为取对人数。恰有 $k$ 人取对，先选这 $k$ 人，再让其余 $n-k$ 人全部配错，得到
\begin{equation}
 \mathbb P(K=k)=\frac{\binom nk D_{n-k}}{n!}
 =\frac1{k!}\sum_{j=0}^{n-k}\frac{(-1)^j}{j!},
 \qquad k=0,\ldots,n.
 \label{eq:study-matching-exact}
\end{equation}
特别地，$D_1=0$，所以不可能恰有 $n-1$ 人取对；$k=n$ 时按空对象约定得到 $1/n!$。不同的 $k$ 对应互斥事件，因而
\[
 \mathbb P(K\ge m)=\sum_{k=m}^{n}\frac1{k!}
       \sum_{j=0}^{n-k}\frac{(-1)^j}{j!}.
\]
对每个固定 $k$，由指数函数的幂级数，$n\to\infty$ 时 $\mathbb P(K=k)\to e^{-1}/k!$。将帽子改为通知与信封，概率模型完全相同；这种对应可以复用已有结论。
\end{solve}

\section{扩展阅读1：贝叶斯公式和垃圾邮件识别}
\label{sec:study-bayes-spam}

邮件分类把“根据已观察到的特征更新类别概率”具体化。令 $S$ 表示垃圾邮件，$W$ 表示邮件中出现某关键词，记 $\pi=\mathbb P(S)$、$\theta_1=\mathbb P(W\mid S)$、$\theta_0=\mathbb P(W\mid S^{\mathrm c})$。当分母为正时，贝叶斯公式给出
\begin{equation}
 \mathbb P(S\mid W)
 =\frac{\pi\theta_1}{\pi\theta_1+(1-\pi)\theta_0}.
 \label{eq:study-spam-single-word}
\end{equation}
教材用假设数值 $\pi=0.5$、$\theta_1=0.05$、$\theta_0=0.001$ 得到 $50/51\approx0.9804$。这里各数字是演示参数，不是关于当前邮件环境的实测结论；特征似然相同时，类别先验不同也会改变后验。

训练时可以使用带类别标记的历史邮件：若 $N_1$ 封垃圾邮件中有 $N_{1,W}$ 封包含该词，就以 $N_{1,W}/N_1$ 估计 $\theta_1$，正常邮件同理。这些频率是对总体条件概率的估计，不是定义上的真值。为了避免有限样本中未出现某词就将其概率置零，可对二元词特征采用平滑，例如给定 $\lambda>0$，取
\[
 \widehat\theta_1=\frac{N_{1,W}+\lambda}{N_1+2\lambda},\qquad
 \widehat\theta_0=\frac{N_{0,W}+\lambda}{N_0+2\lambda}.
\]
这是对教材“避免零概率”做法的一种明确实现，$\lambda$ 是人为选择的平滑参数。若训练集人为各取相同数量的两类邮件，其类别比例不一定等于实际应用中的 $\pi$，应分别估计或指定。

使用 $d$ 个关键词时，令二元特征 $X_j$ 表示第 $j$ 个词是否出现，观测值为 $x_j\in\{0,1\}$。\emph{朴素贝叶斯（naive Bayes）}的核心假设是：\emph{在给定邮件类别后，各词特征条件独立}，即对 $C\in\{S,S^{\mathrm c}\}$，
\begin{equation}
 \mathbb P(X_1=x_1,\ldots,X_d=x_d\mid C)
 =\prod_{j=1}^{d}\mathbb P(X_j=x_j\mid C).
 \label{eq:study-naive-bayes-assumption}
\end{equation}
它不是说词语在所有邮件中无条件独立。记 $\theta_{1j}=\mathbb P(X_j=1\mid S)$、$\theta_{0j}=\mathbb P(X_j=1\mid S^{\mathrm c})$，则两类的联合特征概率分别为
\[
 L_1(x)=\prod_{j=1}^{d}\theta_{1j}^{x_j}(1-\theta_{1j})^{1-x_j},\qquad
 L_0(x)=\prod_{j=1}^{d}\theta_{0j}^{x_j}(1-\theta_{0j})^{1-x_j},
\]
并有 $\mathbb P(S\mid X=x)=\pi L_1(x)/[\pi L_1(x)+(1-\pi)L_0(x)]$。上式按 $0<\theta_{cj}<1$ 书写；端点参数可直接逐个取相应的二元概率，避免 $0^0$ 的记号歧义。若仅以“这几个词全出现”为证据，则只取各出现概率的乘积，不额外引入未观察词的条件。

教材中三个词全出现、每个词的两类条件出现概率分别为 $0.05$ 与 $0.001$、先验各半的示例，在上述条件独立假设下给出
\[
 \mathbb P(S\mid X_1=X_2=X_3=1)
 =\frac{0.05^3}{0.05^3+0.001^3}\approx0.999992.
\]
这是模型内的后验；相近含义的词可能相关，直接相乘可能重复计算同一证据，因而“特征更多”不保证实际分类更可靠。分类还需选阈值：将后验超过 $\tau$ 的邮件判为垃圾邮件时，$\tau$ 应考虑错拦正常邮件与漏放垃圾邮件的代价；模型概率、估计误差与决策规则应分开理解。

\section{扩展阅读2：三门问题}
\label{sec:study-monty-hall}

\emph{三门问题（Monty Hall problem）}的关键是观察信息怎样产生。先明确标准规则：汽车等可能藏在三门之一，参赛者的初选与汽车位置独立；主持人知道位置，总会打开一扇未被选中的羊门，总会允许换到另一扇未开的门；如果有两个羊门可开，主持人等可能选一个。假定初选为 $1$ 号门，令 $C_i$ 表示汽车在 $i$ 号门后，$H_3$ 表示主持人开 $3$ 号门。

在这些规则下，$\mathbb P(C_i)=1/3$，且
\[
 \mathbb P(H_3\mid C_1)=\frac12,\qquad
 \mathbb P(H_3\mid C_2)=1,\qquad
 \mathbb P(H_3\mid C_3)=0.
\]
于是全概率公式给出 $\mathbb P(H_3)=1/2$，贝叶斯公式给出
\begin{equation}
 \mathbb P(C_1\mid H_3)=\frac{(1/3)(1/2)}{1/2}=\frac13,
 \qquad
 \mathbb P(C_2\mid H_3)=\frac{(1/3)\cdot1}{1/2}=\frac23.
 \label{eq:study-monty-hall-standard}
\end{equation}
所以已知开的是 $3$ 号羊门后，换到 $2$ 号门的获胜概率为 $2/3$。从完整策略看也容易解释：若初选正确，换门输；若初选错误，主持人必须排除另一扇羊门，换门赢。因此始终换门的获胜概率等于初选错误的概率 $2/3$。

若主持人不知道汽车位置，并在未选两门中独立等可能开一门，则他可能直接开出汽车。现在额外观察到“开 $3$ 号且是羊”，记为 $E$，则
\[
 \mathbb P(E\mid C_1)=\mathbb P(E\mid C_2)=\frac12,
 \qquad\mathbb P(E\mid C_3)=0.
\]
由贝叶斯公式，$\mathbb P(C_1\mid E)=\mathbb P(C_2\mid E)=1/2$，此时换与不换的条件获胜概率相同。\emph[caution]{“看见一扇羊门被打开”并没有脱离开门机制的固定概率含义。}

还可以考察知情主持人的偏好：若初选正确时，他以概率 $q\in[0,1]$ 开 $3$ 号门，其余规则不变，则
\[
 \mathbb P(C_1\mid H_3)=\frac{q}{1+q},\qquad
 \mathbb P(C_2\mid H_3)=\frac1{1+q}.
\]
因此，给定某扇具体门被打开后的概率会依赖 $q$；标准的 $1/3$ 与 $2/3$ 使用了 $q=1/2$。不过，在主持人始终知情、开羊门并允许换门的前提下，始终换门的事前胜率仍为 $2/3$，因为它恰在初选错误时获胜。这里应区分策略的事前胜率与观察具体开门后的条件胜率。

\section*{本章总结}
\addcontentsline{toc}{section}{本章总结}

一个完整的概率模型包括样本空间、事件类和概率函数。样本空间决定记录什么，事件运算把语言条件转成集合条件，概率公理则保证计数、加法、条件化与极限推理相容。建模时应先说明等可能性、均匀取样或独立性来自哪里；不同记录方式可以描述同一问题，但合并样本点后要保留相应的权重。

计算并事件时，互斥分解可直接相加，一般重叠用容斥，至少一次发生常转向对立事件；计算交事件时，乘法公式总要带上适用的条件，只有独立性成立时才可化为边缘概率的乘积。全概率公式按互斥完备的情形加权，贝叶斯公式再用观察结果更新这些情形的权重。选择划分、确定足够的状态与利用对称性，往往比直接代公式更重要。

本章需特别区分五组概念：样本点与事件、频率与概率、互斥与独立、两两独立与相互独立、条件概率与反向条件概率。概率为零或一不等于逻辑上的不可能或必然；小概率、近似计算和条件独立模型的结论，都应连同其假设一起使用。下一章将用随机变量把事件条件转化为数值条件，并系统研究这些数值的概率分布。

\section*{习题}
\addcontentsline{toc}{section}{习题}

以下选取教材第 1 章的代表性习题，覆盖概率约束、几何模型、条件化、独立性与事件列；题号均指教材本章题号，未逐题抄录全部习题。

% 教材第 1 章习题 4。
\begin{example}
\label{ex:study-exercise-candy-data}
教材习题 4：$1000$ 人中喜欢三种糖果的人数分别为 $811,752,418$，同时喜欢每两种的分别为 $570,356,348$，三种都喜欢的为 $297$。证明数据有误。
\end{example}

\begin{solve}
在这 $1000$ 人组成的有限总体中，按容斥计数，至少喜欢一种的人数应为
\[
 811+752+418-570-356-348+297=1004>1000.
\]
并集不可能超过全集，因此这组统计数不能同时成立。数据检验不只检查各比例是否在 $[0,1]$，还要检查它们之间的集合约束。
\end{solve}

% 教材第 1 章习题 7，同时涵盖习题 6。
\begin{example}
\label{ex:study-exercise-intersection-bounds}
教材习题 7：已知 $\mathbb P(A)=a>0$、$\mathbb P(B)=b>0$，其中 $a,b\le1$，确定 $\mathbb P(A\cap B)$ 的可能范围。
\end{example}

\begin{solve}
由单调性，交概率不超过 $\min\{a,b\}$；由加法公式与并概率不超过 $1$，交概率至少为 $a+b-1$，并且非负。所以
\begin{equation}
 \max\{0,a+b-1\}\le\mathbb P(A\cap B)\le\min\{a,b\}.
 \label{eq:study-frechet-bounds}
\end{equation}
这些界在适当模型中可以达到。事实上，任选区间内的 $t$，给 $A\cap B,A\cap B^{\mathrm c},A^{\mathrm c}\cap B,A^{\mathrm c}\cap B^{\mathrm c}$ 四个区域分别赋概率 $t,a-t,b-t,1-a-b+t$，则均非负且和为 $1$，构成满足条件的模型。对教材习题 6 的 $a=0.7,b=0.8$，交概率范围为 $[0.5,0.7]$；上界对应 $A\setminus B$ 概率为零，下界对应 $A\cup B$ 概率为一。一般不能把这两个概率条件强化为严格的集合包含或全集相等。
\end{solve}

% 教材第 1 章习题 11。
\begin{example}
\label{ex:study-exercise-four-suits}
教材习题 11：从 $52$ 张扑克牌中等可能地不放回抽取 $10$ 张，求四种花色均出现的概率。
\end{example}

\begin{solve}
令 $M_i$ 为第 $i$ 种花色缺席的事件。缺少指定一种、两种或三种花色的选法分别为 $\binom{39}{10},\binom{26}{10},\binom{13}{10}$；四种全缺不可能。由容斥及对立事件法，
\[
 \mathbb P(\text{四种均有})
 =1-\frac{4\binom{39}{10}-6\binom{26}{10}+4\binom{13}{10}}
 {\binom{52}{10}}.
\]
缺席事件有重叠，不能只减去四个单花色缺席概率。
\end{solve}

% 教材第 1 章习题 17。
\begin{example}
\label{ex:study-exercise-same-bus}
教材习题 17：甲、乙在一小时内各自到达公交始发站，每隔 $15$ 分钟发一班车，见车就乘。补充假定两人到达时间相互独立且在这一小时内均匀，求同乘一辆车的概率。
\end{example}

\begin{solve}
以分钟为单位记到达时间为 $X,Y\in[0,60]$，则 $(X,Y)$ 在面积为 $60^2$ 的正方形内均匀。同车对应两人落在同一个 $15$ 分钟区间，共四个面积为 $15^2$ 的小正方形，故概率为 $4\times15^2/60^2=1/4$。恰逢发车时刻的边界概率为零，可以任意一致约定。若只知道每人的边缘到达时间均匀，却没有给定二者关系，就不足以确定同车概率。
\end{solve}

% 教材第 1 章习题 24。
\begin{example}
\label{ex:study-exercise-two-boxes}
教材习题 24：两箱零件分别有 $50$ 只与 $30$ 只，其中一等品分别有 $10$ 只与 $18$ 只。等可能选一箱，再从中不放回抽两只，求第一只是一等品的概率，以及已知第一只一等时第二只也一等的概率。
\end{example}

\begin{solve}
令 $B_i$ 表示选第 $i$ 箱，$A_j$ 表示第 $j$ 只是一等品。全概率公式给出 $\mathbb P(A_1)=\tfrac12\tfrac{10}{50}+\tfrac12\tfrac{18}{30}=2/5$。已知第一只一等后，两箱的后验权重成为 $1/4$ 与 $3/4$，所以
\[
 \mathbb P(A_2\mid A_1)
 =\frac14\frac9{49}+\frac34\frac{17}{29}
 =\frac{690}{1421}\approx0.4856.
\]
也可先计算 $\mathbb P(A_1\cap A_2)$ 再除以 $\mathbb P(A_1)$。条件信息不仅改变箱内组成，也改变选中哪一箱的权重，不能仍以各半平均。
\end{solve}

% 教材第 1 章习题 37。
\begin{example}
\label{ex:study-exercise-conditional-independence}
教材习题 37：$\mathbb P(C)=1/2$，$A,B$ 在给定 $C$ 下以及给定 $C^{\mathrm c}$ 下均条件独立，且
\[
 \mathbb P(A\mid C)=\mathbb P(B\mid C)=0.9,\qquad
 \mathbb P(A\mid C^{\mathrm c})=0.2,\quad
 \mathbb P(B\mid C^{\mathrm c})=0.1.
\]
求 $\mathbb P(A),\mathbb P(B),\mathbb P(A\cap B)$，判断无条件独立性。
\end{example}

\begin{solve}
全概率公式给出 $\mathbb P(A)=0.55$、$\mathbb P(B)=0.5$。在每层中使用条件独立性，再加权，得
\[
 \mathbb P(A\cap B)=\tfrac12(0.9\times0.9)+\tfrac12(0.2\times0.1)
 =0.415\ne0.55\times0.5=0.275.
\]
所以 $A,B$ 无条件不独立。把两个内部独立的层混合起来，可能由于层间成功率不同而产生总体相关性。
\end{solve}

% 教材第 1 章习题 38、40。
\begin{example}
\label{ex:study-exercise-independent-patterns}
教材习题 38：三次射击相互独立，命中率为 $0.5,0.6,0.8$，求恰中一次与至少中一次的概率。教材习题 40：三个继电器断开事件相互独立，概率为 $0.3,0.4,0.6$；第一个断开或第二、三个同时断开时电路断开，求电路断开概率。
\end{example}

\begin{solve}
恰中一次的三个互斥模式分别是中失失、失中失、失失中，由独立性相乘再相加，概率为 $0.5\times0.4\times0.2+0.5\times0.6\times0.2+0.5\times0.4\times0.8=0.26$；至少中一次的概率为 $1-0.5\times0.4\times0.2=0.96$。继电器题中记断开事件为 $F_i$，电路断开为 $F_1\cup(F_2\cap F_3)$，故概率为
\[
 0.3+0.4\times0.6-0.3\times0.4\times0.6=0.468.
\]
两个引发断路的情形可以同时发生，必须扣去重叠。
\end{solve}

% 教材第 1 章习题 19（带星号）。
\begin{example}
\label{ex:study-exercise-borel-cantelli}
教材习题 19：令 $\{A_n\text{ 无穷多次发生}\}=\bigcap_{n=1}^{\infty}\bigcup_{k=n}^{\infty}A_k$。证明\emph{博雷尔--坎泰利引理（Borel--Cantelli lemmas）}的两个结论：若 $\sum_n\mathbb P(A_n)<\infty$，则无穷多次发生的概率为零；若事件列相互独立且 $\sum_n\mathbb P(A_n)=\infty$，则无穷多次发生的概率为一。
\end{example}

\begin{proof}
记 $B_n=\bigcup_{k=n}^{\infty}A_k$，则 $B_n\downarrow\bigcap_n B_n$。第一种情形中，由次可加性，$\mathbb P(B_n)\le\sum_{k=n}^{\infty}\mathbb P(A_k)\to0$，再用上连续性即得结论，整个论证无需独立性。第二种情形中，对固定 $n$ 和 $m\ge n$，由独立性与 $1-x\le e^{-x}$（$0\le x\le1$），
\[
 \mathbb P\!\left(\bigcap_{k=n}^{m}A_k^{\mathrm c}\right)
 =\prod_{k=n}^{m}[1-\mathbb P(A_k)]
 \le\exp\!\left(-\sum_{k=n}^{m}\mathbb P(A_k)\right)\longrightarrow0.
\]
由上连续性，$\mathbb P(B_n^{\mathrm c})=0$，故每个 $B_n$ 的概率均为 $1$；再对递减列 $B_n$ 应用上连续性，得到交事件的概率为 $1$。第二个结论需要独立性假设，不能仅凭概率级数发散得出。
\end{proof}

\chapter{随机变量及其分布}
\label{chap:study-random-variables}

本章把第一章中关于事件的概率转化为实数轴上的分布问题。学习的主线是：先用随机变量记录试验结果，再用分布律、分布函数或密度刻画其概率规律，最后研究确定的数值变换如何改变分布。各种分布的名称并不能代替建模条件，尤其要区分固定次数内的成功数、达到指定成功次数所需的试验数，以及给定时间或空间范围内的事件计数。

\section{随机变量的概念}
\label{sec:study-random-variable-concept}

随机试验的结果可能是硬币的正反面、球的颜色或比赛的胜负记录。研究时，我们常常只关心其中某个数值特征，例如正面出现的次数。随机变量就是按照所研究的问题选定一种数值记录规则，给每个样本点赋予一个实数，从而把关于具体结果的问题转化为关于数值的问题。

\begin{definition}
\label{def:study-random-variable}
设随机试验的样本空间为 $\Omega$，并已给定事件及其概率。若实值函数
\[
  X:\Omega\longrightarrow\mathbb R,\qquad
  \omega\longmapsto X(\omega)
\]
满足：对任意实数 $u$，集合
\begin{equation}
  \{X\le u\}:=\{\omega\in\Omega:X(\omega)\le u\}
  \label{eq:study-random-variable-measurability}
\end{equation}
都是该概率模型中的事件，则称 $X$ 为一个\emph{随机变量（random variable）}。
\end{definition}

上述事件条件称为\emph{可测性（measurability）}要求，保证相应概率有意义；在有限样本空间且所有子集均为事件时，每个实值函数都满足这一要求。赋值规则 $X$ 是事先确定的，随机性来自试验最终出现哪个样本点 $\omega$。

\begin{example}
\label{ex:study-random-variable-two-coins}
独立抛掷一枚公平硬币两次，每次正面和反面出现的概率均为 $1/2$。用 $H$ 表示正面，$T$ 表示反面，记随机变量 $X$ 为两次抛掷中正面出现的次数。写出 $X$ 的赋值规则，说明具体取值与事件的区别，并求事件 $\{X=1\}$ 和 $\{X\le1\}$ 的概率。
\end{example}

\begin{solve}
用两个字母依次记录两次抛掷的结果，样本空间为
\[
  \Omega=\{HH,HT,TH,TT\}.
\]
由公平性与独立性，四个样本点等可能，各自对应的单点事件的概率为 $1/4$。按照“数出正面的个数”这一规则，
\begin{equation}
  X(HH)=2,\qquad X(HT)=X(TH)=1,\qquad X(TT)=0.
  \label{eq:study-random-variable-coin-map}
\end{equation}
所以 $X$ 的可能取值为 $0,1,2$。若实际结果为 $HT$，则本次观察到的\emph{实现值（realization）}为 $x=X(HT)=1$。这里，$HT$ 是样本点，$X$ 是赋值函数，$1$ 是具体取值；通常用大写字母表示随机变量，用小写字母表示其取值。

“恰好出现一次正面”对应事件
\begin{equation}
  \{X=1\}:=\{\omega\in\Omega:X(\omega)=1\}
    =\{HT,TH\}.
  \label{eq:study-random-variable-coin-event}
\end{equation}
因此，仅知道 $X=1$ 不能区分 $HT$ 与 $TH$，随机变量不一定保留结果的全部信息。记号 $\mathbb P(X=1)$ 是 $\mathbb P(\{X=1\})$ 的简写；由两个单点事件的互斥性，
\[
  \mathbb P(X=1)=\mathbb P(\{HT\})+\mathbb P(\{TH\})
    =\frac14+\frac14=\frac12.
\]
同样，“至多出现一次正面”对应事件
\[
  \{X\le1\}=\{\omega\in\Omega:X(\omega)\le1\}
    =\{TT,HT,TH\}.
\]
它包含三个等可能的样本点，故 $\mathbb P(X\le1)=3/4$。这些计算说明，关于随机变量的数值条件，仍然是在描述样本空间中的事件。

若还要区分 $HT$ 与 $TH$，可以在同一样本空间上再定义 $Y$：第一次为正面时取 $1$，否则取 $0$。这样，$HT$ 与 $TH$ 虽然都使 $X=1$，却分别使 $Y=1$ 与 $Y=0$，说明可以用不同随机变量记录同一试验的不同特征。\emph{同一次试验只产生一个样本点，却可以同时确定多个随机变量的取值。}
\end{solve}

\section{离散型随机变量的分布}
\label{sec:study-discrete-distributions}

\begin{definition}
\label{def:study-discrete-mass}
若存在有限或可数的实数集合 $S=\{x_1,x_2,\ldots\}$，使 $\mathbb P(X\in S)=1$，则称 $X$ 为\emph{离散型随机变量（discrete random variable）}。将互不相同的取值与相应概率
\begin{equation}
  p_k=\mathbb P(X=x_k),\qquad p_k\ge0,\qquad \sum_kp_k=1
  \label{eq:study-discrete-mass-conditions}
\end{equation}
合在一起，称为 $X$ 的\emph{分布律}，也称\emph{概率质量函数（probability mass function, pmf）}。
\end{definition}

各事件 $\{X=x_k\}$ 两两互斥，并以概率 $1$ 覆盖全部结果，因此式~\eqref{eq:study-discrete-mass-conditions} 中的归一化条件来自可数可加性。对任意所讨论的可测集合 $A\subseteq\mathbb R$，有
\begin{equation}
  \mathbb P(X\in A)=\sum_{k:x_k\in A}p_k.
  \label{eq:study-discrete-set-probability}
\end{equation}
所以掌握分布律后，计算概率就是确定哪些取值满足条件，再将对应概率相加。反过来，任意满足非负性和总和为 $1$ 的数列，都可在可数样本空间上构成一个分布律；仅列出取值而不给出概率，还没有确定随机变量的分布。

\subsection{0-1 分布与二项分布}
\label{subsec:study-bernoulli-binomial}

把试验的两类结果分别称为“成功”和“失败”。若只记录一次试验是否成功，就得到 0-1 分布；若进行固定次数的独立重复试验，每次成功的概率相同，并记录成功的总次数，就得到二项分布。“成功”和“失败”只是结果的名称，不表示好坏。

\begin{definition}
\label{def:study-bernoulli-distribution}
设 $p\in[0,1]$。若随机变量 $X$ 取值于 $\{0,1\}$，且
\begin{equation}
  \mathbb P(X=1)=p,\qquad \mathbb P(X=0)=1-p,
  \label{eq:study-bernoulli-mass}
\end{equation}
则称 $X$ 服从参数为 $p$ 的\emph{0-1 分布}，也称\emph{伯努利分布（Bernoulli distribution）}，记作 $X\sim\mathrm{Bern}(p)$。
\end{definition}

例如，抛掷一次硬币，正面出现的概率为 $p$；把正面记为 $1$、反面记为 $0$，所得随机变量便服从 $\mathrm{Bern}(p)$。

\begin{definition}
\label{def:study-binomial-distribution}
设 $n$ 为正整数，$p\in[0,1]$。进行 $n$ 次相互独立的试验，每次只有成功与失败两类结果，且每次成功的概率均为 $p$，这样的试验序列称为\emph{伯努利试验序列（Bernoulli trials）}。令 $X_i$ 在第 $i$ 次试验成功时取 $1$，失败时取 $0$，其中 $i=1,\ldots,n$，则成功的总次数为
\begin{equation}
  Y=X_1+X_2+\cdots+X_n.
  \label{eq:study-binomial-sum}
\end{equation}
称 $Y$ 的分布为参数为 $n,p$ 的\emph{二项分布（binomial distribution）}，记作 $Y\sim\mathrm{Bin}(n,p)$。其中 $X_1,\ldots,X_n$ 相互独立，均服从 $\mathrm{Bern}(p)$，而 $Y$ 取值于 $\{0,1,\ldots,n\}$。
\end{definition}

计算二项分布的概率时，应区分“指定哪些次试验成功”和“只指定成功的总次数”：前者确定一个完整的结果序列，后者通常包含多个不同序列。

\begin{example}
\label{ex:study-binomial-three-trials}
独立进行三次试验，每次成功的概率均为 $p\in(0,1)$。记 $Y$ 为成功的总次数，求恰好成功两次的概率，并说明计数系数的来源。
\end{example}

\begin{solve}
用 $S$ 表示成功，$F$ 表示失败，按试验顺序记录结果。恰好成功两次对应 $SSF$、$SFS$ 和 $FSS$ 三种结果，由独立性，每种结果的概率均为 $p^2(1-p)$。它们对应的事件两两互斥，所以
\begin{equation}
  \mathbb P(Y=2)=3p^2(1-p).
  \label{eq:study-binomial-three-trials}
\end{equation}
其中 $3=\binom32$ 是从三个试验位置中选出两个成功位置的方法数。
\end{solve}

\begin{proposition}
\label{prop:study-binomial-mass}
设 $Y\sim\mathrm{Bin}(n,p)$，其中 $n$ 为正整数。当 $0<p<1$ 时，$Y$ 的概率分布为
\begin{equation}
  \mathbb P(Y=k)=\binom nk p^k(1-p)^{n-k},
  \qquad k=0,1,\ldots,n,
  \label{eq:study-binomial-mass}
\end{equation}
其余取值的概率为 $0$。当 $p=0$ 时，$\mathbb P(Y=0)=1$；当 $p=1$ 时，$\mathbb P(Y=n)=1$。
\end{proposition}

\begin{proof}
先设 $0<p<1$，固定整数 $k\in\{0,1,\ldots,n\}$。若事先指定哪 $k$ 次试验成功、其余 $n-k$ 次试验失败，由各次试验的独立性，这一种完整安排的概率为
\[
  \underbrace{p\cdots p}_{k\text{ 个}}
  \underbrace{(1-p)\cdots(1-p)}_{n-k\text{ 个}}
  =p^k(1-p)^{n-k},
\]
其中没有因子时相应的乘积按 $1$ 计算。由于每次试验的成功概率相同，这一概率只取决于成功与失败的次数，不取决于它们出现的具体位置。

从 $n$ 个试验位置中选出 $k$ 个成功位置，共有
\[
  \binom nk=\frac{n!}{k!(n-k)!}
\]
种选法，其中 $0!=1$。选定成功位置后，其余位置全部为失败，所以每个选法恰好对应一种完整安排；这里只需选择位置，无须再给选中的位置另作排列。不同安排对应的事件两两互斥，而成功总次数恰为 $k$ 的所有结果都包含在这些安排中。由概率的有限可加性，将这 $\binom nk$ 个相同的概率相加，得到
\[
  \mathbb P(Y=k)=\binom nk p^k(1-p)^{n-k}.
\]
此外，成功次数必在 $0$ 与 $n$ 之间；若 $p=0$，每次试验都以概率 $1$ 失败，故 $Y=0$ 的概率为 $1$；若 $p=1$，同理可得 $Y=n$ 的概率为 $1$。
\end{proof}

因此，式~\eqref{eq:study-binomial-mass} 中的组合数统计的是成功位置的安排数，$p^k(1-p)^{n-k}$ 则是一种指定安排的概率。\emph[tip]{同一种安排内部，由独立性将各次试验的概率相乘；不同安排之间，由互斥性将各个事件的概率相加。}这里具有相同成功次数的序列等可能，并不要求全部结果序列等可能。

两种分布的关系也可由式~\eqref{eq:study-binomial-sum} 看出：\emph{二项随机变量是若干个相互独立、参数相同的 0-1 随机变量之和。}当 $n=1$ 时，$\mathrm{Bin}(1,p)$ 就是 $\mathrm{Bern}(p)$。应用二项分布时，应核对试验次数固定、各次试验相互独立、成功概率相同这三个条件；若抽样方式不满足这些条件，就需要根据实际模型重新计算分布。

\subsection{离散均匀分布}
\label{subsec:study-discrete-uniform}

设 $x_1,\ldots,x_n$ 为 $n\ge1$ 个不同的实数。若
\begin{equation}
  \mathbb P(X=x_k)=\frac1n,\qquad k=1,\ldots,n,
  \label{eq:study-discrete-uniform}
\end{equation}
则称 $X$ 服从这 $n$ 个取值上的\emph{离散均匀分布（discrete uniform distribution）}。例如，均匀骰子的点数在 $\{1,\ldots,6\}$ 上均匀分布。古典概型中的样本点等可能，但把多个样本点合并为同一个随机变量取值后，各取值未必仍等可能；例~\ref{ex:study-random-variable-two-coins} 中四个样本点等可能，正面次数 $0,1,2$ 的概率却分别为 $1/4,1/2,1/4$。下面的抽样问题也体现了这一差别。

\subsection{标记重捕模型与超几何分布}
\label{subsec:study-capture-recapture}

% 教材例 2.5：标记重捕（Capture-Recapture）模型。
\begin{example}
\label{ex:study-capture-recapture}
池塘中共有 $N$ 条鱼，第一次捕捞 $M$ 条，标记后全部放回，充分混合后再不放回地捕捞 $n$ 条，其中 $N,M,n$ 为整数，$N\ge1$、$0\le M\le N$、$1\le n\le N$。假定鱼群组成与标记保持不变，第二次捕捞时每个包含 $n$ 条不同鱼的组合等可能。求其中恰有 $m$ 条鱼带标记的概率。
\end{example}

\begin{solve}
这称为\emph{标记重捕（capture--recapture）模型}。记 $X$ 为第二次捕捞中有标记鱼的数量。“不放回”指第二次捕捞时同一条鱼不会被重复抽到，与第一次标记后放回并不矛盾。不考虑捕捞顺序，每个具体的 $n$ 条鱼的组合是一个样本点，全部样本点共有 $\binom Nn$ 个，且等可能；这里 $\binom ab$ 表示从 $a$ 个不同对象中选出 $b$ 个的组合数。

要使 $X=m$，需从 $M$ 条有标记鱼中选出 $m$ 条，再从 $N-M$ 条无标记鱼中选出 $n-m$ 条。\emph[tip]{选定有标记部分后，无标记部分的不同选择仍会形成不同的完整结果。}两部分的任意选法可以配对，每个完整结果又唯一确定这两部分，故有利结果数为 $\binom Mm\binom{N-M}{n-m}$。由古典概率公式，
\begin{equation}
  \mathbb P(X=m)=\frac{\binom Mm\binom{N-M}{n-m}}{\binom Nn}.
  \label{eq:study-hypergeometric-mass}
\end{equation}
这称为参数为 $(N,M,n)$ 的\emph{超几何分布（hypergeometric distribution）}。由于两类鱼的数量有限，可能的整数取值满足
\begin{equation}
  \max\{0,n-(N-M)\}\le m\le\min\{M,n\},
  \label{eq:study-hypergeometric-support}
\end{equation}
范围外的概率为 $0$。这里使用的是选法数的乘法原理，不要求额外的独立性假设；等可能的是具体组合，而不是 $X$ 的不同取值。
\end{solve}

\subsection{负二项分布与几何分布}
\label{subsec:study-negative-binomial-geometric}

二项分布先固定试验次数，再数成功次数；如果改为不断进行伯努利试验，直到第 $r$ 次成功才停止，随机的就是停止时的试验次数。设各次试验相互独立，成功概率为 $p\in(0,1)$，记 $q=1-p$，$r\ge1$ 为整数，$T_r$ 表示第 $r$ 次成功发生时的累计试验次数。对于整数 $k\ge r$，事件 $\{T_r=k\}$ 要求前 $k-1$ 次恰好成功 $r-1$ 次，第 $k$ 次成功，所以
\begin{equation}
  \mathbb P(T_r=k)=\binom{k-1}{r-1}p^rq^{k-r},
  \qquad k=r,r+1,\ldots.
  \label{eq:study-negative-binomial-mass}
\end{equation}
这称为参数为 $r,p$ 的\emph{负二项分布（negative binomial distribution）}或\emph{帕斯卡分布（Pascal distribution）}，记为 $T_r\sim\mathrm{NB}(r,p)$。利用负二项级数，在 $0\le q<1$ 时，
\[
  \sum_{k=r}^{\infty}\binom{k-1}{r-1}p^rq^{k-r}
  =p^r\sum_{j=0}^{\infty}\binom{r+j-1}{j}q^j
  =p^r(1-q)^{-r}=1.
\]
这也说明在这些假设下，第 $r$ 次成功以概率 $1$ 在有限次试验内出现。有些文献把第 $r$ 次成功前的失败次数 $W_r=T_r-r$ 称为负二项随机变量，此时取值从 $0$ 开始，分布律为 $\mathbb P(W_r=j)=\binom{r+j-1}{j}p^rq^j$；使用公式前须先核对计数对象。

当 $r=1$ 时，$T=T_1$ 表示首次成功时的试验次数，其分布称为\emph{几何分布（geometric distribution）}，记为 $T\sim\mathrm{Ge}(p)$，且
\begin{equation}
  \mathbb P(T=k)=q^{k-1}p\quad(k=1,2,\ldots),\qquad
  \mathbb P(T>m)=q^m\quad(m=0,1,\ldots).
  \label{eq:study-geometric-mass-tail}
\end{equation}
这里尾概率也可直接由“前 $m$ 次全部失败”得到。几何分布的核心性质是\emph{无记忆性（memoryless property）}：在已经等待 $m$ 次仍未成功的条件下，还要等待超过 $n$ 次的概率，与一开始等待超过 $n$ 次的概率相同。

\begin{theorem}[几何分布的无记忆性及刻画]
\label{thm:study-geometric-memoryless}
若 $T\sim\mathrm{Ge}(p)$，$0<p<1$，则对任意非负整数 $m,n$，
\begin{equation}
  \mathbb P(T>m+n\mid T>m)=\mathbb P(T>n).
  \label{eq:study-geometric-memoryless}
\end{equation}
反之，若 $T$ 以概率 $1$ 取正整数值，所有尾概率 $\mathbb P(T>m)$ 均为正，且满足此等式，则 $T$ 服从某个参数 $p\in(0,1)$ 的几何分布。
\end{theorem}

\begin{proof}
正向结论由 $q^{m+n}/q^m=q^n$ 得到。反向令 $s_m=\mathbb P(T>m)$，则 $s_0=1$，条件概率等式给出 $s_{m+n}=s_ms_n$，从而 $s_m=s_1^m$。记 $q=s_1$，由假设 $q>0$；又因 $T$ 有限且取正整数值，事件 $\{T>m\}$ 下降至空集，故 $s_m\to0$，排除 $q=1$。于是 $0<q<1$，并有
\[
  \mathbb P(T=k)=s_{k-1}-s_k=(1-q)q^{k-1},\qquad k\ge1,
\]
即为参数 $p=1-q$ 的几何分布。
\end{proof}

% 教材例 2.8：巴拿赫火柴问题。
\begin{example}
\label{ex:study-banach-matches}
两盒火柴最初各有 $n\ge1$ 根，每次独立地以相同概率选择其中一盒，若盒中有火柴则取用一根，直到某次选中的盒子已经空了才停止。求停止时另一盒剩有 $r$ 根火柴的概率，其中 $r=0,\ldots,n$。
\end{example}

\begin{solve}
先指定被发现为空的是甲盒。把选择甲盒看作成功，停止时恰为第 $n+1$ 次选中甲盒；此前已从乙盒取走 $n-r$ 根，因此总选择次数为 $2n-r+1$。其概率为
\[
  \binom{2n-r}{n}\left(\frac12\right)^{2n-r+1}.
\]
在这类序列的前 $2n-r$ 次选择中，甲盒恰被选择 $n$ 次，乙盒恰被选择 $n-r\le n$ 次，因此此前不会已经选中空盒。由甲、乙对称性，且“最后选中甲空盒”与“最后选中乙空盒”互斥，所求概率为
\begin{equation}
  \mathbb P(R=r)=\binom{2n-r}{n}2^{-(2n-r)},
  \qquad r=0,\ldots,n.
  \label{eq:study-banach-matches}
\end{equation}
即使 $r=0$，两种最后选择仍是互斥的，可以相加。这里的停止时刻是\emph{发现盒空}，与刚刚用完某盒最后一根火柴的时刻不同。
\end{solve}

\subsection{泊松分布与稀有事件近似}
\label{subsec:study-poisson}

\begin{definition}
\label{def:study-poisson}
设 $\lambda>0$。若随机变量 $X$ 的分布律为
\begin{equation}
  \mathbb P(X=k)=e^{-\lambda}\frac{\lambda^k}{k!},
  \qquad k=0,1,2,\ldots,
  \label{eq:study-poisson-mass}
\end{equation}
则称 $X$ 服从参数为 $\lambda$ 的\emph{泊松分布（Poisson distribution）}，记为 $X\sim\mathrm{Pois}(\lambda)$。
\end{definition}

各项非负，而指数函数的幂级数给出 $\sum_{k=0}^{\infty}e^{-\lambda}\lambda^k/k!=1$。泊松分布常用于给定时间或空间范围内的事件计数，但采用这一模型需要相应的机制或数据支持，不能只因事件发生次数较少就断定其服从泊松分布。教材的微生物模型可这样表述：$N$ 个微生物独立、均匀地分布在体积为 $V$ 的液体中，固定体积 $v$ 内的个数服从 $\mathrm{Bin}(N,v/V)$；若 $N,V\to\infty$，且 $N/V\to d>0$，则可由下述定理得到参数 $dv$ 的泊松极限。独立性和空间均匀性在这里是建模假设。

\begin{theorem}[泊松逼近定理]
\label{thm:study-poisson-limit}
设 $X_n\sim\mathrm{Bin}(n,p_n)$，其中 $np_n\to\lambda\in(0,\infty)$。则对每个固定的非负整数 $k$，
\begin{equation}
  \lim_{n\to\infty}\mathbb P(X_n=k)
  =e^{-\lambda}\frac{\lambda^k}{k!}.
  \label{eq:study-poisson-limit}
\end{equation}
\end{theorem}

\begin{proof}
因为 $p_n\to0$，对固定 $k$ 和充分大的 $n$，
\[
  \mathbb P(X_n=k)
  =\frac{(np_n)^k}{k!}
    \prod_{j=0}^{k-1}\left(1-\frac jn\right)
    (1-p_n)^n(1-p_n)^{-k}.
\]
当 $k=0$ 时乘积取 $1$。有限乘积及最后一项均趋于 $1$，且由 $\log(1-u)=-u+O(u^2)$，
\[
  n\log(1-p_n)=-np_n+O(np_n^2)\longrightarrow-\lambda,
\]
其中 $np_n^2=(np_n)p_n\to0$，故 $(1-p_n)^n\to e^{-\lambda}$，代入即得。
\end{proof}

实际计算中，$n$ 较大、$p$ 较小时可取 $\lambda=np$，用泊松概率近似二项概率；这是近似而非恒等式，精度还取决于参数和所求事件。例如 $X\sim\mathrm{Pois}(2)$ 时，
\[
  \mathbb P(X\ge3)=1-\mathbb P(X\le2)
  =1-e^{-2}\left(1+2+\frac{2^2}{2!}\right)
  =1-5e^{-2}\approx0.3233.
\]
类似地，对于 $T_r\sim\mathrm{NB}(r,p_r)$，若 $r\to\infty$ 且 $r(1-p_r)\to\lambda>0$，则对每个固定的 $j\ge0$，
\[
  \mathbb P(T_r-r=j)
  =\frac{[r(1-p_r)]^j}{j!}
    \prod_{\ell=0}^{j-1}\left(1+\frac\ell r\right)p_r^r
  \longrightarrow e^{-\lambda}\frac{\lambda^j}{j!}.
\]
此时近似为泊松分布的是停止前的\emph{失败次数}，而总试验次数还包含确定的 $r$ 次成功。

% 教材例 2.14：维修人员配备。
\begin{example}
\label{ex:study-poisson-repair}
有 $200$ 台设备，各台在所考察时刻是否故障相互独立，故障概率均为 $0.01$。每名工人同时只能处理一台设备，至少需配备多少人，才能使故障设备数超过工人数的概率小于 $0.01$？
\end{example}

\begin{solve}
令 $X$ 为故障设备数，则 $X\sim\mathrm{Bin}(200,0.01)$，所求为满足 $\mathbb P(X>N)<0.01$ 的最小非负整数 $N$。用 $Z\sim\mathrm{Pois}(2)$ 近似，
\[
  \mathbb P(Z>5)\approx0.01656,\qquad
  \mathbb P(Z>6)\approx0.00453,
\]
由尾概率随 $N$ 非增，近似计算给出 $N=6$。若按精确二项概率核验，相应尾概率约为 $0.01602$ 和 $0.00430$，因而本题精确的最少人数也是 $6$。只凭极限定理本身不能对任意有限参数作出严格的误差保证，本题的最终判断可由有限二项和直接核实。
\end{solve}

% 教材例 2.15：泊松计数的独立记录。
\begin{example}
\label{ex:study-poisson-thinning}
设单位时间内发射的粒子数 $X\sim\mathrm{Pois}(\lambda)$。给定发射的粒子后，各粒子以概率 $p\in(0,1)$ 被独立记录，且这一记录机制不随 $X$ 改变。求记录到的粒子数 $Y$ 的分布。
\end{example}

\begin{solve}
给定 $X=n$ 时，$Y$ 的条件分布为 $\mathrm{Bin}(n,p)$。对整数 $k\ge0$，由全概率公式与非负级数求和，
\begin{align*}
  \mathbb P(Y=k)
  &=\sum_{n=k}^{\infty}\binom nk p^k(1-p)^{n-k}
       e^{-\lambda}\frac{\lambda^n}{n!}\\
  &=e^{-\lambda}\frac{(\lambda p)^k}{k!}
       \sum_{j=0}^{\infty}\frac{[\lambda(1-p)]^j}{j!}
    =e^{-\lambda p}\frac{(\lambda p)^k}{k!}.
\end{align*}
所以 $Y\sim\mathrm{Pois}(\lambda p)$。这里随机的是总发射数和每个粒子是否被记录，先给定总数再平均，可以把这两层随机性分开处理。
\end{solve}

\section{连续型随机变量的分布}
\label{sec:study-continuous-distributions}

\subsection{适用于所有随机变量的分布函数}
\label{subsec:study-cdf}

离散分布可逐点相加，但一般随机变量的取值未必能逐一列举，甚至每个单点的概率都为零。为统一描述各类分布，先记录随机变量落在半直线上的概率。

\begin{definition}
\label{def:study-cdf}
随机变量 $X$ 的\emph{分布函数（cumulative distribution function, cdf）}定义为
\begin{equation}
  F_X(x)=\mathbb P(X\le x),\qquad x\in\mathbb R.
  \label{eq:study-cdf-definition}
\end{equation}
不致混淆时简记为 $F(x)$。自变量 $x$ 是确定的实数，$F_X$ 是由 $X$ 的分布确定的普通函数。
\end{definition}

\begin{proposition}[分布函数的基本性质]
\label{prop:study-cdf-properties}
任意随机变量的分布函数均非减、右连续，并满足
\begin{equation}
  0\le F(x)\le1,\qquad
  \lim_{x\to-\infty}F(x)=0,\qquad
  \lim_{x\to+\infty}F(x)=1.
  \label{eq:study-cdf-properties}
\end{equation}
\end{proposition}

\begin{proof}
若 $x_1<x_2$，则 $\{X\le x_1\}\subseteq\{X\le x_2\}$，由概率的单调性得到非减性，取值界限也直接来自概率的界限。事件 $\{X\le n\}$ 随正整数 $n$ 增加而上升至整个样本空间，$\{X\le-n\}$ 则下降至空集，由概率的连续性得到两端极限，再结合单调性推广到实数趋于无穷的情形。若 $x_n\downarrow x$，则 $\{X\le x_n\}\downarrow\{X\le x\}$，故 $F(x_n)\to F(x)$，这证明了右连续性。
\end{proof}

记 $F(x-)=\lim_{t\uparrow x}F(t)$。由 $\{X\le x-1/n\}\uparrow\{X<x\}$，有
\begin{equation}
  F(x-)=\mathbb P(X<x),\qquad
  \mathbb P(X=x)=F(x)-F(x-).
  \label{eq:study-cdf-atom}
\end{equation}
因此，分布函数在一点的跳跃高度就是该点的概率。对任意 $a<b$，端点的包含方式对应下列不同公式：
\begin{equation}
  \begin{aligned}
    \mathbb P(a<X\le b)&=F(b)-F(a),&
    \mathbb P(a\le X\le b)&=F(b)-F(a-),\\
    \mathbb P(a<X<b)&=F(b-)-F(a),&
    \mathbb P(a\le X<b)&=F(b-)-F(a-).
  \end{aligned}
  \label{eq:study-cdf-intervals}
\end{equation}
例如第一式来自 $\{X\le b\}=\{X\le a\}\cup\{a<X\le b\}$ 的不交分解，其他公式用左右端点的单点概率调整即可。

对离散型随机变量，分布律与分布函数相互确定：
\begin{equation}
  F(x)=\sum_{k:x_k\le x}p_k,\qquad
  p_k=F(x_k)-F(x_k-).
  \label{eq:study-discrete-cdf}
\end{equation}
这一写法不要求所有取值都能按相邻大小顺序排列。有限取值的情形给出常见的阶梯图形；一般可数取值的情形也应使用式~\eqref{eq:study-discrete-cdf}，而不预设两个相邻取值之间必有非空区间。分布相同的随机变量可以定义在不同的样本空间上，也不必在每次试验中取相同的值；分布描述的是概率规律。

% 教材例 2.17：同时含连续部分和跳跃部分的分布函数。
\begin{example}
\label{ex:study-mixed-cdf}
设随机变量 $X$ 的分布函数为
\[
  F(x)=\begin{cases}
    0,&x<0,\\
    x/2,&0\le x<1,\\
    2/3,&1\le x<2,\\
    11/12,&2\le x<3,\\
    1,&x\ge3.
  \end{cases}
\]
求 $\mathbb P(X<3)$、$\mathbb P(X=1)$、$\mathbb P(X>1/2)$ 和 $\mathbb P(2<X\le4)$。
\end{example}

\begin{solve}
分别依据左极限、跳跃、补事件和半开区间公式，得到
\[
  \begin{aligned}
    \mathbb P(X<3)&=F(3-)=\frac{11}{12},&
    \mathbb P(X=1)&=\frac23-\frac12=\frac16,\\
    \mathbb P(X>1/2)&=1-F(1/2)=\frac34,&
    \mathbb P(2<X\le4)&=F(4)-F(2)=\frac1{12}.
  \end{aligned}
\]
这个分布在 $(0,1)$ 上的连续部分共有概率 $1/2$，在 $1,2,3$ 三点的概率分别为 $1/6,1/4,1/12$，合计为 $1$。它属于\emph{混合型分布（mixed distribution）}，既不是纯离散型，也不能仅用一个普通密度函数表示全部概率。
\end{solve}

\subsection{密度、积分与绝对连续性}
\label{subsec:study-density}

\begin{definition}
\label{def:study-density}
若存在非负可积函数 $f_X$，使随机变量 $X$ 的分布函数可写为
\begin{equation}
  F_X(x)=\int_{-\infty}^{x}f_X(t)\,\mathrm dt,
  \qquad x\in\mathbb R,
  \label{eq:study-density-definition}
\end{equation}
则称 $X$ 为\emph{连续型随机变量}，更确切地称为\emph{绝对连续型随机变量（absolutely continuous random variable）}，$f_X$ 称为其\emph{概率密度函数（probability density function, pdf）}。
\end{definition}

密度满足 $f_X(x)\ge0$ 以及 $\int_{\mathbb R}f_X(x)\,\mathrm dx=1$，并由积分计算概率：
\begin{equation}
  \mathbb P(a<X\le b)=\int_a^bf_X(x)\,\mathrm dx,
  \qquad
  \mathbb P(X\in A)=\int_Af_X(x)\,\mathrm dx.
  \label{eq:study-density-probability}
\end{equation}
后一式对可测集合 $A$ 成立，是前一式从区间推广到一般集合的形式。绝对连续的 $F_X$ 必连续，故 $\mathbb P(X=x)=0$；因此有限区间的四种端点写法具有相同概率。\emph[caution]{单点概率为零不表示该点不可能作为观察值，密度值也不是单点概率。}密度可以大于 $1$，只要在整个实轴上的积分等于 $1$；改变有限个点上的密度值不改变分布。更一般地，在长度为零的集合上改变密度值，也不改变各个概率，这就是密度只在\emph{几乎处处（almost everywhere）}意义下唯一。

若 $f_X$ 在 $x$ 处连续，由微积分基本定理有 $F_X'(x)=f_X(x)$，从而当 $h\downarrow0$ 时，
\begin{equation}
  \mathbb P(x<X\le x+h)=f_X(x)h+o(h).
  \label{eq:study-density-local}
\end{equation}
密度描述的是单位长度上的局部概率集中程度，而非不依赖区间长度的概率。将观测数据分组时，直方图应以“组内相对频数除以组距”为高度，使矩形面积等于该组的相对频数；它是根据样本构成的密度近似，不能直接等同于未知总体的真实密度。

连续型随机变量的密度不一定连续；反过来，分布函数连续只保证没有单点概率，不能保证存在密度，奇异连续分布就是反例。因此，离散型和本节定义的连续型没有穷尽所有分布，分布函数才是统一的描述工具。

% 教材例 2.19：分段密度的积分。
\begin{example}
\label{ex:study-piecewise-density}
设随机变量 $X$ 的密度为
\[
  f(x)=\begin{cases}
    2x/3,&0\le x<1,\\
    2(x-2)^2,&1\le x\le2,\\
    0,&\text{其他}.
  \end{cases}
\]
求其分布函数，并说明密度的不连续是否导致分布函数跳跃。
\end{example}

\begin{solve}
两段的积分分别为 $1/3$ 和 $2/3$，总积分为 $1$。分段积分时须计入此前各段已累积的概率，得到
\[
  F(x)=\begin{cases}
    0,&x<0,\\
    x^2/3,&0\le x<1,\\
    1+\dfrac23(x-2)^3,&1\le x\le2,\\
    1,&x>2.
  \end{cases}
\]
在 $x=1$ 处，密度左极限为 $2/3$，右侧值为 $2$，但 $F(1-)=F(1)=1/3$。所以密度的跳变只会影响分布函数的局部斜率，并不产生单点概率。
\end{solve}

\subsection{均匀分布与指数分布}
\label{subsec:study-uniform-exponential}

\begin{definition}
\label{def:study-continuous-uniform}
设 $-\infty<a<b<\infty$。若随机变量 $X$ 的密度为
\begin{equation}
  f_X(x)=\begin{cases}
    1/(b-a),&a<x<b,\\
    0,&\text{其他},
  \end{cases}
  \label{eq:study-uniform-density}
\end{equation}
则称 $X$ 服从区间 $(a,b)$ 上的\emph{均匀分布（uniform distribution）}，记为 $X\sim U(a,b)$。
\end{definition}

积分得到
\begin{equation}
  F_X(x)=\begin{cases}
    0,&x\le a,\\
    (x-a)/(b-a),&a<x<b,\\
    1,&x\ge b.
  \end{cases}
  \label{eq:study-uniform-cdf}
\end{equation}
因此，对于 $a\le c<d\le b$，有 $\mathbb P(c<X<d)=(d-c)/(b-a)$；等长子区间具有相同概率。这里的均匀性是对区间长度而言，不能解释为不可数个点各占有同一个正概率。端点的单点概率为零，所以把 $(a,b)$ 改成相应的闭区间不会改变分布。

\begin{definition}
\label{def:study-exponential}
设 $\lambda>0$。若随机变量 $X$ 的密度为
\begin{equation}
  f_X(x)=\begin{cases}
    \lambda e^{-\lambda x},&x>0,\\
    0,&x\le0,
  \end{cases}
  \label{eq:study-exponential-density}
\end{equation}
则称 $X$ 服从参数为 $\lambda$ 的\emph{指数分布（exponential distribution）}，记为 $X\sim\mathrm{Exp}(\lambda)$。本笔记采用\emph{率参数（rate parameter）}约定，其倒数为另一种常用的尺度参数。
\end{definition}

对 $x\ge0$，积分给出
\begin{equation}
  F_X(x)=1-e^{-\lambda x},\qquad
  S_X(x):=\mathbb P(X>x)=e^{-\lambda x}.
  \label{eq:study-exponential-cdf-tail}
\end{equation}
对 $x<0$，$F_X(x)=0$，$S_X(x)=1$。函数 $S_X$ 称为\emph{生存函数（survival function）}，适合表示寿命超过某一时长的概率。

\begin{theorem}[指数分布的无记忆性及刻画]
\label{thm:study-exponential-memoryless}
若 $X\sim\mathrm{Exp}(\lambda)$，则对任意 $s,t\ge0$，
\begin{equation}
  \mathbb P(X>s+t\mid X>t)=\mathbb P(X>s).
  \label{eq:study-exponential-memoryless}
\end{equation}
反之，若 $X$ 为非负、有限的连续型随机变量，$S_X(t)>0$ 对每个 $t\ge0$ 成立，且满足此等式，则 $X$ 服从某个指数分布。
\end{theorem}

\begin{proof}
正向等式为 $e^{-\lambda(s+t)}/e^{-\lambda t}=e^{-\lambda s}$。反向令 $H(t)=-\log S_X(t)$，无记忆性给出 $H(s+t)=H(s)+H(t)$，且 $H$ 非减。记 $\lambda=H(1)$，由可加性先得 $H(m/n)=\lambda m/n$，再用从两侧逼近实数的非负有理数及单调性，得到 $H(t)=\lambda t$。因为 $S_X(t)\to0$，$\lambda$ 不能为零，故 $S_X(t)=e^{-\lambda t}$，即得结论。
\end{proof}

几何分布刻画离散试验次数中的等待，指数分布刻画连续时间中的等待，两者的无记忆性形式相似，但取值范围和尾概率公式不同。若把 $X$ 解释为寿命，式~\eqref{eq:study-exponential-memoryless} 表明已经正常工作的时长不改变剩余寿命的条件分布；这一假设不适合直接描述明显存在老化效应的对象。

进一步，若非负寿命 $X$ 在 $(0,\infty)$ 上有连续密度，且 $S_X(t)>0$，则在 $t>0$ 处可定义\emph{失效率（hazard rate）}
\begin{equation}
  r(t)=\lim_{h\downarrow0}
    \frac{\mathbb P(t<X\le t+h\mid X>t)}{h}
  =\frac{f_X(t)}{S_X(t)},
  \label{eq:study-hazard-rate}
\end{equation}
失效率是每单位时间的条件失效强度，$r(t)h$ 才是短时间内条件失效概率的近似。由 $S_X'(t)=-f_X(t)$，若 $S_X(0)=1$，则
\begin{equation}
  S_X(t)=\exp\left(-\int_0^tr(u)\,\mathrm du\right),\qquad
  f_X(t)=r(t)\exp\left(-\int_0^tr(u)\,\mathrm du\right).
  \label{eq:study-hazard-survival}
\end{equation}
对一般局部可积的非负 $r$，这些等式在相应的几乎处处意义下成立；要使寿命有限的概率为 $1$，还须有 $\int_0^{\infty}r(u)\,\mathrm du=\infty$。常数 $r(t)=\lambda$ 导出指数分布；若 $r(t)=kt^{k-1}/\eta^k$（$t>0$），其中 $k,\eta>0$，则
\[
  F_X(t)=1-\exp\bigl(-(t/\eta)^k\bigr),\qquad t\ge0,
\]
得到\emph{韦布尔分布（Weibull distribution）}。$k=1$ 时它成为指数分布，$k>1$ 和 $0<k<1$ 分别允许递增和递减的失效率。

\subsection{正态分布与标准化}
\label{subsec:study-normal}

\begin{definition}
\label{def:study-normal}
设 $\mu\in\mathbb R$、$\sigma>0$。若随机变量 $X$ 的密度为
\begin{equation}
  f_X(x)=\frac1{\sqrt{2\pi}\sigma}
    \exp\left(-\frac{(x-\mu)^2}{2\sigma^2}\right),
  \qquad x\in\mathbb R,
  \label{eq:study-normal-density}
\end{equation}
则称 $X$ 服从\emph{正态分布（normal distribution）}，也称\emph{高斯分布（Gaussian distribution）}，记为 $X\sim N(\mu,\sigma^2)$。其中 $\mu$ 是位置参数，$\sigma$ 是尺度参数；记号中的第二个参数是 $\sigma^2$。
\end{definition}

密度关于 $x=\mu$ 对称，在该点达到最大值 $1/(\sqrt{2\pi}\sigma)$，拐点位于 $\mu\pm\sigma$。固定 $\sigma$ 改变 $\mu$ 使曲线平移；固定 $\mu$ 增大 $\sigma$ 则使曲线变宽、峰值降低，但总面积仍为 $1$。取 $\mu=0,\sigma=1$ 得到\emph{标准正态分布（standard normal distribution）}，其密度和分布函数分别记为
\begin{equation}
  \varphi(z)=\frac1{\sqrt{2\pi}}e^{-z^2/2},\qquad
  \Phi(z)=\int_{-\infty}^z\varphi(t)\,\mathrm dt.
  \label{eq:study-standard-normal}
\end{equation}
正态密度的归一化可由高斯积分说明：令 $I=\int_{\mathbb R}e^{-x^2/2}\,\mathrm dx$，由非负被积函数的重积分公式和极坐标变换，
\[
  I^2=\int_{\mathbb R^2}e^{-(x^2+y^2)/2}\,\mathrm dx\,\mathrm dy
  =\int_0^{2\pi}\!\int_0^{\infty}e^{-r^2/2}r\,\mathrm dr\,\mathrm d\theta
  =2\pi.
\]
故 $I=\sqrt{2\pi}$，再经平移和尺度变换即得一般情形。

\begin{proposition}[标准化公式]
\label{prop:study-normal-standardization}
若 $X\sim N(\mu,\sigma^2)$，则 $Z=(X-\mu)/\sigma\sim N(0,1)$，且对 $a<b$，
\begin{equation}
  F_X(x)=\Phi\left(\frac{x-\mu}{\sigma}\right),\qquad
  \mathbb P(a<X\le b)=
    \Phi\left(\frac{b-\mu}{\sigma}\right)
    -\Phi\left(\frac{a-\mu}{\sigma}\right).
  \label{eq:study-normal-standardization}
\end{equation}
\end{proposition}

\begin{proof}
在 $F_X(x)=\int_{-\infty}^xf_X(t)\,\mathrm dt$ 中作代换 $u=(t-\mu)/\sigma$，即得分布函数公式。又由 $\sigma>0$，$\mathbb P(Z\le z)=F_X(\mu+\sigma z)=\Phi(z)$，区间概率由两端分布函数之差得到。
\end{proof}

由对称性，$\Phi(-z)=1-\Phi(z)$，因而
\begin{equation}
  \mathbb P(|X-\mu|\le k\sigma)=2\Phi(k)-1\quad(k>0).
  \label{eq:study-normal-central-probability}
\end{equation}
取 $k=1,2,3$，相应概率约为 $0.6827,0.9545,0.9973$。因此所谓“三倍标准差范围”包含很大概率，但区间外仍有约 $0.0027$ 的概率，不能把正态随机变量的取值范围严格截为 $[\mu-3\sigma,\mu+3\sigma]$。

由于 $\varphi(z)>0$，$\Phi$ 连续且严格递增，并把实轴映到 $(0,1)$，其反函数记为 $\Phi^{-1}$，称为标准正态分布的\emph{分位数函数（quantile function）}。对 $u\in(0,1)$，$\Phi^{-1}(u)$ 是使左侧累积概率等于 $u$ 的唯一实数；因而可把概率约束反过来转化为高度、时间等数值的设计条件。

% 教材例 2.23：车门高度。
\begin{example}
\label{ex:study-normal-door}
假定所研究人群的男子身高 $X\sim N(172,36)$，单位为厘米。以身高超过车门高度作为碰头事件，怎样选择高度 $H$，使碰头概率小于 $0.01$？
\end{example}

\begin{solve}
这里 $\sigma=6$，所需条件为
\[
  1-\Phi\left(\frac{H-172}{6}\right)<0.01,
\]
等价于 $H>172+6\Phi^{-1}(0.99)\approx185.958$ 厘米。因此选 $H=186$ 厘米即可，模型下的超高概率约为 $0.0098$。若身高模型或碰头条件改变，应重新判断；这一计算是给定模型下的概率设计。
\end{solve}

\section{随机变量函数的分布}
\label{sec:study-transformed-distributions}

若已知 $X$ 的分布，而研究对象是 $Y=g(X)$，就需要把 $X$ 的概率经过确定的函数 $g$ 转移到新的数轴上。这里假定 $g$ 为可测函数，从而 $Y$ 仍是随机变量。变换后的类型由 $X$ 与 $g$ 共同决定，连续型随机变量经过函数变换后也可能变成离散型或混合型。

\subsection{离散变换与分布函数法}
\label{subsec:study-transformation-cdf}

若 $X$ 离散，$\mathbb P(X=x_k)=p_k$，则对每个可能的新取值 $y$，
\begin{equation}
  \mathbb P(Y=y)=\sum_{k:g(x_k)=y}p_k.
  \label{eq:study-discrete-transformation}
\end{equation}
\emph[tip]{先列出不同的新取值，再合并映到同一取值的概率。}例如教材例 2.24 中，$X$ 在 $-1,0,1,3,4$ 上的概率依次为 $0.1,0.2,0.2,0.4,0.1$，则 $Y=X^2$ 的分布律为
\[
  \begin{array}{c|cccc}
    y&0&1&9&16\\\hline
    \mathbb P(Y=y)&0.2&0.3&0.4&0.1
  \end{array}
\]
其中 $Y=1$ 同时包含 $X=-1$ 和 $X=1$ 两种情形。若 $g$ 恒等于常数 $c$，则 $\mathbb P(Y=c)=1$，称为\emph{退化分布（degenerate distribution）}。

对任意类型的 $X$，都可先用\emph{分布函数法（distribution function method）}：
\begin{equation}
  F_Y(y)=\mathbb P(g(X)\le y).
  \label{eq:study-general-transformation-cdf}
\end{equation}
其关键是解出集合 $\{x:g(x)\le y\}$。若 $X$ 有密度 $f_X$，则
\begin{equation}
  F_Y(y)=\int_{\{x:g(x)\le y\}}f_X(x)\,\mathrm dx.
  \label{eq:study-continuous-transformation-cdf}
\end{equation}
求得 $F_Y$ 后，应先检查它是否存在跳跃或其他不属于绝对连续部分的概率，再决定能否仅以求导得到完整的密度。

\subsection{单调变换及分段单调变换}
\label{subsec:study-density-transformation}

\begin{theorem}[单调变换的密度公式]
\label{thm:study-monotone-density-transformation}
设 $X$ 具有密度 $f_X$，且 $\mathbb P(X\in I)=1$，其中 $I$ 为开区间。若 $g$ 在 $I$ 上连续可微且 $g'(x)\ne0$，则 $g$ 在 $I$ 上严格单调，反函数 $h:g(I)\to I$ 连续可微，$Y=g(X)$ 具有密度
\begin{equation}
  f_Y(y)=\begin{cases}
    f_X(h(y))|h'(y)|,&y\in g(I),\\
    0,&y\notin g(I).
  \end{cases}
  \label{eq:study-monotone-density-transformation}
\end{equation}
密度等式按几乎处处意义理解，区间端点的密度值可任意指定为非负值。
\end{theorem}

\begin{proof}
连续且不为零的 $g'$ 在区间内符号恒定，故 $g$ 严格单调，反函数定理给出 $h'(y)=1/g'(h(y))$。若 $g$ 递增，则在 $g(I)$ 内有 $F_Y(y)=F_X(h(y))$；若 $g$ 递减，则 $F_Y(y)=1-F_X(h(y))$，因为 $X$ 的单点概率为零。更直接地，在任一子区间 $(c,d)\subseteq g(I)$ 上作积分换元，两种情况均得到
\[
  \mathbb P(c<Y\le d)=\int_c^df_X(h(y))|h'(y)|\,\mathrm dy.
\]
所得函数在整个像区间的积分为 $1$，因此确为 $Y$ 的密度。绝对值来自递减变换时积分方向的反转。
\end{proof}

特别地，当 $Y=aX+b$、$a\ne0$ 时，
\begin{equation}
  f_Y(y)=\frac1{|a|}f_X\left(\frac{y-b}{a}\right).
  \label{eq:study-affine-density}
\end{equation}
若 $X\sim N(\mu,\sigma^2)$，代入正态密度可得
\begin{equation}
  aX+b\sim N(a\mu+b,a^2\sigma^2).
  \label{eq:study-affine-normal}
\end{equation}
这一公式也适用于 $a<0$；若 $a=0$，则新分布退化在 $b$，须单独处理。

若 $X$ 以概率 $1$ 落在有限或可数个互不相交的开区间 $I_j$ 的并集中，$g$ 在每个 $I_j$ 上满足上述条件，反函数记为 $h_j$，则将各分支的概率相加得到
\begin{equation}
  f_Y(y)=\sum_{j:y\in g(I_j)}f_X(h_j(y))|h_j'(y)|.
  \label{eq:study-piecewise-transformation}
\end{equation}
对可数个分支，可由非负函数的求和与积分交换验证此式；分割时漏掉的边界点必须具有总概率零。不能把一个多对一变换任意选一个反函数后直接代入单调公式。

% 教材例 2.26：平方变换。
\begin{example}
\label{ex:study-square-transformation}
设 $X$ 的密度为 $f_X$，分布函数为 $F_X$。求 $Y=X^2$ 的分布。
\end{example}

\begin{solve}
$Y\ge0$ 且 $\mathbb P(Y=0)=0$。当 $y>0$ 时，
\[
  F_Y(y)=\mathbb P(-\sqrt y\le X\le\sqrt y)
  =F_X(\sqrt y)-F_X(-\sqrt y),
\]
而 $y\le0$ 时 $F_Y(y)=0$。在正负两个半轴分别换元，得到密度
\begin{equation}
  f_Y(y)=\begin{cases}
    \dfrac{f_X(\sqrt y)+f_X(-\sqrt y)}{2\sqrt y},&y>0,\\
    0,&y\le0.
  \end{cases}
  \label{eq:study-square-density}
\end{equation}
当 $X\sim N(0,1)$ 时，两分支概率贡献相同，因而 $f_Y(y)=e^{-y/2}/\sqrt{2\pi y}$（$y>0$）。该密度在 $0$ 的右侧趋于无穷，但可积，且不在 $0$ 处产生正概率。
\end{solve}

% 教材例 2.27：正弦变换。
\begin{example}
\label{ex:study-sine-transformation}
设 $X\sim U(0,\pi)$，求 $Y=\sin X$ 的密度。
\end{example}

\begin{solve}
对 $0<y<1$，不等式 $\sin x\le y$ 在 $(0,\pi)$ 中对应两个区间 $(0,\arcsin y]$ 和 $[\pi-\arcsin y,\pi)$，所以
\[
  F_Y(y)=\frac{2\arcsin y}{\pi},\qquad
  f_Y(y)=\frac{2}{\pi\sqrt{1-y^2}}.
\]
在 $y\le0$ 时 $F_Y(y)=0$，在 $y\ge1$ 时 $F_Y(y)=1$；密度在 $(0,1)$ 外可取零。这与式~\eqref{eq:study-piecewise-transformation} 中两个反函数分支的贡献之和相同。
\end{solve}

\begin{notealert}[type=caution,title={变换可能改变分布类型}]
若 $X\sim U(0,1)$，而 $Y=\mathbf 1_{\{X>1/2\}}$，则 $Y\sim\mathrm{Bern}(1/2)$，连续型变量变成了离散型。若令 $Z=\max\{X,1/2\}$，则 $\mathbb P(Z=1/2)=1/2$，其余概率在 $(1/2,1)$ 上连续分布，得到混合型分布。变换在一个有正概率的区域内取常值时，必须保留这一整块概率，不能只对可微部分套用密度公式。
\end{notealert}

\section{扩展阅读：正态分布的由来}
\label{sec:study-normal-origin}

教材从二项分布的近似和测量误差的建模两条线索解释正态分布的来源，分别联系棣莫弗、拉普拉斯关于极限分布的研究和高斯关于误差分析的研究。理解这些线索，重点在于说明钟形密度从什么数学条件产生，而不是把所有观测数据都预设成正态分布。

\subsection{从二项概率到正态极限}
\label{subsec:study-normal-binomial-origin}

先考虑 $S_n\sim\mathrm{Bin}(n,1/2)$。当 $n$ 为偶数，且 $k=n/2+j$ 时，
\[
  \mathbb P(S_n=k)=\binom{n}{n/2+j}2^{-n}.
\]
利用斯特林公式 $m!\sim\sqrt{2\pi m}(m/e)^m$，再对对数作展开，可在 $|j|\le C\sqrt n$（$C>0$ 固定）的中心范围内得到一致近似
\begin{equation}
  \mathbb P(S_n=n/2+j)
  =\sqrt{\frac2{\pi n}}\exp\left(-\frac{2j^2}{n}\right)
    \bigl(1+o(1)\bigr).
  \label{eq:study-binomial-local-normal}
\end{equation}
这里给出的是推导思路：令 $u=2j/n=O(n^{-1/2})$，阶乘比的主要指数项为
\[
  -\frac n2\bigl[(1+u)\log(1+u)+(1-u)\log(1-u)\bigr]
  =-\frac{nu^2}{2}+O(nu^4)
  =-\frac{2j^2}{n}+O(n^{-1}),
\]
平方根因子则给出 $\sqrt{2/(\pi n)}(1+o(1))$，这说明中心尺度为 $\sqrt n$，并解释了平方指数的来源。

令 $z_j=2j/\sqrt n$，相邻标准化取值间距为 $\Delta z=2/\sqrt n$，式~\eqref{eq:study-binomial-local-normal} 的右侧正好为 $\varphi(z_j)\Delta z$。因而对固定的有限区间，将相应的二项概率相加，就是正态密度的黎曼和。这给出了离散概率为何会趋近连续曲线下面积的解释。一般成功概率下的完整结论为\emph{棣莫弗--拉普拉斯定理（de Moivre--Laplace theorem）}：若 $S_n\sim\mathrm{Bin}(n,p)$，$p\in(0,1)$ 固定，则
\begin{equation}
  \lim_{n\to\infty}
  \mathbb P\left(\frac{S_n-np}{\sqrt{np(1-p)}}\le z\right)
  =\Phi(z),\qquad z\in\mathbb R.
  \label{eq:study-de-moivre-laplace}
\end{equation}
一般情形的证明留待第~\ref{sec:study-limit-theorems} 节的极限定理讨论。有限 $n$ 时，这是概率近似；尤其在 $p$ 接近 $0$ 或 $1$ 时，不能仅以 $n$ 较大为由忽略分布偏斜。

\subsection{从算术平均的最优性反推误差密度}
\label{subsec:study-normal-error-origin}

另一条线索从测量误差出发。设同一未知真值为 $\theta\in\mathbb R$，测量值为 $x_1,\ldots,x_n$，各次随机误差相互独立、具有相同的严格正密度 $f$，且 $f$ 在实轴上连续可微。给定观测值后，将其联合密度看成 $\theta$ 的函数，得到\emph{似然函数（likelihood function）}
\[
  L(\theta)=\prod_{i=1}^nf(x_i-\theta).
\]
\emph{最大似然估计（maximum likelihood estimation）}选择使此函数最大的参数值。现在提出一个强的建模要求：对每个样本量及每组实数观测，算术平均 $\bar x=n^{-1}\sum_i x_i$ 都是似然函数的最大点。由于参数范围为整个实轴，最大点须满足可微函数的一阶必要条件。

令 $\psi(u)=f'(u)/f(u)$。将 $\theta=\bar x$ 代入对数似然的一阶导数，得到
\begin{equation}
  \sum_{i=1}^n\psi(x_i-\bar x)=0.
  \label{eq:study-gaussian-score-condition}
\end{equation}
在 $n=2$ 时取观测 $u,-u$，得到 $\psi(-u)=-\psi(u)$；在 $n=3$ 时取观测 $u,v,-u-v$，再用奇性，得到
\[
  \psi(u+v)=\psi(u)+\psi(v).
\]
连续的可加函数必为线性函数：先由可加性在有理数上得到 $\psi(r)=cr$，再用连续性推广到所有实数。因此 $\psi(u)=cu$，积分得到
\[
  \log f(u)=\frac c2u^2+C.
\]
若 $c\ge0$，所得正函数不能在整个实轴上归一化为概率密度；故 $c<0$，令 $c=-1/\sigma^2$，再用高斯积分归一化，得到
\begin{equation}
  f(u)=\frac1{\sqrt{2\pi}\sigma}
    \exp\left(-\frac{u^2}{2\sigma^2}\right).
  \label{eq:study-gaussian-error-characterization}
\end{equation}
反过来，在这一密度下，最大化 $L(\theta)$ 等价于最小化平方误差和，因为
\[
  \sum_{i=1}^n(x_i-\theta)^2
  =\sum_{i=1}^n(x_i-\bar x)^2+n(\theta-\bar x)^2,
\]
其唯一最小点确为 $\bar x$。这说明正态误差模型与算术平均、最小二乘法之间有严格联系；结论依赖独立同分布、正密度的光滑性以及对所有样本都要求平均值最优等条件，不能仅凭某一次采用了平均数便断定误差服从正态分布。

两条线索也有不同的逻辑地位：二项极限描述许多独立随机贡献相加后的近似规律，误差推导则在明确的最优性要求下刻画密度。把复杂误差视为许多微小因素之和，可以为正态近似提供动机，但仍须核对相应极限定理的条件；偏斜、重尾或受共同因素影响的数据，未必适合直接使用正态模型。

\section*{本章总结}
\addcontentsline{toc}{section}{本章总结}

随机变量是样本空间上的数值函数，分布刻画其取值的概率规律。离散型分布由取值与概率质量确定，连续型分布由密度的积分确定，各类分布都能用 $F_X(x)=\mathbb P(X\le x)$ 描述。计算区间概率时，应先判断端点是否有概率质量；只有没有相应单点概率时，才可自由改变端点开闭。密度是概率对长度的集中程度，既不等于单点概率，也不要求逐点小于 $1$。

常用模型按计数方式区分：伯努利分布记录一次试验是否成功，二项分布记录固定次数中的成功数，超几何分布记录有限总体不放回抽样中的某类对象数，负二项与几何分布记录达到成功条件的试验次数，泊松分布描述给定范围内的计数。均匀分布把区间概率与长度联系起来，指数分布对应具有无记忆性的连续等待，正态分布则通过标准化归结为 $\Phi$ 的计算。分布选择依赖条件，泊松和正态极限提供近似依据，不能把近似等同于精确模型。

求 $Y=g(X)$ 的分布时，离散情形将相同像的概率相加；一般情形先写 $F_Y(y)=\mathbb P(g(X)\le y)$，再解出原像集合。只有确认绝对连续性和单调分支条件后，才使用带导数绝对值的密度变换公式。下面选取教材习题练习有限抽样、跳跃端点、模型计算和分布变换，题号均指教材本章习题。

\section*{习题}
\addcontentsline{toc}{section}{习题}

\begin{example}
\label{ex:study-ch2-exercise-6}
教材习题 6。$10$ 件产品中有 $8$ 件正品、$2$ 件次品，每次从剩余产品中等可能地抽取一件，直到首次取到正品。设 $T$ 为抽取次数，求其分布律和分布函数。
\end{example}

\begin{solve}
由于至多先抽到两件次品，$T$ 只能取 $1,2,3$。逐次使用条件概率，得到
\[
  \mathbb P(T=1)=\frac45,\qquad
  \mathbb P(T=2)=\frac2{10}\frac89=\frac8{45},\qquad
  \mathbb P(T=3)=\frac2{10}\frac19=\frac1{45}.
\]
三项之和为 $1$，因而
\[
  F_T(t)=\begin{cases}
    0,&t<1,\\
    4/5,&1\le t<2,\\
    44/45,&2\le t<3,\\
    1,&t\ge3.
  \end{cases}
\]
虽然 $T$ 是首次成功的等待次数，但不放回抽样使后续成功率随已抽结果改变，因此不能直接使用几何分布。
\end{solve}

\begin{example}
\label{ex:study-ch2-exercise-8}
教材习题 8。独立射击 $20$ 次，每次命中率为 $0.2$，以 $X$ 记命中次数。求至少命中一次的概率，并求 $X$ 最有可能的取值。
\end{example}

\begin{solve}
$X\sim\mathrm{Bin}(20,0.2)$，所以 $\mathbb P(X\ge1)=1-0.8^{20}\approx0.9885$。设 $p_k=\mathbb P(X=k)$，对 $0\le k\le19$，
\[
  \frac{p_{k+1}}{p_k}=\frac{20-k}{k+1}\frac{0.2}{0.8}
  =\frac{20-k}{4(k+1)}.
\]
此比值在 $k\le3$ 时大于 $1$，在 $k\ge4$ 时小于 $1$，故 $p_k$ 先增加至 $k=4$，之后下降，唯一最可能的命中次数为 $4$。
\end{solve}

\begin{example}
\label{ex:study-ch2-exercise-18}
教材习题 18。设随机变量 $X$ 的分布函数为
\[
  F(x)=\begin{cases}
    0,&x<0,\\
    x/4,&0\le x<1,\\
    1/2+(x-1)/4,&1\le x<2,\\
    5/6,&2\le x<3,\\
    1,&x\ge3.
  \end{cases}
\]
求 $\mathbb P(X=k)$（$k=1,2,3$）及 $\mathbb P(1/2<X<3/2)$。
\end{example}

\begin{solve}
逐点取跳跃高度，而非代入分布函数后直接当作单点概率，得
\[
  \mathbb P(X=1)=\frac12-\frac14=\frac14,\quad
  \mathbb P(X=2)=\frac56-\frac34=\frac1{12},\quad
  \mathbb P(X=3)=1-\frac56=\frac16.
\]
两个区间端点都是连续点，因而
\[
  \mathbb P(1/2<X<3/2)=F(3/2)-F(1/2)=\frac58-\frac18=\frac12.
\]
这一概率已经包含 $X=1$ 的概率质量，不能只把两个端点间的分段导数积分后就停止计算。
\end{solve}

\begin{example}
\label{ex:study-ch2-exercise-24}
教材习题 24。设 $X\sim U(-5,5)$，求关于实数 $u$ 的方程 $u^2+Xu+1=0$ 有实根的概率。
\end{example}

\begin{solve}
系数 $X$ 随机，未知数 $u$ 是方程中的实数变量。有实根等价于判别式 $X^2-4\ge0$，即 $X\in(-5,-2]\cup[2,5)$。由均匀分布，所求概率是这些区间总长度与 $10$ 之比，即 $6/10=3/5$。边界 $X=\pm2$ 的概率为零，不改变结果。
\end{solve}

\begin{example}
\label{ex:study-ch2-exercise-28}
教材习题 28。机器检修时间 $X\sim\mathrm{Exp}(1)$，单位为小时。求检修时间超过 $2$ 小时的概率；已知检修 $2$ 小时仍未完成，求总检修时间超过 $4$ 小时的条件概率。
\end{example}

\begin{solve}
由生存函数及无记忆性，
\[
  \mathbb P(X>2)=e^{-2},\qquad
  \mathbb P(X>4\mid X>2)=\frac{e^{-4}}{e^{-2}}=e^{-2}.
\]
二者均约为 $0.1353$，但描述的信息条件不同；未给定已等待时间时，总检修时间超过 $4$ 小时的概率应为 $e^{-4}$。
\end{solve}

\begin{example}
\label{ex:study-ch2-exercise-31}
教材习题 31（1）。设 $X\sim N(1,4)$，求 $\mathbb P(0\le X\le4)$、$\mathbb P(X>2.4)$ 和 $\mathbb P(|X|>2)$。
\end{example}

\begin{solve}
这里 $\sigma=2$，依次标准化得到
\[
  \begin{aligned}
    \mathbb P(0\le X\le4)&=\Phi(1.5)-\Phi(-0.5)\approx0.6247,\\
    \mathbb P(X>2.4)&=1-\Phi(0.7)\approx0.2420,\\
    \mathbb P(|X|>2)&=\Phi(-1.5)+1-\Phi(0.5)\approx0.3753.
  \end{aligned}
\]
最后一个事件需分成 $X<-2$ 与 $X>2$，因为分布中心为 $1$，不能直接把两尾视为相等。
\end{solve}

\begin{example}
\label{ex:study-ch2-exercise-42}
教材习题 42。设随机变量 $X$ 的分布函数 $F$ 连续且在实轴上严格递增，证明 $U=F(X)$ 服从 $U(0,1)$。
\end{example}

\begin{proof}
由 $F$ 连续、严格递增及两端极限，对每个 $u\in(0,1)$，存在唯一的 $x_u\in\mathbb R$，使 $F(x_u)=u$。因此
\[
  \mathbb P(U\le u)=\mathbb P(F(X)\le u)
  =\mathbb P(X\le x_u)=F(x_u)=u.
\]
又对有限的 $X$ 有 $0<F(X)<1$，故 $u\le0$ 时分布函数为 $0$，$u\ge1$ 时为 $1$，与均匀分布相同。这一结论称为\emph{概率积分变换（probability integral transform）}。严格递增是此处使用普通反函数的充分条件；如果分布函数有平坦区间，就需要另外处理反函数的定义。
\end{proof}

\begin{example}
\label{ex:study-ch2-exercise-44}
教材习题 44。设 $X$ 的密度为 $f_X(x)=2(1-x)$（$0<x<1$），其余处为零。构造 $(0,1)$ 上的单调递增函数 $g$，使 $g(X)\sim\mathrm{Exp}(1)$。
\end{example}

\begin{solve}
在 $(0,1)$ 上积分得 $F_X(x)=2x-x^2=1-(1-x)^2$，它连续且严格递增。由概率积分变换的同样论证，$U=F_X(X)\sim U(0,1)$。对 $y\ge0$，
\[
  \mathbb P(-\log(1-U)\le y)
  =\mathbb P(U\le1-e^{-y})=1-e^{-y}.
\]
所以可取
\begin{equation}
  g(x)=-\log(1-F_X(x))=-2\log(1-x),\qquad0<x<1.
  \label{eq:study-exercise-44-transform}
\end{equation}
由于 $g'(x)=2/(1-x)>0$，它确为所需的单调递增函数。这也展示了从一种连续分布先变到均匀分布、再变到目标分布的构造方法。
\end{solve}

\chapter{多维随机变量及其分布}
\label{chap:study-multivariate}
% 教材范围：第 3 章，第 88--122 页；正文六节及章末代表习题。

只知道身高的分布和体重的分布，还不能回答“身高较高的人是否更可能体重较大”。这类问题既涉及各个量的取值规律，也涉及它们之间的依赖关系，需要把同一次随机试验中的多个观测量合在一起研究。本章沿教材第 3 章的顺序，先用联合分布描述整体，再从中导出边缘分布与条件分布，随后讨论独立性以及多个随机变量的函数。联合、边缘、条件这三个层次贯穿全章：联合分布保留共同变化的信息，边缘分布只观察部分分量，条件分布则描述给定信息后的概率规律。

\section{多维随机变量及其分布}
\label{sec:study-joint-distributions}

\subsection{随机向量、联合分布函数与矩形概率}

\begin{definition}
\label{def:study-ch3-random-vector}
设 $X_1,\ldots,X_n$ 是定义在同一概率空间 $(\Omega,\mathcal F,\mathbb P)$ 上的实随机变量，称
\[
  \boldsymbol X=(X_1,\ldots,X_n)
\]
为 $n$ 维\emph{随机向量（random vector）}，也称 $n$ 维随机变量。其\emph{联合分布函数（joint cumulative distribution function）}定义为
\begin{equation}
  F_{\boldsymbol X}(x_1,\ldots,x_n)
  =\mathbb P(X_1\le x_1,\ldots,X_n\le x_n),
  \qquad (x_1,\ldots,x_n)\in\mathbb R^n.
  \label{eq:study-ch3-joint-cdf-n}
\end{equation}
\end{definition}

“同一概率空间”表示这些分量来自同一个随机结果 $\omega$，但不预设分量之间独立。以下主要讨论二维情形，记
\[
  F(x,y)=F_{X,Y}(x,y)=\mathbb P(X\le x,Y\le y).
\]
逗号表示两个条件同时成立。与一维分布函数相似，$F$ 也通过对概率累加来描述分布；与一维不同的是，要分别控制两个坐标，并保留它们共同取值的信息。

\begin{proposition}
\label{prop:study-ch3-joint-cdf-rectangle}
若 $a<b$、$c<d$，则
\begin{equation}
  \mathbb P(a<X\le b,c<Y\le d)
  =F(b,d)-F(a,d)-F(b,c)+F(a,c).
  \label{eq:study-ch3-joint-cdf-rectangle}
\end{equation}
因此，联合分布函数对每个坐标分别单调不减、分别右连续，并且满足
\[
  0\le F(x,y)\le1,\qquad
  \lim_{x\to-\infty}F(x,y)=\lim_{y\to-\infty}F(x,y)=0,
  \qquad \lim_{t\to+\infty}F(t,t)=1,
\]
其中前两个极限分别固定另一个坐标。它还满足\emph{矩形增量非负}：式~\eqref{eq:study-ch3-joint-cdf-rectangle} 的右端对任意 $a<b,c<d$ 都不小于 $0$。
\end{proposition}

\begin{proof}
从事件 $\{X\le b,Y\le d\}$ 中去掉 $\{X\le a,Y\le d\}$ 与 $\{X\le b,Y\le c\}$，再加回被重复去掉的 $\{X\le a,Y\le c\}$，便得到矩形概率公式。单调性来自事件的包含关系。若 $x_k\downarrow x$，则 $\{X\le x_k,Y\le y\}$ 递减至 $\{X\le x,Y\le y\}$，由概率对递减事件列的连续性得到关于 $x$ 的右连续性，关于 $y$ 的证明相同。坐标趋于 $-\infty$ 时，对应事件递减至空集；$\{X\le t,Y\le t\}$ 随 $t\to+\infty$ 递增至 $\Omega$，由此得到极限。矩形增量非负则直接来自概率非负。
\end{proof}

矩形端点的开闭与分布函数采用的“$\le$”配套，离散情形不能随意改变。\emph[caution]{分别单调不减不能代替矩形增量非负。}例如，教材例 3.1 所用的反例可写成
\[
  H(x,y)=
  \begin{cases}
    1,&x\ge0,\ y\ge0,\ x+y\ge1,\\
    0,&\text{其他}.
  \end{cases}
\]
它分别单调不减、分别右连续，并满足上述边界极限；但取 $a=c=0,b=d=1$，矩形增量为 $1-1-1+0=-1$，所以它不可能是联合分布函数。二维分布的合法性还必须检查不同坐标之间的相容性。

\subsection{离散型联合分布与多项分布}

\begin{definition}
\label{def:study-ch3-joint-pmf}
若 $X,Y$ 分别只取有限或可数个值 $x_i,y_j$，则称 $(X,Y)$ 为二维离散型随机向量。其\emph{联合概率质量函数（joint probability mass function）}或联合分布律为
\[
  p_{ij}=\mathbb P(X=x_i,Y=y_j),\qquad i,j=1,2,\ldots.
\]
这些数满足 $p_{ij}\ge0$、$\sum_i\sum_jp_{ij}=1$；反过来，满足这两个条件的一组数给出一个离散型联合分布。对任意集合 $A\subseteq\mathbb R^2$，
\begin{equation}
  \mathbb P((X,Y)\in A)=\sum_{(x_i,y_j)\in A}p_{ij},\qquad
  F(x,y)=\sum_{x_i\le x}\sum_{y_j\le y}p_{ij}.
  \label{eq:study-ch3-joint-pmf-probability}
\end{equation}
\end{definition}

取值表中的零概率格点也有意义，它们说明某些分量组合不能出现。将每个分量的所有可能值作笛卡尔积，只是便于列表，并不表示每种组合都有正概率。

\begin{example}
\label{ex:study-ch3-joint-exam}
教材例 3.2：某学生数学成绩优秀的概率为 $0.6$；数学优秀时，物理优秀的条件概率为 $0.8$；数学未达优秀时，物理也未达优秀的条件概率为 $0.7$。令 $X,Y$ 分别表示数学、物理是否优秀，优秀记为 $1$，否则记为 $0$。求联合分布。
\end{example}

\begin{solve}
用乘法公式逐项计算，例如
\[
  \mathbb P(X=1,Y=1)
  =\mathbb P(X=1)\mathbb P(Y=1\mid X=1)=0.48.
\]
其余三项同理，得到
\[
  \begin{array}{c|cc}
    X\backslash Y&0&1\\\hline
    0&0.28&0.12\\
    1&0.12&0.48
  \end{array}
\]
四项之和为 $1$。表中概率来自题设的条件规律，不能在没有独立性依据时直接把两个边缘概率相乘。
\end{solve}

\begin{definition}
\label{def:study-ch3-multinomial}
一次试验恰有 $m$ 种互斥且穷尽的结果，其概率为 $p_1,\ldots,p_m$，其中 $p_i\ge0$、$\sum_{i=1}^mp_i=1$。将试验独立重复 $N$ 次，令 $X_i$ 表示第 $i$ 种结果出现的次数，则 $\boldsymbol X=(X_1,\ldots,X_m)$ 服从\emph{多项分布（multinomial distribution）}，记作 $M(N;p_1,\ldots,p_m)$，其联合分布律为
\begin{equation}
  \mathbb P(X_1=k_1,\ldots,X_m=k_m)
  =\frac{N!}{k_1!\cdots k_m!}\prod_{i=1}^mp_i^{k_i},
  \qquad k_i\in\{0,1,\ldots,N\},\quad \sum_{i=1}^mk_i=N,
  \label{eq:study-ch3-multinomial-pmf}
\end{equation}
其他取值的概率为 $0$，并约定出现 $p_i=0,k_i=0$ 时 $0^0=1$。
\end{definition}

\begin{proof}
对满足次数约束的一组 $(k_1,\ldots,k_m)$，先选择各类结果在 $N$ 次试验中的位置，共有 $N!/(k_1!\cdots k_m!)$ 种排列。每种排列由独立性具有相同概率 $\prod_i p_i^{k_i}$，对这些互斥排列求和即得式~\eqref{eq:study-ch3-multinomial-pmf}。归一性也可由多项式展开 $(p_1+\cdots+p_m)^N=1$ 验证。
\end{proof}

当 $m=2$ 时，只要知道 $X_1$ 就知道 $X_2=N-X_1$，因而多项分布化为二项分布。\emph[caution]{各次试验独立，不表示各类结果的总次数相互独立。}等式 $X_1+\cdots+X_m=N$ 已经给出了分量之间的约束。

\subsection{联合密度、区域概率与常见模型}

\begin{definition}
\label{def:study-ch3-joint-density}
若存在非负可积函数 $f(x,y)$，使任意 $(x,y)\in\mathbb R^2$ 都满足
\begin{equation}
  F(x,y)=\int_{-\infty}^{x}\int_{-\infty}^{y}f(u,v)\,\mathrm dv\,\mathrm du,
  \label{eq:study-ch3-joint-density-cdf}
\end{equation}
则称 $(X,Y)$ 为具有联合密度的二维连续型随机向量，$f=f_{X,Y}$ 称为\emph{联合概率密度函数（joint probability density function）}。它满足
\begin{equation}
  \int_{\mathbb R^2}f(x,y)\,\mathrm dx\,\mathrm dy=1,\qquad
  \mathbb P((X,Y)\in A)=\iint_A f(x,y)\,\mathrm dx\,\mathrm dy
  \label{eq:study-ch3-joint-density-region}
\end{equation}
对任意 Borel 区域 $A$ 成立。
\end{definition}

联合密度是单位面积上的概率强度，$f(x,y)$ 本身不是 $\mathbb P(X=x,Y=y)$，也不必小于 $1$。密度存在时，单点、直线及其他二维面积为零的集合都具有概率 $0$；因此，积分区域边界上的开闭通常不影响结果。密度在零面积集合上的改动不改变分布，本章涉及密度相等的结论均按\emph{几乎处处（almost everywhere）}理解。若 $F$ 在某开区域内二阶连续可微，则由矩形增量公式及积分表示，该区域内的密度可取为
\[
  f(x,y)=\frac{\partial^2F}{\partial x\,\partial y}(x,y).
\]
在密度的连续点附近，对边长很小的矩形，其概率近似为 $f(x,y)\,\Delta x\,\Delta y$；这只是局部近似，精确概率仍由积分给出。上述定义与积分公式直接推广到 $n$ 维，将面积和二重积分分别换成体积与 $n$ 重积分。

\begin{example}
\label{ex:study-ch3-continuous-marginals-singular}
教材例 3.4：设 $X$ 在 $(0,1)$ 上均匀分布，令 $Y=X$。说明虽然两个分量都有连续型分布，$(X,Y)$ 却没有二维联合密度。
\end{example}

\begin{solve}
两个边缘分布均为 $U(0,1)$，但 $\mathbb P(Y=X)=1$。若存在二维联合密度，对直线 $y=x$ 积分应为 $0$，与这一概率等于 $1$ 矛盾。实际上，其联合分布函数为
\[
  F(x,y)=\max\{0,\min\{x,y,1\}\},
\]
它是连续函数，概率却全部集中在对角线段上。因此，联合分布函数连续、每个分量有密度，都不能单独保证存在联合密度。
\end{solve}

若平面区域 $D$ 的面积 $|D|$ 满足 $0<|D|<\infty$，在 $D$ 上的\emph{二维均匀分布（bivariate uniform distribution）}定义为
\begin{equation}
  f(x,y)=\frac1{|D|}\boldsymbol 1_D(x,y),\qquad
  \mathbb P((X,Y)\in A)=\frac{|A\cap D|}{|D|},
  \label{eq:study-ch3-uniform-region}
\end{equation}
其中 $\boldsymbol 1_D$ 为示性函数，点在 $D$ 内时取 $1$，否则取 $0$。均匀指单位面积上的概率强度相同，分量的边缘密度却可能并不均匀，因为不同位置的截面长度可能不同。

本章另一常见模型是\emph{二元正态分布（bivariate normal distribution）}。给定 $\mu_1,\mu_2\in\mathbb R$、$\sigma_1,\sigma_2>0$、$-1<\rho<1$，若联合密度为
\begin{equation}
  f(x,y)=\frac1{2\pi\sigma_1\sigma_2\sqrt{1-\rho^2}}
  \exp\!\left\{-\frac1{2(1-\rho^2)}
    \left[\frac{(x-\mu_1)^2}{\sigma_1^2}
    -\frac{2\rho(x-\mu_1)(y-\mu_2)}{\sigma_1\sigma_2}
    +\frac{(y-\mu_2)^2}{\sigma_2^2}\right]\right\},
  \label{eq:study-ch3-bivariate-normal-density}
\end{equation}
则记作 $(X,Y)\sim N(\mu_1,\mu_2,\sigma_1^2,\sigma_2^2,\rho)$。这里的 $\rho$ 控制两个分量的依赖结构，其与相关系数的关系在后续数字特征章节中说明；本式要求 $|\rho|<1$，不包含概率集中在直线上的退化情形。令
\[
  \boldsymbol\mu=\begin{pmatrix}\mu_1\\\mu_2\end{pmatrix},\qquad
  \Sigma=\begin{pmatrix}\sigma_1^2&\rho\sigma_1\sigma_2\\
    \rho\sigma_1\sigma_2&\sigma_2^2\end{pmatrix},
\]
则 $\det\Sigma=\sigma_1^2\sigma_2^2(1-\rho^2)>0$，密度也可写成
\[
  f(\boldsymbol x)=\frac{\exp\{-\tfrac12(\boldsymbol x-\boldsymbol\mu)^{\mathsf T}
    \Sigma^{-1}(\boldsymbol x-\boldsymbol\mu)\}}
    {2\pi\sqrt{\det\Sigma}}.
\]
对 $n$ 维非退化正态分布，将 $\Sigma$ 换成 $n$ 阶对称正定矩阵，分母中的 $2\pi$ 换成 $(2\pi)^{n/2}$，即得到相应联合密度；转置记号 $\mathsf T$ 用来区分列向量及其转置。

\begin{example}
\label{ex:study-ch3-joint-exponential-region}
教材例 3.7：设 $f(x,y)=e^{-(x+y)}$，$x>0,y>0$，其他位置取 $0$。求联合分布函数以及 $\mathbb P(X+Y\le1)$。
\end{example}

\begin{solve}
对正半轴上的变量分别积分，得到
\[
  F(x,y)=
  \begin{cases}
    (1-e^{-x})(1-e^{-y}),&x>0,y>0,\\
    0,&\text{其他}.
  \end{cases}
\]
事件 $X+Y\le1$ 在密度非零区域内对应三角形 $0<x<1,0<y<1-x$，故
\[
  \mathbb P(X+Y\le1)
  =\int_0^1\int_0^{1-x}e^{-(x+y)}\,\mathrm dy\,\mathrm dx
  =\int_0^1(e^{-x}-e^{-1})\,\mathrm dx
  =1-2e^{-1}\approx0.2642.
\]
积分限必须同时反映事件条件和密度的非零区域。
\end{solve}

\section{边缘（际）分布}
\label{sec:study-marginal-distributions}

\subsection{从联合分布消去其余分量}

\begin{definition}
\label{def:study-ch3-marginal-distributions}
联合分布中部分分量自身的分布，称为相应的\emph{边缘分布（marginal distribution）}或边际分布。对于 $(X,Y)$，由事件列的连续性有
\begin{equation}
  F_X(x)=\lim_{y\to+\infty}F(x,y),\qquad
  F_Y(y)=\lim_{x\to+\infty}F(x,y).
  \label{eq:study-ch3-marginal-cdf}
\end{equation}
若联合分布律为 $p_{ij}$，则
\begin{equation}
  p_{i\cdot}=\mathbb P(X=x_i)=\sum_jp_{ij},\qquad
  p_{\cdot j}=\mathbb P(Y=y_j)=\sum_ip_{ij}.
  \label{eq:study-ch3-marginal-pmf}
\end{equation}
若联合密度为 $f(x,y)$，则边缘密度可取为
\begin{equation}
  f_X(x)=\int_{-\infty}^{+\infty}f(x,y)\,\mathrm dy,\qquad
  f_Y(y)=\int_{-\infty}^{+\infty}f(x,y)\,\mathrm dx.
  \label{eq:study-ch3-marginal-density}
\end{equation}
\end{definition}

边缘化保留所考察的分量，对其余分量的全部可能值求和或积分。对密度公式，可以先用非负函数积分的 Tonelli 定理交换积分次序，得到
\[
  F_X(x)=\int_{-\infty}^x
    \left(\int_{-\infty}^{+\infty}f(u,y)\,\mathrm dy\right)\mathrm du,
\]
由一维密度的定义即可识别括号内的函数，无须额外假定可以逐点在积分号下求导。多维情形也完全相同：若 $\boldsymbol U\in\mathbb R^k$、$\boldsymbol V\in\mathbb R^{n-k}$ 具有联合密度 $f(\boldsymbol u,\boldsymbol v)$，则
\[
  f_{\boldsymbol U}(\boldsymbol u)
  =\int_{\mathbb R^{n-k}}f(\boldsymbol u,\boldsymbol v)\,\mathrm d\boldsymbol v.
\]
从 $n$ 个坐标选取非空真子集共有 $2^n-2$ 种，因此按坐标子集计数有 $2^n-2$ 个边缘分布；这不保证这些分布作为函数都互不相同。

\begin{example}
\label{ex:study-ch3-triangle-marginals}
设 $(X,Y)$ 在三角形 $D=\{(x,y):0<x<y<1\}$ 上均匀分布。求两个边缘密度，并解释其形状。
\end{example}

\begin{solve}
三角形面积为 $1/2$，所以联合密度为 $2\boldsymbol 1_D(x,y)$。固定 $x\in(0,1)$ 时，$y$ 从 $x$ 变化至 $1$；固定 $y\in(0,1)$ 时，$x$ 从 $0$ 变化至 $y$。由式~\eqref{eq:study-ch3-marginal-density}，
\[
  f_X(x)=2(1-x)\boldsymbol 1_{(0,1)}(x),\qquad
  f_Y(y)=2y\boldsymbol 1_{(0,1)}(y).
\]
两个边缘密度分别随截面长度递减、递增，均非均匀分布；它们的积分都为 $1$。这说明区域内的联合均匀性不能直接传给各个坐标。
\end{solve}

\subsection{边缘分布通常不能恢复联合分布}

\begin{example}
\label{ex:study-ch3-cards-marginals}
教材例 3.8：袋中有 $3$ 张标为 $0$、$2$ 张标为 $1$ 的卡片。依次随机抽取两张，用 $X,Y$ 表示两次抽到的数字。比较有放回与不放回抽取的联合分布及边缘分布。
\end{example}

\begin{solve}
有放回时，第二次仍从原来的五张卡片中抽取；不放回时，第二次只剩四张。因此两种联合分布分别为
\[
  \begin{array}{c|cc|c}
    X\backslash Y&0&1&p_{i\cdot}\\\hline
    0&9/25&6/25&3/5\\
    1&6/25&4/25&2/5\\\hline
    p_{\cdot j}&3/5&2/5&1
  \end{array}
  \qquad
  \begin{array}{c|cc|c}
    X\backslash Y&0&1&p_{i\cdot}\\\hline
    0&6/20&6/20&3/5\\
    1&6/20&2/20&2/5\\\hline
    p_{\cdot j}&3/5&2/5&1
  \end{array}
\]
两个模型的边缘分布完全相同，联合分布却不同；例如 $\mathbb P(X=1,Y=1)$ 分别为 $4/25$ 与 $1/10$。第二次抽取的边缘分布相同，并不表示第二次的结果不受第一次的影响。
\end{solve}

对于二元正态分布，令 $u=(x-\mu_1)/\sigma_1$、$v=(y-\mu_2)/\sigma_2$，利用
\[
  u^2-2\rho uv+v^2=(1-\rho^2)u^2+(v-\rho u)^2
\]
对 $y$ 积分，第二项对应的正态密度积分为 $1$，于是式~\eqref{eq:study-ch3-bivariate-normal-density} 的边缘分布为
\begin{equation}
  X\sim N(\mu_1,\sigma_1^2),\qquad
  Y\sim N(\mu_2,\sigma_2^2).
  \label{eq:study-ch3-bivariate-normal-marginals}
\end{equation}
这也验证了联合正态密度的归一性。两个边缘密度均不含 $\rho$，所以即使限制在二元正态分布族内，边缘分布也不能确定联合分布。

更进一步，两个分量各自正态不保证它们联合正态。固定 $0<r<1$，将标准边缘、参数分别为 $r$ 与 $-r$ 的两种二元正态密度记为 $f_r,f_{-r}$，则混合密度 $\tfrac12f_r+\tfrac12f_{-r}$ 的两个边缘仍为标准正态。若它等于某个标准边缘的二元正态密度 $f_\rho$，由于该混合密度关于 $x$、$y$ 各自对称，式~\eqref{eq:study-ch3-bivariate-normal-density} 中的交叉项要求 $\rho=0$；但它在原点的值为 $1/(2\pi\sqrt{1-r^2})$，不同于 $f_0(0,0)=1/(2\pi)$。这些密度连续，因此不能仅靠零面积集合上的取值差异消除矛盾。这给出了教材例 3.10 所用的正态混合反例。

\section{条件分布}
\label{sec:study-conditional-distributions}

\subsection{条件分布、联合分布与贝叶斯公式}

条件分布描述已知某一分量取值后，另一分量的概率规律。对离散型随机向量，若 $p_{\cdot j}=\mathbb P(Y=y_j)>0$，直接应用事件的条件概率定义，有
\begin{equation}
  \mathbb P(X=x_i\mid Y=y_j)=\frac{p_{ij}}{p_{\cdot j}}.
  \label{eq:study-ch3-conditional-pmf}
\end{equation}
固定条件 $Y=y_j$ 后，对 $i$ 求和为 $1$，因此它确实是一个概率分布。相应地，给定 $X=x_i$ 时，$Y$ 的条件分布律为 $p_{ij}/p_{i\cdot}$，要求 $p_{i\cdot}>0$。在联合分布表中，条件化就是取定一行或一列，再除以该行或该列的总概率。

连续型情形需另作处理：有密度时 $\mathbb P(Y=y)=0$，不能把 $\mathbb P(X\in A,Y=y)/\mathbb P(Y=y)$ 当作定义。教材从窄条事件 $\{y<Y\le y+h\}$ 出发说明条件密度的含义；严格使用这种极限时，还需相应密度与积分随 $y$ 变化的正则性。更直接的处理是用联合密度与边缘密度的比值给出条件分布，并核对它能恢复原来的联合概率。

\begin{definition}
\label{def:study-ch3-conditional-density}
设 $(X,Y)$ 具有联合密度 $f(x,y)$。对满足 $0<f_Y(y)<\infty$ 的 $y$，定义给定 $Y=y$ 时 $X$ 的\emph{条件概率密度函数（conditional probability density function）}为
\begin{equation}
  f_{X\mid Y}(x\mid y)=\frac{f(x,y)}{f_Y(y)},\qquad
  F_{X\mid Y}(x\mid y)=\int_{-\infty}^x f_{X\mid Y}(u\mid y)\,\mathrm du.
  \label{eq:study-ch3-conditional-density}
\end{equation}
给定 $X=x$ 时，同样有 $f_{Y\mid X}(y\mid x)=f(x,y)/f_X(x)$，要求 $0<f_X(x)<\infty$。
\end{definition}

固定合法的 $y$ 后，$f_{X\mid Y}(x\mid y)$ 非负且对 $x$ 的积分为 $1$。这里 $y$ 是条件参数，归一化时不能对 $y$ 积分代替对 $x$ 积分。边缘密度为零或不有限的取值构成对相应边缘分布概率为零的集合，条件分布在这些取值上可另行约定而不改变联合概率；因此，连续型条件分布由联合分布确定到几乎处处相同，并非在每个单点上天然唯一。

由定义得到\emph{密度的乘法公式}与\emph{全概率公式}：
\begin{equation}
  f(x,y)=f_{Y\mid X}(y\mid x)f_X(x)
        =f_{X\mid Y}(x\mid y)f_Y(y),\qquad
  f_Y(y)=\int_{\mathbb R}f_{Y\mid X}(y\mid x)f_X(x)\,\mathrm dx.
  \label{eq:study-ch3-density-product-total}
\end{equation}
例如，对 Borel 集 $A,B$，将第一种分解代入区域积分便得
\[
  \mathbb P(X\in A,Y\in B)
  =\int_B\left(\int_A f_{X\mid Y}(x\mid y)\,\mathrm dx\right)f_Y(y)\,\mathrm dy,
\]
这正是条件分布与联合分布相容的含义。交换已知条件与所研究的变量，还得到密度形式的\emph{贝叶斯公式（Bayes' formula）}：
\begin{equation}
  f_{X\mid Y}(x\mid y)
  =\frac{f_{Y\mid X}(y\mid x)f_X(x)}
    {\displaystyle\int_{\mathbb R}f_{Y\mid X}(y\mid u)f_X(u)\,\mathrm du},
  \label{eq:study-ch3-density-bayes}
\end{equation}
其中分母必须为有限正数。以上公式也适用于向量条件，只需对被消去的整个向量积分。

\begin{example}
\label{ex:study-ch3-two-stage-uniform}
教材例 3.12：先在 $(0,1)$ 中均匀选取 $X$，给定 $X=x$ 后，再在 $(0,x)$ 中均匀选取 $Y$。求 $Y$ 的密度和 $F_{X\mid Y}(1/2\mid1/4)$。
\end{example}

\begin{solve}
由抽样机制可知 $f_X(x)=1$，$0<x<1$，且 $f_{Y\mid X}(y\mid x)=1/x$，$0<y<x$。于是
\[
  f(x,y)=\frac1x\boldsymbol 1_{\{0<y<x<1\}}(x,y),\qquad
  f_Y(y)=\int_y^1\frac{\mathrm dx}{x}
       =-\ln y,\quad0<y<1,
\]
而 $f_Y$ 在其他位置取 $0$。给定 $0<y<1$，
\[
  f_{X\mid Y}(x\mid y)=\frac1{x(-\ln y)}\boldsymbol 1_{(y,1)}(x),
  \qquad
  F_{X\mid Y}(1/2\mid1/4)
  =\frac{\ln(1/2)-\ln(1/4)}{-\ln(1/4)}=\frac12.
\]
联合取值区域是三角形，但密度为 $1/x$，所以这个模型并非三角形上的均匀分布。给定 $X=x$ 时 $Y$ 均匀，也不表示给定 $Y=y$ 时 $X$ 均匀。
\end{solve}

\subsection{正态模型中的条件分布}

\begin{proposition}
\label{prop:study-ch3-normal-conditional}
若 $(X,Y)\sim N(\mu_1,\mu_2,\sigma_1^2,\sigma_2^2,\rho)$，其中 $|\rho|<1$，则
\begin{equation}
  Y\mid X=x\sim N\!\left(
    \mu_2+\rho\frac{\sigma_2}{\sigma_1}(x-\mu_1),
    \sigma_2^2(1-\rho^2)\right).
  \label{eq:study-ch3-normal-conditional}
\end{equation}
交换两组下标即可得到 $X\mid Y=y$ 的分布。
\end{proposition}

\begin{proof}
用式~\eqref{eq:study-ch3-bivariate-normal-density} 除以 $X$ 的边缘密度，再对 $y$ 配平方，得到
\[
  f_{Y\mid X}(y\mid x)
  =\frac1{\sqrt{2\pi}\sigma_2\sqrt{1-\rho^2}}
  \exp\!\left\{-\frac{
    [y-\mu_2-\rho\sigma_2(x-\mu_1)/\sigma_1]^2}
    {2\sigma_2^2(1-\rho^2)}\right\},
\]
识别其正态密度参数即可。分母在每个 $x\in\mathbb R$ 都为正，因此这里可选取对所有实数 $x$ 都有定义的这一版本。
\end{proof}

条件均值随 $x$ 线性变化，条件方差不随 $x$ 变化，并小于或等于边缘方差 $\sigma_2^2$。这是联合正态模型的具体性质，不能只凭两个边缘正态就使用式~\eqref{eq:study-ch3-normal-conditional}。

教材例 3.13 给出另一个按条件规律建模的例子：以厘米为单位的身高 $H\sim N(172,36)$，给定 $H=h$ 后，以千克为单位的体重满足 $W\mid H=h\sim N(h-105,49)$。定义残差 $E=W-(H-105)$，其条件分布恒为 $N(0,49)$，所以对任意 Borel 集 $A,B$，
\[
  \mathbb P(E\in A,H\in B)
  =\int_B\mathbb P(E\in A\mid H=h)f_H(h)\,\mathrm dh
  =\mathbb P(E\in A)\mathbb P(H\in B).
\]
这说明 $E$ 与 $H$ 独立；应用第~\ref{sec:study-vector-transformations} 节将推导的独立正态变量相加公式，得 $W=H-105+E\sim N(67,85)$。也可直接对乘积 $f_{W\mid H}(w\mid h)f_H(h)$ 关于 $h$ 积分并配平方得到相同结论。这里 $H,W$ 分别表示以厘米、千克计的数值，线性关系中的系数和常数均以这一单位约定为前提。

\section{相互独立的随机变量}
\label{sec:study-independent-variables}

\subsection{独立性的定义与等价判据}

\begin{definition}
\label{def:study-ch3-independent-variables}
若对任意 $x,y\in\mathbb R$ 都有
\begin{equation}
  F_{X,Y}(x,y)=F_X(x)F_Y(y),
  \label{eq:study-ch3-independence-cdf}
\end{equation}
则称随机变量 $X,Y$ \emph{相互独立（independent）}。等价地，对任意 Borel 集 $A,B\subseteq\mathbb R$，
\begin{equation}
  \mathbb P(X\in A,Y\in B)=\mathbb P(X\in A)\mathbb P(Y\in B).
  \label{eq:study-ch3-independence-borel}
\end{equation}
\end{definition}

式~\eqref{eq:study-ch3-independence-borel} 表明，关于 $X$ 的任意可测事件与关于 $Y$ 的任意可测事件都独立。取 $A=(-\infty,x],B=(-\infty,y]$ 就得到式~\eqref{eq:study-ch3-independence-cdf}；反向结论可先用矩形增量公式推广到半开区间，再由概率测度在这类生成集合上的唯一性推广到 Borel 集。因此，独立性要求整个分布的乘积分解，验证某一对事件独立通常不够。

对离散型随机向量，独立性等价于 $p_{ij}=p_{i\cdot}p_{\cdot j}$ 对所有 $i,j$ 成立。必要性由单点事件直接得到；充分性由对式~\eqref{eq:study-ch3-joint-pmf-probability} 求和得到。对具有联合密度的随机向量，独立性等价于
\begin{equation}
  f(x,y)=f_X(x)f_Y(y)\quad\text{几乎处处成立}.
  \label{eq:study-ch3-independence-density}
\end{equation}
充分性由矩形积分分离得到，必要性则由乘积密度与原联合分布在所有矩形上相同，以及密度的几乎处处唯一性得到。独立时，在合法的条件取值上，可以选取等于边缘分布的条件分布版本；例如 $f_{Y\mid X}(y\mid x)=f_Y(y)$ 对几乎所有 $x$ 成立。

判断可分离性时须包括整个取值区域。若 $f(x,y)=a(x)b(y)$ 几乎处处成立，其中 $a,b$ 为非负可积函数，且其积分 $A,B$ 为正，则联合归一性给出 $AB=1$，两个边缘分别为 $Ba(x)$、$Ab(y)$，故 $X,Y$ 独立。但只在一个三角形内把非零表达式写成乘积，并不能完成这一分解，因为区域的示性函数还将 $x,y$ 耦合在一起。

\begin{example}
\label{ex:study-ch3-disk-dependence}
教材例 3.16：设 $(X,Y)$ 在单位圆盘 $x^2+y^2<1$ 内均匀分布，证明 $X,Y$ 不独立。
\end{example}

\begin{solve}
令 $A=\{X>3/4\}$、$B=\{Y>3/4\}$。圆盘内分别有正面积的部分满足这两个条件，所以 $\mathbb P(A)>0$、$\mathbb P(B)>0$；但两条件同时成立要求 $x^2+y^2>18/16>1$，因此 $\mathbb P(A\cap B)=0$。这与独立时应有 $\mathbb P(A\cap B)=\mathbb P(A)\mathbb P(B)>0$ 矛盾。该证明直接使用正概率事件，避免只在某个密度取值点上判断独立性。
\end{solve}

对二元正态分布，$X,Y$ 独立当且仅当 $\rho=0$。若 $\rho=0$，式~\eqref{eq:study-ch3-bivariate-normal-density} 直接分解为两个边缘密度的乘积。反之，独立要求这两个连续密度处处相等；在 $(\mu_1,\mu_2)$ 比较其值，得到 $1/\sqrt{1-\rho^2}=1$，所以 $\rho=0$。这里用到了联合正态密度的具体形式；对一般分布，后续章节定义的相关系数为零并不足以推出独立。

\subsection{相互独立、两两独立与分组函数}

随机变量 $X_1,\ldots,X_n$ \emph{相互独立（mutually independent）}的定义是：对任意 $x_1,\ldots,x_n\in\mathbb R$，
\[
  F_{\boldsymbol X}(x_1,\ldots,x_n)=\prod_{i=1}^nF_{X_i}(x_i).
\]
对离散型随机向量，这等价于所有联合点概率等于各边缘点概率的乘积；对具有联合密度的随机向量，等价于联合密度几乎处处等于 $\prod_i f_{X_i}(x_i)$。相互独立必然推出每一子组相互独立，特别推出两两独立；其逆命题一般不成立。

\begin{example}
\label{ex:study-ch3-pairwise-not-mutual}
教材例 3.18：设 $X,Y$ 独立，且都以概率 $1/2$ 取 $-1,1$，令 $Z=XY$。证明 $X,Y,Z$ 两两独立但不相互独立。
\end{example}

\begin{solve}
$Z$ 也等概率取 $-1,1$。对任意 $x,z\in\{-1,1\}$，有
\[
  \mathbb P(X=x,Z=z)=\mathbb P(X=x,Y=xz)=\frac14
  =\mathbb P(X=x)\mathbb P(Z=z),
\]
所以 $X,Z$ 独立；交换 $X,Y$ 得 $Y,Z$ 独立，而 $X,Y$ 已由假设独立。然而
\[
  \mathbb P(X=1,Y=1,Z=-1)=0
  \ne\mathbb P(X=1)\mathbb P(Y=1)\mathbb P(Z=-1)=\frac18.
\]
第三个变量被前两个的乘积完全确定，因此三者不相互独立。
\end{solve}

若 $X_1,\ldots,X_n$ 相互独立，从两个不相交指标组分别构造的 Borel 可测函数
\[
  U=g(X_1,\ldots,X_k),\qquad
  V=h(X_{k+1},\ldots,X_n),\qquad1\le k<n,
\]
仍然独立。其依据是 $\{U\in A\}$ 与 $\{V\in B\}$ 分别属于两组独立坐标生成的事件族，因而概率可分解。特别地，$X,Y$ 独立时 $g(X),h(Y)$ 独立；函数内部不要求单调。但若两个函数共用了某个原变量，就不能仅凭原坐标独立而直接断言它们独立。

\begin{example}
\label{ex:study-ch3-meeting-probability}
教材例 3.15：两人分别在一小时内独立、均匀地到达，先到者等候 $20$ 分钟后离开。求能会面的概率。
\end{example}

\begin{solve}
以小时为单位，把起始时刻记为 $0$，则 $X,Y$ 独立且均服从 $U(0,1)$，联合密度在单位正方形内为 $1$。会面事件为 $|X-Y|\le1/3$，其补集由两个直角边均为 $2/3$ 的直角三角形组成，所以
\[
  \mathbb P(|X-Y|\le1/3)=1-2\cdot\frac12\left(\frac23\right)^2=\frac59.
\]
独立、均匀的到达机制是概率模型的假设；仅说“约在一小时内到达”并不能自动推出这些假设。
\end{solve}

\section{随机向量函数的分布}
\label{sec:study-vector-transformations}

\subsection{分布函数法与变量变换公式}

求 $Z=g(X,Y)$ 的分布，首先应把关于 $Z$ 的事件还原成关于 $(X,Y)$ 的事件。对 Borel 可测函数 $g$，离散情形为
\[
  \mathbb P(Z=z)=\sum_{g(x_i,y_j)=z}p_{ij},
\]
具有联合密度时则为
\begin{equation}
  F_Z(z)=\iint_{g(x,y)\le z}f(x,y)\,\mathrm dx\,\mathrm dy.
  \label{eq:study-ch3-vector-function-cdf}
\end{equation}
这种\emph{分布函数法（distribution function method）}不要求 $g$ 一一对应。先找出 $Z$ 的取值范围，再计算相应区域概率，可以避免分段条件遗漏；若所得分布绝对连续，再对 $F_Z$ 求导得到密度。向量函数 $(U,V)=(g_1(X,Y),g_2(X,Y))$ 也按同样原则求逆像区域的概率。

\begin{theorem}
\label{thm:study-ch3-jacobian-transform}
设 $(X,Y)$ 具有联合密度 $f$，且以概率 $1$ 落在开集 $D\subseteq\mathbb R^2$ 内。设变换 $T:D\to G=T(D)$ 为一一对应的连续可微映射，具有连续可微逆映射
\[
  x=x(u,v),\qquad y=y(u,v),
\]
并且雅可比行列式在 $D$ 内非零。则 $(U,V)=T(X,Y)$ 具有联合密度
\begin{equation}
  f_{U,V}(u,v)
  =f(x(u,v),y(u,v))
    \left|\frac{\partial(x,y)}{\partial(u,v)}\right|
    \boldsymbol 1_G(u,v),
  \label{eq:study-ch3-jacobian-density}
\end{equation}
其中在 $G$ 上计算右侧的密度与行列式，在 $G$ 外令密度为 $0$，且
\[
  \frac{\partial(x,y)}{\partial(u,v)}
  =\det\begin{pmatrix}
    \partial x/\partial u&\partial x/\partial v\\
    \partial y/\partial u&\partial y/\partial v
  \end{pmatrix}.
\]
\end{theorem}

\begin{proof}
对任意 Borel 区域 $A\subseteq G$，积分换元给出
\[
  \mathbb P((U,V)\in A)
  =\iint_{T^{-1}(A)}f(x,y)\,\mathrm dx\,\mathrm dy
  =\iint_A f(x(u,v),y(u,v))
    \left|\frac{\partial(x,y)}{\partial(u,v)}\right|
    \,\mathrm du\,\mathrm dv.
\]
由密度的定义即得结论。绝对值来自面积伸缩，不随变换是否改变方向而改变概率的正负。
\end{proof}

该公式同样适用于 $n$ 维可逆光滑变换，行列式相应变为 $n$ 阶。若变换在整个取值区域上不一一对应，可在有限个互不重叠、覆盖全部概率的分支上分别换元，再把同一点各个原像的密度贡献相加；不能只保留某一个逆分支。若原变量具有联合密度，概率为零的分支边界可以单独忽略。只求一个函数 $Z=g(X,Y)$ 时，常补选另一个函数作为辅助变量，先求二维联合密度，再边缘化；雅可比因子和新取值区域必须同时计算。

\subsection{和、差的分布与卷积}

\begin{proposition}
\label{prop:study-ch3-convolution}
设 $(X,Y)$ 具有联合密度 $f$。则 $S=X+Y$ 具有密度
\begin{equation}
  f_S(s)=\int_{\mathbb R}f(x,s-x)\,\mathrm dx.
  \label{eq:study-ch3-sum-density-general}
\end{equation}
若 $X,Y$ 独立且密度分别为 $f_X,f_Y$，则
\begin{equation}
  f_S(s)=\int_{\mathbb R}f_X(x)f_Y(s-x)\,\mathrm dx
        =(f_X*f_Y)(s),
  \label{eq:study-ch3-convolution}
\end{equation}
右侧称为两个密度的\emph{卷积（convolution）}。同理，$D=X-Y$ 的密度为 $f_D(d)=\int_{\mathbb R}f(d+y,y)\,\mathrm dy$。
\end{proposition}

\begin{proof}
令 $U=X,S=X+Y$，逆变换为 $x=u,y=s-u$，雅可比行列式绝对值为 $1$。由定理~\ref{thm:study-ch3-jacobian-transform}，$f_{U,S}(u,s)=f(u,s-u)$，再对 $u$ 积分即可。独立时将联合密度分解成 $f_Xf_Y$，得到卷积公式；令 $D=X-Y,V=Y$ 可得差的公式。这一推导直接使用换元和边缘化，不需要在积分号下求导。
\end{proof}

\begin{example}
\label{ex:study-ch3-exponential-sum-difference}
教材例 3.20：设 $X,Y$ 独立，均服从率参数为 $\lambda>0$ 的指数分布，即密度为 $\lambda e^{-\lambda x}\boldsymbol 1_{(0,\infty)}(x)$。求和与差的分布。
\end{example}

\begin{solve}
对 $S=X+Y$，密度非零要求 $x>0,s-x>0$，故 $s>0$ 时
\[
  f_S(s)=\int_0^s\lambda^2e^{-\lambda x}e^{-\lambda(s-x)}\,\mathrm dx
        =\lambda^2s e^{-\lambda s},
\]
而 $s\le0$ 时取 $0$。重复卷积可得：若 $X_1,\ldots,X_n$ 独立同分布于此指数分布，则
\begin{equation}
  f_{X_1+\cdots+X_n}(s)
  =\frac{\lambda^n}{(n-1)!}s^{n-1}e^{-\lambda s}
    \boldsymbol 1_{(0,\infty)}(s).
  \label{eq:study-ch3-integer-gamma-density}
\end{equation}
归纳时，把 $n$ 项和的密度与指数密度卷积，剩下的积分是 $\int_0^s x^{n-1}\,\mathrm dx=s^n/n$，由此得到 $n+1$ 项的公式。这是形状参数为 $n$、率参数为 $\lambda$ 的\emph{伽马分布（gamma distribution）}，记为 $\operatorname{Ga}(n,\lambda)$；若改用尺度参数，则尺度是 $1/\lambda$。

对 $D=X-Y$，$y>0$ 与 $d+y>0$ 要求 $y>\max\{0,-d\}$，所以
\[
  f_D(d)=\lambda^2e^{-\lambda d}
    \int_{\max\{0,-d\}}^\infty e^{-2\lambda y}\,\mathrm dy
  =\frac\lambda2e^{-\lambda|d|},\qquad d\in\mathbb R.
\]
这是以 $0$ 为位置参数、以 $1/\lambda$ 为尺度参数的\emph{拉普拉斯分布（Laplace distribution）}。和只取正值，差可取任意实值，因此两者的积分限不能混用。
\end{solve}

教材例 3.21 中，若 $X\sim U(0,2)$、$Y\sim U(-1,1)$ 且独立，则联合密度为矩形内的 $1/4$。卷积积分的长度为区间 $(0,2)\cap(s-1,s+1)$ 的长度，因而
\[
  f_{X+Y}(s)=
  \begin{cases}
    (s+1)/4,&-1<s<1,\\
    (3-s)/4,&1\le s<3,\\
    0,&\text{其他}.
  \end{cases}
\]
三角形密度来自两区间的重叠长度；分段位置由取值区间端点相遇的位置决定。

\begin{proposition}
\label{prop:study-ch3-normal-sum}
若 $X\sim N(\mu_1,\sigma_1^2)$、$Y\sim N(\mu_2,\sigma_2^2)$，且 $X,Y$ 独立、$\sigma_1,\sigma_2>0$，则
\begin{equation}
  X+Y\sim N(\mu_1+\mu_2,\sigma_1^2+\sigma_2^2).
  \label{eq:study-ch3-normal-sum}
\end{equation}
\end{proposition}

\begin{proof}
令 $a=\sigma_1^2$、$b=\sigma_2^2$、$t=s-\mu_1-\mu_2$、$r=x-\mu_1$。卷积指数中的平方和可分解为
\[
  \frac{r^2}{a}+\frac{(t-r)^2}{b}
  =\frac{a+b}{ab}\left(r-\frac{at}{a+b}\right)^2
   +\frac{t^2}{a+b}.
\]
对 $r$ 积分，第一项给出的高斯积分为 $\sqrt{2\pi ab/(a+b)}$，结合原密度的系数 $1/(2\pi\sqrt{ab})$，得到
\[
  f_{X+Y}(s)=\frac1{\sqrt{2\pi(a+b)}}
    \exp\!\left\{-\frac{(s-\mu_1-\mu_2)^2}{2(a+b)}\right\},
\]
即为所述正态密度。
\end{proof}

独立的同类随机变量相加后仍属于同一分布族的性质，称为该分布族的\emph{再生性（reproductive property）}。正态分布的均值、方差参数在独立相加时分别相加；同率参数的整数形状伽马分布，其形状参数也在独立相加时相加。这里的独立性或参数一致条件均属于结论的一部分。一般的两个正态边缘既不能保证式~\eqref{eq:study-ch3-normal-sum} 的方差公式，也不能保证和仍为正态。

\subsection{商、乘积与极坐标变换}

设 $(X,Y)$ 具有联合密度，则 $\mathbb P(Y=0)=0$，商 $Z=X/Y$ 可在 $Y=0$ 上任意补充定义。令 $Z=X/Y,V=Y$，逆变换为 $x=zv,y=v$，其雅可比行列式绝对值为 $|v|$。在 $v>0$、$v<0$ 两个分支分别换元，得到
\begin{equation}
  f_{X/Y}(z)=\int_{\mathbb R}|v|f(zv,v)\,\mathrm dv;
  \qquad
  f_{X/Y}(z)=\int_{\mathbb R}|v|f_X(zv)f_Y(v)\,\mathrm dv
  \quad\text{当 $X,Y$ 独立时}.
  \label{eq:study-ch3-ratio-density}
\end{equation}
绝对值不能去掉，特别是分母允许取负值时。同样令 $T=XY,V=Y$，逆变换为 $x=t/v,y=v$，得到乘积密度
\begin{equation}
  f_{XY}(t)=\int_{\mathbb R\setminus\{0\}}
    \frac1{|v|}f(t/v,v)\,\mathrm dv.
  \label{eq:study-ch3-product-density}
\end{equation}
这些积分公式按几乎处处成立理解，某个零测集上密度取无穷大不妨碍它作为可积密度。

\begin{example}
\label{ex:study-ch3-standard-cauchy-ratio}
教材例 3.23：若 $X,Y$ 独立且均服从 $N(0,1)$，求 $Z=X/Y$ 的密度。
\end{example}

\begin{solve}
由式~\eqref{eq:study-ch3-ratio-density}，
\[
  f_Z(z)=\frac1{2\pi}\int_{\mathbb R}|v|
    e^{-(1+z^2)v^2/2}\,\mathrm dv
  =\frac1\pi\int_0^\infty v e^{-(1+z^2)v^2/2}\,\mathrm dv
  =\frac1{\pi(1+z^2)},\quad z\in\mathbb R.
\]
这是\emph{标准柯西分布（standard Cauchy distribution）}。同样，两个独立、均值为零、方差相同的非退化正态变量之商也服从标准柯西分布，因为共同的尺度会在商中约去；不能把这一结论扩展成任意两个独立正态变量的商都服从标准柯西分布。
\end{solve}

\begin{example}
\label{ex:study-ch3-normal-polar}
教材例 3.24：若 $X,Y$ 独立且均服从 $N(0,1)$，令 $R=\sqrt{X^2+Y^2}$，极角 $\Theta\in[0,2\pi)$ 满足 $X=R\cos\Theta,Y=R\sin\Theta$。求 $(R,\Theta)$ 的分布。
\end{example}

\begin{solve}
去掉原点及选定的极角切割射线后，极坐标变换一一对应；被去掉的集合概率为零。逆变换的雅可比行列式绝对值为 $r$，所以
\[
  f_{R,\Theta}(r,\theta)
  =\frac1{2\pi}r e^{-r^2/2}
    \boldsymbol 1_{(0,\infty)}(r)\boldsymbol 1_{(0,2\pi)}(\theta).
\]
它分解为两个归一化密度的乘积，故 $R,\Theta$ 独立，且
\[
  f_R(r)=r e^{-r^2/2}\boldsymbol 1_{(0,\infty)}(r),\qquad
  \Theta\sim U(0,2\pi).
\]
$R$ 的分布称为尺度为 $1$ 的\emph{瑞利分布（Rayleigh distribution）}。联合密度本身的旋转对称性与极坐标的面积因子 $r$ 共同决定这一结果，不能漏掉 $r$ 后直接对原正态密度积分。
\end{solve}

\begin{example}
\label{ex:study-ch3-uniform-transformations}
教材例 3.25：设 $X_1,X_2,X_3$ 独立且均服从 $U(0,1)$，求 $Y=-a\ln X_1$（$a>0$）、$Z=X_2/X_1$ 和 $U=1/(X_1X_2X_3)$ 的分布。
\end{example}

\begin{solve}
对 $y>0$，$\mathbb P(Y\le y)=\mathbb P(X_1\ge e^{-y/a})=1-e^{-y/a}$，所以 $Y$ 为率参数 $1/a$ 的指数变量。对商 $Z$，在单位正方形内积分得
\[
  F_Z(z)=
  \begin{cases}
    0,&z\le0,\\
    z/2,&0<z<1,\\
    1-1/(2z),&z\ge1.
  \end{cases}
\]
其中 $0<z<1$ 时，事件为 $0<x_2\le zx_1$；$z\ge1$ 时，按 $x_1=1/z$ 分段，面积为 $\int_0^{1/z}zx_1\,\mathrm dx_1+\int_{1/z}^1\mathrm dx_1$。

对乘积的倒数取对数：$V=\ln U=\sum_{i=1}^3(-\ln X_i)$。由于独立变量各自的可测函数仍独立，各项均服从率参数为 $1$ 的指数分布，所以由式~\eqref{eq:study-ch3-integer-gamma-density}，$f_V(v)=v^2e^{-v}/2$，$v>0$。再用 $U=e^V$ 及 $\mathrm dv/\mathrm du=1/u$，得到
\[
  f_U(u)=\frac{(\ln u)^2}{2u^2}\boldsymbol 1_{(1,\infty)}(u).
\]
乘积问题经对数化后转成独立和的问题，适用于这里各因子都为正的情形。
\end{solve}

\subsection{最大值、最小值与系统可靠性}

\begin{proposition}
\label{prop:study-ch3-extreme-distributions}
设 $X_1,\ldots,X_n$ 相互独立，分布函数分别为 $F_1,\ldots,F_n$，令 $M=\max_iX_i$、$L=\min_iX_i$，则
\begin{equation}
  F_M(t)=\prod_{i=1}^nF_i(t),\qquad
  F_L(t)=1-\prod_{i=1}^n[1-F_i(t)].
  \label{eq:study-ch3-extreme-cdf}
\end{equation}
若各分量有密度 $f_i$，相应密度为
\begin{equation}
  f_M(t)=\sum_{i=1}^nf_i(t)\prod_{j\ne i}F_j(t),\qquad
  f_L(t)=\sum_{i=1}^nf_i(t)\prod_{j\ne i}[1-F_j(t)].
  \label{eq:study-ch3-extreme-density}
\end{equation}
\end{proposition}

\begin{proof}
$\{M\le t\}=\bigcap_i\{X_i\le t\}$，而 $\{L>t\}=\bigcap_i\{X_i>t\}$；分别利用相互独立并对后一事件取补集，得到分布函数公式。若各分量有密度，$F_i$ 绝对连续，有限乘积也绝对连续，按乘积法则几乎处处求导便得密度公式。
\end{proof}

若各分量同分布于 $F$、密度为 $f$，则 $f_M(t)=nf(t)F(t)^{n-1}$、$f_L(t)=nf(t)[1-F(t)]^{n-1}$。分布函数公式不要求原变量连续，离散情形只需用跳跃大小求点概率。例如，教材例 3.27 中两枚公平骰子独立掷出点数 $X,Y$，则
\[
  \mathbb P(\max\{X,Y\}=5)=\left(\frac56\right)^2-\left(\frac46\right)^2=\frac14,
\]
\[
  \mathbb P(\min\{X,Y\}=3)=\left(\frac46\right)^2-\left(\frac36\right)^2=\frac7{36}.
\]
求离散极值分布时，相邻的两个分布函数值或生存概率值相减，不能使用密度求导公式。

\begin{example}
\label{ex:study-ch3-system-reliability}
教材例 3.28：两个子系统的寿命 $X,Y$ 独立，分别服从率参数为 $\alpha>0$、$\beta>0$ 的指数分布。求串联、并联与备用三种连接方式下系统寿命的分布。
\end{example}

\begin{solve}
串联系统有一个子系统失效就停止工作，所以寿命为 $L=\min\{X,Y\}$，且
\[
  \mathbb P(L>t)=e^{-(\alpha+\beta)t},\qquad t\ge0,
\]
即 $L$ 仍为指数变量，率参数为 $\alpha+\beta$。并联系统两个子系统全部失效才停止工作，因此寿命为 $M=\max\{X,Y\}$，对 $t>0$ 有
\[
  F_M(t)=(1-e^{-\alpha t})(1-e^{-\beta t}),\qquad
  f_M(t)=\alpha e^{-\alpha t}+\beta e^{-\beta t}
        -(\alpha+\beta)e^{-(\alpha+\beta)t}.
\]
这里假设并联工作不会改变各自原定的寿命规律。

备用系统在第一个失效后启用第二个；在备用期间不老化、切换瞬时且可靠、两段工作寿命独立的模型下，总寿命为 $S=X+Y$。卷积给出，当 $t>0$ 且 $\alpha\ne\beta$ 时，
\[
  f_S(t)=\int_0^t\alpha e^{-\alpha x}\beta e^{-\beta(t-x)}\,\mathrm dx
        =\frac{\alpha\beta}{\beta-\alpha}
          (e^{-\alpha t}-e^{-\beta t}).
\]
若 $\alpha=\beta$，则 $f_S(t)=\alpha^2t e^{-\alpha t}$，不能把相等参数直接代入含 $\beta-\alpha$ 的分母。上述密度在 $t\le0$ 时均取 $0$。三个模型分别对应最小值、最大值和独立和，关键是先根据失效机制写出寿命函数，再选择分布公式。
\end{solve}

\section{扩展阅读：辛普森悖论}
\label{sec:study-simpson-paradox}

\emph{辛普森悖论（Simpson's paradox）}指分组比较与汇总比较可能出现相反方向：某一群体在每个分组中的成功比例都较高，汇总后成功比例却较低。教材用一个明确标为虚构的招生例子说明这一现象，数据如下。这里的“录取率”只是录取人数占申请人数的比例，不能仅凭这些比例判断性别是否造成录取差异。

\begin{table}
  \centering
  \caption{教材中的虚构招生数据}
  \label{tab:study-ch3-simpson-admissions}
  \begin{tabular}{ccrrc}
    \hline
    学院&性别&录取人数&未录取人数&录取率\\\hline
    法学院&男&8&45&$8/53\approx15.1\%$\\
    法学院&女&51&101&$51/152\approx33.6\%$\\
    商学院&男&201&50&$201/251\approx80.1\%$\\
    商学院&女&92&9&$92/101\approx91.1\%$\\\hline
    合计&男&209&95&$209/304\approx68.8\%$\\
    合计&女&143&110&$143/253\approx56.5\%$\\\hline
  \end{tabular}
\end{table}

由表~\ref{tab:study-ch3-simpson-admissions} 可见，每个学院的女生录取率都高于男生，但总体录取率低于男生。为看清其数学原因，从这些申请记录中等概率抽取一人，以 $A$ 表示“被录取”，$G$ 表示性别，$C$ 表示所申请的学院。条件全概率公式给出
\begin{equation}
  \mathbb P(A\mid G=g)
  =\sum_c\mathbb P(A\mid G=g,C=c)\mathbb P(C=c\mid G=g).
  \label{eq:study-ch3-simpson-total-probability}
\end{equation}
各组录取率的比较使用 $\mathbb P(A\mid G=g,C=c)$，总体则是以各自申请学院的比例为权重的加权平均。具体地，男生的法学院申请比例为 $53/304$，女生为 $152/253$，女生更多集中在录取率较低的法学院；男生更多集中在录取率较高的商学院。因此
\[
  \frac8{53}\frac{53}{304}+\frac{201}{251}\frac{251}{304}
  =\frac{209}{304}
  >\frac{143}{253}
  =\frac{51}{152}\frac{152}{253}+\frac{92}{101}\frac{101}{253}.
\]
分组优势不能跨越不同的权重直接传递到总体。如果对男女使用同一组非负且和为 $1$ 的学院权重，逐组均较高的一方，其加权平均也较高；所以这里的反转来自分组比例与权重共同作用，不是概率运算矛盾。

汇总比例和条件比例回答不同问题，均可正确描述各自对象。将男女按相同的学院权重比较，是一种明确了参照结构的标准化比较，但选择何种权重、是否应控制学院，还取决于要回答的实际问题及数据生成机制。\emph[caution]{条件概率的差异不直接等于因果效应。}是否存在不公平待遇，需结合申请选择、资格条件、录取规则等信息进一步研究，不能单凭汇总比例或单凭分组比例下结论。

\section*{本章总结}
\addcontentsline{toc}{section}{本章总结}

本章从联合分布出发形成三条联系。第一，联合分布经过求和、积分或让其余坐标趋于 $+\infty$ 得到边缘分布，因而联合决定边缘；边缘通常丢失依赖信息，不能反推联合。第二，条件分布通过对已知条件归一化得到，联合密度、边缘密度与条件密度满足式~\eqref{eq:study-ch3-density-product-total}；连续型条件取值往往是零概率事件，必须在密度意义下处理。第三，独立性使联合分布分解为边缘分布的乘积，从而既能由边缘恢复联合，也能简化函数分布的计算。

求函数分布时，先把事件写回原变量的取值区域，再选择分布函数法或可逆变量变换。变换后的密度必须同时包含原像处的联合密度、逆变换雅可比行列式的绝对值和新取值区域；辅助变量最后通过边缘化消去。和使用卷积，商使用式~\eqref{eq:study-ch3-ratio-density}，极值优先使用事件的交与补集。最后检查非负性、归一性、参数范围与取值范围，并核对独立性是否确由假设或证明得到。辛普森悖论则说明，从联合信息退回边缘信息时，比较对象和权重也随之改变。

\section*{习题}
\addcontentsline{toc}{section}{习题}

以下选取教材第 3 章习题中的代表题，分别练习联合计数、条件化、混合模型、函数分布与独立性。题号沿用教材；解答补明计算所需的随机抽样或独立性假设。

\begin{example}
\label{ex:study-ch3-exercise-1}
教材第 3 章习题 1：箱中有 $6$ 个球，红、白、黑球分别为 $1,2,3$ 个，不放回地随机取出 $2$ 个。以 $X,Y$ 分别表示取出的红球数与白球数，求联合分布。
\end{example}

\begin{solve}
每个由两个球组成的无序子集等可能，共有 $\binom62=15$ 个。给定 $X=r,Y=s$，还需取 $2-r-s$ 个黑球，因此
\[
  p_{rs}=\frac{\binom1r\binom2s\binom3{2-r-s}}{\binom62},
  \qquad r=0,1,\quad s=0,1,2,
\]
约定组合数下标不合法时取 $0$。具体联合分布律为
\[
  \begin{array}{c|ccc}
    X\backslash Y&0&1&2\\\hline
    0&3/15&6/15&1/15\\
    1&3/15&2/15&0
  \end{array}
\]
总和为 $1$。教材所求的联合分布函数可统一写成
\[
  F_{X,Y}(x,y)=\sum_{r=0}^1\sum_{s=0}^2
    p_{rs}\boldsymbol 1_{\{r\le x\}}\boldsymbol 1_{\{s\le y\}}.
\]
这个形式保留了所有跳跃位置，无须把平面分成大量小区域逐块列式。
\end{solve}

\begin{example}
\label{ex:study-ch3-exercise-6}
教材第 3 章习题 6：设射手各次射击独立，命中概率均为 $p\in(0,1)$，直到第二次命中为止。用 $X,Y$ 分别表示第一次、第二次命中时累计射击的次数，求联合分布及边缘分布。
\end{example}

\begin{solve}
记 $q=1-p$。$X=i,Y=j$ 要求第 $i,j$ 次命中，前 $j$ 次中其余 $j-2$ 次均未命中，所以
\[
  \mathbb P(X=i,Y=j)=p^2q^{j-2},\qquad1\le i<j,
\]
其余取值概率为 $0$。对另一个指标求和，得到
\[
  \mathbb P(X=i)=\sum_{j=i+1}^\infty p^2q^{j-2}=pq^{i-1},\quad i\ge1,
  \qquad
  \mathbb P(Y=j)=(j-1)p^2q^{j-2},\quad j\ge2.
\]
第一个是首次命中等待次数的几何分布，第二个是第二次命中等待次数的负二项分布。在给定 $Y=j$ 时，$X$ 在 $1,\ldots,j-1$ 上均匀分布，概率均为 $1/(j-1)$。尽管射击结果逐次独立，$X,Y$ 却不独立；例如 $\mathbb P(X=2)>0$、$\mathbb P(Y=2)>0$，但二者不可能同时等于 $2$。
\end{solve}

\begin{example}
\label{ex:study-ch3-exercise-19}
教材第 3 章习题 19：设 $X$ 服从率参数为 $1$ 的指数分布，且给定 $X=x$ 时，$Y$ 服从参数为 $x$ 的泊松分布。求 $Y$ 的分布律。
\end{example}

\begin{solve}
条件概率为 $\mathbb P(Y=k\mid X=x)=e^{-x}x^k/k!$，$k=0,1,\ldots$。这里 $X$ 连续、$Y$ 离散，不能把联合模型当作二维连续密度。对条件分布按 $X$ 的密度积分得
\[
  \mathbb P(Y=k)=\int_0^\infty\frac{x^ke^{-x}}{k!}e^{-x}\,\mathrm dx
  =\frac1{2^{k+1}},\qquad k=0,1,\ldots.
\]
计算使用 $\int_0^\infty x^ke^{-2x}\,\mathrm dx=k!/2^{k+1}$，可由分部积分对 $k$ 归纳证明。所得概率之和为 $1$，即 $Y$ 服从从 $0$ 开始计数、成功概率为 $1/2$ 的几何分布。随机化泊松参数后，边缘分布一般不再是泊松分布。
\end{solve}

\begin{example}
\label{ex:study-ch3-exercise-20}
教材第 3 章习题 20：设 $(X,Y)$ 在正方形 $[0,2]^2$ 内均匀分布，求 $Z=|X-Y|$ 的密度。
\end{example}

\begin{solve}
$Z$ 的取值在 $[0,2]$ 内。对 $0\le z\le2$，事件 $|X-Y|>z$ 对应正方形两角的两个直角三角形，直角边长为 $2-z$，总面积为 $(2-z)^2$。于是
\[
  F_Z(z)=1-\frac{(2-z)^2}{4},\quad0\le z\le2,
  \qquad
  f_Z(z)=\left(1-\frac z2\right)\boldsymbol 1_{(0,2)}(z).
\]
在 $z<0$ 时分布函数为 $0$，在 $z>2$ 时为 $1$。密度积分为 $1$，两个端点均无点概率。
\end{solve}

\begin{example}
\label{ex:study-ch3-exercise-24}
教材第 3 章习题 24：从一副 $52$ 张扑克牌中随机取 $13$ 张，以 $X,Y$ 分别表示黑桃和红桃的张数。求联合分布，以及给定 $X=1$ 时 $Y$ 的条件分布。
\end{example}

\begin{solve}
把牌分成 $13$ 张黑桃、$13$ 张红桃和 $26$ 张其他牌，每个 $13$ 张牌的子集等可能。因此
\[
  \mathbb P(X=x,Y=y)
  =\frac{\binom{13}x\binom{13}y\binom{26}{13-x-y}}{\binom{52}{13}},
  \qquad x,y\in\{0,1,\ldots,13\},\quad x+y\le13.
\]
其他取值的概率为 $0$。在恰有一张黑桃的条件下，其余 $12$ 张是从 $39$ 张非黑桃中等可能取出的，而其中有 $13$ 张红桃、$26$ 张其他牌，故
\[
  \mathbb P(Y=y\mid X=1)
  =\frac{\binom{13}y\binom{26}{12-y}}{\binom{39}{12}},
  \qquad y=0,1,\ldots,12.
\]
这是超几何分布。条件化后抽取张数与总体构成都已改变，不能继续使用原来的 $52$ 张总体。
\end{solve}

\begin{example}
\label{ex:study-ch3-exercise-28}
教材第 3 章习题 28：班车起点站的上客人数 $X\sim\operatorname{Pois}(\lambda)$，$\lambda>0$。给定上客人数，每位乘客独立地以概率 $p\in(0,1)$ 在中途下车。令 $Y$ 为中途下车人数，求条件分布与联合分布。
\end{example}

\begin{solve}
给定 $X=n$ 时，$Y$ 为 $n$ 个独立的下车指示变量之和，所以
\[
  \mathbb P(Y=m\mid X=n)=\binom nm p^m(1-p)^{n-m},
  \qquad0\le m\le n.
\]
乘以 $X$ 的点概率，得到
\[
  \mathbb P(X=n,Y=m)
  =e^{-\lambda}\frac{\lambda^n}{n!}\binom nm p^m(1-p)^{n-m},
  \qquad n=0,1,\ldots,\quad m=0,\ldots,n.
\]
进一步令 $W=X-Y$ 为未在中途下车的人数，则对非负整数 $m,k$，
\[
  \mathbb P(Y=m,W=k)
  =e^{-\lambda p}\frac{(\lambda p)^m}{m!}
   e^{-\lambda(1-p)}\frac{[\lambda(1-p)]^k}{k!}.
\]
因此 $Y\sim\operatorname{Pois}(\lambda p)$、$W\sim\operatorname{Pois}(\lambda(1-p))$，并且 $Y,W$ 独立。这一分解也验证了所得联合分布的归一性。它并不表示 $X,Y$ 独立，因为 $X=Y+W$ 且始终有 $Y\le X$。
\end{solve}

\begin{example}
\label{ex:study-ch3-exercise-33}
教材第 3 章习题 33：独立地重复掷一枚公平骰子，直到出现大于 $2$ 的点数。令 $X$ 为掷骰子的总次数，$Y$ 为最后一次的点数。求联合分布与边缘分布，并判断独立性。
\end{example}

\begin{solve}
$X=k,Y=j$ 表示前 $k-1$ 次均出现 $1$ 或 $2$，第 $k$ 次出现指定的 $j\in\{3,4,5,6\}$。因此
\[
  \mathbb P(X=k,Y=j)=\left(\frac13\right)^{k-1}\frac16,
  \qquad k=1,2,\ldots,\quad j=3,4,5,6.
\]
边缘分布为
\[
  \mathbb P(X=k)=\frac23\left(\frac13\right)^{k-1},\qquad
  \mathbb P(Y=j)=\sum_{k=1}^\infty\frac16\left(\frac13\right)^{k-1}
               =\frac14.
\]
每个联合点概率都等于相应的边缘点概率乘积，故 $X,Y$ 独立。等待多久与最后出现的是哪个合格点数，在这个模型中是独立的；停止规则本身并不必然造成依赖。
\end{solve}

\begin{example}
\label{ex:study-ch3-exercise-45}
教材第 3 章习题 45（带星号）：在长度为 $1$ 的线段上独立、均匀地选取两点，将线段切成三段，求三段能组成非退化三角形的概率。
\end{example}

\begin{solve}
以 $X,Y$ 表示两个切点坐标，联合密度在单位正方形内为 $1$。令 $U=\min\{X,Y\}$、$V=\max\{X,Y\}$；在 $0<u<v<1$ 上，两个原像 $(u,v)$ 与 $(v,u)$ 各贡献密度 $1$，故 $f_{U,V}(u,v)=2$。三段长为 $U,V-U,1-V$，总和为 $1$；能组成非退化三角形，当且仅当每段都小于 $1/2$，也就是
\[
  0<U<\frac12,\qquad \frac12<V<U+\frac12.
\]
所以所求概率为
\[
  \int_0^{1/2}\int_{1/2}^{u+1/2}2\,\mathrm dv\,\mathrm du
  =\int_0^{1/2}2u\,\mathrm du=\frac14.
\]
边界对应退化三角形或重合切点，概率为零。结论依赖“两切点独立、均匀”的模型，不能把三段长度直接当作独立均匀变量。
\end{solve}

\chapter{随机变量的数字特征和极限定理}
\label{chap:study-moments-limits}

% 依据缪柏其、张伟平《概率论与数理统计》（2022）第 4 章，教材第 123--179 页。

分布描述随机变量的全部取值规律，数字特征则提取其中便于比较和计算的部分信息：期望、中位数与众数刻画位置，方差与矩刻画波动及形状，协方差与相关系数刻画两个变量之间的线性联系，熵从另一角度度量不确定性。这些量一般不能代替完整分布，但在分布难以直接求出时，仍可能通过分解、条件化或对称性计算。极限定理进一步说明，在适当条件下，大量随机变量的平均会集中在期望附近，而经过适当放大的随机波动具有近似正态分布。

\section{数学期望和中位数}
\label{sec:study-expectation-median}

\subsection{数学期望的定义、存在条件与运算}
\label{subsec:study-expectation-properties}

例~\ref{ex:points-original} 中，若比赛继续，甲以概率 $3/4$ 得到 $300$ 元，以概率 $1/4$ 得到 $0$ 元，因而按当前信息计算的平均所得为 $300\times3/4+0\times1/4=225$ 元。这种按概率加权的平均就是\emph{数学期望（expectation）}，也称\emph{均值（mean）}。期望是由分布决定的常数，不必是随机变量实际能够取到的值；例如公平骰子的点数期望为 $3.5$，却不可能掷出 $3.5$ 点。

\begin{definition}
\label{def:study-expectation}
设离散型随机变量 $X$ 的不同取值为 $x_k$，相应概率为 $p_k$。若
\[
  \mathbb E|X|=\sum_k|x_k|p_k<\infty,
\]
则称 $X$ \emph{可积（integrable）}，定义其期望为
\begin{equation}
  \mathbb E X=\sum_k x_kp_k.
  \label{eq:study-discrete-expectation}
\end{equation}
若 $X$ 为具有密度 $f_X$ 的连续型随机变量，则可积条件及期望分别为
\begin{equation}
  \mathbb E|X|=\int_{-\infty}^{\infty}|x|f_X(x)\,dx<\infty,
  \qquad
  \mathbb E X=\int_{-\infty}^{\infty}xf_X(x)\,dx.
  \label{eq:study-continuous-expectation}
\end{equation}
本章所说“期望存在”均指上述有限期望存在。
\end{definition}

绝对收敛保证离散求和不依赖取值的排列次序；连续情形中，绝对可积排除了正、负部分都发散时的无意义抵消。对一般实值随机变量，写 $X^+=\max\{X,0\}$、$X^-=\max\{-X,0\}$，则 $X=X^+-X^-$、$|X|=X^++X^-$。可积等价于这两部分的期望都有限。非负随机变量的期望还可以按扩展实数理解为 $+\infty$，但不能把 $\infty-\infty$ 当作一个期望值。同分布随机变量的期望相同；这不要求它们相互独立。

\begin{example}
\label{ex:study-cauchy-no-mean}
设 $X$ 具有标准柯西密度 $f_X(x)=1/[\pi(1+x^2)]$。判断其期望是否存在。
\end{example}

\begin{solve}
由
\[
  \mathbb E|X|=\frac2\pi\int_0^\infty\frac{x}{1+x^2}\,dx=\infty,
\]
知 $X$ 不可积。尽管对任意 $R>0$，对称截断积分 $\int_{-R}^Rxf_X(x)\,dx=0$，这只给出了柯西主值；正半轴与负半轴上的积分分别发散，不能据此写 $\mathbb E X=0$。\emph[caution]{分布对称只有在期望存在时，才保证期望等于对称中心。}
\end{solve}

\begin{theorem}
\label{thm:study-expectation-properties}
设 $X_1,\ldots,X_n$ 可积，$a_1,\ldots,a_n,b$ 为实常数，则
\begin{equation}
  \mathbb E\left(b+\sum_{i=1}^na_iX_i\right)
    =b+\sum_{i=1}^na_i\mathbb E X_i.
  \label{eq:study-expectation-linearity}
\end{equation}
若 $X\le Y$ 以概率 $1$ 成立且二者可积，则 $\mathbb E X\le\mathbb E Y$。若 $X_1,\ldots,X_n$ 还相互独立，则它们的乘积可积，且
\begin{equation}
  \mathbb E\left(\prod_{i=1}^nX_i\right)
    =\prod_{i=1}^n\mathbb E X_i.
  \label{eq:study-independent-product-mean}
\end{equation}
\end{theorem}

\begin{proof}
线性组合的绝对值不超过 $|b|+\sum_i|a_i||X_i|$，故其期望存在。以两个离散型变量为例，记联合概率为 $p_{ij}=\mathbb P(X=x_i,Y=y_j)$，利用绝对收敛重排求和得
\[
  \mathbb E(aX+cY)=\sum_{i,j}(ax_i+cy_j)p_{ij}
    =a\sum_i x_i\sum_jp_{ij}+c\sum_jy_j\sum_ip_{ij}
    =a\mathbb E X+c\mathbb E Y.
\]
连续型情形用绝对可积时的积分换序作相同推导。一般情形可用有限值随机变量在平均绝对误差意义下逼近各个可积变量，再用 $|\mathbb E(U-V)|\le\mathbb E|U-V|$ 将有限值情形的线性等式取极限。单调性来自 $\mathbb E(Y-X)\ge0$。独立时联合分布可分解为边缘分布的乘积，先对非负函数 $\prod_i|X_i|$ 求和或积分，得其期望等于 $\prod_i\mathbb E|X_i|<\infty$，再分解乘积的期望，即得式~\eqref{eq:study-independent-product-mean}。
\end{proof}

\emph[caution]{有限和的期望可直接相加，不需要独立性；乘积的期望一般不能直接相乘。}例如 $X\sim\mathrm{Bern}(p)$，$0<p<1$，则 $\mathbb E(X^2)=p\ne p^2=(\mathbb E X)^2$。同样，一般不能把 $\mathbb E[g(X)]$ 写成 $g(\mathbb E X)$。

\begin{proposition}
\label{prop:study-function-expectation}
设 $g$ 为可测实值函数，且 $\mathbb E|g(X)|<\infty$，则
\begin{equation}
  \mathbb E[g(X)]=
  \begin{cases}
    \displaystyle\sum_k g(x_k)p_k,&X\text{ 为离散型},\\
    \displaystyle\int_{-\infty}^{\infty}g(x)f_X(x)\,dx,
      &X\text{ 具有密度 }f_X.
  \end{cases}
  \label{eq:study-function-expectation}
\end{equation}
对随机向量 $(X,Y)$ 的可积函数 $g(X,Y)$，可相应使用联合概率或联合密度作二重求和或积分。
\end{proposition}

离散情形可把产生相同 $g(x_k)$ 的取值归组，由绝对收敛得到公式；连续情形的公式是分布变换下积分的基本性质，对一般可测变换的严格证明需要更一般的积分理论。应用时无须先求 $g(X)$ 的分布。例如 $X\sim U(-1,1)$ 时，$\mathbb E(2|X|)=\int_{-1}^1|x|\,dx=1$。

\begin{example}
\label{ex:study-poisson-binomial-means}
求二项分布和泊松分布的期望，并说明如何利用随机变量的分解简化计算。
\end{example}

\begin{solve}
若 $X\sim\mathrm{Bin}(n,p)$，写成独立伯努利变量之和 $X=\sum_{i=1}^nI_i$，由 $\mathbb E I_i=p$ 得 $\mathbb E X=np$。若 $X\sim\mathrm{Pois}(\lambda)$，$\lambda>0$，则
\[
  \mathbb E X=\sum_{k=1}^{\infty}k e^{-\lambda}\frac{\lambda^k}{k!}
    =\lambda e^{-\lambda}\sum_{j=0}^{\infty}\frac{\lambda^j}{j!}
    =\lambda.
\]
若每次成功概率为 $p\in(0,1)$，各次试验独立，以 $G$ 表示首次成功所需的试验次数，则 $\mathbb P(G=k)=p(1-p)^{k-1}$，$k\ge1$，几何级数求导给出 $\mathbb E G=1/p$。第 $r$ 次成功所需的总次数是 $r$ 个独立同分布的这种等待时间之和，因此其负二项分布的期望为 $r/p$。这里计数包含成功的试验；若改为成功前的失败次数，期望相应减少 $r$。
\end{solve}

\subsection{条件期望与全期望公式}
\label{subsec:study-conditional-expectation}

\emph{条件期望（conditional expectation）}是在已知信息下，按条件分布计算的平均值。设 $Y$ 可积。当 $X$ 离散且 $\mathbb P(X=x)>0$ 时，
\[
  h(x):=\mathbb E(Y\mid X=x)=\sum_y y\,\mathbb P(Y=y\mid X=x),
\]
其中右侧按 $Y$ 的可能取值求和，且在相应条件绝对期望有限时定义。若 $(X,Y)$ 有联合密度 $f_{X,Y}$，则在 $f_X(x)>0$ 且条件绝对积分有限的点，可取
\begin{equation}
  h(x)=\mathbb E(Y\mid X=x)
    =\int_{-\infty}^{\infty}y\frac{f_{X,Y}(x,y)}{f_X(x)}\,dy.
  \label{eq:study-conditional-mean-density}
\end{equation}
连续情形中 $\mathbb P(X=x)=0$，因此这不是把单点事件的概率直接代入条件概率商式。$Y$ 可积保证上述条件平均对 $X$ 的分布几乎处处有定义；在剩余的零概率部分任意赋值，不改变后续结论。固定 $x$ 时 $h(x)$ 是数，而
\[
  \mathbb E(Y\mid X):=h(X)
\]
是 $X$ 的函数，通常仍为随机变量。它作为 $x$ 的函数也称 $Y$ 关于 $X$ 的\emph{回归函数（regression function）}。

\begin{theorem}
\label{thm:study-total-expectation}
若 $Y$ 可积，则 $\mathbb E(Y\mid X)$ 也可积，并且
\begin{equation}
  \mathbb E[\mathbb E(Y\mid X)]=\mathbb E Y.
  \label{eq:study-total-expectation}
\end{equation}
这称为\emph{全期望公式（law of total expectation）}或平滑公式。
\end{theorem}

\begin{proof}
条件平均满足 $|h(X)|\le\mathbb E(|Y|\mid X)$。离散情形按联合概率求和即得 $\mathbb E|h(X)|\le\mathbb E|Y|<\infty$，并有
\[
  \sum_xh(x)\mathbb P(X=x)
    =\sum_{x,y}y\,\mathbb P(X=x,Y=y)=\mathbb E Y.
\]
有联合密度时，由 $\iint|y|f_{X,Y}(x,y)\,dx\,dy=\mathbb E|Y|<\infty$ 可合法换序，故
\[
  \int h(x)f_X(x)\,dx
    =\iint y f_{X,Y}(x,y)\,dy\,dx
    =\int y f_Y(y)\,dy=\mathbb E Y.
\]
一般情形可由条件期望的积分定义得到相同结论。
\end{proof}

\begin{example}
\label{ex:study-normal-conditional-mean}
设 $(X,Y)$ 服从二元正态分布，边缘均值为 $\mu_X,\mu_Y$，标准差为 $\sigma_X,\sigma_Y>0$，相关参数为 $\rho\in(-1,1)$。求 $\mathbb E(Y\mid X=x)$。
\end{example}

\begin{solve}
由联合正态密度除以 $X$ 的边缘密度，并对关于 $y$ 的二次式配方，可得
\[
  Y\mid X=x\sim N\left(\mu_Y+\rho\frac{\sigma_Y}{\sigma_X}(x-\mu_X),
                       \sigma_Y^2(1-\rho^2)\right).
\]
因此 $\mathbb E(Y\mid X=x)=\mu_Y+\rho(\sigma_Y/\sigma_X)(x-\mu_X)$。全期望公式给出 $\mathbb E[\mathbb E(Y\mid X)]=\mu_Y$，因为 $\mathbb E(X-\mu_X)=0$。条件平均随 $x$ 线性变化是联合正态模型的性质，一般条件期望可以是非线性函数。
\end{solve}

\subsection{中位数、众数与分位数}
\label{subsec:study-location-quantiles}

均值用全部取值按概率加权，因此少量极端值也可能明显影响它。中位数按两侧的概率划分位置，众数则关注概率或密度最大的地方，三者回答的问题不同。

\begin{definition}
\label{def:study-median-quantile}
设 $X$ 的分布函数为 $F$。若实数 $m$ 满足
\begin{equation}
  \mathbb P(X\le m)\ge\frac12,\qquad
  \mathbb P(X\ge m)\ge\frac12,
  \label{eq:study-median}
\end{equation}
则称 $m$ 为 $X$ 的\emph{中位数（median）}。更一般地，对 $p\in(0,1)$，若 $q$ 满足 $F(q-)\le p\le F(q)$，即 $\mathbb P(X\le q)\ge p$ 且 $\mathbb P(X\ge q)\ge1-p$，则称 $q$ 为一个 $p$ \emph{分位数（quantile）}。为统一选取一个代表值，通常约定
\begin{equation}
  Q(p)=\inf\{x\in\mathbb R:F(x)\ge p\}.
  \label{eq:study-generalized-quantile}
\end{equation}
\end{definition}

由分布函数在两端的极限，式~\eqref{eq:study-generalized-quantile} 的下确界有限；由单调性和右连续性，$F(Q(p)-)\le p\le F(Q(p))$，故任意实值随机变量都有中位数和上述分位数，即使其均值不存在。满足分位数不等式的点可能不唯一，但选定下确界的约定后 $Q(p)$ 是确定的。若 $F$ 连续且严格递增，则 $F(Q(p))=p$ 且解唯一。只假设 $F$ 连续时，水平区间仍可能造成不唯一。

\begin{example}
\label{ex:study-bernoulli-medians}
若 $X\sim\mathrm{Bern}(1/2)$，求全部中位数。
\end{example}

\begin{solve}
当 $0<m<1$ 时，两侧概率均为 $1/2$；在 $m=0$ 或 $m=1$ 时，两侧概率仍均不小于 $1/2$。区间外则至少一侧概率为 $0$。所以全部中位数组成闭区间 $[0,1]$，其中按式~\eqref{eq:study-generalized-quantile} 选取的 $Q(1/2)=0$。中位数不一定是具有正概率的实际取值。
\end{solve}

\begin{definition}
\label{def:study-mode-iqr}
离散分布中使概率质量函数达到全局最大值的取值，称为\emph{众数（mode）}；对选定的连续密度，若它达到最大值，相应位置称为众数。众数可以不唯一，也可能不存在；连续分布通常针对其自然的连续密度讨论众数，以免单点修改密度造成歧义。称 $Q(1/4),Q(1/2),Q(3/4)$ 为四分位数，称
\[
  \operatorname{IQR}(X)=Q(3/4)-Q(1/4)
\]
为\emph{四分位距（interquartile range）}。
\end{definition}

\begin{example}
\label{ex:study-lognormal-location}
设 $\ln X\sim N(\mu,\sigma^2)$，$\sigma>0$，求对数正态分布的均值、中位数与众数。
\end{example}

\begin{solve}
作变换 $z=(\ln x-\mu)/\sigma$，得到密度
\[
  f_X(x)=\frac1{x\sigma\sqrt{2\pi}}
    \exp\left(-\frac{(\ln x-\mu)^2}{2\sigma^2}\right),\quad x>0.
\]
对任意实数 $k$，利用 $k\sigma z-z^2/2=-(z-k\sigma)^2/2+k^2\sigma^2/2$ 可得
\begin{equation}
  \mathbb E(X^k)
    =e^{k\mu}\int_{-\infty}^{\infty}e^{k\sigma z}
       \frac{e^{-z^2/2}}{\sqrt{2\pi}}\,dz
    =\exp\left(k\mu+\frac{k^2\sigma^2}{2}\right).
  \label{eq:study-lognormal-moments}
\end{equation}
于是均值为 $e^{\mu+\sigma^2/2}$；由 $\mathbb P(X\le e^\mu)=1/2$，中位数为 $e^\mu$。又
\[
  \frac{d}{dx}\ln f_X(x)
    =-\frac1x\left(1+\frac{\ln x-\mu}{\sigma^2}\right),
\]
其符号在 $x=e^{\mu-\sigma^2}$ 处由正变负，故唯一众数为 $e^{\mu-\sigma^2}$。三者依次为众数、中位数、均值，说明右侧长尾会使均值向较大取值偏移。
\end{solve}

若 $X\sim N(\mu,\sigma^2)$，用 $\Phi$ 表示标准正态分布函数，则 $Q(p)=\mu+\sigma\Phi^{-1}(p)$，所以 $\operatorname{IQR}(X)=2\Phi^{-1}(3/4)\sigma\approx1.349\sigma$。这体现了分位距也可描述散布程度，但它与标准差的关系依赖分布，不能把这个系数用于任意总体。

\section{方差和矩}
\label{sec:study-variance-moments}

\subsection{方差、标准差与概率界}
\label{subsec:study-variance-inequalities}

若 $\mathbb E X=\mu$，则 $\mathbb E(X-\mu)=0$，正负偏差会相互抵消，不能用它衡量波动。把偏差平方后再平均，就得到方差。它便于展开计算，但量纲是原变量量纲的平方；标准差恢复了原量纲。

\begin{definition}
\label{def:study-variance}
若 $X$ \emph{平方可积（square-integrable）}，即 $\mathbb E(X^2)<\infty$，定义其\emph{方差（variance）}和\emph{标准差（standard deviation）}为
\begin{equation}
  \operatorname{Var}(X)=\mathbb E[(X-\mathbb E X)^2],\qquad
  \sigma_X=\sqrt{\operatorname{Var}(X)}.
  \label{eq:study-variance-definition}
\end{equation}
\end{definition}

平方可积蕴含可积，因为 $|x|\le1+x^2$。展开定义并用期望的线性性，可得
\begin{equation}
  \operatorname{Var}(X)=\mathbb E(X^2)-(\mathbb E X)^2,
  \qquad
  \operatorname{Var}(aX+b)=a^2\operatorname{Var}(X).
  \label{eq:study-variance-calculation}
\end{equation}
因此平移不改变方差，尺度变化使方差乘以尺度的平方；特别是 $\operatorname{Var}(-X)=\operatorname{Var}(X)$。

\begin{proposition}
\label{prop:study-variance-minimum}
设 $X$ 平方可积，$\mu=\mathbb E X$。对任意常数 $c\in\mathbb R$，
\begin{equation}
  \mathbb E[(X-c)^2]=\operatorname{Var}(X)+(\mu-c)^2.
  \label{eq:study-mean-square-decomposition}
\end{equation}
故均值是使平均平方误差最小的唯一常数，最小值为方差；且 $\operatorname{Var}(X)=0$ 当且仅当 $\mathbb P(X=\mu)=1$。
\end{proposition}

\begin{proof}
将 $X-c=(X-\mu)+(\mu-c)$ 平方并取期望，交叉项因 $\mathbb E(X-\mu)=0$ 消失，得到分解式。若 $Z\ge0$ 且 $\mathbb E Z=0$，则对每个正整数 $k$，$0\ge k^{-1}\mathbb P(Z\ge k^{-1})$，故 $\mathbb P(Z>0)=0$。应用于 $Z=(X-\mu)^2$，即得方差为零的必要性；充分性直接来自定义。
\end{proof}

“以概率 $1$ 成立”也称\emph{几乎必然（almost surely）}成立，允许在一个零概率集合上例外。若 $0<\sigma_X<\infty$，则
\[
  Z=\frac{X-\mathbb E X}{\sigma_X}
\]
称为 $X$ 的\emph{标准化（standardization）}，满足 $\mathbb E Z=0$、$\operatorname{Var}(Z)=1$。标准化消除了位置与尺度，但一般不会把原分布变成正态分布。若原变量确为正态分布，则 $Z\sim N(0,1)$，并有 $\mathbb P(|Z|>3)\approx0.0027$；这一数值不是任意分布都有的性质。

\begin{example}
\label{ex:study-common-variances}
求伯努利、二项与泊松分布的方差，并整理常见分布的均值和方差。
\end{example}

\begin{solve}
若 $I\sim\mathrm{Bern}(p)$，则 $I^2=I$，故 $\operatorname{Var}(I)=p-p^2$。将二项变量分解为 $n$ 个独立伯努利变量之和，利用下文证明的独立和方差公式得 $\operatorname{Var}(X)=np(1-p)$。若 $X\sim\mathrm{Pois}(\lambda)$，则
\[
  \mathbb E[X(X-1)]
    =\sum_{k=2}^{\infty}k(k-1)e^{-\lambda}\frac{\lambda^k}{k!}
    =\lambda^2,
\]
因此 $\mathbb E(X^2)=\lambda^2+\lambda$，$\operatorname{Var}(X)=\lambda$。常见结果可合写为
\[
\begin{array}{c|c|c}
  \text{分布及参数条件}&\mathbb E X&\operatorname{Var}(X)\\ \hline
  \mathrm{Bern}(p),\ 0\le p\le1&p&p(1-p)\\
  \mathrm{Bin}(n,p),\ n\ge1&np&np(1-p)\\
  \mathrm{Pois}(\lambda),\ \lambda>0&\lambda&\lambda\\
  \mathrm{Ge}(p),\ 0<p<1&1/p&(1-p)/p^2\\
  U(a,b),\ a<b&(a+b)/2&(b-a)^2/12\\
  \mathrm{Exp}(\lambda),\ \lambda>0&1/\lambda&1/\lambda^2\\
  N(\mu,\sigma^2),\ \sigma>0&\mu&\sigma^2
\end{array}
\]
其中 $n$ 为正整数，二项分布参数 $p\in[0,1]$；几何分布的支持为 $1,2,\ldots$，指数分布采用速率参数，即密度为 $\lambda e^{-\lambda x}$，$x>0$。均匀分布可对 $X-(a+b)/2$ 直接积分；几何分布的二阶矩可对几何级数再次求导；指数分布的二阶矩由两次分部积分得 $2/\lambda^2$。对标准正态密度 $\varphi$，$\varphi'(z)=-z\varphi(z)$，分部积分得 $\int z^2\varphi(z)\,dz=1$，再作位置尺度变换即得正态分布的结果。
\end{solve}

\begin{theorem}
\label{thm:study-markov-chebyshev}
若 $Y\ge0$ 且 $\mathbb E Y<\infty$，则对任意 $a>0$，有\emph{马尔可夫不等式（Markov's inequality）}
\begin{equation}
  \mathbb P(Y\ge a)\le\frac{\mathbb E Y}{a}.
  \label{eq:study-markov-inequality}
\end{equation}
若 $X$ 的均值为 $\mu$、方差为 $\sigma^2<\infty$，则对任意 $\varepsilon>0$，有\emph{切比雪夫不等式（Chebyshev's inequality）}
\begin{equation}
  \mathbb P(|X-\mu|\ge\varepsilon)
    \le\frac{\sigma^2}{\varepsilon^2}.
  \label{eq:study-chebyshev-inequality}
\end{equation}
\end{theorem}

\begin{proof}
以 $\mathbf1_A$ 表示事件 $A$ 的示性变量，由逐点不等式 $Y\ge a\mathbf1_{\{Y\ge a\}}$ 和期望单调性得 $\mathbb E Y\ge a\mathbb P(Y\ge a)$。将非负变量 $Y$ 取为 $(X-\mu)^2$、阈值取为 $\varepsilon^2$，即得第二式。
\end{proof}

这些界只利用少量矩信息，不依赖具体分布，因而应用范围广，但可能比较保守；右端大于 $1$ 时还应与平凡上界 $1$ 取较小值。切比雪夫不等式是用已知方差控制偏离概率，不表示任意分布的尾部都像正态分布。

\subsection{矩、偏度、峰度与矩母函数}
\label{subsec:study-moments-mgf}

\begin{definition}
\label{def:study-moments}
设正整数 $k$ 满足 $\mathbb E|X|^k<\infty$。称 $\alpha_k=\mathbb E(X^k)$ 为 $k$ 阶\emph{原点矩（raw moment）}，称 $\mu_k=\mathbb E[(X-\mathbb E X)^k]$ 为 $k$ 阶\emph{中心矩（central moment）}，更一般地，$\mathbb E[(X-c)^k]$ 为关于常数 $c$ 的 $k$ 阶矩。若 $\sigma_X>0$，并且相应矩存在，则
\begin{equation}
  \gamma_1=\frac{\mu_3}{\sigma_X^3},\qquad
  \beta_2=\frac{\mu_4}{\sigma_X^4},\qquad
  \gamma_2=\beta_2-3
  \label{eq:study-skewness-kurtosis}
\end{equation}
分别称为\emph{偏度（skewness）}、\emph{峰度（kurtosis）}与\emph{超额峰度（excess kurtosis）}。
\end{definition}

一阶原点矩是均值，二阶中心矩是方差。偏度的正负通常反映较大或较小取值一侧的尾部对三次偏差的影响，不能机械地解释为密度最高点在均值哪一侧。关于均值对称且三阶绝对矩有限时，偏度为零；反向结论一般不成立。峰度更适合解释为标准化后极端偏差的敏感指标，不能只凭它判断密度曲线是否尖峭。正态分布满足 $\gamma_1=0$、$\beta_2=3$、$\gamma_2=0$，使用“峰度”一词时应先确认是否减去了 $3$。

计算矩时可用\emph{伽马函数（gamma function）}
\[
  \Gamma(t)=\int_0^\infty x^{t-1}e^{-x}\,dx,\quad t>0.
\]
分部积分给出 $\Gamma(t+1)=t\Gamma(t)$，并有 $\Gamma(1)=1$、$\Gamma(1/2)=\sqrt\pi$。对 $\alpha>-1$、$\beta>0$，换元 $u=\beta x^2$ 得
\begin{equation}
  \int_0^\infty x^\alpha e^{-\beta x^2}\,dx
    =\frac{\Gamma((\alpha+1)/2)}{2\beta^{(\alpha+1)/2}}.
  \label{eq:study-gamma-integral}
\end{equation}
例如标准正态变量 $Z$ 的奇数阶矩为零，偶数阶矩为 $\mathbb E(Z^{2r})=(2r)!/(2^rr!)$，$r\ge1$，由此特别得到四阶矩为 $3$。

\begin{definition}
\label{def:study-mgf}
称 $M_X(t)=\mathbb E(e^{tX})$ 为 $X$ 的\emph{矩母函数（moment generating function，MGF）}。若存在 $a>0$，使 $M_X(t)$ 对所有 $|t|<a$ 都有限，则称其在零点附近存在。
\end{definition}

零点附近存在是实质条件，因为 $M_X(0)=1$ 对任何随机变量都成立。在该条件下，取 $0<b<a$，有 $e^{b|X|}\le e^{bX}+e^{-bX}$，因此 $\mathbb E e^{b|X|}<\infty$。对 $|t|<b$，可用稍大的指数函数控制 $|X|^k e^{tX}$，从而合法地交换微分与期望，得到
\begin{equation}
  M_X^{(k)}(0)=\mathbb E(X^k),\qquad
  M_X(t)=\sum_{k=0}^{\infty}\frac{\mathbb E(X^k)}{k!}t^k,
  \quad |t|<a.
  \label{eq:study-mgf-moments}
\end{equation}
级数交换也可由 $\mathbb E e^{|t||X|}<\infty$ 直接保证。独立变量之和满足
\begin{equation}
  M_{X_1+\cdots+X_n}(t)=\prod_{i=1}^nM_{X_i}(t),
  \label{eq:study-mgf-independent-sum}
\end{equation}
其中 $t$ 取在各矩母函数共同有限的邻域内。常见例子为 $M_X(t)=(1-p+pe^t)^n$（二项分布）、$\exp[\lambda(e^t-1)]$（泊松分布）、$\lambda/(\lambda-t)$（速率为 $\lambda$ 的指数分布，$t<\lambda$），以及 $\exp(\mu t+\sigma^2t^2/2)$（正态分布）。这些式子也给出了计算矩和识别独立和分布的方法。

\begin{theorem}
\label{thm:study-mgf-uniqueness}
若 $X,Y$ 的矩母函数在零点的某个共同开邻域内有限且相等，则 $X,Y$ 同分布。
\end{theorem}

这里引用矩母函数的唯一性定理而不证明。例如 $M_X(t)=2/(2-t)$ 在零点附近成立，便可识别 $X\sim\mathrm{Exp}(2)$。\emph[caution]{所有阶矩都存在，不保证矩母函数在零点附近存在。}例~\ref{ex:study-lognormal-location} 的对数正态分布就具有所有正整数阶矩，但对任意 $t>0$，其矩母函数发散：在密度积分中，$tx$ 最终压过 $(\ln x-\mu)^2/(2\sigma^2)+\ln x$，使被积函数不能在无穷远处可积。因此不能不核对条件便把矩的形式级数当成矩母函数，也不能把一般分布的全部矩必然唯一确定分布作为本定理的结论。

\subsection{协方差、相关系数与独立性}
\label{subsec:study-covariance-correlation}

\begin{definition}
\label{def:study-covariance-correlation}
设 $X,Y$ 平方可积，定义\emph{协方差（covariance）}
\begin{equation}
  \operatorname{Cov}(X,Y)
    =\mathbb E[(X-\mathbb E X)(Y-\mathbb E Y)]
    =\mathbb E(XY)-\mathbb E X\,\mathbb E Y.
  \label{eq:study-covariance}
\end{equation}
若二者方差还均为正，定义\emph{相关系数（correlation coefficient）}
\begin{equation}
  \rho_{X,Y}=\frac{\operatorname{Cov}(X,Y)}{\sigma_X\sigma_Y}.
  \label{eq:study-correlation}
\end{equation}
若协方差为零，称二者\emph{不相关（uncorrelated）}。方差为零的变量仍可讨论协方差，但上述相关系数无定义。
\end{definition}

由 $2|XY|\le X^2+Y^2$ 可知协方差中的期望存在。协方差关于两个变量分别线性，且 $\operatorname{Cov}(X,Y)=\operatorname{Cov}(Y,X)$、$\operatorname{Cov}(X,X)=\operatorname{Var}(X)$，所以
\begin{equation}
  \operatorname{Var}\left(\sum_{i=1}^na_iX_i\right)
    =\sum_{i=1}^na_i^2\operatorname{Var}(X_i)
      +2\sum_{i<j}a_ia_j\operatorname{Cov}(X_i,X_j).
  \label{eq:study-variance-linear-combination}
\end{equation}
这是将中心化的和平方、再逐项取期望得到的。两两不相关时交叉项消失，独立是保证这一点的充分条件。因此独立变量相减时，$\operatorname{Var}(X-Y)=\operatorname{Var}(X)+\operatorname{Var}(Y)$，方差并不相减。对独立同分布、公共方差为 $\sigma^2$ 的 $X_1,\ldots,X_n$，样本均值 $\overline X_n=n^{-1}\sum_iX_i$ 的方差为 $\sigma^2/n$。

\begin{proposition}
\label{prop:study-correlation-bound}
若 $X,Y$ 平方可积，则
\begin{equation}
  |\operatorname{Cov}(X,Y)|^2
    \le\operatorname{Var}(X)\operatorname{Var}(Y).
  \label{eq:study-covariance-cauchy-schwarz}
\end{equation}
若二者方差均为正，则 $|\rho_{X,Y}|\le1$，且等号成立当且仅当存在常数 $a\ne0,b$，使 $Y=aX+b$ 几乎必然成立；此时相关系数的符号与 $a$ 相同。
\end{proposition}

\begin{proof}
记 $U=X-\mathbb E X$、$V=Y-\mathbb E Y$。若 $\operatorname{Var}(X)=0$，则 $U=0$ 几乎必然成立，结论直接成立。否则对任意实数 $t$，
\[
  0\le\mathbb E(V-tU)^2
    =\operatorname{Var}(Y)-2t\operatorname{Cov}(X,Y)
       +t^2\operatorname{Var}(X).
\]
在 $t=\operatorname{Cov}(X,Y)/\operatorname{Var}(X)$ 处取最小值，得到不等式。两方差均正时，等号成立等价于该最小值为零，由命题~\ref{prop:study-variance-minimum} 证明中的非负变量性质，等价于 $V=tU$ 几乎必然成立。此时 $t\ne0$，反向代入也立即得到 $|\rho|=1$。
\end{proof}

相关系数正是两个标准化变量乘积的期望。在把几乎必然相等的变量视为同一元素后，中心化的平方可积变量构成内积空间，内积为 $\langle U,V\rangle=\mathbb E(UV)$；相关系数就是其中的夹角余弦，零相关对应正交。必须先中心化：在全部平方可积变量上直接用协方差作内积，会使非零常数的“长度”为零，因而不满足内积的正定性。

\begin{example}
\label{ex:study-uncorrelated-dependent}
设 $X\sim U(-1,1)$，$Y=X^2$。证明 $X,Y$ 不相关但不独立。
\end{example}

\begin{solve}
由对称性 $\mathbb E X=\mathbb E(X^3)=0$，故 $\operatorname{Cov}(X,Y)=\mathbb E(X^3)-\mathbb E X\,\mathbb E(X^2)=0$。但令 $A=\{|X|\le1/2\}$、$B=\{Y\le1/4\}$，则 $A=B$，所以
\[
  \mathbb P(A\cap B)=\frac12\ne\frac14=\mathbb P(A)\mathbb P(B).
\]
因此不独立。零相关仅排除了这种二阶线性联系，并未排除非线性依赖。
\end{solve}

独立且平方可积蕴含不相关，反向一般不成立。若 $(X,Y)$ \emph{联合服从二元正态分布}，则不相关与独立等价：当相关参数 $\rho=0$ 时，联合正态密度分解为两个边缘正态密度的乘积。这里只知道两个边缘都正态还不够。例如 $X\sim N(0,1)$，$S$ 独立于 $X$ 且等概率取 $\pm1$，令 $Y=SX$，则 $Y\sim N(0,1)$ 且 $\operatorname{Cov}(X,Y)=0$，但 $|Y|=|X|$ 使二者不独立。

条件期望还能简化协方差与方差的计算。平方可积条件下，条件分布积分给出 $\mathbb E[Xg(Y)]=\mathbb E[g(Y)\mathbb E(X\mid Y)]$，只要两边可积；特别地，
\begin{equation}
  \operatorname{Cov}(X,Y)=\operatorname{Cov}(\mathbb E(X\mid Y),Y).
  \label{eq:study-covariance-conditional-mean}
\end{equation}
对 $\mathbb E(Y^2)<\infty$，定义 $\operatorname{Var}(Y\mid X)=\mathbb E[(Y-\mathbb E(Y\mid X))^2\mid X]$。将 $Y-\mathbb E Y$ 分解为 $Y-\mathbb E(Y\mid X)$ 与 $\mathbb E(Y\mid X)-\mathbb E Y$，交叉项的条件期望为零，得到\emph{全方差公式（law of total variance）}
\begin{equation}
  \operatorname{Var}(Y)=\mathbb E[\operatorname{Var}(Y\mid X)]
    +\operatorname{Var}(\mathbb E(Y\mid X)).
  \label{eq:study-total-variance}
\end{equation}
这把总波动分为条件内的平均波动和不同条件平均值之间的波动。式~\eqref{eq:study-covariance-conditional-mean} 应用于例~\ref{ex:study-normal-conditional-mean} 的线性条件均值，还可直接得到二元正态分布的协方差为 $\rho\sigma_X\sigma_Y$，从而验证其参数 $\rho$ 就是相关系数。

\section{熵的基本概念}
\label{sec:study-entropy}

\subsection{离散熵与微分熵}
\label{subsec:study-entropy-definitions}

均值和方差涉及取值的大小，熵则从概率分配的角度描述不确定性。若一项结果的概率为 $p>0$，可用 $-\log_2p$ 表示观察到它时的信息量：越罕见的结果，所带来的信息量越大。对这个量取平均得到离散熵。

\begin{definition}
\label{def:study-entropy}
设离散型随机变量 $X$ 的不同取值的概率为 $p_k$，定义其\emph{熵（entropy）}
\begin{equation}
  H(X)=-\sum_kp_k\log_2p_k,
  \label{eq:study-discrete-entropy}
\end{equation}
约定 $0\log_20=0$。它可以为 $+\infty$，以 $2$ 为底时单位为比特（bit）。若 $X$ 有密度 $f$，则在积分有定义时，定义\emph{微分熵（differential entropy）}
\begin{equation}
  h(X)=-\int_{-\infty}^{\infty}f(x)\ln f(x)\,dx,
  \label{eq:study-differential-entropy}
\end{equation}
约定 $0\ln0=0$，使用自然对数时单位为奈特（nat）。
\end{definition}

这里分别用 $H$ 和 $h$ 区分两类量。离散熵非负，且与各个取值的数值大小无关；只要重新标号不合并不同取值，熵就不变。微分熵则不是离散熵的直接替代品，它可以为负，并且依赖坐标尺度。由密度变换公式，对 $a\ne0$，在有关积分存在时，
\begin{equation}
  h(aX+b)=h(X)+\ln|a|.
  \label{eq:study-entropy-scale}
\end{equation}
因此二者都平移不变，但只有离散熵在非零线性缩放下也不变。例如 $X\sim U(0,c)$ 时，$h(X)=\ln c$，$0<c<1$ 时为负。比较微分熵时必须固定单位和坐标约定，也不能在任意分布间把较大方差直接等同于较大熵。

\begin{proposition}
\label{prop:study-discrete-max-entropy}
若 $X$ 至多取 $n$ 个值，则 $0\le H(X)\le\log_2n$。上界等号成立当且仅当这 $n$ 个值的概率全为 $1/n$；下界等号成立当且仅当分布集中于一个值。
\end{proposition}

\begin{proof}
因为 $0\le p_k\le1$，各项 $-p_k\log_2p_k$ 非负，且在 $0<p_k<1$ 时严格为正，得到下界及其等号条件。补齐零概率项，使共有 $n$ 项；由 $\ln u\le u-1$，
\[
  H(X)-\log_2n
    =\frac1{\ln2}\sum_{p_k>0}p_k\ln\frac1{np_k}
    \le\frac1{\ln2}\left(\frac{\#\{k:p_k>0\}}n-1\right)
    \le0.
\]
等号要求所有 $n$ 项均为正，并且每项 $1/(np_k)=1$，即均匀分布。
\end{proof}

\begin{example}
\label{ex:study-bernoulli-entropy}
若 $X\sim\mathrm{Bern}(p)$，则
\[
  H(X)=-p\log_2p-(1-p)\log_2(1-p).
\]
在 $0<p<1$ 上，其导数为 $\log_2[(1-p)/p]$，因此 $p=1/2$ 时达到最大值 $1$；$p=0,1$ 时熵为零。均匀的两种结果最难在观察前预测，而确定结果不提供这种随机信息量。
\end{example}

\subsection{约束下的最大熵分布}
\label{subsec:study-maximum-entropy}

已知的支持、均值与方差等条件限制了可选分布；最大熵问题是在这些约束内寻找熵最大的密度。它提供一种建模原则，不能据此认定真实数据必然服从所得分布。下面用积分不等式证明最优性，避免只求出变分方程的驻点便断言获得全局最大值。

\begin{lemma}
\label{lem:study-gibbs-inequality}
设 $f,g$ 为概率密度，且在 $f>0$ 处 $g>0$ 几乎处处成立。则\emph{相对熵（relative entropy）}满足
\begin{equation}
  D(f\Vert g):=\int f(x)\ln\frac{f(x)}{g(x)}\,dx\ge0,
  \label{eq:study-gibbs-inequality}
\end{equation}
其中允许 $D=+\infty$。等号成立当且仅当 $f=g$ 几乎处处成立。
\end{lemma}

\begin{proof}
在集合 $A=\{x:f(x)>0\}$ 上，$-\ln(g/f)\ge1-g/f$。乘以 $f$ 再积分，得到 $D(f\Vert g)\ge1-\int_Ag\ge0$。该积分的负部分有限，因为令 $r=f/g$，在 $r<1$ 处有 $-r\ln r\le1/e$，故负部分积分不超过 $\int g/e=1/e$。等号要求 $g/f=1$ 在 $A$ 上几乎处处成立，且 $g$ 在 $A$ 外没有质量，等价于两密度几乎处处相等。
\end{proof}

\begin{theorem}
\label{thm:study-maximum-entropy}
在具有密度的实值随机变量中，有以下最大微分熵结论：若支持限制在有限区间 $[a,b]$，$a<b$，则
\[
  h(X)\le\ln(b-a),
\]
等号当且仅当 $X\sim U(a,b)$。若 $X\ge0$ 且 $\mathbb E X=1/\lambda$，$\lambda>0$，则
\[
  h(X)\le1-\ln\lambda,
\]
等号当且仅当 $X\sim\mathrm{Exp}(\lambda)$。若 $\mathbb E X=\mu$、$\operatorname{Var}(X)=\sigma^2>0$，则
\begin{equation}
  h(X)\le\frac12\ln(2\pi e\sigma^2),
  \label{eq:study-normal-max-entropy}
\end{equation}
等号当且仅当 $X\sim N(\mu,\sigma^2)$。
\end{theorem}

\begin{proof}
分别以相应的均匀、指数和正态密度作为 $g$。在每一组约束下，$-\int f\ln g$ 都是同一个有限常数，依次为
\[
  \ln(b-a),\qquad -\ln\lambda+\lambda\mathbb E X=1-\ln\lambda,
  \qquad
  \frac12\ln(2\pi\sigma^2)+\frac{\mathbb E[(X-\mu)^2]}{2\sigma^2}
    =\frac12\ln(2\pi e\sigma^2).
\]
由引理~\ref{lem:study-gibbs-inequality}，$D(f\Vert g)=-h(X)-\int f\ln g\ge0$，从而得到三条上界和等号条件。若 $D=+\infty$，则对应 $h(X)=-\infty$，也不可能达到最大值。
\end{proof}

\section{大数定律和中心极限定理}
\label{sec:study-limit-theorems}

\subsection{依概率收敛与依分布收敛}
\label{subsec:study-probability-distribution-convergence}

随机变量序列的取值随样本点而变，因此必须说明“趋近”是哪一种意义。依概率收敛关注同一样本点上两个变量差得很远的概率；依分布收敛只比较分布函数，不要求变量取值逐个接近。

\begin{definition}
\label{def:study-convergence-probability}
设 $X_1,X_2,\ldots,X$ 定义在同一概率空间上。若对每个 $\varepsilon>0$，
\begin{equation}
  \mathbb P(|X_n-X|>\varepsilon)\longrightarrow0,
  \label{eq:study-convergence-probability}
\end{equation}
则称 $X_n$ \emph{依概率收敛（convergence in probability）}于 $X$，记作 $X_n\xrightarrow{\mathbb P}X$。把事件中的 $>$ 改成 $\ge$ 得到等价定义，可由阈值 $\varepsilon$ 与 $\varepsilon/2$ 比较验证。
\end{definition}

例如 $X_n\sim\mathrm{Exp}(n)$ 时，对 $\varepsilon>0$ 有 $\mathbb P(|X_n|>\varepsilon)=e^{-n\varepsilon}\to0$，故 $X_n\xrightarrow{\mathbb P}0$。这里不需要假设不同 $n$ 的变量独立。

\begin{definition}
\label{def:study-convergence-distribution}
设 $F_n,F$ 分别为 $X_n,X$ 的分布函数。若在 $F$ 的每一个连续点 $x$ 都有
\begin{equation}
  F_n(x)\longrightarrow F(x),
  \label{eq:study-convergence-distribution}
\end{equation}
则称 $X_n$ \emph{依分布收敛（convergence in distribution）}或\emph{弱收敛（weak convergence）}于 $X$，记作 $X_n\xrightarrow{\mathrm d}X$。此定义不要求这些变量定义在同一概率空间上。
\end{definition}

只要求极限分布函数的连续点是必要的。例如 $X_n\sim U(0,1/n)$，极限变量恒等于 $0$。对每个固定 $x<0$，$F_n(x)=0$；对每个固定 $x>0$，充分大的 $n$ 使 $F_n(x)=1$，所以依分布收敛成立。但 $F_n(0)=0$，而极限分布函数在 $0$ 处为 $1$，该不连续点不需要收敛。

\begin{proposition}
\label{prop:study-convergence-relations}
若 $X_n\xrightarrow{\mathbb P}X$，则 $X_n\xrightarrow{\mathrm d}X$。若 $c$ 为常数，则 $X_n\xrightarrow{\mathrm d}c$ 与 $X_n\xrightarrow{\mathbb P}c$ 等价。
\end{proposition}

\begin{proof}
若 $X_n\xrightarrow{\mathbb P}X$，由事件包含关系，对任意 $\delta>0$ 有
\[
  F(x-\delta)-\mathbb P(|X_n-X|>\delta)
    \le F_n(x)
    \le F(x+\delta)+\mathbb P(|X_n-X|>\delta).
\]
先令 $n\to\infty$，再在 $F$ 的连续点令 $\delta\downarrow0$，便得 $F_n(x)\to F(x)$。若 $X_n\xrightarrow{\mathrm d}c$，则对任意 $\varepsilon>0$，点 $c\pm\varepsilon$ 都是极限分布函数的连续点，且
\[
  \mathbb P(|X_n-c|>\varepsilon)
    \le F_n(c-\varepsilon)+1-F_n(c+\varepsilon)\longrightarrow0.
\]
反向由第一条结论得到。
\end{proof}

一般极限变量不是常数时，依分布收敛不能推出依概率收敛。例如令 $X,X_1,X_2,\ldots$ 为独立的 $\mathrm{Bern}(1/2)$ 变量，各个 $X_n$ 与 $X$ 同分布，当然依分布收敛于 $X$，但 $\mathbb P(|X_n-X|>1/2)=1/2$ 恒成立。

\subsection{大数定律：平均值的集中}
\label{subsec:study-law-large-numbers}

\begin{theorem}
\label{thm:study-weak-law}
设 $X_1,X_2,\ldots$ \emph{独立同分布（independent and identically distributed，i.i.d.）}，公共均值为 $\mu$、公共方差为 $\sigma^2<\infty$。令 $S_n=\sum_{i=1}^nX_i$、$\overline X_n=S_n/n$，则
\begin{equation}
  \overline X_n\xrightarrow{\mathbb P}\mu.
  \label{eq:study-weak-law}
\end{equation}
这称为有限方差情形的\emph{弱大数定律（weak law of large numbers）}。
\end{theorem}

\begin{proof}
由期望线性性和独立和的方差公式，$\mathbb E\overline X_n=\mu$、$\operatorname{Var}(\overline X_n)=\sigma^2/n$。对每个 $\varepsilon>0$，切比雪夫不等式给出
\begin{equation}
  \mathbb P(|\overline X_n-\mu|\ge\varepsilon)
    \le\frac{\sigma^2}{n\varepsilon^2}\longrightarrow0.
  \label{eq:study-weak-law-bound}
\end{equation}
因此依概率收敛成立；若 $\sigma^2=0$，各变量几乎必然等于 $\mu$，结论仍成立。
\end{proof}

上述证明实际只需各变量有共同均值、两两不相关且方差有统一有限上界。独立同分布并不是唯一的充分条件，但不能在应用时随意删去控制依赖和波动的假设。另有\emph{强大数定律（strong law of large numbers）}，其结论是 $\mathbb P(\lim_n\overline X_n=\mu)=1$；独立同分布且 $\mathbb E|X_1|<\infty$ 时它成立，这里仅引用而不证明。它描述几乎每一条样本路径上的收敛，比本节直接证明的依概率收敛更强；两种收敛均不意味着误差在每次增加样本时都会减小。

若 $I_i$ 是独立重复试验中事件 $A$ 的示性变量，$\mathbb P(A)=p$，则 $\overline I_n$ 是事件发生的频率，弱大数定律给出 $\overline I_n\xrightarrow{\mathbb P}p$，这就是\emph{伯努利大数定律（Bernoulli's law of large numbers）}。它为以频率估计概率、以样本均值估计总体均值提供依据，前提是实际抽样与模型条件相符。

\begin{example}
\label{ex:study-survey-sample-size}
设 $I_1,\ldots,I_n$ 独立且均服从 $\mathrm{Bern}(p)$，以 $\widehat p_n=n^{-1}\sum_iI_i$ 估计未知比例 $p$。要使 $\mathbb P(|\widehat p_n-p|<0.005)\ge0.95$，用切比雪夫不等式给出统一适用于 $p\in[0,1]$ 的充分样本量。
\end{example}

\begin{solve}
由 $p(1-p)\le1/4$，
\[
  \mathbb P(|\widehat p_n-p|\ge0.005)
    \le\frac{p(1-p)}{n(0.005)^2}
    \le\frac1{4n(0.005)^2}.
\]
要求右端不超过 $0.05$，得到 $n\ge200000$。这是模型下的严格充分界，通常较保守，不是实际所需最小样本量。随机抽样也未必自动意味着独立；若从有限总体不放回抽样，须按相应抽样模型另行处理。
\end{solve}

\subsection{中心极限定理与正态近似}
\label{subsec:study-central-limit-theorem}

\begin{theorem}
\label{thm:study-iid-clt}
设 $X_1,X_2,\ldots$ 独立同分布，$\mathbb E X_i=\mu$，$\operatorname{Var}(X_i)=\sigma^2\in(0,\infty)$。则
\begin{equation}
  Z_n=\frac{S_n-n\mu}{\sigma\sqrt n}
    =\frac{\sqrt n(\overline X_n-\mu)}\sigma
    \xrightarrow{\mathrm d}N(0,1),
  \label{eq:study-iid-clt}
\end{equation}
即对每个 $x\in\mathbb R$，$\mathbb P(Z_n\le x)\to\Phi(x)$。这称为\emph{林德伯格--莱维中心极限定理（Lindeberg--Lévy central limit theorem）}。
\end{theorem}

\begin{proof}[在矩母函数存在这一附加条件下的证明]
这里额外假设 $X_1$ 的矩母函数在零点附近存在；一般有限方差版本引用上述定理，本证明只覆盖这个较强条件。令 $W_i=(X_i-\mu)/\sigma$，则 $\mathbb E W_i=0$、$\mathbb E(W_i^2)=1$，其矩母函数在零点处满足
\[
  M_W(u)=1+\frac{u^2}{2}+o(u^2),\qquad u\to0.
\]
由独立性，对每个固定 $t$，当 $n$ 充分大时，
\[
  M_{Z_n}(t)=\left[M_W\left(\frac{t}{\sqrt n}\right)\right]^n
    =\left[1+\frac{t^2}{2n}+o(n^{-1})\right]^n
    \longrightarrow e^{t^2/2}.
\]
最后引用矩母函数的收敛定理：若矩母函数在零点某邻域内逐点趋于一个分布在该邻域内有限的矩母函数，则对应变量依分布收敛于该分布。此处极限 $e^{t^2/2}$ 正是标准正态矩母函数，故结论成立。这个步骤用的是收敛定理，不能仅由矩母函数的唯一性直接替代。
\end{proof}

大数定律说明 $\overline X_n$ 集中于常数 $\mu$，中心极限定理则描述放大后的误差 $\sqrt n(\overline X_n-\mu)$ 的极限分布。有限 $n$ 时常近似写作 $S_n\approx N(n\mu,n\sigma^2)$、$\overline X_n\approx N(\mu,\sigma^2/n)$，这不是一般总体下的精确分布等式。它也不是说未经放大的样本均值收敛到一个固定的非退化正态分布。

\begin{notealert}[type=caution,title={核对条件与近似精度}]
不存在仅凭 $n\ge30$ 就对所有有限方差分布保证良好正态近似的规则；偏斜程度、尾部及所求是否为极端概率都会影响精度。把结论推广到独立但不同分布的变量时，还需要控制单个大偏差对总和的影响，不能只由各方差夹在两个正数之间便断言中心极限定理成立。本章采用定理~\ref{thm:study-iid-clt} 的独立同分布条件。
\end{notealert}

若 $X_i\sim\mathrm{Bern}(p)$，$0<p<1$，则 $S_n\sim\mathrm{Bin}(n,p)$，代入定理得到\emph{棣莫弗--拉普拉斯中心极限定理（de Moivre--Laplace central limit theorem）}
\begin{equation}
  \frac{S_n-np}{\sqrt{np(1-p)}}\xrightarrow{\mathrm d}N(0,1).
  \label{eq:study-binomial-clt}
\end{equation}
在有限样本中，$np$ 与 $n(1-p)$ 都不宜太小；概率很靠近 $0$ 或 $1$ 时，仅有 $n$ 大仍可能不足。小 $p$ 且适当控制 $np$ 时，可另考虑泊松近似，而不是仅由 $np$ 小就断定泊松近似适用。

二项变量取整数，正态变量连续。对整数 $a\le b$，把每个整数看作宽为 $1$ 的小区间，有\emph{连续性修正（continuity correction）}
\begin{equation}
  \mathbb P(a\le S_n\le b)
    \approx\Phi\left(\frac{b+1/2-np}{\sqrt{np(1-p)}}\right)
      -\Phi\left(\frac{a-1/2-np}{\sqrt{np(1-p)}}\right).
  \label{eq:study-continuity-correction}
\end{equation}
特别地，$\mathbb P(S_n\ge a)\approx1-\Phi((a-1/2-np)/\sqrt{np(1-p)})$。连续性修正通常改善整数边界带来的误差，但不提供对所有参数的统一精度保证。

\begin{example}
\label{ex:study-binomial-normal-approximation}
设 $X\sim\mathrm{Bin}(20,0.3)$，比较精确值与正态近似所得的 $\mathbb P(1\le X\le3)$。
\end{example}

\begin{solve}
精确值为 $\sum_{k=1}^3\binom{20}k0.3^k0.7^{20-k}\approx0.10629$。均值为 $6$、方差为 $4.2$，用连续性修正得
\[
  \mathbb P(1\le X\le3)
    \approx\Phi\left(\frac{3.5-6}{\sqrt{4.2}}\right)
      -\Phi\left(\frac{0.5-6}{\sqrt{4.2}}\right)
    \approx0.10762.
\]
若直接用 $1,3$ 作正态积分边界，得到约 $0.0643$，遗漏的边界宽度在这个例子中影响较大。这是一个具体参数下的数值比较，不表明 $n=20$ 时对任意二项分布都准确。
\end{solve}

例~\ref{ex:study-survey-sample-size} 中，对固定 $p\in(0,1)$，中心极限定理给出近似覆盖概率
\[
  \mathbb P(|\widehat p_n-p|<0.005)
    \approx2\Phi\left(\frac{0.005\sqrt n}{\sqrt{p(1-p)}}\right)-1.
\]
用 $p(1-p)\le1/4$ 和 $\Phi^{-1}(0.975)\approx1.96$ 作保守的正态样本量规划，得到 $n\gtrsim1.96^2/[4(0.005)^2]=38416$。这个数来自渐近近似，不能与切比雪夫给出的严格有限样本保证混同；若需精确保证，应直接检查二项尾概率或使用适当的有限样本界。

\section{扩展阅读：数学期望的计算}
\label{sec:study-expectation-calculation}

\subsection{分解、示性变量与对称性}
\label{subsec:study-indicator-expectation-method}

计算期望不一定需要先求目标变量的分布。对计数问题，优先把总数写成示性变量之和，再用 $\mathbb E\mathbf1_A=\mathbb P(A)$。对等待问题，可以拆成若干阶段；对对称模型，可以先证明若干部分同分布，再利用总和恒定求出各自的均值。这里使用的有限和线性性均不要求这些部分相互独立。

\begin{example}
\label{ex:study-hypergeometric-mean}
在例~\ref{ex:study-capture-recapture} 的模型中，从 $N$ 条鱼中不放回均匀抽取 $n$ 条，其中 $M$ 条带标记。求抽到的带标记鱼数 $X$ 的期望。
\end{example}

\begin{solve}
把均匀抽样等价地看成先将全部鱼随机排列，再取前 $n$ 条。令 $I_i$ 表示第 $i$ 条抽到的鱼是否带标记，则各个位置的边缘概率均为 $M/N$，而 $X=\sum_{i=1}^nI_i$。所以
\[
  \mathbb E X=\sum_{i=1}^n\mathbb E I_i=\frac{nM}{N}.
\]
抽到带标记鱼后，余下样本的组成会改变，故这些示性变量一般不独立；这不影响上述期望计算。
\end{solve}

\begin{example}
\label{ex:study-matching-runs}
均匀随机排列 $n$ 顶各有主人的帽子后分给 $n$ 个人，求拿回自己帽子的人数期望。再设含 $m$ 个 $0$、$n$ 个 $1$ 的序列在所有不同排列中均匀选取，其中 $m,n\ge1$，求游程数 $R$ 的期望；游程是一个极大的连续同值片段。
\end{example}

\begin{solve}
对配对问题，令 $I_i$ 表示第 $i$ 人是否拿回自己的帽子。每人拿到每顶帽子的概率为 $1/n$，因此总人数 $X=\sum_iI_i$ 满足 $\mathbb E X=n\times1/n=1$。这个结论与总人数无关，并不意味着每一次恰有一个人拿对。

对游程问题，令 $L=m+n$，依次记序列元素为 $Z_1,\ldots,Z_L$。第一个位置开始一个游程，此后每发生一次相邻值变化就开始一个新游程，因此
\[
  R=1+\sum_{i=2}^L\mathbf1_{\{Z_{i-1}\ne Z_i\}}.
\]
任意两个相邻位置出现 $01$ 的概率为 $mn/[L(L-1)]$，出现 $10$ 的概率相同，故
\begin{equation}
  \mathbb E R=1+(L-1)\frac{2mn}{L(L-1)}
    =1+\frac{2mn}{m+n}.
  \label{eq:study-expected-runs}
\end{equation}
关键是选择“一个游程从哪里开始”作为计数单位，而不是先枚举整个游程数分布。
\end{solve}

\begin{example}
\label{ex:study-lottery-order-means}
从 $1,2,\ldots,33$ 中等可能地抽出 $6$ 个不同号码，按升序记为 $X_1<\cdots<X_6$，求各 $\mathbb E X_i$。
\end{example}

\begin{solve}
定义七个空隙数 $G_0=X_1-1$、$G_i=X_{i+1}-X_i-1$（$1\le i\le5$）、$G_6=33-X_6$。它们均为非负整数，且总和为 $27$。每个六号码子集与一个满足总和为 $27$ 的七元非负整数向量一一对应；子集等可能，所以这些向量也等可能。交换任意两个分量仍在同一集合中，故七个空隙同分布，$\mathbb E G_i=27/7$。于是
\[
  X_i=i+\sum_{j=0}^{i-1}G_j,
  \qquad \mathbb E X_i=i+i\frac{27}{7}=\frac{34i}{7}.
\]
这里需要的是交换空隙后等可能性的论证，不能仅凭“看起来对称”便假定各部分独立；它们的总和固定，实际并不独立。
\end{solve}

\subsection{条件化、递推与尾概率公式}
\label{subsec:study-conditioning-tail-method}

全期望公式的使用通常分为两步：选择一个分布容易掌握且能简化目标问题的条件变量 $Y$，先求 $\mathbb E(X\mid Y=y)$，再对 $Y$ 平均。若问题回到起点后重新开始，这种条件化常给出关于期望的递推方程；在解方程前，必须说明所求期望有限，否则形如 $a=b+ca$ 的式子在 $a=\infty$ 时不能直接移项。

\begin{example}
\label{ex:study-baghdad-thief}
窃贼每次独立等概率选择三扇门：第一扇经过 $3$ 小时到达地面；第二、第三扇分别经过 $5$、$7$ 小时回到起点，随后按相同规则重新选门。求到达地面所需总时间 $T$ 的期望。
\end{example}

\begin{solve}
以 $N$ 表示首次选到第一扇门所需的选择次数，则 $N\sim\mathrm{Ge}(1/3)$，$\mathbb E N=3$。因每次行走时间至多为 $7$，有 $T\le7N$，故 $\mathbb E T<\infty$。再按第一次选择 $J$ 条件化；返回起点后未来过程独立且与原过程同分布，所以
\[
  \mathbb E(T\mid J=1)=3,\quad
  \mathbb E(T\mid J=2)=5+\mathbb E T,\quad
  \mathbb E(T\mid J=3)=7+\mathbb E T.
\]
全期望公式给出 $\mathbb E T=[3+5+\mathbb E T+7+\mathbb E T]/3$，解得 $\mathbb E T=15$ 小时。独立且重新按原规则选择是建立该递推式的模型条件。
\end{solve}

\begin{proposition}
\label{prop:study-tail-expectation}
若 $X\ge0$，则
\begin{equation}
  \mathbb E X=\int_0^\infty\mathbb P(X>t)\,dt.
  \label{eq:study-tail-expectation}
\end{equation}
若 $X$ 还为整数值，则
\begin{equation}
  \mathbb E X=\sum_{n=0}^{\infty}\mathbb P(X>n)
    =\sum_{n=1}^{\infty}\mathbb P(X\ge n).
  \label{eq:study-integer-tail-expectation}
\end{equation}
这些等式允许两端同时为 $+\infty$，也可用于判断有限期望是否存在。
\end{proposition}

\begin{proof}
对每个非负实数 $x$，有 $x=\int_0^\infty\mathbf1_{\{x>t\}}\,dt$。先在有限区间积分，再利用非负函数的积分换序及单调收敛，得到式~\eqref{eq:study-tail-expectation}。整数值情形也可直接从 $X=\sum_{n=0}^{\infty}\mathbf1_{\{X>n\}}$ 出发，对有限部分和取期望后用单调收敛。这是教材习题 2 的尾概率表达，并不只适用于有密度的随机变量。
\end{proof}

\begin{example}
\label{ex:study-uniform-sum-stopping}
设 $U_1,U_2,\ldots$ 独立且均服从 $U(0,1)$，记
\[
  T=\min\left\{n\ge1:\sum_{i=1}^nU_i>1\right\},
\]
求 $\mathbb E T$。
\end{example}

\begin{solve}
对 $n\ge1$，$\{T>n\}$ 等价于 $\sum_{i=1}^nU_i\le1$，故其概率是单位立方体中单纯形 $\{u_i\ge0,\sum_i u_i\le1\}$ 的体积 $1/n!$。这一体积可归纳验证：若 $n-1$ 维、总和上限为 $s$ 的体积为 $s^{n-1}/(n-1)!$，再对第 $n$ 个坐标积分，得到 $\int_0^1(1-u)^{n-1}/(n-1)!\,du=1/n!$。又 $\mathbb P(T>0)=1$，因此
\[
  \mathbb E T=\sum_{n=0}^{\infty}\mathbb P(T>n)
    =\sum_{n=0}^{\infty}\frac1{n!}=e.
\]
上述尾概率趋于零，说明 $T$ 几乎必然有限；级数有限进一步证明它可积，无须先求 $T$ 的完整分布。

还可按第一次取值建立条件期望方程。对 $0\le t\le1$ 定义首次使部分和超过 $t$ 的次数为 $T_t$，并记 $h(t)=\mathbb E T_t$。由于 $T_t\le T$，各期望有限。给定 $U_1=u$，若 $u>t$ 则立即停止，否则仍需一个与 $T_{t-u}$ 同分布的等待过程，因而
\[
  h(t)=1+\int_0^t h(t-u)\,du=1+\int_0^t h(s)\,ds.
\]
有界性保证右侧先连续，再可微，于是 $h'(t)=h(t)$、$h(0)=1$，得到 $h(t)=e^t$，特别有 $h(1)=e$。这里先由尾概率法保证可积，再对递推方程求导，避免把未核实的无穷期望代入微分方程。
\end{solve}

\section*{本章总结}
\addcontentsline{toc}{section}{本章总结}

期望、中位数和众数分别由加权平均、两侧概率和最大概率或密度定义，不能互相替代。计算有限期望先核对绝对可积；计算函数的期望时可直接在原分布下积分；计算计数或等待时间时，分解为示性变量、利用对称性、按第一步条件化以及尾概率公式，往往比先求完整分布更有效。条件期望 $\mathbb E(Y\mid X=x)$ 是固定条件下的数，$\mathbb E(Y\mid X)$ 是随机变量；对后者取平均才还原总体期望。

方差是中心化平方的期望，均值使常数预测的平均平方误差最小；中位数则使平均绝对误差最小，见后面的习题选解。和的方差必须保留协方差项，只有两两不相关时才简化为方差之和。独立且平方可积蕴含不相关，反向结论一般不成立；联合正态分布是等价成立的重要特例。偏度、峰度与矩母函数均有各自的存在条件，不能由形式运算跳过可积性要求。

离散熵描述概率分配的不确定性，微分熵还受尺度影响；最大熵结论依赖明确的支持与矩约束。大数定律回答平均值是否集中，中心极限定理回答标准化波动的极限形状；前者给出依概率收敛，后者给出依分布收敛。实际使用中要分别检查独立同分布、有限矩与正方差等条件，并区分严格的概率界、渐近结论和有限样本的数值近似。

\section*{习题}
\addcontentsline{toc}{section}{习题}

以下选解教材习题 8、22、27、34、39、45、55、56，分别练习等待时间分解、条件期望与全方差、矩母函数、不相关与独立的区别，以及概率界和位置特征的最优化解释；习题 2 已在命题~\ref{prop:study-tail-expectation} 中证明。题号在正文中保留，例题编号仍由共享环境统一生成。

\begin{example}
\label{ex:study-exercise4-8}
教材习题 8。每袋零食独立等概率地含有 $n$ 种卡片之一，记集齐全部种类所需袋数为 $T_n$。求 $\mathbb E T_n$ 以及 $\lim_{n\to\infty}\mathbb E T_n/(n\ln n)$；并求 $n=12$ 时的均值。
\end{example}

\begin{solve}
已有 $k$ 种不同卡片后，下一袋出现新种类的概率为 $(n-k)/n$。令 $G_k$ 为从已有 $k$ 种到获得第 $k+1$ 种所需的袋数，则 $G_k$ 的条件等待分布总是参数为 $(n-k)/n$、支持从 $1$ 开始的几何分布，因而 $\mathbb E G_k=n/(n-k)$。由 $T_n=\sum_{k=0}^{n-1}G_k$ 和期望线性性，
\[
  \mathbb E T_n=n\sum_{j=1}^n\frac1j=nH_n,
\]
其中 $H_n$ 为调和和。积分比较给出 $\ln(n+1)\le H_n\le1+\ln n$，所以所求极限为 $1$。当 $n=12$ 时，均值为 $12H_{12}\approx37.24$ 袋。均值不是集齐的确定保证；接近集齐时获得最后少数新卡的概率小，是等待时间明显增加的原因。
\end{solve}

\begin{example}
\label{ex:study-exercise4-22}
教材习题 22。设 $X$ 等概率取 $0,1,2$，且 $Y\mid X=k\sim N(k,1)$。求 $Y$ 的密度与期望、$X+Y$ 的分布函数及 $\operatorname{Cov}(X,Y)$，并补充计算 $\operatorname{Var}(Y)$。
\end{example}

\begin{solve}
用 $\varphi,\Phi$ 分别表示标准正态密度与分布函数。按 $X$ 条件化可得
\[
  f_Y(y)=\frac13\sum_{k=0}^2\varphi(y-k),\qquad
  \mathbb E Y=\mathbb E[\mathbb E(Y\mid X)]=\mathbb E X=1.
\]
给定 $X=k$ 时，$X+Y\sim N(2k,1)$，因此其分布函数为 $F_{X+Y}(s)=\frac13\sum_{k=0}^2\Phi(s-2k)$。由 $\mathbb E(Y\mid X)=X$、$\mathbb E(X^2)=5/3$，
\[
  \operatorname{Cov}(X,Y)
    =\operatorname{Cov}(X,\mathbb E(Y\mid X))
    =\operatorname{Var}(X)=\frac23.
\]
全方差公式又给出 $\operatorname{Var}(Y)=1+2/3=5/3$。这是混合正态分布，一般不能把它仅因每个条件分布都正态便当作一个正态分布。
\end{solve}

\begin{example}
\label{ex:study-exercise4-27}
教材习题 27。分别求二项、泊松、指数和正态分布的矩母函数，并核对存在区间。
\end{example}

\begin{solve}
二项分布的有限和由二项式定理给出
\[
  M_X(t)=\sum_{k=0}^n\binom nk(pe^t)^k(1-p)^{n-k}
    =(1-p+pe^t)^n,\quad t\in\mathbb R.
\]
泊松分布用指数级数得 $M_X(t)=e^{-\lambda}\sum_{k=0}^{\infty}(\lambda e^t)^k/k!=\exp[\lambda(e^t-1)]$，对全部实数 $t$ 有限。速率为 $\lambda$ 的指数分布满足
\[
  M_X(t)=\int_0^\infty\lambda e^{-(\lambda-t)x}\,dx
    =\frac{\lambda}{\lambda-t},\quad t<\lambda,
\]
而 $t\ge\lambda$ 时发散。若 $X\sim N(\mu,\sigma^2)$，先写 $X=\mu+\sigma Z$（分布意义下），再用 $t\sigma z-z^2/2=-(z-t\sigma)^2/2+t^2\sigma^2/2$，得 $M_X(t)=\exp(\mu t+\sigma^2t^2/2)$，对全部实数 $t$ 有限。四个矩母函数均在零点附近存在，故可在零点求导得到矩。
\end{solve}

\begin{example}
\label{ex:study-exercise4-34}
教材习题 34。设 $(X,Y)$ 在菱形 $G=\{(x,y):|x|+|y|\le1\}$ 内均匀分布。求协方差并判断独立性。
\end{example}

\begin{solve}
$G$ 的面积为 $2$，故联合密度在 $G$ 内为 $1/2$。由关于两坐标轴的对称性，$\mathbb E X=\mathbb E Y=\mathbb E(XY)=0$，所以协方差为零。边缘密度为
\[
  f_X(x)=1-|x|,\quad |x|<1,
  \qquad f_Y(y)=1-|y|,\quad |y|<1.
\]
取事件 $A=\{X>3/4\}$、$B=\{Y>3/4\}$，则 $\mathbb P(A)=\mathbb P(B)=\int_{3/4}^1(1-x)\,dx=1/32$，但两事件不可能同时发生。因此 $\mathbb P(A\cap B)=0\ne\mathbb P(A)\mathbb P(B)$，二者不独立。
\end{solve}

\begin{example}
\label{ex:study-exercise4-39}
教材习题 39。设 $X_1,X_2,\ldots$ 独立且均服从支持为 $1,2,\ldots$ 的几何分布 $\mathrm{Ge}(p)$，$p\in(0,1)$；$N\sim\mathrm{Bin}(m,q)$ 独立于整个序列，其中 $m\ge1$、$q\in(0,1)$。令 $S=\sum_{i=1}^NX_i$，$N=0$ 时取 $S=0$，求条件均值、条件方差与 $\operatorname{Var}(S)$。
\end{example}

\begin{solve}
给定 $N=n$ 后，独立性保证被加变量仍具有原来的独立几何分布，因此
\[
  \mathbb E(S\mid N=n)=\frac np,\qquad
  \operatorname{Var}(S\mid N=n)=\frac{n(1-p)}{p^2}.
\]
全期望给出 $\mathbb E S=mq/p$，全方差公式给出
\[
  \operatorname{Var}(S)
    =\frac{(1-p)\mathbb E N}{p^2}
      +\frac{\operatorname{Var}(N)}{p^2}
    =\frac{mq(2-p-q)}{p^2}.
\]
随机项数自身的波动产生第二项，不能把 $N$ 直接换成其均值后就当作固定项数之和。
\end{solve}

\begin{example}
\label{ex:study-exercise4-45}
教材习题 45。独立重复试验 $500$ 次，每次事件发生的概率为 $0.2$，记总次数为 $K$。用切比雪夫不等式与正态近似估计 $\mathbb P(80\le K\le120)$；这里将题中的“在 80 到 120 之间”按包含两端解释。
\end{example}

\begin{solve}
$K\sim\mathrm{Bin}(500,0.2)$，均值为 $100$、方差为 $80$。因 $K$ 为整数，所求事件的补集是 $\{|K-100|\ge21\}$，故
\[
  \mathbb P(80\le K\le120)\ge1-\frac{80}{21^2}
    =\frac{361}{441}\approx0.81859.
\]
若不利用整数边界而只控制 $|K-100|<20$，也能得到较粗的下界 $0.8$。连续性修正给出
\[
  \mathbb P(80\le K\le120)
    \approx\Phi\left(\frac{120.5-100}{\sqrt{80}}\right)
      -\Phi\left(\frac{79.5-100}{\sqrt{80}}\right)
    =2\Phi\left(\frac{20.5}{\sqrt{80}}\right)-1
    \approx0.97809.
\]
第一种结果是下界，第二种是近似值，两者具有不同含义。
\end{solve}

\begin{example}
\label{ex:study-exercise4-55}
教材习题 55。若 $a\le X\le b$ 几乎必然成立，$a<b$，证明 $\operatorname{Var}(X)\le(b-a)^2/4$ 并给出等号条件。
\end{example}

\begin{solve}
取区间中点 $c=(a+b)/2$，由式~\eqref{eq:study-mean-square-decomposition}，
\[
  \operatorname{Var}(X)\le\mathbb E[(X-c)^2]\le\frac{(b-a)^2}{4}.
\]
第一处等号要求 $\mathbb E X=c$，第二处等号要求 $X$ 只在端点 $a,b$ 取值，概率为 $1$。二者合起来等价于 $\mathbb P(X=a)=\mathbb P(X=b)=1/2$。特别地，$0\le X\le1$ 时方差至多为 $1/4$，等概率取 $0,1$ 时达到上界。
\end{solve}

\begin{example}
\label{ex:study-exercise4-56}
教材习题 56。设 $X$ 可积，中位数为 $m$。证明对任意 $a\in\mathbb R$，$\mathbb E|X-m|\le\mathbb E|X-a|$；若方差为 $\sigma^2<\infty$，再证 $|m-\mathbb E X|\le\sigma$。
\end{example}

\begin{solve}
记 $F$ 为分布函数。对任意 $a,m$，逐点考察绝对值函数的分段斜率，并对有限区间换序积分，可得
\[
  \mathbb E|X-a|-\mathbb E|X-m|
    =\int_m^a(2F(t)-1)\,dt.
\]
因为积分区间有限、被积函数有界，换序合法；$\mathbb P(X=t)$ 的非零点至多可数，不影响关于 $t$ 的积分。若 $a>m$，对 $t>m$ 有 $F(t)\ge F(m)\ge1/2$，右端非负；若 $a<m$，对 $t<m$ 有 $F(t)\le F(m-)\le1/2$，反向积分也非负，得到第一条结论。

记 $\mu=\mathbb E X$。若 $m>\mu$，令 $d=m-\mu>0$、$W=X-\mu$。对任意 $s\ge0$，事件 $\{W\ge d\}$ 包含于 $\{(W+s)^2\ge(d+s)^2\}$，故由马尔可夫不等式，
\[
  \mathbb P(W\ge d)\le\frac{\sigma^2+s^2}{(d+s)^2}.
\]
取 $s=\sigma^2/d$，得到单侧界 $\mathbb P(W\ge d)\le\sigma^2/(\sigma^2+d^2)$。中位数条件使左端至少为 $1/2$，所以 $d^2\le\sigma^2$。若 $m<\mu$，对 $-W$ 作同样论证；$m=\mu$ 时结论直接成立。于是 $|m-\mu|\le\sigma$。这也说明只用双侧切比雪夫界会损失信息，利用单侧事件才能直接得到此处的系数 $1$。
\end{solve}

\chapter{统计学基本概念}
\label{chap:study-statistics-basics}

% 依据教材第 5 章，正文第 180--206 页（PDF 第 192--218 页）整理。
前四章从给定的概率模型推导随机现象的规律，本章转向统计学的出发点：模型中的某些因素未知，需要利用观测数据认识它们。统计推断的基本链条是\emph{明确总体及研究目标，规定获取样本的方式，构造统计量，再利用统计量的分布评价结论的误差}。本章首先区分总体、样本和统计量，继而研究常见统计量的抽样分布，为后续的参数估计和假设检验准备工具；民意调查与双盲对照试验则说明，数据的取得方式也是统计结论成立的条件。

\section{统计学发展简史}
\label{sec:study-statistics-history}

教材从两条相互汇合的线索介绍统计学的发展。一条源于人口、经济和社会状况的记录与整理，逐渐形成描述数据和开展抽样调查的方法；另一条以概率论为基础，围绕测量误差、参数估计和科学假设的判断发展数学理论。前者强调如何取得和概括数据，后者强调如何从有限且有随机波动的数据中推断未知规律。两者共同说明，统计学既依赖实际问题和数据，也依赖对不确定性的数学刻画。

沿教材的叙述，最小二乘法和正态误差理论为处理重复测量提供了工具，相关与回归把研究对象扩展到变量之间的关系，Student 的 $t$ 分布使小样本下的均值推断有了精确依据；此后，估计、假设检验、方差分析和试验设计逐渐形成系统方法。计算工具的发展又推动了复杂模型和大规模数据分析。理解这条路线时，重点是看到\emph{统计方法来自具体的推断需求，其有效性同时依赖模型、数据获取方式和数学条件}，而不只是记忆人名与公式。

\section{基本概念}
\label{sec:study-statistical-concepts}

描述性统计（descriptive statistics）利用均值、分位数、方差及图形概括已有数据；探索性数据分析（exploratory data analysis）帮助发现结构、异常和可能的关系；预测分析（predictive analysis）利用已知变量预测尚未观测的结果；统计推断（statistical inference）则由样本对总体中的未知因素作出判断。本笔记后续主要讨论参数估计、区间估计与假设检验。它们不仅要给出一个答案，还要说明答案如何受抽样波动影响。

例如，沿用教材的面料瑕疵问题，假定单位长度面料的瑕疵点数 $X$ 服从参数为 $\lambda>0$ 的泊松分布，则总体平均瑕疵点数为 $\mathbb E(X)=\lambda$。观测 $n$ 段面料后，可以用样本均值估计 $\lambda$，也可以判断平均瑕疵点数是否超过给定标准 $\lambda_0$。前者是估计问题，后者是检验问题。即使真实的 $\lambda\le\lambda_0$，某次样本均值也可能超过标准；反过来，真实不合格的产品也可能在一次抽检中表现良好。因此，样本中的差异和关于总体的结论之间，需要由抽样分布建立联系。

\subsection{总体与统计模型}

\begin{definition}[统计总体]
\label{def:study-population}
针对明确的研究对象及其数量指标，该指标所有可能取值及其概率分布构成\emph{统计总体（statistical population）}，简称总体。用随机变量 $X$ 表示这一指标，用 $F(x)=\mathbb P(X\le x)$ 表示总体分布函数。若总体分布属于一个由未知参数 $\theta$ 标记的分布族，则记为 $\{F_\theta:\theta\in\Theta\}$，其中 $\Theta$ 为参数空间。
\end{definition}

总体首先由研究问题确定：研究灯管寿命时指标是寿命，研究废品率时指标是产品是否为废品，不能仅列出一批物品而不说明关心什么性质。对含 $N$ 件产品、其中有 $M$ 件废品的有限总体，均匀随机抽取一件，令废品对应 $X=1$、正品对应 $X=0$，便有
\[
  \mathbb P(X=1)=p=\frac MN,\qquad
  \mathbb P(X=0)=1-p.
\]
这里 $N$ 是总体容量，$p$ 是总体废品率；在通常的频率学派模型中，$p$ 是固定而未知的量，随机性来自抽样，并不因为 $p$ 未知就把它视为随机变量。

\emph{统计模型（statistical model）}还要规定全部观测的联合分布。例如，重复称量某物件时，可写为 $X_i=\mu+\varepsilon_i$，其中 $\mu$ 是真实质量，$\varepsilon_i$ 是第 $i$ 次测量误差。若假定各次误差独立、条件稳定、没有系统偏差，并采用 $\varepsilon_i\sim N(0,\sigma^2)$ 的误差模型，则
\begin{equation}
  f_{\mu,\sigma^2}(x_1,\ldots,x_n)
  =(2\pi\sigma^2)^{-n/2}
    \exp\left\{-\frac1{2\sigma^2}\sum_{i=1}^{n}(x_i-\mu)^2\right\},
  \qquad \mu\in\mathbb R,\quad \sigma^2>0.
  \label{eq:study-normal-sample-density}
\end{equation}
大量微小误差叠加可以为正态近似提供动机，但不能仅凭“重复测量”就断言独立或精确正态。式~\eqref{eq:study-normal-sample-density} 是在这些假设下选定的模型；若仪器有漂移、共同环境扰动或系统偏差，联合分布便需要相应改变。

\subsection{样本及其二重性}

\begin{definition}[样本与简单随机样本]
\label{def:study-iid-sample}
按一定抽样或试验方案得到的 $n$ 个观测指标组成随机向量 $(X_1,\ldots,X_n)$，称为\emph{样本（sample）}，$n$ 称为\emph{样本量（sample size）}。若 $X_1,\ldots,X_n$ 相互独立，且都具有总体分布 $F$，则称其为\emph{简单随机样本（simple random sample）}，也称独立同分布样本，简记为 $X_1,\ldots,X_n\overset{\mathrm{i.i.d.}}{\sim}F$。
\end{definition}

抽样前，每个 $X_i$ 的取值尚未确定，样本是随机向量；抽样后，得到具体的\emph{样本值（sample realization）} $(x_1,\ldots,x_n)$。研究某种方法在重复抽样中的误差时使用大写随机变量，计算这一次观测的结果时使用小写数据。这称为样本的二重性，并不意味着某一已记录数值仍会自行变化。一次抽取 $n$ 件产品得到的是一个样本量为 $n$ 的样本，而不是 $n$ 个样本。

对简单随机样本，独立性给出联合分布函数
\begin{equation}
  \mathbb P(X_1\le x_1,\ldots,X_n\le x_n)
  =\prod_{i=1}^{n}F(x_i).
  \label{eq:study-iid-joint-cdf}
\end{equation}
若总体有密度 $f$，联合密度为 $\prod_{i=1}^{n}f(x_i)$；若总体为离散型，联合概率质量同样按各分量相乘。\emph[caution]{同分布只规定单次观测的规律，独立性才允许将联合分布写为乘积。}

\begin{example}
\label{ex:study-replacement-sampling}
教材例 5.3 的抽样模型。设一批产品共有 $N\ge2$ 件，其中 $M$ 件为废品，$p=M/N$。用 $X_i=1$ 表示第 $i$ 次抽到废品，$X_i=0$ 表示抽到正品，比较均匀有放回抽样与均匀无放回抽样得到的样本分布。
\end{example}

\begin{solve}
有放回且每次独立均匀抽取时，$X_i$ 独立同分布于 $0$--$1$ 分布。对任意 $x_i\in\{0,1\}$，记 $a=\sum_{i=1}^{n}x_i$，则
\begin{equation}
  \mathbb P(X_1=x_1,\ldots,X_n=x_n)=p^a(1-p)^{n-a}.
  \label{eq:study-replacement-sample-law}
\end{equation}
边界情形 $p=0,1$ 按相应的退化分布理解。此时 $K=\sum_{i=1}^{n}X_i\sim\operatorname{Bin}(n,p)$。

无放回时设 $1\le n\le N$，并用 $(b)_r=b(b-1)\cdots(b-r+1)$ 表示下降阶乘，约定 $(b)_0=1$。每一有序的不同产品序列等可能；指定一个含 $a$ 个 $1$ 的指标序列，可供相应位置选取的废品及正品序列分别有 $(M)_a$ 和 $(N-M)_{n-a}$ 种，故
\begin{equation}
  \mathbb P(X_1=x_1,\ldots,X_n=x_n)
  =\frac{(M)_a(N-M)_{n-a}}{(N)_n},
  \label{eq:study-without-replacement-law}
\end{equation}
其中 $0\le a\le M$、$0\le n-a\le N-M$，不可行序列的概率为 $0$。每个位置仍满足 $\mathbb P(X_i=1)=p$，但在 $0<M<N$ 时
\[
  \mathbb P(X_2=1\mid X_1=1)=\frac{M-1}{N-1}\ne\frac MN.
\]
因此指标同分布而不独立，不能使用式~\eqref{eq:study-replacement-sample-law}。共有 $\binom na$ 个含 $a$ 个 $1$ 的指标序列，求和得到超几何分布
\[
  \mathbb P(K=a)=\frac{\binom Ma\binom{N-M}{n-a}}{\binom Nn}.
\]
进一步，$i\ne j$ 时 $\mathbb E(X_iX_j)=M(M-1)/[N(N-1)]$，故 $\operatorname{Cov}(X_i,X_j)=-p(1-p)/(N-1)$，从而
\begin{equation}
  \operatorname{Var}(\overline X)
  =\frac{p(1-p)}n\frac{N-n}{N-1},
  \qquad \overline X=\frac1n\sum_{i=1}^{n}X_i.
  \label{eq:study-finite-population-variance}
\end{equation}
当抽样比例 $n/N$ 很小时，在样本均值的方差等计算中忽略有限总体修正通常是合理近似，但原始样本并未因此变成严格独立。
\end{solve}

有限总体抽样中，“简单随机抽样”也常用于每个固定容量子集等可能的无放回设计；本章沿用教材的概率论约定，以独立同分布定义简单随机样本。使用结论时应先看抽样设计及独立性，不能只看名称。随机抽样也不保证每一次样本都与总体完全相似；它提供的是可以分析抽样误差的概率机制。

\subsection{统计量与常用的数据概括}

\begin{definition}[统计量]
\label{def:study-statistic}
设样本为 $(X_1,\ldots,X_n)$。若一个可测函数 $T=t(X_1,\ldots,X_n)$ 的计算规则完全由样本和已知常数确定，不含未知参数，则称 $T$ 为\emph{统计量（statistic）}。对应的观测值为 $t(x_1,\ldots,x_n)$。
\end{definition}

例如，$\overline X=n^{-1}\sum_iX_i$ 是统计量；当 $\mu$ 未知时，$\overline X-\mu$ 不是统计量，而 $\overline X-\mu_0$ 在 $\mu_0$ 已知时是统计量。\emph[caution]{计算规则不含未知参数，并不意味着统计量的分布不依赖未知参数。}例如正态样本的 $\overline X$ 的分布仍含 $\mu$ 和 $\sigma^2$。研究含未知参数的标准化随机变量也很有用，只是要与统计量的定义区分。

对实值样本，常用的\emph{样本均值（sample mean）}、\emph{样本方差（sample variance）}与\emph{样本标准差（sample standard deviation）}定义为
\begin{equation}
  \overline X=\frac1n\sum_{i=1}^{n}X_i,\qquad
  S^2=\frac1{n-1}\sum_{i=1}^{n}(X_i-\overline X)^2,\qquad
  S=\sqrt{S^2},\quad n\ge2.
  \label{eq:study-sample-mean-variance}
\end{equation}
若样本独立同分布，$\mathbb E(X_i)=\mu$、$\operatorname{Var}(X_i)=\sigma^2<\infty$，则
\begin{equation}
  \mathbb E(\overline X)=\mu,\qquad
  \operatorname{Var}(\overline X)=\frac{\sigma^2}{n},\qquad
  \mathbb E(S^2)=\sigma^2.
  \label{eq:study-sample-mean-variance-moments}
\end{equation}
前两式来自期望的线性性与独立性；第三式可由恒等式
\begin{equation}
  \sum_{i=1}^{n}(X_i-\overline X)^2
  =\sum_{i=1}^{n}(X_i-\mu)^2-n(\overline X-\mu)^2
  =\sum_{i=1}^{n}X_i^2-n\overline X^{,2}
  \label{eq:study-residual-squares}
\end{equation}
取期望得到：右端前两项的期望分别为 $n\sigma^2$ 和 $\sigma^2$，所以除以 $n-1$ 后期望恰为 $\sigma^2$。这解释了样本方差分母的选择，并不要求总体正态；分母为 $n$ 的样本二阶中心矩则有期望 $(n-1)\sigma^2/n$。

\emph{样本矩（sample moment）}包括 $k$ 阶原点矩与中心矩
\[
  A_k=\frac1n\sum_{i=1}^{n}X_i^k\quad(k\ge1),\qquad
  B_k=\frac1n\sum_{i=1}^{n}(X_i-\overline X)^k\quad(k\ge2).
\]
特别地，$A_1=\overline X$，$B_2=(n-1)S^2/n$。对独立同分布样本，若 $\mathbb E|X|^k<\infty$，由大数定律及中心矩的多项式展开，$A_k$、$B_k$ 分别依概率收敛到 $\mathbb E(X^k)$、$\mathbb E[(X-\mu)^k]$。在 $B_2>0$ 时，样本偏度（sample skewness）可定义为 $g_1=B_3/B_2^{3/2}$，样本峰度（sample kurtosis）可用 $B_4/B_2^2$ 表示；若采用使正态总体峰度为 $0$ 的约定，则用超额峰度（excess kurtosis）$g_2=B_4/B_2^2-3$。两种峰度约定应明确区分，不能直接混用。

对于二维总体的样本 $(X_i,Y_i)$，\emph{样本相关系数（sample correlation coefficient）}为
\[
  r=\frac{\sum_{i=1}^{n}(X_i-\overline X)(Y_i-\overline Y)}
  {\sqrt{\sum_{i=1}^{n}(X_i-\overline X)^2
               \sum_{i=1}^{n}(Y_i-\overline Y)^2}},
  \qquad \overline Y=\frac1n\sum_{i=1}^{n}Y_i,
\]
其中两组数据的离差平方和都须为正。柯西--施瓦茨不等式给出 $|r|\le1$。这些量分别概括位置、离散程度、形状及线性关联，但不能由一个统计量恢复样本的全部特征。

将样本按大小排列为 $X_{(1)}\le\cdots\le X_{(n)}$，得到\emph{次序统计量（order statistics）}。样本中位数（sample median）通常取为
\[
  M_n=
  \begin{cases}
    X_{((n+1)/2)},&n\text{ 为奇数},\\
    \dfrac{X_{(n/2)}+X_{(n/2+1)}}2,&n\text{ 为偶数}.
  \end{cases}
\]
极小值、极大值及极差分别为 $X_{(1)}$、$X_{(n)}$、$R_n=X_{(n)}-X_{(1)}$。对称总体的对称中心可以作为中位数；只有在均值存在时，均值才也等于该中心，例如柯西总体就没有均值。

\begin{definition}[经验分布函数]
\label{def:study-empirical-cdf}
对样本 $(X_1,\ldots,X_n)$，定义\emph{经验分布函数（empirical distribution function）}
\begin{equation}
  F_n(x)=\frac1n\sum_{i=1}^{n}\mathbf 1_{\{X_i\le x\}},
  \qquad x\in\mathbb R,
  \label{eq:study-empirical-cdf}
\end{equation}
其中示性变量（indicator variable）$\mathbf 1_{\{X_i\le x\}}$ 在 $X_i\le x$ 时取 $1$，否则取 $0$。它表示样本中不大于 $x$ 的观测所占的比例。
\end{definition}

$F_n$ 单调不减、右连续，在左右无穷远处分别趋于 $0,1$，因而确为分布函数；一个不同观测值出现 $r$ 次，就在该点跳跃 $r/n$，不能在重复值处只记一次 $1/n$ 的跳跃。对固定 $x$，若样本独立同分布于 $F$，则各示性变量独立且服从参数为 $F(x)$ 的 $0$--$1$ 分布，故
\begin{equation}
  nF_n(x)\sim\operatorname{Bin}(n,F(x)),\qquad
  \mathbb E[F_n(x)]=F(x),\qquad
  \operatorname{Var}[F_n(x)]=\frac{F(x)[1-F(x)]}{n}.
  \label{eq:study-empirical-cdf-law}
\end{equation}
切比雪夫不等式进而给出 $\mathbb P(|F_n(x)-F(x)|>\varepsilon)\le1/(4n\varepsilon^2)$，所以对每个固定的 $x$，$F_n(x)$ 依概率收敛到 $F(x)$。这是逐点结论，不能不加论证就换成对全部 $x$ 同时成立的一致误差界。

样本分位数（sample quantile）存在多种有限样本约定。为使定义覆盖全部 $0<p<1$，本笔记用经验分布的广义逆定义
\[
  Q_n(p)=\inf\{x:F_n(x)\ge p\}=X_{(\lceil np\rceil)}.
\]
其中 $\lceil np\rceil$ 表示向上取整。教材采用基于 $\lfloor(n+1)p\rfloor$ 的次序位置约定，使用时须另行处理位置为 $0$ 的边界；采用插值的其他约定也很常见。应说明所用规则，特别是偶数样本量时 $Q_n(1/2)$ 与上述取两点平均的 $M_n$ 未必相同。以同一规则取 $p=1/4,3/4$ 的分位数，相减得到样本四分位距。

\section{抽样分布}
\label{sec:study-sampling-distributions}

\begin{definition}[抽样分布]
\label{def:study-sampling-distribution}
在给定总体及抽样方案下，统计量 $T=t(X_1,\ldots,X_n)$ 的概率分布称为它的\emph{抽样分布（sampling distribution）}。样本分布描述整个随机向量 $(X_1,\ldots,X_n)$ 的联合规律，抽样分布描述由该向量构成的某个统计量的规律。
\end{definition}

例如，$0$--$1$ 总体的独立样本具有式~\eqref{eq:study-replacement-sample-law} 的联合概率质量，而统计量 $K=\sum_iX_i$ 的抽样分布是二项分布。这两个分布的对象和样本空间不同。抽样分布回答的是“如果按相同规则反复抽样，所得统计量怎样变化”，它可能依赖未知总体参数，也可能是精确分布或仅有大样本近似。

\begin{example}
\label{ex:study-uniform-maximum}
教材例 5.6。设 $X_1,\ldots,X_n$ 独立同分布于均匀总体 $U(0,\theta)$，其中 $\theta>0$ 未知，求样本最大值 $T=X_{(n)}$ 的抽样分布。
\end{example}

\begin{solve}
对任意 $t$，事件 $\{T\le t\}$ 等价于所有 $X_i\le t$，故由独立性
\begin{equation}
  F_T(t)=F(t)^n=
  \begin{cases}
    0,&t\le0,\\
    (t/\theta)^n,&0<t<\theta,\\
    1,&t\ge\theta.
  \end{cases}
  \label{eq:study-uniform-maximum-cdf}
\end{equation}
在 $(0,\theta)$ 内求导，得到 $f_T(t)=nt^{n-1}/\theta^n$，区间外密度取 $0$。因此 $T$ 的分布依赖未知的 $\theta$，但 $T$ 本身只由样本计算，仍为统计量。积分还给出 $\mathbb E(T)=n\theta/(n+1)$，说明最大值通常低估上端点；这一性质将在点估计中再次用到。
\end{solve}

\subsection{统计三大分布及其构造}

正态总体的均值与方差推断需要处理标准正态变量的平方和，以及标准正态变量与独立平方和构成的比值。\emph{卡方分布（chi-square distribution）}、\emph{Student 的 $t$ 分布（Student's $t$ distribution）}和 $F$ 分布（$F$ distribution）正是从这些结构产生的，通常合称统计三大分布。

\begin{definition}[卡方分布]
\label{def:study-chi-square}
设 $Z_1,\ldots,Z_\nu$ 相互独立且都服从 $N(0,1)$，其中 $\nu\ge1$ 为整数。则
\[
  U=\sum_{i=1}^{\nu}Z_i^2
\]
服从自由度（degrees of freedom）为 $\nu$ 的卡方分布，记为 $U\sim\chi^2_\nu$。
\end{definition}

记伽马函数 $\Gamma(a)=\int_0^\infty u^{a-1}e^{-u}\,\mathrm du$，$a>0$。由平方变换可得 $Z_i^2$ 的密度为 $x^{-1/2}e^{-x/2}/\sqrt{2\pi}$，$x>0$；独立求和的卷积给出
\begin{equation}
  f_{\chi^2_\nu}(x)
  =\frac{x^{\nu/2-1}e^{-x/2}}{2^{\nu/2}\Gamma(\nu/2)},
  \quad x>0,
  \label{eq:study-chi-square-density}
\end{equation}
在 $x\le0$ 处可取密度为 $0$。若以形状参数 $a>0$、\emph{率参数} $\lambda>0$ 定义伽马分布（gamma distribution）
\[
  f_{\operatorname{Ga}(a,\lambda)}(x)
  =\frac{\lambda^a}{\Gamma(a)}x^{a-1}e^{-\lambda x},\qquad x>0,
\]
则 $\chi^2_\nu=\operatorname{Ga}(\nu/2,1/2)$；这里第二参数不是尺度参数。由 $\mathbb E(Z_i^2)=1$、$\mathbb E(Z_i^4)=3$ 得
\[
  \mathbb E(U)=\nu,\qquad \operatorname{Var}(U)=2\nu.
\]
若 $U\sim\chi^2_m$、$V\sim\chi^2_\nu$ 且相互独立，则 $U+V\sim\chi^2_{m+\nu}$，可将它们表示为相互独立的标准正态平方和来理解。密度在自由度 $1,2$ 时递减，在 $\nu>2$ 时于 $x=\nu-2$ 处达到最大。

\begin{definition}[$t$ 分布]
\label{def:study-student-t}
若 $Z\sim N(0,1)$、$U\sim\chi^2_\nu$，其中 $\nu\ge1$ 为整数，且 $Z,U$ 相互独立，则
\begin{equation}
  T=\frac{Z}{\sqrt{U/\nu}}\sim t_\nu.
  \label{eq:study-t-construction}
\end{equation}
称 $T$ 服从自由度为 $\nu$ 的 $t$ 分布。
\end{definition}

对联合密度作变换 $z=t\sqrt{u/\nu}$，其雅可比因子为 $\sqrt{u/\nu}$，积分消去 $u$ 后得到
\begin{equation}
  f_{t_\nu}(t)
  =\frac{\Gamma((\nu+1)/2)}{\sqrt{\nu\pi}\,\Gamma(\nu/2)}
     \left(1+\frac{t^2}{\nu}\right)^{-(\nu+1)/2},
  \qquad t\in\mathbb R.
  \label{eq:study-t-density}
\end{equation}
它关于 $0$ 对称，但尾部比标准正态分布更厚。自由度为 $1$ 时是标准柯西分布，没有期望；自由度 $\nu>1$ 时 $\mathbb E(T)=0$，$\nu>2$ 时 $\operatorname{Var}(T)=\nu/(\nu-2)$。这些条件可由密度尾部为 $|t|^{-(\nu+1)}$ 的阶及相应积分检验，再用伽马积分计算得到。\emph[caution]{对称性只有结合绝对可积性才能推出期望为零。}又因 $\mathbb E(U/\nu)=1$、$\operatorname{Var}(U/\nu)=2/\nu$，当 $\nu\to\infty$ 时 $U/\nu$ 依概率趋于 $1$，式~\eqref{eq:study-t-construction} 的分布因而趋于 $N(0,1)$；有限自由度时二者一般不相同。

\begin{definition}[$F$ 分布]
\label{def:study-f-distribution}
若 $U\sim\chi^2_m$、$V\sim\chi^2_\nu$，其中 $m,\nu\ge1$ 为整数，且 $U,V$ 相互独立，则
\begin{equation}
  W=\frac{U/m}{V/\nu}\sim F_{m,\nu}.
  \label{eq:study-f-construction}
\end{equation}
称 $W$ 服从第一自由度为 $m$、第二自由度为 $\nu$ 的 $F$ 分布。
\end{definition}

令 $u=mwv/\nu$，用雅可比因子 $mv/\nu$ 对联合密度积分，得到
\begin{equation}
  f_{F_{m,\nu}}(w)
  =\frac{\Gamma((m+\nu)/2)}{\Gamma(m/2)\Gamma(\nu/2)}
    \left(\frac m\nu\right)^{m/2}
    w^{m/2-1}\left(1+\frac m\nu w\right)^{-(m+\nu)/2},
  \quad w>0.
  \label{eq:study-f-density}
\end{equation}
在 $w\le0$ 处密度取 $0$。由构造直接得到
\begin{equation}
  W\sim F_{m,\nu}\ \Longrightarrow\ W^{-1}\sim F_{\nu,m},
  \qquad
  T\sim t_\nu\ \Longrightarrow\ T^2\sim F_{1,\nu}.
  \label{eq:study-f-relations}
\end{equation}
第一式须同时交换两个自由度；第二式使用了 $Z^2\sim\chi^2_1$。三大分布的定义均要求分子、分母所用的随机量独立，只有边缘分布正确还不足以识别比值的分布。

\subsection{上分位数及尾概率的转换}

对本章的连续分布，\emph{上 $\alpha$ 分位数（upper-tail quantile）} $q(\alpha)$ 由 $\mathbb P(W>q(\alpha))=\alpha$ 定义，其中 $0<\alpha<1$。本笔记分别使用 $z(\alpha)$、$\chi^2_\nu(\alpha)$、$t_\nu(\alpha)$、$F_{m,\nu}(\alpha)$ 表示标准正态、卡方、$t$ 和 $F$ 分布的上分位数。这里的 $\alpha$ 是右侧尾概率，若软件采用累计概率 $u$，应输入 $u=1-\alpha$。

由对称性及式~\eqref{eq:study-f-relations}，有
\begin{equation}
  t_\nu(1-\alpha)=-t_\nu(\alpha),\qquad
  F_{m,\nu}(1-\alpha)=\frac1{F_{\nu,m}(\alpha)}.
  \label{eq:study-upper-quantile-conversion}
\end{equation}
第二式的依据是 $\mathbb P(W>c)=\mathbb P(W^{-1}<1/c)$，再利用连续性转换尾概率。双侧对称区域则应各留 $\alpha/2$ 的尾概率，例如
\[
  \mathbb P\{-t_\nu(\alpha/2)\le T\le t_\nu(\alpha/2)\}=1-\alpha.
\]
卡方和 $F$ 分布不关于某个中心对称，不能照搬这组正负对称端点；其等尾区间的下、上端点分别使用上 $1-\alpha/2$ 与上 $\alpha/2$ 分位数。

\subsection{正态总体的样本均值与样本方差}

\begin{theorem}[正态样本的基本抽样分布]
\label{thm:study-normal-sampling}
设 $X_1,\ldots,X_n\overset{\mathrm{i.i.d.}}{\sim}N(\mu,\sigma^2)$，其中 $\sigma^2>0$、$n\ge2$，$\overline X,S^2$ 由式~\eqref{eq:study-sample-mean-variance} 定义。对不全为零的实常数 $c_1,\ldots,c_n$，
\begin{align}
  \sum_{i=1}^{n}c_iX_i
    &\sim N\left(\mu\sum_{i=1}^{n}c_i,
                  \sigma^2\sum_{i=1}^{n}c_i^2\right),
       \label{eq:study-normal-linear-combination}\\
  \overline X&\sim N(\mu,\sigma^2/n),
  \qquad Z=\frac{\sqrt n(\overline X-\mu)}{\sigma}\sim N(0,1),
       \label{eq:study-normal-mean-law}\\
  U=\frac{(n-1)S^2}{\sigma^2}&\sim\chi^2_{n-1}.
       \label{eq:study-normal-variance-law}
\end{align}
并且 $\overline X$ 与 $S^2$ 相互独立，从而
\begin{equation}
  \frac{\sqrt n(\overline X-\mu)}{S}\sim t_{n-1}.
  \label{eq:study-normal-t-law}
\end{equation}
\end{theorem}

\begin{proof}
记 $L=\sum_i c_iX_i$。由独立性，$L$ 的矩母函数为
\[
  \mathbb E(e^{tL})
  =\prod_{i=1}^{n}\exp\left(\mu c_it+\frac{\sigma^2c_i^2t^2}{2}\right)
  =\exp\left(\mu t\sum_i c_i+\frac{\sigma^2t^2}{2}\sum_i c_i^2\right),
  \qquad t\in\mathbb R.
\]
这正是式~\eqref{eq:study-normal-linear-combination} 中正态分布的矩母函数，由矩母函数的唯一性得到其分布；取 $c_i=1/n$ 即得式~\eqref{eq:study-normal-mean-law}。为同时证明方差的分布与独立性，令 $Z_i=(X_i-\mu)/\sigma$，$\boldsymbol Z=(Z_1,\ldots,Z_n)^{\mathsf T}$。将单位向量 $(1,\ldots,1)/\sqrt n$ 扩充为一组正交规范基，得到以它为第一行的正交矩阵 $A$，再令 $\boldsymbol Y=A\boldsymbol Z$。

正交变换保持长度，且雅可比绝对值为 $1$，所以变换后联合密度仍为
\[
  (2\pi)^{-n/2}\exp\left(-\frac12\sum_{i=1}^{n}y_i^2\right)
  =\prod_{i=1}^{n}\frac1{\sqrt{2\pi}}e^{-y_i^2/2}.
\]
因此 $Y_1,\ldots,Y_n$ 独立且均为标准正态变量，而
\[
  Y_1=\frac{\sqrt n(\overline X-\mu)}\sigma,\qquad
  \frac{(n-1)S^2}{\sigma^2}
    =\sum_{i=1}^{n}Z_i^2-Y_1^2
    =\sum_{i=2}^{n}Y_i^2.
\]
后一式是 $n-1$ 个独立标准正态变量的平方和，故服从 $\chi^2_{n-1}$。$\overline X$ 只依赖 $Y_1$，$S^2$ 只依赖 $Y_2,\ldots,Y_n$，故二者独立。将 $Y_1$ 除以 $\sqrt{U/(n-1)}=S/\sigma$，由 $t$ 分布定义得到式~\eqref{eq:study-normal-t-law}。
\end{proof}

这里的自由度 $n-1$ 来自中心化后 $\sum_i(X_i-\overline X)=0$ 的一个线性约束；正交变换把这个约束具体化为只剩 $n-1$ 个独立正态方向。\emph[caution]{样本方差的无偏性只需独立同分布和有限方差，但它的卡方分布、与样本均值独立及上述精确 $t$ 分布还依赖正态假设。}对非正态总体，不能只把 $\sigma$ 换为 $S$ 就断言得到精确的 $t$ 分布。

\subsection{两个独立正态样本}

\begin{corollary}
\label{cor:study-two-normal-samples}
设 $X_1,\ldots,X_m$ 来自 $N(\mu_1,\sigma_1^2)$，$Y_1,\ldots,Y_n$ 来自 $N(\mu_2,\sigma_2^2)$，两组均为简单随机样本且相互独立，$m,n\ge2$、$\sigma_1^2,\sigma_2^2>0$。记各自样本均值为 $\overline X,\overline Y$，样本方差为 $S_X^2,S_Y^2$，则
\begin{equation}
  \frac{S_X^2/\sigma_1^2}{S_Y^2/\sigma_2^2}
  \sim F_{m-1,n-1}.
  \label{eq:study-two-normal-f}
\end{equation}
若还假定 $\sigma_1^2=\sigma_2^2=\sigma^2$，定义合并样本方差
\begin{equation}
  S_p^2=\frac{(m-1)S_X^2+(n-1)S_Y^2}{m+n-2},
  \label{eq:study-ch5-pooled-variance}
\end{equation}
则
\begin{equation}
  \frac{(\overline X-\overline Y)-(\mu_1-\mu_2)}
       {S_p\sqrt{1/m+1/n}}
  \sim t_{m+n-2}.
  \label{eq:study-two-normal-t}
\end{equation}
\end{corollary}

\begin{proof}
由定理~\ref{thm:study-normal-sampling}，$(m-1)S_X^2/\sigma_1^2$ 与 $(n-1)S_Y^2/\sigma_2^2$ 分别服从自由度为 $m-1,n-1$ 的独立卡方分布，除以相应自由度后作比即得式~\eqref{eq:study-two-normal-f}。若方差相等，这两个卡方变量的和为 $(m+n-2)S_p^2/\sigma^2\sim\chi^2_{m+n-2}$。又由两组样本独立及组内均值与方差独立，$S_p^2$ 与 $\overline X-\overline Y$ 独立，而
\[
  \frac{(\overline X-\overline Y)-(\mu_1-\mu_2)}
       {\sigma\sqrt{1/m+1/n}}\sim N(0,1).
\]
由 $t$ 分布的定义得到式~\eqref{eq:study-two-normal-t}。
\end{proof}

合并方差是以各自自由度加权，不能一般地写成 $(S_X^2+S_Y^2)/2$。精确的合并方差 $t$ 结论要求两个总体方差相等；方差比的 $F$ 结论则无此要求。它们分别用于后续的均值差和方差比推断。

\subsection{其他常见总体的抽样分布}

若 $X_1,\ldots,X_n$ 为独立同分布样本，样本和 $K=\sum_iX_i=n\overline X$ 的分布常可利用总体分布的再生性求出。对 $0$--$1$ 总体，$K\sim\operatorname{Bin}(n,p)$，因而
\[
  \mathbb P\left(\overline X=\frac kn\right)
    =\binom nk p^k(1-p)^{n-k},\qquad k=0,\ldots,n.
\]
对泊松总体 $X_i\sim\operatorname{Pois}(\lambda)$，则 $K\sim\operatorname{Pois}(n\lambda)$，故
\begin{equation}
  \mathbb P\left(\overline X=\frac kn\right)
    =e^{-n\lambda}\frac{(n\lambda)^k}{k!},\qquad k=0,1,\ldots.
  \label{eq:study-poisson-mean-law}
\end{equation}
注意样本均值的取值为 $k/n$；样本和是泊松变量，不意味着样本均值仍服从某个泊松分布。

对率参数为 $\lambda>0$ 的指数总体 $f(x)=\lambda e^{-\lambda x}$，$x>0$，变换 $u=2\lambda x$ 给出 $2\lambda X_i$ 的密度为 $e^{-u/2}/2$，这正是 $\chi^2_2$ 的密度。因此由独立性和卡方分布的可加性，
\begin{equation}
  2\lambda\sum_{i=1}^{n}X_i=2n\lambda\overline X\sim\chi^2_{2n},
  \qquad
  \sum_{i=1}^{n}X_i\sim\operatorname{Ga}(n,\lambda),\qquad
  \overline X\sim\operatorname{Ga}(n,n\lambda).
  \label{eq:study-exponential-sampling-law}
\end{equation}
这些都是有限样本下的精确结果，不需要先使用中心极限定理。一般总体若仅知 $\mathbb E(X)=\mu$、$0<\operatorname{Var}(X)=\sigma^2<\infty$，则独立同分布条件下只有渐近结论 $\sqrt n(\overline X-\mu)/\sigma$ 依分布收敛到 $N(0,1)$；应将这种近似与上述精确分布分开。

\section{扩展阅读1：民意调查}
\label{sec:study-opinion-polls}

民意调查用部分受访者的数据认识目标人群对某一明确问题的态度。首先应规定目标人群、调查时间、问题措辞和抽样框，即实际可以被抽取的对象范围。若抽样框遗漏某些群体，或者被抽中者是否回应与其态度有关，即使收集大量问卷也可能产生系统偏差。教材的选举民调例子旨在说明：\emph{样本量控制随机波动，抽样与回应机制决定数据能否有效反映目标人群，两者不能互相替代。}

以支持某一观点的比例为例，令 $X_i=1$ 表示第 $i$ 位受访者支持，$X_i=0$ 表示不支持。暂以独立同分布的 $0$--$1$ 模型描述数据，设总体支持率为 $0<p<1$，样本支持率为
\[
  \widehat p=\overline X=\frac1n\sum_{i=1}^{n}X_i,
  \qquad \mathbb E(\widehat p)=p,\qquad
  \operatorname{Var}(\widehat p)=\frac{p(1-p)}n.
\]
对固定的 $p\in(0,1)$，由中心极限定理及 $\widehat p$ 依概率趋于 $p$，当 $n$ 增大时可得常用的近似置信区间
\begin{equation}
  \left[\widehat p-z(\alpha/2)
                   \sqrt{\frac{\widehat p(1-\widehat p)}n},\quad
        \widehat p+z(\alpha/2)
                   \sqrt{\frac{\widehat p(1-\widehat p)}n}\right],
  \label{eq:study-poll-wald-interval}
\end{equation}
名义置信水平为 $1-\alpha$。其含义是，在模型成立并反复抽样时，按此规则构造的随机区间包含固定参数 $p$ 的比例近似为 $1-\alpha$；观测后区间端点成为具体数值。这个近似通常要求支持和不支持的预期人数都足够大，$p$ 接近 $0,1$ 或样本很小时可能表现较差，不能仅因算出了端点就认为覆盖率得到保证。置信区间的严格定义与其他构造将在后续章节展开。

若希望在近似置信水平 $1-\alpha$ 下将绝对误差控制在 $d>0$ 内，以正态近似的半宽为准，需要计划
\begin{equation}
  n\ge\frac{z(\alpha/2)^2p(1-p)}{d^2}.
  \label{eq:study-poll-sample-size}
\end{equation}
抽样前 $p$ 未知，可用可靠的先验调查结果规划，或利用 $p(1-p)\le1/4$ 取保守样本量
\[
  n_0=\left\lceil\frac{z(\alpha/2)^2}{4d^2}\right\rceil.
\]
例如取 $1-\alpha=0.95$、$d=0.03$、$z(0.025)\approx1.96$，计算得到 $1067.11\ldots$，按整数样本量向上取整为 $1068$。这是正态近似下的设计计算，不能解释为“每次误差必定不超过 $3\%$”，也没有包含无回应、覆盖不足等误差。

当有限总体容量为 $N$ 且采用均匀无放回抽样时，应以式~\eqref{eq:study-finite-population-variance} 代替独立样本的方差。若记未取整的无限总体样本量为 $c=z(\alpha/2)^2/(4d^2)$，相同的近似半宽条件给出 $n\ge Nc/(N-1+c)$。所以“所需样本量与总人口数量无关”只是在总体足够大、抽样比例很小并忽略有限总体修正时的近似说法。分层、整群等设计也要使用对应的方差，不能直接套用独立样本公式。此外，调查反映的是规定时点与问题下的态度，不能将某次调查比例当作永久不变的总体属性。

\section{扩展阅读2：双盲对照试验}
\label{sec:study-double-blind-trials}

观察某组对象在采取措施后有所改善，还不能仅凭前后变化把原因归于该措施：自然变化、对象选择、共同环境，以及参与者或评价者的预期都可能产生影响。\emph{对照试验（controlled experiment）}通过设置可比较的处理组与对照组，使这些因素能够得到区分；\emph{随机分组（random assignment）}用随机机制决定对象接受哪种处理，以减少处理选择与原有特征之间的系统联系。

\emph{双盲（double blinding）}通常指试验参与者与直接实施或评价结果的相关人员不知道具体分组，以减少预期对行为、测量和评价的影响；实施时应明确到底对哪些人员设盲。对照措施应与试验目的相匹配，例如在适用时采用安慰剂，使两组在目标处理以外尽量可比。教材的胃冷冻案例所要强调的是：没有对照的早期观察难以排除非目标因素，随机、对照与设盲分别解决不同环节的问题，不能只凭“有改善”就确认处理效果。

随机抽样与随机分组还应区分：前者主要服务于由样本向目标总体的推广，后者主要服务于所比较处理之间的因果解释。随机分组使组间差异具有可分析的随机化机制，却不保证每一次分组都在所有特征上完全平衡，也不自动使参与者代表全部人群。试验仍需关注失访、处理执行和结果测量等条件。最后，\emph[caution]{某次试验未发现显著差异，不能直接推出两种处理完全相同。}是否能排除有实际意义的差异，还取决于效应估计的精度和研究设计，这将在区间估计与假设检验中进一步讨论。

\section*{本章总结}
\addcontentsline{toc}{section}{本章总结}

统计总体由研究指标及其分布刻画，统计模型还规定全部观测的联合规律；样本是按给定方案取得的随机向量，样本值是它的一次实现，统计量则是由样本和已知量计算的函数。独立同分布条件允许将样本的联合分布写成乘积，但有限总体无放回抽样通常不独立。使用均值、方差、矩、次序统计量和经验分布函数概括数据时，应始终分清总体量、随机统计量与具体计算值。

抽样分布把一次观测和重复抽样中的误差联系起来。正态样本的均值为正态变量，残差平方和经标准化后服从卡方分布，且均值与方差独立，从而产生精确的 $t$ 和 $F$ 分布；二项、泊松与指数总体也有各自的精确分布关系。识别这些分布时，要核对独立性、标准化常数、自由度及参数约定；处理尾概率时，要核对分位数的方向。民意调查与对照试验进一步提醒我们，数学公式对误差的刻画建立在数据获取机制和模型条件之上。

\section*{习题}
\addcontentsline{toc}{section}{习题}

以下选取教材本章的代表性习题，练习概念辨析、经验分布和抽样分布的推导；题目中注明原题号，节选或改写处另作说明。

\begin{example}
\label{ex:study-sample-duality-exercise}
教材习题 1、2。说明样本的二重性。某人独立射击 $5$ 次，结果为 $8,9,7,10,6$ 环，如何描述总体、样本与样本值？
\end{example}

\begin{solve}
研究指标是该人在规定条件下一次射击的环数，故总体用环数随机变量 $X$ 的分布表示。仅由这 $5$ 次结果不能确定整个总体分布；若再假定射击条件稳定、各次环数同分布，则 $(X_1,\ldots,X_5)$ 为来自该总体的简单随机样本，$(8,9,7,10,6)$ 为其一次实现。理论中的 $\overline X$ 随重复抽样变化，本次样本均值则是确定的 $\overline x=8$。独立性本身不足以推出同分布，故是否采用简单随机样本模型仍要看相同试验条件的假定。
\end{solve}

\begin{example}
\label{ex:study-bernoulli-statistic-exercise}
教材习题 8 节选。设 $X_1,\ldots,X_5$ 为参数 $p\in(0,1)$ 未知的 $0$--$1$ 总体的简单样本，写出样本空间及样本分布，并判断 $X_1+X_2$、$\min_{1\le i\le5}X_i$、$X_5+2p$ 是否为统计量。
\end{example}

\begin{solve}
样本空间为 $\{0,1\}^5$；对 $x_i\in\{0,1\}$，联合概率为 $p^{\sum_i x_i}(1-p)^{5-\sum_i x_i}$。前两个量只由样本计算，都是统计量；最后一个量需要未知的 $p$，因而不是统计量。若进一步求前两个量的抽样分布，则 $X_1+X_2\sim\operatorname{Bin}(2,p)$，而最小值只有在全部五个观测均为 $1$ 时才取 $1$，所以它服从参数为 $p^5$ 的 $0$--$1$ 分布。这也区分了样本的联合分布与特定统计量的抽样分布。
\end{solve}

\begin{example}
\label{ex:study-empirical-cdf-exercise}
教材习题 10。样本值为 $0.4,0.3,-0.3,-0.1,1.7,0.6,-0.1,0.9,2.6,0.5$，求经验分布函数。
\end{example}

\begin{solve}
由小到大排列为 $-0.3,-0.1,-0.1,0.3,0.4,0.5,0.6,0.9,1.7,2.6$，按不大于 $x$ 的个数计数，得到
\[
  F_{10}(x)=
  \begin{cases}
    0,&x<-0.3,\\
    1/10,&-0.3\le x<-0.1,\\
    3/10,&-0.1\le x<0.3,\\
    4/10,&0.3\le x<0.4,\\
    5/10,&0.4\le x<0.5,\\
    6/10,&0.5\le x<0.6,\\
    7/10,&0.6\le x<0.9,\\
    8/10,&0.9\le x<1.7,\\
    9/10,&1.7\le x<2.6,\\
    1,&x\ge2.6.
  \end{cases}
\]
在 $-0.1$ 处有两个重复观测，跳跃大小为 $2/10$；每个区间包含左端点，体现了 $F_{10}$ 的右连续性。
\end{solve}

\begin{example}
\label{ex:study-normal-variance-moments-exercise}
教材习题 11。设 $X_1,\ldots,X_n$ 为正态总体 $N(\mu,\sigma^2)$ 的简单随机样本，$n\ge2$、$\sigma>0$。求样本方差 $S^2$ 的方差和样本标准差 $S$ 的期望。
\end{example}

\begin{solve}
令 $\nu=n-1$，则 $U=\nu S^2/\sigma^2\sim\chi^2_\nu$，所以
\[
  \operatorname{Var}(S^2)
  =\frac{\sigma^4}{\nu^2}\operatorname{Var}(U)
  =\frac{2\sigma^4}{n-1}.
\]
为求 $\mathbb E(S)$，对式~\eqref{eq:study-chi-square-density} 积分：当 $r>-\nu/2$ 时，令 $v=u/2$ 可得
\[
  \mathbb E(U^r)
   =2^r\frac{\Gamma(\nu/2+r)}{\Gamma(\nu/2)}.
\]
取 $r=1/2$，利用 $S=\sigma\sqrt{U/\nu}$，得到
\[
  \mathbb E(S)
  =\sigma\sqrt{\frac2{n-1}}\,
       \frac{\Gamma(n/2)}{\Gamma((n-1)/2)}.
\]
因此 $S^2$ 对总体方差无偏，并不意味着开平方后的 $S$ 对总体标准差也无偏。事实上 $S^2$ 非退化，平方根严格凹，由詹森不等式 $\mathbb E(S)<\sqrt{\mathbb E(S^2)}=\sigma$。
\end{solve}

\begin{example}
\label{ex:study-reciprocal-t-square-exercise}
教材习题 12。设 $X\sim t_\nu$，求 $Y=1/X^2$ 的分布，其中 $\nu>1$ 为整数。
\end{example}

\begin{solve}
由式~\eqref{eq:study-f-relations}，$X^2\sim F_{1,\nu}$，故 $Y\sim F_{\nu,1}$。由于 $X$ 有连续密度，$\mathbb P(X=0)=0$，倒数几乎处处有定义；此分布关系实际上对 $\nu=1$ 也成立，不依赖 $t$ 分布的均值或方差是否存在。
\end{solve}

\begin{example}
\label{ex:study-sample-mean-laws-exercise}
教材习题 18。样本为简单随机样本，分别来自 $N(\mu,\sigma^2)$、参数为 $\lambda$ 的泊松总体，以及率参数为 $\lambda$ 的指数总体，求样本均值的分布。
\end{example}

\begin{solve}
正态总体有 $\overline X\sim N(\mu,\sigma^2/n)$；泊松总体有 $n\overline X\sim\operatorname{Pois}(n\lambda)$，其样本均值的概率质量由式~\eqref{eq:study-poisson-mean-law} 给出；指数总体有 $\overline X\sim\operatorname{Ga}(n,n\lambda)$，等价地 $2n\lambda\overline X\sim\chi^2_{2n}$。关键是先求和的分布，再作除以 $n$ 的尺度变换，不能只把总体分布中的参数机械地除以 $n$。
\end{solve}

\begin{example}
\label{ex:study-bernoulli-central-moment-exercise}
教材习题 19。设 $X_1,\ldots,X_n$ 为参数 $0<p<1$ 的 $0$--$1$ 总体的简单样本，求 $B_2=n^{-1}\sum_i(X_i-\overline X)^2$ 的期望。
\end{example}

\begin{solve}
由 $X_i^2=X_i$ 及式~\eqref{eq:study-residual-squares}，有 $B_2=\overline X-\overline X^{,2}$。又 $\mathbb E(\overline X)=p$、$\mathbb E(\overline X^{,2})=p^2+p(1-p)/n$，因此
\[
  \mathbb E(B_2)=\frac{n-1}{n}p(1-p).
\]
这与 $n\ge2$ 时 $B_2=(n-1)S^2/n$ 及 $S^2$ 的无偏性一致，也表明样本二阶中心矩与样本方差的区别。
\end{solve}

\begin{example}
\label{ex:study-normal-prediction-exercise}
教材习题 20。设 $X_1,\ldots,X_n$ 为 $N(\mu,\sigma^2)$ 的简单样本，$n\ge2$，其样本均值与标准差为 $\overline X,S$。若 $X_{n+1}$ 与前 $n$ 个观测独立且服从相同总体，求
\[
  T=\sqrt{\frac n{n+1}}\,\frac{X_{n+1}-\overline X}{S}
\]
的分布。
\end{example}

\begin{solve}
由独立正态变量的线性组合性质，$X_{n+1}-\overline X\sim N(0,\sigma^2(1+1/n))$。$X_{n+1}$ 与整个原样本独立，且原样本中 $\overline X$ 与 $S$ 独立，因此向量 $(X_{n+1},\overline X)$ 与 $S$ 独立，进而差值也与 $S$ 独立。于是
\[
  Z=\frac{X_{n+1}-\overline X}{\sigma\sqrt{1+1/n}}\sim N(0,1),
  \qquad U=\frac{(n-1)S^2}{\sigma^2}\sim\chi^2_{n-1},
\]
且 $Z,U$ 独立，所以 $T=Z/\sqrt{U/(n-1)}\sim t_{n-1}$。新观测与样本均值之差的方差包含两部分，不能沿用仅标准化 $\overline X-\mu$ 时的 $\sigma^2/n$。
\end{solve}

\begin{example}
\label{ex:study-uniform-range-exercise}
教材习题 22。设 $X_{(1)},\ldots,X_{(n)}$ 来自均匀总体 $U(0,1)$ 的简单样本。求使 $\mathbb P(X_{(n)}>0.99)\ge0.95$ 成立的最小整数 $n$，以及极差 $R_n=X_{(n)}-X_{(1)}$ 的期望。
\end{example}

\begin{solve}
由式~\eqref{eq:study-uniform-maximum-cdf}，$\mathbb P(X_{(n)}>0.99)=1-0.99^n$，故条件等价于 $0.99^n\le0.05$。因 $\ln0.99<0$，取对数并相除时不等号反向，得到
\[
  n\ge\frac{\ln0.05}{\ln0.99}=298.07\ldots,
\]
最小整数为 $299$。又对 $0<x<1$，最小值满足 $\mathbb P(X_{(1)}>x)=(1-x)^n$，所以其密度为 $n(1-x)^{n-1}$；最大值密度为 $nx^{n-1}$。积分得到 $\mathbb E(X_{(1)})=1/(n+1)$、$\mathbb E(X_{(n)})=n/(n+1)$，由期望的线性性
\[
  \mathbb E(R_n)=\frac{n-1}{n+1}.
\]
这里计算期望之差不要求最大值与最小值独立。
\end{solve}

\chapter{参数点估计}
\label{chap:study-point-estimation}

从已知总体分布推导样本的性质，是概率论的基本方向；从样本反推分布中的未知参数，则进入统计推断。本章以教材第六章为线索，讨论如何构造点估计、如何比较估计量，以及样本量增大后估计误差怎样变化。构造方法与评价准则是两个不同问题：矩估计或最大似然估计说明估计量怎样产生，无偏性、均方误差和相合性等说明它在重复抽样中有什么性质。
% 教材依据：缪柏其、张伟平《概率论与数理统计》（2022），第六章，正文第 207--229 页，习题第 230--239 页。

\section{参数点估计的概念}
\label{sec:study-point-estimation-concept}

设总体分布属于分布族 $\{F_\theta:\theta\in\Theta\}$，其中 $\theta=(\theta_1,\ldots,\theta_k)$ 是未知参数，$\Theta\subseteq\mathbb R^k$ 称为\emph{参数空间（parameter space）}。本章除特别说明外，$X_1,\ldots,X_n$ 均为来自该总体的简单随机样本，即独立且同分布；$\boldsymbol X=(X_1,\ldots,X_n)$ 表示随机样本，$\boldsymbol x=(x_1,\ldots,x_n)$ 表示它的一次观测值。记 $\mathbb E_\theta$、$\operatorname{Var}_\theta$ 和 $\mathbb P_\theta$ 为参数固定在 $\theta$ 时的期望、方差和概率。\emph[caution]{频率学派的点估计把参数视为固定的未知常数，随机性来自抽样。}

\begin{definition}[点估计量与估计值]
\label{def:study-point-estimator}
设 $g:\Theta\to\mathbb R$ 为已知函数。用统计量
\[
  T_n=T_n(\boldsymbol X)
\]
估计 $g(\theta)$，称 $T_n$ 为 $g(\theta)$ 的\emph{点估计量（point estimator）}；观测到 $\boldsymbol x$ 后所得的数 $T_n(\boldsymbol x)$ 称为相应的\emph{估计值（estimate）}。统计量的计算规则不能含有未知参数。取 $g(\theta)=\theta_j$ 就是在估计某个参数分量，向量参数也可逐分量估计。
\end{definition}

例如，总体具有有限均值 $\mu$ 时，样本均值 $\overline X=n^{-1}\sum_{i=1}^nX_i$ 是 $\mu$ 的一个估计量；测得三个寿命值 $5,6.4,5.9$ 后，估计值为 $\overline x=17.3/3\approx5.77$。前者具有抽样分布，后者是当前数据算出的数。点估计提供一个代表值，并不同时给出误差范围；误差的量化要结合估计量的分布，进一步引出下一章的区间估计。

\begin{example}
\label{ex:study-voltage-estimators}
教材例 6.1 的 $20$ 个击穿电压观测值之和为 $555.86$，最小值、最大值分别为 $24.46,30.88$，中间两个次序值为 $27.94,27.98$，单位均为 $\mathrm{kV}$。若用正态总体建模，怎样估计总体均值 $\mu$？
\end{example}

\begin{solve}
正态分布的均值、中位数和对称中心相同，因此可考虑样本均值、样本中位数（sample median）和中程数（midrange）。令 $X_{(1)}\le\cdots\le X_{(n)}$ 为由小到大排列的次序统计量，三种估计值分别为
\[
  \overline x=\frac{555.86}{20}=27.793,\qquad
  \widetilde x=\frac{27.94+27.98}{2}=27.960,\qquad
  \frac{x_{(1)}+x_{(20)}}2=27.670.
\]
第三种使用最大值与最小值的平均，不能与极差的一半 $(x_{(20)}-x_{(1)})/2$ 混淆。不同估计值不完全相同并不构成矛盾；比较估计方法，需要研究它们在重复抽样中的整体误差，而不能仅凭这一次数据判断谁离未知的 $\mu$ 更近。
\end{solve}

\section{矩估计法}
\label{sec:study-method-of-moments}

\subsection{以样本矩代替总体矩}

对整数 $r\ge1$，若 $\mathbb E_\theta|X|^r<\infty$，总体的 $r$ 阶原点矩和中心矩分别记为
\[
  \alpha_r(\theta)=\mathbb E_\theta(X^r),\qquad
  \mu_r(\theta)=\mathbb E_\theta[(X-\alpha_1(\theta))^r],
\]
相应的样本矩记为
\begin{equation}
  A_r=\frac1n\sum_{i=1}^nX_i^r,\qquad
  M_r=\frac1n\sum_{i=1}^n(X_i-\overline X)^r.
  \label{eq:study-sample-moments}
\end{equation}
原点矩可用大数定律得到 $A_r\xrightarrow{\mathbb P}\alpha_r(\theta)$；中心矩展开为样本原点矩的多项式后，也可得到 $M_r\xrightarrow{\mathbb P}\mu_r(\theta)$。这说明在相应矩存在时，样本矩能够近似总体矩。

\begin{definition}[矩估计]
\label{def:study-moment-estimator}
设选定的 $k$ 个总体矩为 $h_1(\theta),\ldots,h_k(\theta)$，对应样本矩为 $\widehat h_1,\ldots,\widehat h_k$。将矩之间的近似关系写成方程组
\begin{equation}
  h_j(\widehat\theta_{\mathrm M})=\widehat h_j,
  \qquad j=1,\ldots,k,
  \label{eq:study-moment-equations}
\end{equation}
若能在参数空间内选出一个由样本确定的解，就称它为参数的\emph{矩估计量（method-of-moments estimator）}。相应的参数函数估计为 $g(\widehat\theta_{\mathrm M})$。
\end{definition}

方程个数与参数个数相同并不保证解存在或唯一，求解后还须核对参数约束；矩存在也不保证所选矩能识别参数，例如对称分布的二阶矩可能无法区分某个位置参数的正负。矩估计无需写出完整密度，但必须知道选定总体矩与未知参数的关系。通常先使用低阶矩，因为高阶矩对尾部观测更敏感；这是一条方法选择的经验，不能脱离分布与误差准则断言低阶矩估计总是更优。

\subsection{常见总体与参数函数}

\begin{example}
\label{ex:study-normal-moments}
设 $X_i\sim N(\mu,\sigma^2)$，样本量 $n\ge2$，$\mu\in\mathbb R,\sigma^2>0$ 均未知，求二者的矩估计。
\end{example}

\begin{solve}
由 $\mathbb E(X)=\mu$、$\mathbb E(X^2)=\mu^2+\sigma^2$，矩方程给出
\begin{equation}
  \widehat\mu_{\mathrm M}=\overline X,\qquad
  \widehat{\sigma^2}_{\mathrm M}=A_2-\overline X^2
    =M_2=\frac1n\sum_{i=1}^n(X_i-\overline X)^2.
  \label{eq:study-normal-moment-estimates}
\end{equation}
当 $n\ge2$ 时，通常另记
\begin{equation}
  S^2=\frac1{n-1}\sum_{i=1}^n(X_i-\overline X)^2
     =\frac n{n-1}M_2,\qquad S=\sqrt{S^2}.
  \label{eq:study-unbiased-sample-variance}
\end{equation}
$S^2$ 是对 $M_2$ 作无偏修正后的估计，二者不能在推导中不加说明地互换。标准差 $\sigma$ 的直接矩估计为 $\sqrt{M_2}$，使用 $S$ 则是另一种选择。
\end{solve}

\begin{example}
\label{ex:study-rate-moments}
分别求指数分布、泊松分布和二项分布参数的低阶矩估计。
\end{example}

\begin{solve}
本章以速率参数定义指数分布：
\begin{equation}
  X\sim\operatorname{Exp}(\lambda)
  \quad\Longleftrightarrow\quad
  f(x;\lambda)=\lambda e^{-\lambda x}\mathbf1_{\{x\ge0\}},
  \qquad \lambda>0.
  \label{eq:study-exponential-rate}
\end{equation}
它的均值为 $1/\lambda$，所以 $\widehat\lambda_{\mathrm M}=1/\overline X$，而平均寿命 $1/\lambda$ 的矩估计为 $\overline X$。若 $X\sim\operatorname{Pois}(\lambda)$，由 $\mathbb E(X)=\lambda$ 得 $\widehat\lambda_{\mathrm M}=\overline X$。若 $X\sim\operatorname{Bin}(N,p)$，其中试验次数 $N\ge1$ 已知，则 $\mathbb E(X)=Np$，故 $\widehat p_{\mathrm M}=\overline X/N$；只有 $N=1$ 时才直接等于 $\overline X$。
\end{solve}

\begin{example}
\label{ex:study-uniform-moments}
设 $X_i\sim U(a,b)$，样本量 $n\ge2$，其中 $a<b$ 均未知，求矩估计。
\end{example}

\begin{solve}
由 $\mathbb E(X)=(a+b)/2$、$\operatorname{Var}(X)=(b-a)^2/12$，解方程
\[
  \frac{a+b}{2}=\overline X,\qquad
  \frac{(b-a)^2}{12}=M_2
\]
得到
\begin{equation}
  \widehat a_{\mathrm M}=\overline X-\sqrt{3M_2},\qquad
  \widehat b_{\mathrm M}=\overline X+\sqrt{3M_2}.
  \label{eq:study-uniform-moment-estimates}
\end{equation}
取正平方根是因为 $b-a>0$。当 $n\ge2$ 时，连续抽样下 $M_2>0$ 几乎必然成立。若按教材中的修正做法用 $S^2$ 替代 $M_2$，会得到 $\overline X\mp\sqrt3S$，这是相近但不同的估计。矩方程只匹配矩，并不保证估计出的区间包含全部观测值。
\end{solve}

参数函数可在得到参数估计后代入。例如，对于正态总体，设 $\Phi$ 为标准正态分布函数，则固定阈值 $c$ 对应的尾概率
\[
  q=\mathbb P(X>c)=1-\Phi\!\left(\frac{c-\mu}{\sigma}\right)
\]
的直接矩估计为 $1-\Phi((c-\overline X)/\sqrt{M_2})$，其中要求 $M_2>0$。若总体方差非零且四阶绝对矩有限，在 $M_2>0$ 时，偏度（skewness）与峰度（kurtosis）的矩估计分别为
\[
  \widehat\gamma_1=\frac{M_3}{M_2^{3/2}},\qquad
  \widehat\beta_2=\frac{M_4}{M_2^2};
\]
这里峰度采用未减去 $3$ 的约定，超额峰度的估计为 $\widehat\beta_2-3$。代入构造一般不保持无偏性。

\section{最大似然估计}
\label{sec:study-maximum-likelihood}

\subsection{固定样本，比较参数}

矩估计使用分布的若干数字特征，最大似然法则使用样本的联合分布。若总体概率密度或概率质量函数为 $f(x;\theta)$，独立样本在给定参数下的联合密度或联合概率质量为 $\prod_{i=1}^nf(x_i;\theta)$。

\begin{definition}[似然函数与最大似然估计]
\label{def:study-likelihood-mle}
固定观测值 $\boldsymbol x$，把联合密度或联合概率质量视为参数 $\theta$ 的函数，
\begin{equation}
  L(\theta;\boldsymbol x)=\prod_{i=1}^nf(x_i;\theta),
  \qquad\theta\in\Theta,
  \label{eq:study-likelihood}
\end{equation}
称为\emph{似然函数（likelihood function）}。若存在最大值点
\begin{equation}
  \widehat\theta_{\mathrm{ML}}(\boldsymbol x)
    \in\operatorname*{arg\,max}_{\theta\in\Theta}
       L(\theta;\boldsymbol x),
  \label{eq:study-mle-definition}
\end{equation}
则称其为\emph{最大似然估计（maximum likelihood estimate，MLE）}；把样本值换成随机样本便得到最大似然估计量。
\end{definition}

离散情形中，似然是不同参数下出现当前样本的概率；连续情形中，它是密度值，不能解释为取到某个精确样本点的正概率。\emph[caution]{似然函数不是参数的概率分布，也不要求对参数积分为 $1$。}它只在给定模型内比较参数对当前数据的支持程度，并不保证最大值对应真实参数。

当 $L>0$ 时，因对数严格递增，最大化 $L$ 等价于最大化\emph{对数似然（log-likelihood）}
\[
  \ell(\theta;\boldsymbol x)=\log L(\theta;\boldsymbol x)
      =\sum_{i=1}^n\log f(x_i;\theta).
\]
可导且位于参数空间内部的最大值点必须满足似然方程 $\partial\ell/\partial\theta_j=0$，但方程的根只是候选点，仍须检查整体最大性、边界和参数约束。若上确界只在参数空间外的边界趋近而未达到，则通常意义的 MLE 不存在；若存在多个最大值点，还须给出不含未知参数的选取规则。数值优化得到的驻点也不能自动当作全局最大值。

\begin{proposition}[最大似然估计的不变性]
\label{prop:study-mle-invariance}
若 $\widehat\theta_{\mathrm{ML}}$ 是 $\theta$ 的最大似然估计，则 $g(\widehat\theta_{\mathrm{ML}})$ 是参数函数 $g(\theta)$ 的最大似然估计。这里对 $\eta=g(\theta)$，以
\[
  L_g(\eta;\boldsymbol x)
    =\sup_{\theta:\,g(\theta)=\eta}L(\theta;\boldsymbol x),
  \qquad\eta\in g(\Theta),
\]
作为消去其余参数后的剖面似然（profile likelihood）。
\end{proposition}

\begin{proof}
对任意 $\eta$，都有 $L_g(\eta;\boldsymbol x)\le L(\widehat\theta_{\mathrm{ML}};\boldsymbol x)$；在 $\eta=g(\widehat\theta_{\mathrm{ML}})$ 处，取 $\theta=\widehat\theta_{\mathrm{ML}}$ 可达到右端。因此该 $\eta$ 为最大值点，不要求 $g$ 为一一映射。
\end{proof}

\subsection{光滑似然与参数边界}

\begin{example}
\label{ex:study-normal-mle}
设 $X_i\sim N(\mu,\sigma^2)$，$\mu,\sigma^2$ 均未知，求最大似然估计。
\end{example}

\begin{solve}
对给定样本，令 $\overline x=n^{-1}\sum x_i$、$Q=\sum(x_i-\overline x)^2$。对数似然为
\[
  \ell(\mu,\sigma^2)
   =-\frac n2\log(2\pi)-\frac n2\log\sigma^2
      -\frac1{2\sigma^2}\sum_{i=1}^n(x_i-\mu)^2.
\]
利用 $\sum(x_i-\mu)^2=Q+n(\overline x-\mu)^2$，固定 $\sigma^2$ 时唯一的最优位置为 $\mu=\overline x$。代入后令 $v=\sigma^2>0$，有
\[
  \frac{d}{dv}\ell(\overline x,v)=\frac{Q-nv}{2v^2}.
\]
若 $Q>0$，导数在 $v=Q/n$ 左侧为正、右侧为负，故得到全局最大值
\begin{equation}
  \widehat\mu_{\mathrm{ML}}=\overline X,\qquad
  \widehat{\sigma^2}_{\mathrm{ML}}=M_2.
  \label{eq:study-normal-mle}
\end{equation}
当 $n\ge2$ 时，$Q>0$ 在正态抽样下几乎必然成立。若所有观测完全相同，$Q=0$，则令 $\mu=\overline x,\sigma^2\downarrow0$ 会使似然无界，在 $\sigma^2>0$ 内没有最大值；$n=1$ 必属此情形。若均值 $\mu_0$ 已知，则方差的 MLE 改为 $n^{-1}\sum(X_i-\mu_0)^2$，不能再减去 $\overline X$。
\end{solve}

对于指数速率，$\ell(\lambda)=n\log\lambda-\lambda\sum x_i$，在 $\sum x_i>0$ 时唯一最大值点为 $\widehat\lambda_{\mathrm{ML}}=1/\overline X$。对于泊松参数，$\ell(\lambda)=-n\lambda+(\sum x_i)\log\lambda-\sum\log(x_i!)$，在 $\lambda\ge0$ 的参数空间内得到 $\widehat\lambda_{\mathrm{ML}}=\overline X$，全零样本对应 $\lambda=0$。若模型要求 $\lambda>0$，全零样本只有在 $\lambda\downarrow0$ 时达到上确界，MLE 不存在。二项总体 $\operatorname{Bin}(N,p)$ 在 $p\in[0,1]$ 内的 MLE 为 $\overline X/N$，全失败和全成功样本分别在 $0,1$ 达到最大；若把参数空间限定为 $(0,1)$，同样要单独处理这两种情形。

\begin{example}
\label{ex:study-uniform-mle}
设 $X_i\sim U(0,\theta)$，$\theta>0$，求最大似然估计。
\end{example}

\begin{solve}
采用端点包含在支撑中的密度版本 $f(x;\theta)=\theta^{-1}\mathbf1_{\{0\le x\le\theta\}}$。对正的观测值，有
\begin{equation}
  L(\theta;\boldsymbol x)=\theta^{-n}
      \mathbf1_{\{\theta\ge x_{(n)}\}},
  \qquad
  \widehat\theta_{\mathrm{ML}}=X_{(n)}.
  \label{eq:study-uniform-mle}
\end{equation}
当 $\theta<x_{(n)}$ 时，参数不能解释当前样本；在 $\theta\ge x_{(n)}$ 的可行域内，$\theta^{-n}$ 严格递减，因此最大值位于边界。若直接删除示性函数再求导，会丢失决定答案的约束。连续分布的端点取值不改变其概率分布，但会影响逐点似然是否在边界取到最大；这里明确端点约定，以得到教材采用的 MLE。
\end{solve}

\begin{example}
\label{ex:study-uniform-location-mle}
设 $X_i\sim U(\theta-\tfrac12,\theta+\tfrac12)$，求位置参数的最大似然估计。
\end{example}

\begin{solve}
仍取包含端点的密度版本，则
\[
  L(\theta;\boldsymbol x)
    =\mathbf1_{\{x_{(n)}-1/2\le\theta\le x_{(1)}+1/2\}}.
\]
因此区间 $[x_{(n)}-1/2,x_{(1)}+1/2]$ 中每一点都是最大值点；若另有 $\theta>0$ 的约束，则与 $(0,\infty)$ 取交集。在来自该模型的有限样本下，交集几乎必然非空。中程数 $(X_{(1)}+X_{(n)})/2$ 是无额外约束时的一种选取。这说明最大似然估计不一定唯一，也不一定来自似然方程。
\end{solve}

类似地，当 $U(a,b)$ 的两个端点均未知且样本极差为正时，似然为 $(b-a)^{-n}\mathbf1_{\{a\le x_{(1)},\,b\ge x_{(n)}\}}$，最小的可行区间给出 $\widehat a_{\mathrm{ML}}=X_{(1)}$、$\widehat b_{\mathrm{ML}}=X_{(n)}$。若极差为零，则可以令区间长度趋于零，似然无界，不能使用这一结论。

\subsection{离散参数与不存在矩的总体}

\begin{example}
\label{ex:study-capture-mle}
教材例 6.7 中，先标记并放回 $500$ 尾鱼，再随机捕获 $400$ 尾，其中标记鱼有 $60$ 尾。假定总体封闭、标记不脱落，标记鱼充分混合且与其他鱼具有相同捕获机会，求总数 $N$ 的最大似然估计。
\end{example}

\begin{solve}
第二次捕获数中的标记鱼数服从超几何分布。固定观测结果后，似然为
\[
  L(N)=\frac{\binom{500}{60}\binom{N-500}{340}}{\binom N{400}},
  \qquad N=840,841,\ldots.
\]
对 $N\ge841$，相邻似然比为
\[
  \frac{L(N)}{L(N-1)}
    =\frac{(N-500)(N-400)}{N(N-840)}.
\]
分子减去分母等于 $200000-60N$，所以似然先增后减，在 $N=\lfloor200000/60\rfloor=3333$ 达到最大。这里的 $N$ 是整数参数，不需要对它求导。一般地，若对应阈值恰为可行域内部的整数 $r$，相邻比在 $N=r$ 等于 $1$，则 $r-1,r$ 都是最大值点，不能忽略并列最大值。
\end{solve}

\begin{example}
\label{ex:study-cauchy-location}
设总体具有柯西位置分布（Cauchy location distribution）
\[
  f(x;\theta)=\frac1{\pi[1+(x-\theta)^2]},
  \qquad x,\theta\in\mathbb R.
\]
讨论位置参数的估计。
\end{example}

\begin{solve}
柯西分布没有有限的一阶绝对矩，因而不能用 $\mathbb E(X)=\theta$ 建立通常的一阶矩方程。最大似然法仍可使用，其对数似然和似然方程为
\[
  \ell(\theta)=-n\log\pi-\sum_{i=1}^n\log[1+(x_i-\theta)^2],
  \qquad
  \ell'(\theta)=2\sum_{i=1}^n\frac{x_i-\theta}{1+(x_i-\theta)^2}=0.
\]
由于 $\ell$ 连续且在 $|\theta|\to\infty$ 时趋于 $-\infty$，最大值存在，但似然方程可能有多个根，求得一个根后仍须比较似然值。另一方面，$\theta$ 是该分布唯一的中位数，可直接用样本中位数
\[
  \widetilde X=
  \begin{cases}
    X_{((n+1)/2)},&n\text{ 为奇数},\\
    (X_{(n/2)}+X_{(n/2+1)})/2,&n\text{ 为偶数}
  \end{cases}
\]
来估计。它是利用分位数构造的估计，不应仅因形式简单就称为最大似然估计；其优劣仍需在给定准则下研究。
\end{solve}

\section{优良性准则}
\label{sec:study-estimator-criteria}

评价估计量时应把样本重新视为随机向量。某一次估计碰巧接近真值，并不能说明该方法在重复抽样中更准确；一种方法在某个参数值下较好，也不意味着它在整个参数空间内都较好。以下先考虑标量目标 $g(\theta)$，所有涉及期望或方差的结论均要求相应量存在。

\subsection{偏差、无偏性及修正}

\begin{definition}[偏差与无偏性]
\label{def:study-estimator-bias}
估计量 $T$ 的\emph{偏差（bias）}为
\begin{equation}
  b_\theta(T)=\mathbb E_\theta(T)-g(\theta).
  \label{eq:study-estimator-bias}
\end{equation}
若对每个 $\theta\in\Theta$ 都有 $b_\theta(T)=0$，则称 $T$ 为 $g(\theta)$ 的\emph{无偏估计量（unbiased estimator）}。若估计量序列 $T_n$ 满足对每个 $\theta$ 都有 $\mathbb E_\theta(T_n)\to g(\theta)$，则称其为渐近无偏估计。
\end{definition}

无偏性是抽样分布的均值性质，不是每次观测都估计正确。固定样本量 $n$，独立重复整个抽样过程，若 $\mathbb E_\theta|T|<\infty$，则这些重复估计值的平均依强大数定律趋于 $\mathbb E_\theta T$；无偏时这个极限恰好是 $g(\theta)$。这里增加的是独立重复抽样的次数，不能与单次样本量 $n\to\infty$ 混为一谈。

\begin{proposition}[均值和方差的无偏估计]
\label{prop:study-unbiased-mean-variance}
若 $X_i$ 独立同分布，$\mathbb E(X_i)=\mu$、$\operatorname{Var}(X_i)=\sigma^2<\infty$，则
\begin{equation}
  \mathbb E(\overline X)=\mu,\qquad
  \operatorname{Var}(\overline X)=\frac{\sigma^2}{n},\qquad
  \mathbb E(M_2)=\frac{n-1}{n}\sigma^2.
  \label{eq:study-mean-variance-expectations}
\end{equation}
因此 $n\ge2$ 时，式~\eqref{eq:study-unbiased-sample-variance} 中的 $S^2$ 是总体方差的无偏估计；结论不要求总体正态。
\end{proposition}

\begin{proof}
期望的线性性与样本独立性给出前两式。再由恒等式
\[
  \sum_{i=1}^n(X_i-\overline X)^2
    =\sum_{i=1}^n(X_i-\mu)^2-n(\overline X-\mu)^2,
\]
取期望得到 $n\mathbb E(M_2)=n\sigma^2-n(\sigma^2/n)=(n-1)\sigma^2$。乘以 $1/(n-1)$ 即得 $\mathbb E(S^2)=\sigma^2$。
\end{proof}

残差满足 $\sum(X_i-\overline X)=0$，因而落在一个 $n-1$ 维线性子空间中，这给出自由度减少一的几何解释，但不能据此声称一般总体的各个残差相互独立。无偏性也不在非线性变换下保留：由柯西--施瓦茨不等式，$\mathbb E(S)\le\sqrt{\mathbb E(S^2)}=\sigma$，当 $S$ 非退化时严格小于 $\sigma$。所以 $S^2$ 无偏并不意味着 $S$ 无偏；同理，$\overline X$ 对 $1/\lambda$ 无偏也不意味着 $1/\overline X$ 对 $\lambda$ 无偏。

\begin{example}
\label{ex:study-uniform-unbiased-correction}
设 $X_i\sim U(0,\theta)$，计算最大值的偏差并作无偏修正。
\end{example}

\begin{solve}
对 $0\le x\le\theta$，独立性给出
\[
  \mathbb P_\theta(X_{(n)}\le x)=\left(\frac x\theta\right)^n,
  \qquad
  f_{X_{(n)}}(x)=\frac{nx^{n-1}}{\theta^n},\quad0<x<\theta.
\]
因此，对 $r>0$，
\begin{equation}
  \mathbb E_\theta(X_{(n)}^r)
    =\int_0^\theta x^r\frac{nx^{n-1}}{\theta^n}\,dx
    =\frac n{n+r}\theta^r.
  \label{eq:study-uniform-maximum-moments}
\end{equation}
特别地，$\mathbb E_\theta(X_{(n)})=n\theta/(n+1)$，偏差为 $-\theta/(n+1)$，故
\begin{equation}
  \widetilde\theta=\frac{n+1}{n}X_{(n)},\qquad
  \operatorname{Var}_\theta(\widetilde\theta)
     =\frac{\theta^2}{n(n+2)}
  \label{eq:study-uniform-corrected-maximum}
\end{equation}
是无偏修正及其方差。修正后的统计量不再是原来的 MLE；原来的 MLE 虽有偏，却是渐近无偏的。
\end{solve}

\begin{example}
\label{ex:study-no-unbiased-reciprocal}
设只观测 $Y\sim\operatorname{Bin}(N,p)$，其中正整数 $N$ 已知、$0<p<1$。证明 $1/p$ 不存在实值无偏估计量。
\end{example}

\begin{solve}
任一实值统计量由有限个数 $t_0,\ldots,t_N$ 决定，其期望
\[
  P(p)=\sum_{j=0}^Nt_j\binom Njp^j(1-p)^{N-j}
\]
是次数至多为 $N$ 的多项式。若它恒等于 $1/p$，则多项式 $pP(p)-1$ 在 $(0,1)$ 的每一点为零，故必须为零多项式；但其在 $p=0$ 处的值为 $-1$，矛盾。无偏估计不一定存在，因此无偏性不能成为不加区别的唯一要求。
\end{solve}

\subsection{均方误差、有效性与最小方差无偏估计}

\begin{definition}[均方误差]
\label{def:study-estimator-mse}
若 $\mathbb E_\theta(T^2)<\infty$，定义估计 $g(\theta)$ 的\emph{均方误差（mean squared error，MSE）}为
\begin{equation}
  \operatorname{MSE}_\theta(T)
     =\mathbb E_\theta[(T-g(\theta))^2]
     =\operatorname{Var}_\theta(T)+b_\theta(T)^2.
  \label{eq:study-mse-decomposition}
\end{equation}
\end{definition}

最后一个等式由 $T-g=(T-\mathbb E_\theta T)+b_\theta(T)$ 展开平方并取期望得到，其中交叉项为零。它把随机波动与系统偏差纳入同一个准则。对无偏估计，MSE 等于方差；对有偏估计，只比较方差可能得出错误结论。也可以采用平均绝对误差 $\mathbb E_\theta|T-g(\theta)|$ 等准则，不同损失函数对应不同的评价目标。

\begin{definition}[有效性与最小方差无偏估计]
\label{def:study-mvue}
若 $T_1,T_2$ 都是 $g(\theta)$ 的无偏估计，且对所有 $\theta\in\Theta$ 都有
\[
  \operatorname{Var}_\theta(T_1)\le\operatorname{Var}_\theta(T_2),
\]
并至少在一个参数值处严格小于，则称 $T_1$ 比 $T_2$ 更有效。若无偏估计 $T_*$ 对每个具有有限方差的无偏估计 $T$ 都满足
\begin{equation}
  \operatorname{Var}_\theta(T_*)\le\operatorname{Var}_\theta(T),
  \qquad\theta\in\Theta,
  \label{eq:study-umvue}
\end{equation}
则称 $T_*$ 为\emph{一致最小方差无偏估计量（uniformly minimum variance unbiased estimator，UMVUE）}，教材简称 MVUE。“一致”指同一个估计量在整个参数空间逐点最优，不是只在某个未知真值处最优。
\end{definition}

将一个估计量与少数候选量作比较，只能得到这些候选量之间的结论，不能据此宣布它在所有无偏估计中最优。例如，设总体均值为 $\mu$、方差为 $\sigma^2>0$，取不含未知参数的常数权重 $\sum_{i=1}^nw_i=1$，则 $T_w=\sum w_iX_i$ 无偏，且
\[
  \operatorname{Var}(T_w)=\sigma^2\sum_{i=1}^nw_i^2
    \ge\frac{\sigma^2}{n}
    =\operatorname{Var}(\overline X).
\]
这里使用 $(\sum w_i)^2\le n\sum w_i^2$，等号当且仅当 $w_i=1/n$。这证明样本均值在常系数线性无偏估计类内最优，尚未覆盖所有非线性无偏估计。

在均匀总体 $U(0,\theta)$ 中，矩估计 $2\overline X$ 和式~\eqref{eq:study-uniform-corrected-maximum} 中的 $\widetilde\theta$ 都无偏，而
\[
  \operatorname{Var}_\theta(2\overline X)=\frac{\theta^2}{3n},
  \qquad
  \operatorname{Var}_\theta(\widetilde\theta)
       =\frac{\theta^2}{n(n+2)}.
\]
当 $n>1$ 时后者严格较小，说明即使样本均值在线性类内最优，使用次序统计量仍可能进一步降低方差。无偏修正也不必降低 MSE：方差与偏差平方的权衡见本章习题中教材第 14 题的解答。

\subsection{克拉默--拉奥方差下界}

要证明某个估计量在所有无偏估计中最优，可以先寻找共同的方差下界。下面只讨论一维参数 $\theta$，并明确正则条件：参数空间为开区间，样本支撑与参数无关，密度或概率质量函数关于参数可微，且对归一化条件和下述无偏期望求导时均可交换微分与积分（离散情形为求和）。例如，在每个参数邻域内，相关导数被可积函数控制，是保证这些交换的一组充分条件。此外要求信息量正且有限、估计量方差有限。

\begin{definition}[得分与费希尔信息]
\label{def:study-fisher-information}
单个观测的\emph{得分（score）}和\emph{费希尔信息（Fisher information）}定义为
\begin{equation}
  u_\theta(X)=\frac{\partial}{\partial\theta}\log f(X;\theta),
  \qquad
  I_1(\theta)=\mathbb E_\theta[u_\theta(X)^2].
  \label{eq:study-fisher-information}
\end{equation}
样本得分为 $U_n(\theta)=\sum_{i=1}^nu_\theta(X_i)$。在上述正则条件下，$\mathbb E_\theta u_\theta=0$，由独立性有
\[
  I_n(\theta):=\operatorname{Var}_\theta(U_n(\theta))
      =nI_1(\theta).
\]
\end{definition}

\begin{theorem}[克拉默--拉奥不等式]
\label{thm:study-cramer-rao}
在上述正则条件下，设 $T$ 是可微参数函数 $g(\theta)$ 的无偏估计，且允许在 $\mathbb E_\theta T=g(\theta)$ 中交换微分与积分，则
\begin{equation}
  \operatorname{Var}_\theta(T)
    \ge\frac{[g'(\theta)]^2}{nI_1(\theta)}.
  \label{eq:study-cramer-rao}
\end{equation}
在给定 $\theta$ 处等号成立，当且仅当
\begin{equation}
  T-g(\theta)
     =\frac{g'(\theta)}{nI_1(\theta)}U_n(\theta)
     \quad\mathbb P_\theta\text{-几乎处处}.
  \label{eq:study-cramer-rao-equality}
\end{equation}
\end{theorem}

\begin{proof}
设联合密度为 $f_n(\boldsymbol x;\theta)$。对 $\int f_n\,d\boldsymbol x=1$ 和 $\int T f_n\,d\boldsymbol x=g(\theta)$ 求导，利用 $\partial_\theta f_n=U_n f_n$，得到
\[
  \mathbb E_\theta U_n=0,\qquad
  \mathbb E_\theta(TU_n)=g'(\theta).
\]
故 $\operatorname{Cov}_\theta(T,U_n)=g'(\theta)$。协方差的柯西--施瓦茨不等式给出
\[
  [g'(\theta)]^2
    \le\operatorname{Var}_\theta(T)
       \operatorname{Var}_\theta(U_n)
    =\operatorname{Var}_\theta(T)\,nI_1(\theta),
\]
除以正数 $nI_1(\theta)$ 即得下界。等号条件也可由
\[
  \mathbb E_\theta\!\left[
    \left(T-g(\theta)-\frac{g'(\theta)}{nI_1(\theta)}U_n(\theta)\right)^2
  \right]
   =\operatorname{Var}_\theta(T)
     -\frac{[g'(\theta)]^2}{nI_1(\theta)}
\]
看出：右端为零，当且仅当括号内的随机变量几乎处处为零。
\end{proof}

从几何上看，应在零均值、平方可积的随机变量空间中，以 $\mathbb E_\theta(ZY)$ 为内积，将 $T-g(\theta)$ 投影到得分方向；方差下界就是投影长度的平方。这样可避免把所有随机变量上的协方差直接当作内积，因为非零常数的协方差范数为零。下界不一定能够达到；若某个无偏估计对所有参数都达到下界，它必为 UMVUE，反过来 UMVUE 不一定达到该下界。

\begin{example}
\label{ex:study-cramer-rao-applications}
用克拉默--拉奥下界判断正态均值、指数均值和二项比例的样本均值型估计是否为 UMVUE。
\end{example}

\begin{solve}
正态总体 $N(\mu,\sigma^2)$ 中先设 $\sigma^2>0$ 已知，以 $\mu$ 为参数，则
\[
  u_\mu(X)=\frac{X-\mu}{\sigma^2},\qquad
  I_1(\mu)=\frac1{\sigma^2}.
\]
$\mu$ 的无偏估计方差不小于 $\sigma^2/n$，而 $\overline X$ 恰好达到，故为 UMVUE。即使 $\sigma^2$ 也未知，任何对所有 $(\mu,\sigma^2)$ 无偏的估计，在固定任意一个 $\sigma^2$ 的子模型中仍须满足该下界；$\overline X$ 因而仍为 $\mu$ 的 UMVUE。若反而是 $\mu=\mu_0$ 已知，以 $v=\sigma^2$ 为参数，则 $I_1(v)=1/(2v^2)$，无偏估计 $n^{-1}\sum(X_i-\mu_0)^2$ 的方差为 $2v^2/n$，也达到下界。

指数分布 $\operatorname{Exp}(\lambda)$ 中，$u_\lambda(X)=1/\lambda-X$、$I_1(\lambda)=1/\lambda^2$。估计的是均值 $g(\lambda)=1/\lambda$，其导数为 $-1/\lambda^2$，故下界为
\[
  \frac{[g'(\lambda)]^2}{nI_1(\lambda)}
    =\frac1{n\lambda^2}
    =\operatorname{Var}_\lambda(\overline X).
\]
因此 $\overline X$ 是 $1/\lambda$ 的 UMVUE。这一结论不能直接改写成 $1/\overline X$ 是 $\lambda$ 的无偏或最优估计。

对于 $\operatorname{Bin}(N,p)$，$N$ 已知、$0<p<1$，有
\[
  u_p(X)=\frac{X-Np}{p(1-p)},\qquad
  I_1(p)=\frac{N}{p(1-p)}.
\]
估计 $p$ 的方差下界为 $p(1-p)/(nN)$，正好等于无偏估计 $\overline X/N$ 的方差，故后者为 UMVUE。类似地，泊松总体有 $I_1(\lambda)=1/\lambda$，$\overline X$ 以方差 $\lambda/n$ 达到下界；这也证明它优于其他无偏估计，包括用 $S^2$ 估计同一个参数 $\lambda$ 的做法。
\end{solve}

\begin{notealert}[type=caution,title={支撑依赖参数时不能机械套用下界}]
对于 $U(0,\theta)$，若只在支撑内部计算 $\partial_\theta\log f=-1/\theta$，会得到一个期望不为零的“得分”，说明正则推导已失效。积分上限依赖 $\theta$，归一化条件求导时存在边界项。因此不能把 $I_1(\theta)=1/\theta^2$ 代入式~\eqref{eq:study-cramer-rao}，再断言方差至少为 $\theta^2/n$；式~\eqref{eq:study-uniform-corrected-maximum} 的方差实际更小，并不与定理矛盾。
\end{notealert}

\begin{proposition}[均匀总体中修正最大值的最优性]
\label{prop:study-uniform-maximum-umvue}
在 $U(0,\theta)$ 的简单随机样本中，$\widetilde\theta=(n+1)X_{(n)}/n$ 是 $\theta$ 的 UMVUE。
\end{proposition}

\begin{proof}
当 $n=1$ 时也可使用以下论证。记 $M=X_{(n)}$。在给定 $M=m$ 后，最大值所在位置等可能，其余 $n-1$ 个观测的条件分布为独立的 $U(0,m)$，不含 $\theta$。因此，对任意无偏估计 $T$，条件期望 $H(M)=\mathbb E_\theta(T\mid M)$ 可取为不含未知参数的统计量，并且 $\mathbb E_\theta H(M)=\theta$。由全方差公式，$\operatorname{Var}_\theta H(M)\le\operatorname{Var}_\theta T$。

只需证明作为 $M$ 的函数的无偏估计唯一。若 $h(M)$ 在所有 $\theta>0$ 下期望为零，则
\[
  \int_0^\theta h(m)m^{n-1}\,dm=0,\qquad\theta>0.
\]
因 $h(M)$ 可积，被积函数局部可积，对上式按积分上限求导，得到 $h(m)m^{n-1}=0$ 对几乎所有 $m>0$ 成立。于是 $h(M)=0$ 在每个 $\theta$ 下几乎必然成立。取 $h(M)=H(M)-(n+1)M/n$，便知 $H(M)=\widetilde\theta$，再结合方差不等式得到结论。
\end{proof}

\section{点估计量的大样本理论}
\label{sec:study-estimator-asymptotics}

有限样本性质是在固定 $n$ 下研究估计量的抽样分布；大样本性质则研究估计量序列在 $n\to\infty$ 时的行为。两者的区别是研究方式，而不是某个统一的样本量门槛。无偏性控制期望，相合性控制误差概率，渐近正态性进一步描述适当缩放后的误差分布，三者不能互相替代。

\subsection{相合性}

\begin{definition}[相合估计]
\label{def:study-consistent-estimator}
若对每个 $\theta\in\Theta$、每个 $\varepsilon>0$，都有
\begin{equation}
  \mathbb P_\theta(|T_n-g(\theta)|>\varepsilon)\longrightarrow0,
  \qquad n\to\infty,
  \label{eq:study-estimator-consistency}
\end{equation}
则称 $T_n$ 为 $g(\theta)$ 的\emph{弱相合估计量（weakly consistent estimator）}，记为 $T_n\xrightarrow{\mathbb P_\theta}g(\theta)$，简称相合估计。
\end{definition}

定义中的“对每个 $\theta$”是逐点要求，不等于更强的 $\sup_{\theta\in\Theta}\mathbb P_\theta(|T_n-g(\theta)|>\varepsilon)\to0$。依概率收敛也不表示每一条样本路径的误差都单调减小；它表示固定精度以外的误差越来越少见。

\begin{proposition}[由均方误差趋于零判断相合性]
\label{prop:study-mse-consistency}
若对每个 $\theta$，$\operatorname{MSE}_\theta(T_n)\to0$，则 $T_n$ 相合。特别地，渐近无偏且方差趋于零是相合性的一组充分条件。
\end{proposition}

\begin{proof}
由非负随机变量的马尔可夫不等式，
\[
  \mathbb P_\theta(|T_n-g(\theta)|>\varepsilon)
    \le\frac{\mathbb E_\theta[(T_n-g(\theta))^2]}{\varepsilon^2}
    \longrightarrow0.
\]
后一结论再用式~\eqref{eq:study-mse-decomposition} 即得。
\end{proof}

相合性不要求每个有限样本量下都无偏。反过来，无偏也不保证相合：若总体非退化，始终使用 $T_n=X_1$ 估计均值，则虽对每个 $n$ 都无偏，却存在 $\varepsilon>0$ 使 $\mathbb P(|X_1-\mu|>\varepsilon)>0$ 不随 $n$ 改变。

\begin{example}
\label{ex:study-moment-consistency}
设 $X_i$ 独立同分布且 $\mathbb E(X_i^2)<\infty$，证明 $\overline X$、$M_2$ 和 $S^2$ 分别相合地估计均值与方差。
\end{example}

\begin{solve}
由大数定律，
\[
  \overline X\xrightarrow{\mathbb P}\mu,\qquad
  A_2\xrightarrow{\mathbb P}\mathbb E(X^2).
\]
连续函数保持依概率收敛，故
\[
  M_2=A_2-\overline X^2
      \xrightarrow{\mathbb P}\mathbb E(X^2)-\mu^2=\sigma^2.
\]
再由 $S^2=nM_2/(n-1)$ 得 $S^2\xrightarrow{\mathbb P}\sigma^2$。这里不需要四阶矩存在：证明使用二阶原点矩的大数定律，而不是先计算样本方差自身的方差。更一般地，若矩方程在真值附近有连续的逆，且所选样本解落在对应分支上，则由样本矩的收敛可推出该矩估计相合；不能略去解的存在、分支选择等条件。
\end{solve}

\begin{example}
\label{ex:study-uniform-maximum-consistency}
设 $X_i\sim U(0,\theta)$，证明 MLE $X_{(n)}$ 相合，并考察其误差缩放。
\end{example}

\begin{solve}
对 $0<\varepsilon<\theta$，因 $X_{(n)}\le\theta$ 几乎必然成立，
\begin{equation}
  \mathbb P_\theta(|X_{(n)}-\theta|>\varepsilon)
    =\left(1-\frac{\varepsilon}{\theta}\right)^n\longrightarrow0.
  \label{eq:study-uniform-maximum-consistency}
\end{equation}
当 $\varepsilon\ge\theta$ 时该概率为零，故相合性成立。进一步，对固定 $t\ge0$、充分大的 $n$，
\[
  \mathbb P_\theta\!\left(
      \frac{n(\theta-X_{(n)})}{\theta}>t
  \right)=\left(1-\frac tn\right)^n\longrightarrow e^{-t}.
\]
因而 $n(\theta-X_{(n)})/\theta$ 依分布趋于速率为 $1$ 的指数分布。它的自然误差尺度为 $1/n$，极限分布不是正态分布；这也是不能对支撑依赖参数的模型套用一般正则 MLE 结论的例子。
\end{solve}

\subsection{渐近正态性与误差近似}

\begin{definition}[标准化误差的渐近正态性]
\label{def:study-asymptotic-normality}
沿用教材的标准化形式，设 $v_n(\theta)=\operatorname{Var}_\theta(T_n)\in(0,\infty)$。若
\begin{equation}
  \frac{T_n-g(\theta)}{\sqrt{v_n(\theta)}}
      \xrightarrow{\mathcal D}N(0,1),
  \label{eq:study-asymptotic-normality}
\end{equation}
则称该标准化误差具有\emph{渐近正态性（asymptotic normality）}，其中 $\xrightarrow{\mathcal D}$ 表示依分布收敛。
\end{definition}

常用的另一种表达是 $\sqrt n(T_n-g(\theta))\xrightarrow{\mathcal D}N(0,v(\theta))$，其中 $0<v(\theta)<\infty$。若另有 $nv_n(\theta)\to v(\theta)$，它与式~\eqref{eq:study-asymptotic-normality} 等价；若尚未证明二阶矩收敛，就不能仅凭分布极限把 $v(\theta)$ 当作 $\lim n\operatorname{Var}_\theta(T_n)$。正态近似通常用于相合估计，在上述标准化形式下还应有 $v_n(\theta)\to0$ 才能据此推出误差趋于零。

\begin{theorem}[样本均值及其可行标准化]
\label{thm:study-studentized-mean-limit}
若 $X_i$ 独立同分布，均值为 $\mu$，且 $0<\sigma^2<\infty$，则
\begin{equation}
  \frac{\sqrt n(\overline X-\mu)}{\sigma}
       \xrightarrow{\mathcal D}N(0,1),\qquad
  \frac{\sqrt n(\overline X-\mu)}{S}
       \xrightarrow{\mathcal D}N(0,1).
  \label{eq:study-studentized-mean-limit}
\end{equation}
第二个统计量在 $S=0$ 时可任意约定其值，该事件的概率随 $n$ 增大趋于零。
\end{theorem}

\begin{proof}
第一式是独立同分布中心极限定理。由例~\ref{ex:study-moment-consistency}，$S^2\xrightarrow{\mathbb P}\sigma^2>0$，故 $S/\sigma\xrightarrow{\mathbb P}1$，且 $\mathbb P(S=0)\to0$。将第一式的左端乘以 $\sigma/S$，由斯卢茨基定理（Slutsky's theorem）得到第二式：依分布收敛的随机变量与依概率收敛到常数的随机变量相乘，其极限为相应极限之积。
\end{proof}

因此，在大样本下 $\overline X$ 的标准误差可估为 $S/\sqrt n$。总体本身不必正态；只有总体正态时，固定有限 $n$ 的 $(\overline X-\mu)/(S/\sqrt n)$ 才具有精确的 $t_{n-1}$ 分布。大样本近似的精度取决于总体分布，不能只凭一个固定样本量阈值保证。

参数函数的极限可用\emph{德尔塔方法（delta method）}得到。若 $\widehat\theta_n\xrightarrow{\mathbb P}\theta$、$\sqrt n(\widehat\theta_n-\theta)\xrightarrow{\mathcal D}N(0,v)$，且 $g$ 在 $\theta$ 处可微，则一阶展开给出
\begin{equation}
  \sqrt n\,[g(\widehat\theta_n)-g(\theta)]
      \xrightarrow{\mathcal D}
          N\!\left(0,[g'(\theta)]^2v\right).
  \label{eq:study-delta-method}
\end{equation}
其依据是可微性将差值写成 $[g'(\theta)+r_n](\widehat\theta_n-\theta)$，其中 $r_n\xrightarrow{\mathbb P}0$，再用斯卢茨基定理。若 $g'(\theta)=0$，右端退化为零，不能继续按非退化正态分布标准化，往往需要更高阶展开。

例如，指数分布下 $\mathbb E X=1/\lambda$、$\operatorname{Var}X=1/\lambda^2$，对 $g(t)=1/t$ 使用式~\eqref{eq:study-delta-method}，得到
\[
  \sqrt n\left(\frac1{\overline X}-\lambda\right)
      \xrightarrow{\mathcal D}N(0,\lambda^2).
\]
泊松总体下，零次事件概率 $q=e^{-\lambda}$ 的 MLE 为 $\widehat q=e^{-\overline X}$；当 $\lambda>0$ 时，同理有
\[
  \sqrt n(\widehat q-q)
      \xrightarrow{\mathcal D}N(0,\lambda e^{-2\lambda}).
\]
这些是对具体模型的直接验证。一般 MLE 的相合性与渐近正态性还涉及参数可识别、真值为内点、似然的光滑性与信息量等条件，不是最大化定义本身的无条件推论。

\section{扩展阅读：德军坦克问题}
\label{sec:study-german-tank}

将教材中的序列号问题抽象为：总体编号为 $1,2,\ldots,N$，其中 $N$ 为未知正整数；从中等可能、不放回地抽取 $k$ 个不同编号，$k$ 已知且 $1\le k\le N$。记抽到的编号集合为 $A$，最大编号为 $M=\max A$，其观测值为 $m$。关键假设是每个大小为 $k$ 的编号子集等可能；若某些编号更容易被观测，以下估计的性质便可能改变。这是有限总体的不放回抽样，不能把编号视为独立的连续均匀样本。

\subsection{最大编号的分布及无偏修正}

在 $\binom Nk$ 个等可能子集中，最大值等于 $m$ 时必须包含 $m$，其余 $k-1$ 个编号从 $1,\ldots,m-1$ 中选择，故
\begin{equation}
  \mathbb P_N(M=m)
    =\frac{\binom{m-1}{k-1}}{\binom Nk},
  \qquad m=k,k+1,\ldots,N.
  \label{eq:study-tank-maximum-mass}
\end{equation}
利用组合恒等式 $m\binom{m-1}{k-1}=k\binom mk$ 及 $\sum_{m=k}^N\binom mk=\binom{N+1}{k+1}$，有
\begin{equation}
  \mathbb E_N M
    =\frac{k\binom{N+1}{k+1}}{\binom Nk}
    =\frac{k(N+1)}{k+1}.
  \label{eq:study-tank-maximum-mean}
\end{equation}
因此，直接用 $M$ 会低估 $N$，偏差为 $-(N-k)/(k+1)$；解出期望关系并用 $M$ 替换 $\mathbb E_NM$，得到
\begin{equation}
  \widehat N_{\mathrm U}=\frac{k+1}{k}M-1,
  \qquad \mathbb E_N\widehat N_{\mathrm U}=N.
  \label{eq:study-tank-unbiased-estimator}
\end{equation}
它可解释为在最大观测编号之后补上一个估计的平均间隔。该估计量可能不是整数；若再四舍五入，应另行评价舍入后的偏差，不能沿用原来的无偏结论。

为了评价精度，进一步利用
\[
  m(m+1)\binom{m-1}{k-1}
      =k(k+1)\binom{m+1}{k+1}
\]
求和，得到
\[
  \mathbb E_N[M(M+1)]=\frac{k(N+1)(N+2)}{k+2}.
\]
减去 $\mathbb E_N M$ 得到二阶矩，再减去均值平方，可得
\begin{equation}
  \operatorname{Var}_N(M)
    =\frac{k(N+1)(N-k)}{(k+1)^2(k+2)},\qquad
  \operatorname{Var}_N(\widehat N_{\mathrm U})
    =\frac{(N+1)(N-k)}{k(k+2)}.
  \label{eq:study-tank-estimator-variance}
\end{equation}
当 $N\gg k$ 时，后者约为 $N^2/[k(k+2)]$，标准差约为 $N/\sqrt{k(k+2)}$；若 $k$ 还较大，可进一步近似为 $N/k$。这里 $N/k$ 是标准差的近似量级，不能与方差混淆。

\subsection{最大似然估计与最小方差性质}

固定观测集合 $A$ 后，
\[
  L(N;A)=\binom Nk^{-1}\mathbf1_{\{N\ge m\}},
  \qquad N=k,k+1,\ldots.
\]
因为 $\binom Nk$ 随 $N\ge k$ 严格递增，故 $\widehat N_{\mathrm{ML}}=M$。最大似然估计与无偏修正不同：前者使当前样本的似然最大，后者使重复抽样的平均值等于总体大小，这两个目标没有理由产生同一个统计量。

\begin{proposition}[坦克总数的最小方差无偏估计]
\label{prop:study-tank-umvue}
在上述等可能不放回抽样模型下，式~\eqref{eq:study-tank-unbiased-estimator} 是 $N$ 的 UMVUE。
\end{proposition}

\begin{proof}
给定 $M=m$ 后，样本必含 $m$，其余元素在 $1,\ldots,m-1$ 的所有 $k-1$ 元子集中等可能，条件分布不依赖 $N$。因此，对任意无偏估计 $T(A)$，$H(M)=\mathbb E_N(T\mid M)$ 是不含 $N$ 的统计量，仍无偏，且由全方差公式有 $\operatorname{Var}_N H(M)\le\operatorname{Var}_N T$。

若一个函数 $h(M)$ 对每个 $N\ge k$ 都满足 $\mathbb E_Nh(M)=0$，则
\[
  \sum_{m=k}^Nh(m)\binom{m-1}{k-1}=0,\qquad N=k,k+1,\ldots.
\]
取 $N=k$ 得 $h(k)=0$；若已知 $h(k),\ldots,h(N-1)$ 都为零，等式便给出 $h(N)\binom{N-1}{k-1}=0$，故 $h(N)=0$。由归纳法，$h$ 在所有可能取值上均为零。因此 $M$ 的函数中无偏估计唯一，必为 $\widehat N_{\mathrm U}$，再结合条件平均的方差不等式即得最优性。
\end{proof}

若从总体一阶矩 $\mathbb E_N X=(N+1)/2$ 出发，也能构造无偏估计 $2\overline X-1$；这与利用最大编号的期望关系作修正是不同的构造。坦克问题说明，次序统计量能有效使用参数决定的支撑信息，但这种优势以抽样模型成立为前提，并不能仅凭某次估计值接近真值来验证模型。

\section*{本章总结}
\addcontentsline{toc}{section}{本章总结}

点估计首先要区分固定未知参数、随机估计量和观测后的估计值。矩估计通过样本矩与总体矩匹配构造参数估计，须检查矩的存在、方程的可解性和参数约束；最大似然估计固定数据后比较参数，须保留支撑条件、检查边界与全局最大值，并允许不唯一或不存在的情形。对参数函数作代入估计很方便，但非线性变换一般不保持无偏性。

无偏性是期望正确，MSE 同时计算方差与偏差平方；有效性是在给定估计类与整个参数空间内的比较。证明 UMVUE 需要与所有无偏估计比较，正则模型中可使用克拉默--拉奥下界，支撑依赖参数的模型则要另作论证。相合性说明误差概率趋于零，渐近正态性说明标准化误差的分布近似；无偏、相合和渐近正态之间均不能无条件互推。均匀总体最大值与坦克编号最大值都能作无偏修正，但前者来自独立连续抽样，后者来自有限总体不放回抽样，其分布和修正公式必须分别推导。

\section*{习题}
\addcontentsline{toc}{section}{习题}

以下选解教材第六章第 12、14、20、27、44、54、56、58 题，分别检验次序统计量的修正、偏差与方差的权衡、加权估计、参数函数估计及相合性；题目中的端点密度沿用本章包含边界的版本，不改变原分布。

\begin{example}
\label{ex:study-ch6-ex12}
教材第 12 题。设 $X_1,X_2,X_3$ 为 $U(0,\theta)$ 的简单随机样本。证明 $\tfrac43X_{(3)}$ 和 $4X_{(1)}$ 都是 $\theta$ 的无偏估计，并比较有效性。
\end{example}

\begin{solve}
由式~\eqref{eq:study-uniform-maximum-moments}，$\mathbb E X_{(3)}=3\theta/4$，$\operatorname{Var}(X_{(3)})=3\theta^2/80$。由于 $\theta-X_i$ 仍为独立的 $U(0,\theta)$ 变量，$X_{(1)}$ 与 $\theta-X_{(3)}$ 同分布，故 $\mathbb E X_{(1)}=\theta/4$，方差同样为 $3\theta^2/80$。于是两个估计均无偏，而
\[
  \operatorname{Var}\!\left(\frac43X_{(3)}\right)
     =\frac{\theta^2}{15},\qquad
  \operatorname{Var}(4X_{(1)})=\frac{3\theta^2}{5}.
\]
前者的方差为后者的 $1/9$，因此更有效；虽然最大值与最小值本身的方差相同，无偏修正所需倍数的不同改变了最终精度。
\end{solve}

\begin{example}
\label{ex:study-ch6-ex14}
教材第 14 题。设 $X_i\sim N(\mu,\sigma^2)$，$n\ge2$，记 $Q=\sum(X_i-\overline X)^2$。比较无偏估计 $Q/(n-1)$ 与有偏估计 $Q/(n+1)$ 对 $\sigma^2$ 的均方误差。
\end{example}

\begin{solve}
对正态样本作正交变换，残差平方和 $Q/\sigma^2$ 可写成 $n-1$ 个独立标准正态变量的平方和，故服从 $\chi^2_{n-1}$，从而
\[
  \mathbb E Q=(n-1)\sigma^2,\qquad
  \operatorname{Var}(Q)=2(n-1)\sigma^4.
\]
在常数倍估计类 $cQ$ 中，
\[
  \frac{\operatorname{MSE}(cQ)}{\sigma^4}
     =2(n-1)c^2+[(n-1)c-1]^2
     =(n-1)(n+1)c^2-2(n-1)c+1.
\]
此二次式在 $c=1/(n+1)$ 处最小，因此
\[
  \operatorname{MSE}\!\left(\frac Q{n+1}\right)
     =\frac{2\sigma^4}{n+1}
     <\frac{2\sigma^4}{n-1}
     =\operatorname{MSE}\!\left(\frac Q{n-1}\right).
\]
$Q/(n+1)$ 的偏差为 $-2\sigma^2/(n+1)$，但方差的下降超过偏差平方的增加。这证明有偏估计可能具有更小的 MSE；这里只在 $cQ$ 这一估计类内求得最优值。
\end{solve}

\begin{example}
\label{ex:study-ch6-ex20}
教材第 20 题。$k$ 台仪器独立测量同一常数 $\theta$，所得 $X_i$ 满足 $\mathbb E X_i=\theta$、$\operatorname{Var}(X_i)=\sigma_i^2$，其中 $\sigma_i^2>0$ 已知。求常系数线性无偏估计 $T=\sum_{i=1}^ka_iX_i$ 中方差最小的权重。
\end{example}

\begin{solve}
在以所有实数 $\theta$ 为可能值的模型中，无偏条件为 $\sum a_i=1$，独立性给出 $\operatorname{Var}(T)=\sum a_i^2\sigma_i^2$。由
\[
  1=\left(\sum_{i=1}^k(a_i\sigma_i)\frac1{\sigma_i}\right)^2
    \le\left(\sum_{i=1}^ka_i^2\sigma_i^2\right)
       \left(\sum_{i=1}^k\sigma_i^{-2}\right),
\]
最小方差及其达到条件为
\[
  a_i=\frac{\sigma_i^{-2}}{\sum_{j=1}^k\sigma_j^{-2}},
  \qquad
  \min\operatorname{Var}(T)
       =\frac1{\sum_{j=1}^k\sigma_j^{-2}}.
\]
即按方差倒数加权，精度越高的仪器权重越大；方差相同时退化为等权平均。结论限于给定的线性无偏类，且依赖测量独立与方差已知。
\end{solve}

\begin{example}
\label{ex:study-ch6-ex27}
教材第 27 题。设 $X_i\sim\operatorname{Pois}(\lambda)$，求 $q=\mathbb P_\lambda(X=0)$ 的 MLE。$122$ 个观测中，取值 $0,1,2,3,4,5$ 的频数依次为 $44,42,21,9,4,2$，没有更大值，求相应估计值。
\end{example}

\begin{solve}
由 $q=e^{-\lambda}$ 及最大似然估计的不变性，
\[
  \widehat q_{\mathrm{ML}}=e^{-\overline X},\qquad
  \overline x=\frac{42+2\cdot21+3\cdot9+4\cdot4+5\cdot2}{122}
     =\frac{137}{122},
\]
故估计值为 $e^{-137/122}\approx0.3253$。它利用泊松模型中的全部频数，不同于直接使用零值频率 $44/122$。而且
\[
  \mathbb E_\lambda(e^{-\overline X})
      =\exp\!\left[n\lambda(e^{-1/n}-1)\right]
\]
一般不等于 $e^{-\lambda}$，所以不变性不意味着变换后仍无偏。
\end{solve}

\begin{example}
\label{ex:study-ch6-ex44}
教材第 44 题。设总体密度为
\[
  f(x;\theta)=\frac1{\sigma}e^{-(x-\theta)/\sigma}
       \mathbf1_{\{x\ge\theta\}},
  \qquad\theta\in\mathbb R,
\]
其中 $\sigma>0$ 已知。求位置参数 $\theta$ 的矩估计、最大似然估计及无偏修正，并比较修正后的有效性。
\end{example}

\begin{solve}
写成 $X=\theta+Y$，其中 $Y$ 是均值为 $\sigma$ 的指数变量，则 $\mathbb E X=\theta+\sigma$、$\operatorname{Var}X=\sigma^2$，故
\[
  \widehat\theta_{\mathrm M}=\overline X-\sigma,\qquad
  \mathbb E\widehat\theta_{\mathrm M}=\theta,\qquad
  \operatorname{Var}(\widehat\theta_{\mathrm M})=\frac{\sigma^2}{n}.
\]
固定样本后的似然为
\[
  L(\theta)=\sigma^{-n}
     \exp\!\left[-\frac{\sum x_i-n\theta}{\sigma}\right]
     \mathbf1_{\{\theta\le x_{(1)}\}}.
\]
它在可行域内随 $\theta$ 严格递增，故 $\widehat\theta_{\mathrm{ML}}=X_{(1)}$。由
\[
  \mathbb P(X_{(1)}-\theta>t)
      =\prod_{i=1}^n\mathbb P(Y_i>t)
      =e^{-nt/\sigma},\qquad t\ge0,
\]
可知 $X_{(1)}-\theta$ 的均值和方差分别为 $\sigma/n,\sigma^2/n^2$。因此无偏修正为 $\widetilde\theta=X_{(1)}-\sigma/n$，方差为 $\sigma^2/n^2$，当 $n>1$ 时比 $\overline X-\sigma$ 更有效。该模型仍属于支撑依赖参数的情形。
\end{solve}

\begin{example}
\label{ex:study-ch6-ex54}
教材第 54 题。设 $X_i\sim\operatorname{Exp}(\lambda)$，采用式~\eqref{eq:study-exponential-rate} 的速率参数化。$\overline X$ 对 $1/\lambda$ 无偏，$1/\overline X$ 是否对 $\lambda$ 无偏？
\end{example}

\begin{solve}
记 $W=\sum_{i=1}^nX_i$，则其密度为
\[
  f_W(w)=\frac{\lambda^n}{\Gamma(n)}
              w^{n-1}e^{-\lambda w}\mathbf1_{\{w>0\}},
\]
即形状为 $n$、速率为 $\lambda$ 的伽马分布。对 $n>1$，直接积分得
\[
  \mathbb E(W^{-1})
    =\frac{\lambda^n}{\Gamma(n)}
       \int_0^\infty w^{n-2}e^{-\lambda w}\,dw
    =\lambda\frac{\Gamma(n-1)}{\Gamma(n)}
    =\frac{\lambda}{n-1}.
\]
因此 $\mathbb E(1/\overline X)=n\lambda/(n-1)>\lambda$，其无偏修正为 $(n-1)/W$。$n=1$ 时倒数的期望发散，不能使用这个修正公式宣称无偏。类似地，当 $n>2$ 时，
\[
  \mathbb E(W^{-2})=\frac{\lambda^2}{(n-1)(n-2)},\qquad
  \operatorname{Var}\!\left(\frac{n-1}{W}\right)
      =\frac{\lambda^2}{n-2}.
\]
在 $n=2$ 时该无偏估计虽有有限期望，方差却无限，进一步说明讨论无偏性与讨论精度所需的矩条件不同。
\end{solve}

\begin{example}
\label{ex:study-ch6-ex56}
教材第 56 题。总体取值 $0,1,2,3$ 的概率依次为 $\theta/2,\theta,3\theta/2,1-3\theta$，取 $\Theta=[0,1/3]$。观测样本为 $(0,3,1,1,0,2,0,0,3,0)$。求矩方程的解与 MLE，讨论无偏性、有效性及参数约束。
\end{example}

\begin{solve}
由 $\mathbb E_\theta X=3-5\theta$，矩方程的形式解为
\[
  T_{\mathrm M}=\frac{3-\overline X}{5}.
\]
记 $N_3=\sum_{i=1}^n\mathbf1_{\{X_i=3\}}$，则忽略与参数无关的正因子后，似然为 $\theta^{n-N_3}(1-3\theta)^{N_3}$。求导并检查两端，可得在闭参数空间上的 MLE
\[
  \widehat\theta_{\mathrm{ML}}=\frac{n-N_3}{3n}.
\]
给定数据中 $n=10$、$\overline x=1$、$N_3=2$，故形式矩解为 $0.4$，MLE 为 $4/15$。\emph[caution]{$0.4>1/3$，所以这组数据的矩方程在参数空间内没有解。}$T_{\mathrm M}$ 仍可作为一个实值点估计量研究，但不能把这个形式解说成满足参数约束的矩估计值。截断到参数空间会改变估计量的性质。

对于未截断的统计量，$\mathbb E_\theta T_{\mathrm M}=\theta$；又因 $n-N_3\sim\operatorname{Bin}(n,3\theta)$，MLE 也无偏。由
\[
  \operatorname{Var}_\theta X
    =(9-20\theta)-(3-5\theta)^2
    =5\theta(2-5\theta)
\]
得到
\[
  \operatorname{Var}_\theta(T_{\mathrm M})
      =\frac{\theta(2-5\theta)}{5n},\qquad
  \operatorname{Var}_\theta(\widehat\theta_{\mathrm{ML}})
      =\frac{\theta(1-3\theta)}{3n}.
\]
两者之差为 $\theta/(15n)$，在 $\theta>0$ 时严格为正，所以 MLE 比未截断的形式矩估计更有效。此例同时说明，无偏性本身不保证估计值始终落在参数空间内。
\end{solve}

\begin{example}
\label{ex:study-ch6-ex58}
教材第 58 题。设 $T_n$ 是 $\theta$ 的无偏估计，$0<\operatorname{Var}_\theta(T_n)<\infty$ 且 $\operatorname{Var}_\theta(T_n)\to0$。证明 $T_n^2$ 是 $\theta^2$ 的相合估计，却不是无偏估计。
\end{example}

\begin{solve}
由无偏性与切比雪夫不等式，
\[
  \mathbb P_\theta(|T_n-\theta|>\varepsilon)
     \le\frac{\operatorname{Var}_\theta(T_n)}{\varepsilon^2}
     \longrightarrow0.
\]
平方函数连续，故 $T_n^2\xrightarrow{\mathbb P}\theta^2$。另一方面，
\[
  \mathbb E_\theta(T_n^2)
      =\theta^2+\operatorname{Var}_\theta(T_n)>\theta^2,
\]
所以每个有限 $n$ 下都不是无偏估计，但其偏差趋于零，因而是渐近无偏的。这里能得到渐近无偏，来自给定的方差条件，并不是依概率收敛本身保证期望收敛。
\end{solve}

\chapter{区间估计}
\label{chap:study-interval-estimation}

点估计用一个数概括未知参数，却不能单独说明抽样误差的大小。本章在点估计的基础上，用由样本决定的区间表达估计的不确定性：先规定重复抽样时的可靠程度，再研究区间的长度和计算方法。学习的主线是\emph{抽样分布或其近似 $\longrightarrow$ 概率不等式 $\longrightarrow$ 对参数反解 $\longrightarrow$ 检查覆盖率与适用条件}。以下按教材第七章第 240--265 页的结构整理；涉及概率的区间端点使用大写样本，代入观测数据后的端点使用小写样本。

\section{基本概念}
\label{sec:study-interval-concepts}
% 教材第 240--242 页。

设 $X_1,\ldots,X_n$ 是来自某总体的独立同分布样本，$\theta\in\Theta\subseteq\mathbb R$ 是要估计的固定未知参数。希望寻找两个统计量 $L(X_1,\ldots,X_n)$、$U(X_1,\ldots,X_n)$，使 $L\le U$，并用随机区间 $[L,U]$ 估计 $\theta$。端点可以依赖样本、已知常数和预先指定的置信水平，但不能包含仍然未知的参数。例如，在总体均值估计中，$\bar X\pm1.96\sigma/\sqrt n$ 只有在 $\sigma$ 已知时才是可直接计算的区间。

\begin{definition}
\label{def:study-confidence-interval}
给定 $\alpha\in(0,1)$。若对每个允许的参数值都有
\begin{equation}
  \mathbb P_\theta\{L(X_1,\ldots,X_n)\le\theta
                         \le U(X_1,\ldots,X_n)\}\ge1-\alpha,
  \label{eq:study-confidence-coverage}
\end{equation}
则称 $[L,U]$ 为 $\theta$ 的\emph{置信水平（confidence level）}至少为 $1-\alpha$ 的\emph{置信区间（confidence interval）}。左边称为该方法在参数 $\theta$ 处的\emph{覆盖率（coverage probability）}。若对所有 $\theta$ 都恰好等于 $1-\alpha$，便得到教材定义中的等式形式；一般地，覆盖率在参数空间上的下确界称为\emph{置信系数（confidence coefficient）}。
\end{definition}

下标 $\theta$ 表示按真实参数对应的抽样分布计算概率，并不是把 $\theta$ 当作随机变量。若模型中还有其他未知参数，式~\eqref{eq:study-confidence-coverage} 的要求必须对这些参数的所有允许取值同时成立。例如，正态总体中估计 $\mu$ 且 $\sigma^2$ 未知时，覆盖率应在每个 $\mu\in\mathbb R$、$\sigma^2>0$ 下得到保证。

\emph[caution]{频率学派置信区间的随机性来自样本和区间，参数是固定的。}观测前可以说“该随机区间以 $95\%$ 的概率覆盖真实参数”；观测后得到具体数值区间，它对固定真值只有包含或不包含两种情况，不能由置信区间的定义推出“参数以 $95\%$ 的概率落在这个已经算出的区间内”。后一句若作概率陈述，需要另外建立参数的概率模型。

重复抽样有助于理解这一点。假设某方法在真参数处的覆盖率恰为 $0.95$，独立抽取 $B$ 组样本，每组都用同一规则计算区间，则覆盖真值的区间个数服从二项分布 $\operatorname{Bin}(B,0.95)$；覆盖比例依概率趋于 $0.95$。当 $B=100$ 时，覆盖个数的均值为 $95$、标准差为 $\sqrt{100\times0.95\times0.05}\approx2.18$，所以一次模拟得到 $93$ 个或 $97$ 个覆盖区间并不反常，也不能要求每 $100$ 次恰有 $95$ 次覆盖。

可靠性与精度需要同时考虑。可靠性由覆盖率表达，精度可由随机长度 $U-L$、其期望 $\mathbb E_\theta(U-L)$ 或某种预先规定的误差界限表达。在同一种构造中，提高置信水平通常要扩大区间，增加样本量则通常能缩短区间；但比较不同方法时，还要核对它们是否达到相同的覆盖要求，不能只看哪个区间更短。对称区间中的\emph{误差界限（margin of error）}是半宽度，不等于观测后已经知道的真实估计误差。

\begin{example}
\label{ex:study-confidence-interpretation}
用同一组数据估计某总体均值，两个方法分别给出 $[150,170]$ 和 $[159,161]$。能否仅凭第二个区间较短，就认为它更好？
\end{example}
\begin{solve}
不能。区间长度只表达精度的一方面，还必须说明抽样假设、置信水平及覆盖性质。如果第二个方法没有可靠的覆盖保证，其较短区间可能只是低估了抽样波动。相反，即使两个方法都有相同的名义置信水平，有限样本实际覆盖率也可能不同，仍需区分精确校准与近似校准。
\end{solve}

\section{枢轴变量法}
\label{sec:study-pivotal-quantity}
% 教材第 243--248 页。

\subsection{构造原理与分位数约定}
\label{subsec:study-pivot-construction}

\begin{definition}
\label{def:study-pivotal-quantity}
由样本和待估参数构成的随机量 $Q(X_1,\ldots,X_n;\theta)$，若其分布完全已知，且不依赖任何未知参数，就称为\emph{枢轴变量或枢轴量（pivotal quantity）}。它可以显式含有 $\theta$，因此通常不是统计量；关键是其\emph{分布}不含未知参数。
\end{definition}

若能选取已知常数 $a<b$，使
\[
  \mathbb P_\theta\{a\le Q(X;\theta)\le b\}=1-\alpha,
\]
并将事件 $\{a\le Q(X;\theta)\le b\}$ 等价地改写成 $\{L(X)\le\theta\le U(X)\}$，就得到了精确置信区间。反解时必须同时考虑参数空间、分母的正负和函数的单调性；如果反解得到不相连的集合，则首先得到的是置信集合，不能未经说明直接写成一个区间。寻找良好点估计常能启发枢轴量的构造，但“无偏”或“最大似然”本身并不保证相应的区间有正确覆盖率。

本章沿用教材的\emph{上尾分位数}约定。对这里使用的连续严格递增分布，$q(a)$ 满足 $\mathbb P\{Q>q(a)\}=a$。具体记
\begin{equation}
  \begin{aligned}
    \mathbb P\{Z>u_a\}&=a, && Z\sim N(0,1),\\
    \mathbb P\{T>t_\nu(a)\}&=a, && T\sim t_\nu,\\
    \mathbb P\{V>\chi_\nu^2(a)\}&=a, && V\sim\chi_\nu^2,\\
    \mathbb P\{W>F_{r,s}(a)\}&=a, && W\sim F_{r,s}.
  \end{aligned}
  \label{eq:study-upper-quantiles}
\end{equation}
因此 $a$ 越小，上尾分位数越大。正态和 $t$ 分布关于零对称，故 $u_{1-a}=-u_a$、$t_\nu(1-a)=-t_\nu(a)$；卡方和 $F$ 分布不能套用这种变号关系。常用数值为 $u_{0.025}\approx1.9600$、$u_{0.05}\approx1.6449$。选择上下两尾各留 $\alpha/2$ 的概率称为等尾构造，它易于计算，但对不对称分布不必给出最短区间。

\subsection{正态总体均值与样本量}
\label{subsec:study-normal-mean-intervals}

设 $X_1,\ldots,X_n$ 独立同分布于 $N(\mu,\sigma^2)$，其中 $\sigma>0$。记
\begin{equation}
  \bar X=\frac1n\sum_{i=1}^nX_i,\qquad
  S^2=\frac1{n-1}\sum_{i=1}^n(X_i-\bar X)^2\quad(n\ge2).
  \label{eq:study-interval-sample-statistics}
\end{equation}
$S$ 度量单次观测的波动，$S/\sqrt n$ 则估计样本均值的标准差，后者称为\emph{标准误（standard error）}。总体标准差与均值的标准误不可混用。

\begin{theorem}
\label{thm:study-normal-mean-intervals}
在上述正态简单随机样本模型下，$\mu$ 的置信系数为 $1-\alpha$ 的等尾区间为
\begin{equation}
  \left[\bar X-u_{\alpha/2}\frac\sigma{\sqrt n},\,
        \bar X+u_{\alpha/2}\frac\sigma{\sqrt n}\right]
  \quad(\sigma^2\text{ 已知}),
  \label{eq:study-normal-known-variance-ci}
\end{equation}
以及
\begin{equation}
  \left[\bar X-t_{n-1}(\alpha/2)\frac S{\sqrt n},\,
        \bar X+t_{n-1}(\alpha/2)\frac S{\sqrt n}\right]
  \quad(\sigma^2\text{ 未知},\ n\ge2).
  \label{eq:study-normal-unknown-variance-ci}
\end{equation}
\end{theorem}
\begin{proof}
方差已知时，$Z=\sqrt n(\bar X-\mu)/\sigma\sim N(0,1)$。将 $-u_{\alpha/2}\le Z\le u_{\alpha/2}$ 对 $\mu$ 反解，即得式~\eqref{eq:study-normal-known-variance-ci}。方差未知时，正态样本的抽样分布给出
\[
  V=\frac{(n-1)S^2}{\sigma^2}\sim\chi_{n-1}^2,
  \qquad Z\text{ 与 }V\text{ 独立}.
\]
所以
\[
  T=\frac{Z}{\sqrt{V/(n-1)}}
   =\frac{\sqrt n(\bar X-\mu)}S\sim t_{n-1}.
\]
再利用 $t$ 分布的对称性反解，得到式~\eqref{eq:study-normal-unknown-variance-ci}。这里的 $t$ 分布结论依赖总体正态性，不能只把已知方差公式中的 $\sigma$ 换成 $S$，同时仍声称标准化后精确服从正态分布。
\end{proof}

\begin{example}
\label{ex:study-normal-mean-residents}
教材例 7.3：假设某项居民月支出服从正态分布，独立抽取 $20$ 户，得到 $\bar x=160$ 元、$s=20$ 元，求总体平均支出 $\mu$ 的 $95\%$ 置信区间。
\end{example}
\begin{solve}
总体方差未知，使用式~\eqref{eq:study-normal-unknown-variance-ci}，自由度为 $19$。由 $t_{19}(0.025)\approx2.0930$，得
\[
  160\pm2.0930\frac{20}{\sqrt{20}}
  \approx[150.64,169.36]\text{ 元}.
\]
区间估计的对象是总体平均支出，并不表示 $95\%$ 的居民支出在该范围内，也不表示下一户的支出有 $95\%$ 的概率在其中。
\end{solve}

在正态总体方差已知时，若事先要求半宽度不超过 $d_0>0$，则
\begin{equation}
  u_{\alpha/2}\frac\sigma{\sqrt n}\le d_0
  \quad\Longleftrightarrow\quad
  n\ge\left(\frac{u_{\alpha/2}\sigma}{d_0}\right)^2.
  \label{eq:study-normal-ci-sample-size}
\end{equation}
因此最小正整数样本量为右端向上取整，必要时至少取 $1$。教材例 7.4 中，$90\%$ 置信水平、$\sigma=200$ 元、$d_0=100$ 元给出 $n\ge10.82$，至少调查 $11$ 人。若 $\sigma$ 未知，可以用先导样本或可信经验作规划，但该规划的可靠性取决于所用标准差估计；仅知道一个经验取值范围，不能无条件推出严格的样本量保证。方差未知的 $t$ 区间宽度还依赖随机的 $S$，因而需要区分“规划的预期宽度”与“每次抽样都满足的宽度”。

\subsection{正态总体方差}
\label{subsec:study-normal-variance-intervals}

\begin{theorem}
\label{thm:study-normal-variance-interval}
设 $X_1,\ldots,X_n$ 独立同分布于 $N(\mu,\sigma^2)$，均值未知且 $n\ge2$，则 $\sigma^2$ 的置信系数为 $1-\alpha$ 的等尾区间为
\begin{equation}
  \left[
    \frac{(n-1)S^2}{\chi_{n-1}^2(\alpha/2)},\,
    \frac{(n-1)S^2}{\chi_{n-1}^2(1-\alpha/2)}
  \right].
  \label{eq:study-normal-variance-ci}
\end{equation}
\end{theorem}
\begin{proof}
由正态样本方差的抽样分布，
\[
  \mathbb P_{\mu,\sigma^2}\left\{
    \chi_{n-1}^2(1-\alpha/2)
    \le\frac{(n-1)S^2}{\sigma^2}
    \le\chi_{n-1}^2(\alpha/2)
  \right\}=1-\alpha.
\]
两个分位数都为正，对 $\sigma^2$ 反解时不等号方向随取倒数改变，便得到式~\eqref{eq:study-normal-variance-ci}。
\end{proof}

上下端点分别除以较大和较小的卡方分位数，这是核对公式方向的直接办法。该区间通常不以 $S^2$ 为中心；对其两个端点开平方，就得到总体标准差 $\sigma$ 的同一置信水平区间，因为平方根在正数上严格递增。若总体均值 $\mu_0$ 已知，则应改用
\begin{equation}
  Q_0=\frac{\sum_{i=1}^n(X_i-\mu_0)^2}{\sigma^2}
      \sim\chi_n^2,
  \label{eq:study-normal-known-mean-variance-pivot}
\end{equation}
相应地以 $\sum(X_i-\mu_0)^2$ 替换分子、以 $n$ 替换自由度。\emph[caution]{已知均值与用样本均值估计均值，消耗的自由度不同。}

\subsection{两个独立正态总体的均值差与方差比}
\label{subsec:study-two-normal-intervals}

设 $X_1,\ldots,X_m$ 来自 $N(\mu_1,\sigma_1^2)$，$Y_1,\ldots,Y_n$ 来自 $N(\mu_2,\sigma_2^2)$，两个样本相互独立。分别记样本均值为 $\bar X,\bar Y$，样本方差为 $S_1^2,S_2^2$，并统一估计
\[
  \delta=\mu_2-\mu_1,\qquad D=\bar Y-\bar X.
\]
若总体方差都已知，由独立正态变量之和的分布，
\[
  \frac{D-\delta}{\sqrt{\sigma_1^2/m+\sigma_2^2/n}}
  \sim N(0,1),
\]
所以精确区间为
\begin{equation}
  D\pm u_{\alpha/2}\sqrt{\frac{\sigma_1^2}{m}
                              +\frac{\sigma_2^2}{n}}.
  \label{eq:study-two-normal-known-ci}
\end{equation}
这里及下文的 $A\pm d$ 均表示闭区间 $[A-d,A+d]$。

若方差未知但可假设 $\sigma_1^2=\sigma_2^2=\sigma^2$，定义\emph{合并方差（pooled variance）}
\begin{equation}
  S_p^2=\frac{(m-1)S_1^2+(n-1)S_2^2}{m+n-2},
  \qquad m,n\ge2.
  \label{eq:study-interval-pooled-variance}
\end{equation}
两组样本方差对应的独立卡方量相加，给出 $(m+n-2)S_p^2/\sigma^2\sim\chi_{m+n-2}^2$；它又与两个样本均值独立。因此
\[
  \frac{D-\delta}{S_p\sqrt{1/m+1/n}}\sim t_{m+n-2},
\]
得到精确区间
\begin{equation}
  D\pm t_{m+n-2}(\alpha/2)S_p\sqrt{\frac1m+\frac1n}.
  \label{eq:study-two-normal-pooled-ci}
\end{equation}
合并的依据是两个总体有共同方差，而不是单纯为了增加自由度；不能因为两个样本方差数值接近就把共同方差假设当作已经证实。

\begin{example}
\label{ex:study-two-production-lines}
教材例 7.6：两条生产线满足上述独立正态、共同方差假设。数据为 $m=5$、$n=7$，$\bar x=62$、$\bar y=48$，$s_1^2=25$、$s_2^2=16$，求 $\mu_2-\mu_1$ 的 $90\%$ 置信区间。
\end{example}
\begin{solve}
由式~\eqref{eq:study-interval-pooled-variance}，$s_p^2=(4\times25+6\times16)/10=19.6$。利用 $t_{10}(0.05)\approx1.8125$，得
\[
  -14\pm1.8125\sqrt{19.6}\sqrt{\frac15+\frac17}
  \approx[-18.70,-9.30].
\]
整个区间为负，表明在模型及该区间方法下，数据支持第二条生产线的总体平均产量较低；这仍然需要两条生产线的抽样具有可比性。
\end{solve}

若不假设共同方差，式~\eqref{eq:study-two-normal-pooled-ci} 不再有上述精确覆盖保证。作为教材所述近似解的补充，常用\emph{韦尔奇近似（Welch approximation）}为
\begin{equation}
  D\pm t_{\widehat\nu}(\alpha/2)
             \sqrt{\frac{S_1^2}{m}+\frac{S_2^2}{n}},\qquad
  \widehat\nu=
  \frac{(S_1^2/m+S_2^2/n)^2}
       {(S_1^2/m)^2/(m-1)+(S_2^2/n)^2/(n-1)}.
  \label{eq:study-welch-ci}
\end{equation}
这里自由度可以不是整数，所得区间是近似区间；独立性仍然必需。若数据按同一对象的前后测量配成对，应先构造每一对的差值，再对差值总体作单样本推断，不能直接使用独立两样本公式。

方差比 $r=\sigma_1^2/\sigma_2^2$ 则有另一个精确枢轴量：
\[
  \frac{S_1^2/\sigma_1^2}{S_2^2/\sigma_2^2}
   =\frac{S_1^2/S_2^2}{r}\sim F_{m-1,n-1}.
\]
故 $r$ 的等尾区间为
\begin{equation}
  \left[
    \frac{S_1^2/S_2^2}{F_{m-1,n-1}(\alpha/2)},\,
    \frac{S_1^2/S_2^2}{F_{m-1,n-1}(1-\alpha/2)}
  \right].
  \label{eq:study-normal-variance-ratio-ci}
\end{equation}
对 $a\in(0,1)$，利用 $F_{m-1,n-1}(1-a)=1/F_{n-1,m-1}(a)$，也可将其写为
$[(S_1^2/S_2^2)F_{n-1,m-1}(1-\alpha/2),\,(S_1^2/S_2^2)F_{n-1,m-1}(\alpha/2)]$。交换两组自由度时必须同时取倒数，否则端点就会出错。

\section{大样本方法}
\label{sec:study-large-sample-intervals}
% 教材第 248--253 页。

当精确抽样分布难以得到时，可以用极限定理构造\emph{渐近置信区间（asymptotic confidence interval）}。若一列区间满足
\[
  \mathbb P_\theta\{L_n\le\theta\le U_n\}\longrightarrow1-\alpha
  \qquad(n\to\infty),
\]
则有限样本覆盖率只是近似达到名义水平。这里若只证明了每个固定参数处的收敛，不能自动推出在整个参数空间上一致准确；参数接近边界时尤其需要谨慎。以下的 $\xrightarrow{d}$ 表示依分布收敛，$\xrightarrow{\mathbb P}$ 表示依概率收敛。

\subsection{比例的得分区间与瓦尔德区间}
\label{subsec:study-proportion-intervals}

设 $X_i\in\{0,1\}$ 独立同分布，$\mathbb P_p(X_i=1)=p$，其中 $0<p<1$；令 $K=\sum_{i=1}^nX_i\sim\operatorname{Bin}(n,p)$、$\widehat p=K/n$。由中心极限定理，
\begin{equation}
  \frac{\sqrt n(\widehat p-p)}{\sqrt{p(1-p)}}
  \xrightarrow{d}N(0,1).
  \label{eq:study-proportion-clt}
\end{equation}
令 $z=u_{\alpha/2}$。保留分母中的未知参数，先考虑近似概率为 $1-\alpha$ 的事件
\[
  \left|\frac{\widehat p-p}{\sqrt{p(1-p)/n}}\right|\le z.
\]
在 $0<p<1$ 下平方并移项，得到
\[
  (n+z^2)p^2-(2n\widehat p+z^2)p+n\widehat p^{\,2}\le0.
\]
解这个二次不等式，便得到\emph{得分区间（score interval），亦称威尔逊区间（Wilson interval）}：
\begin{equation}
  \left[
    \frac{\widehat p+z^2/(2n)
          -z\sqrt{\widehat p(1-\widehat p)/n+z^2/(4n^2)}}{1+z^2/n},\,
    \frac{\widehat p+z^2/(2n)
          +z\sqrt{\widehat p(1-\widehat p)/n+z^2/(4n^2)}}{1+z^2/n}
  \right].
  \label{eq:study-wilson-ci}
\end{equation}
端点在 $[0,1]$ 内，中心一般不是 $\widehat p$；若 $K=0$，区间仍为 $[0,z^2/(n+z^2)]$，不会退化为单点。该公式仍由正态近似导出，不能因为端点形式较好就称为精确区间。

另一种做法是先用 $\widehat p(1-\widehat p)$ 替换 $p(1-p)$。对固定的内部参数 $p\in(0,1)$，由大数定律和斯卢茨基定理得到\emph{瓦尔德区间（Wald interval）}
\begin{equation}
  \widehat p\pm u_{\alpha/2}
       \sqrt{\frac{\widehat p(1-\widehat p)}n}.
  \label{eq:study-proportion-wald-ci}
\end{equation}
它形式简单，却可能越出 $[0,1]$，并在 $K=0$ 或 $K=n$ 时退化。把端点截到 $[0,1]$ 不会修复其覆盖不足的问题，因为真参数本来就在该参数空间中。得分区间通常能改善边界附近的表现，但两者有限样本覆盖率都应结合 $np$、$n(1-p)$ 等实际信息判断；“$n>100$”或“两个观测频数都大于 $10$”只能作为经验检查，不能代替覆盖率分析。

\begin{example}
\label{ex:study-phishing-proportion}
沿用教材例 7.7 的计数，在符合上述二项抽样模型的 $320$ 次观测中，有 $157$ 次发生目标事件，求事件概率 $p$ 的近似 $95\%$ 置信区间。
\end{example}
\begin{solve}
取 $\widehat p=157/320=0.490625$、$z=u_{0.025}$，代入式~\eqref{eq:study-wilson-ci}，得到约
\[
  [0.4363,0.5452].
\]
计算时保留比例的完整数值，最后才舍入端点。这是在给定抽样与试验条件下对 $p$ 的估计，不能只凭这组计数推出改变另一项条件的因果效果。
\end{solve}

若关心百分数 $100p\%$，直接把比例区间的两个端点乘以 $100$ 后附上百分号即可，不应在含 $p(1-p)$ 的公式中仅把 $p$ 替换成百分数而保留其他项不变。对于稀少事件，可以补充使用二项分布尾概率的精确反演：当观测计数 $0<k<n$ 时，令 $p_L,p_U$ 分别满足
\begin{equation}
  \mathbb P_{p_L}\{K\ge k\}=\alpha/2,\qquad
  \mathbb P_{p_U}\{K\le k\}=\alpha/2.
  \label{eq:study-binomial-exact-inversion}
\end{equation}
取 $k=0$ 时的下端点为 $0$，$k=n$ 时的上端点为 $1$。由二项分布尾概率的单调性可得覆盖率至少为 $1-\alpha$ 的区间；由于计数离散，其覆盖率通常不处处恰等于名义水平。这也说明离散模型虽未必方便使用连续枢轴量，仍可构造精确覆盖的置信区间。

\subsection{比例估计的样本量规划}
\label{subsec:study-proportion-sample-size}

令目标总宽度为 $w\in(0,1)$，注意它是半宽度的两倍。若采用瓦尔德区间作规划，并以先验规划值 $p_0$ 替代尚未观测到的比例，则要求
\begin{equation}
  n\ge\frac{4u_{\alpha/2}^2p_0(1-p_0)}{w^2}.
  \label{eq:study-proportion-wald-size}
\end{equation}
这里的“先验规划值”只表示在本次调查前已有的参考信息，不意味着已经采用贝叶斯推断。没有可靠参考值时，由 $p(1-p)\le1/4$ 得保守规划 $n\ge u_{\alpha/2}^2/w^2$。

得分区间的总宽度为
\[
  W_n(\widehat p)
  =\frac{2z\sqrt{n\widehat p(1-\widehat p)+z^2/4}}{n+z^2},
  \qquad z=u_{\alpha/2}.
\]
将 $\widehat p$ 替换成规划值 $p_0$，记 $r_0=p_0(1-p_0)$，由 $W_n(p_0)\le w$ 对 $n$ 解二次不等式，得到
\begin{equation}
  n\ge\frac{z^2}{w^2}
       \left\{2r_0-w^2+
        \sqrt{4r_0^2+(1-4r_0)w^2}\right\}.
  \label{eq:study-proportion-wilson-size}
\end{equation}
由于宽度在 $\widehat p=1/2$ 处达到最大，直接要求所有观测比例对应的区间宽度都不超过 $w$，可用
\begin{equation}
  \sup_{0\le q\le1}W_n(q)=\frac z{\sqrt{n+z^2}}\le w,
  \qquad n\ge z^2\left(\frac1{w^2}-1\right).
  \label{eq:study-proportion-wilson-conservative-size}
\end{equation}
教材例 7.8 中，$1-\alpha=0.95$、$w=0.08$，得分区间所需样本量向上取整为 $597$，瓦尔德区间的保守规划为 $601$。这里“保守”指宽度规划采用最不利比例，并不把原有的近似覆盖变成有限样本精确保证。

\subsection{一般总体均值及具体模型}
\label{subsec:study-general-mean-intervals}

设 $X_1,\ldots,X_n$ 独立同分布，$\mathbb E(X_i)=\mu$、$0<\operatorname{Var}(X_i)=\sigma^2<\infty$。由中心极限定理及样本方差的相合性，
\begin{equation}
  \frac{\sqrt n(\bar X-\mu)}\sigma\xrightarrow{d}N(0,1),
  \qquad S^2\xrightarrow{\mathbb P}\sigma^2,
  \qquad
  \frac{\sqrt n(\bar X-\mu)}S\xrightarrow{d}N(0,1).
  \label{eq:study-studentized-mean-clt}
\end{equation}
最后一步由斯卢茨基定理得到，在 $S>0$ 上使用该比值；在上述条件下 $\mathbb P(S>0)\to1$。因此均值的渐近区间为
\begin{equation}
  \bar X\pm u_{\alpha/2}\frac S{\sqrt n}.
  \label{eq:study-general-mean-ci}
\end{equation}
它不要求总体正态，却要求抽样模型和矩条件成立。总体偏斜、重尾或存在强依赖时，某个固定样本量未必足以保证正态近似；尤其不能把“$n>30$”当作对一切总体都有效的数学定理。若总体为正态，式~\eqref{eq:study-normal-unknown-variance-ci} 已提供精确 $t$ 区间，无须先退回正态近似。

两个独立总体的均值差也有类似构造。若两总体方差有限且为正，$m,n\to\infty$ 且 $m/n$ 趋于某个正的有限常数，则用各自样本方差替换总体方差，得到
\[
  (\bar Y-\bar X)\pm u_{\alpha/2}
             \sqrt{\frac{S_1^2}{m}+\frac{S_2^2}{n}}
\]
这一渐近区间，不要求两个总体方差相同。该结论与小样本正态共同方差模型下的精确 $t$ 区间具有不同的假设与覆盖含义。

\begin{example}
\label{ex:study-poisson-mean-interval}
教材例 7.9：$X_1,\ldots,X_n$ 独立同分布于泊松分布 $\operatorname{Pois}(\lambda)$，$\lambda>0$，求 $\lambda$ 的渐近区间。
\end{example}
\begin{solve}
由于 $\mathbb E(X_i)=\operatorname{Var}(X_i)=\lambda$，有 $\sqrt n(\bar X-\lambda)/\sqrt\lambda\xrightarrow{d}N(0,1)$。令 $z=u_{\alpha/2}$，对事件 $n(\bar X-\lambda)^2\le z^2\lambda$ 反解，得到
\begin{equation}
  \left[
    \bar X+\frac{z^2}{2n}
      -\frac z{\sqrt n}\sqrt{\bar X+\frac{z^2}{4n}},\,
    \bar X+\frac{z^2}{2n}
      +\frac z{\sqrt n}\sqrt{\bar X+\frac{z^2}{4n}}
  \right].
  \label{eq:study-poisson-score-ci}
\end{equation}
若改用一般均值公式，则得到 $\bar X\pm zS/\sqrt n$；也可利用 $\lambda$ 的相合估计 $\bar X$，得到 $\bar X\pm z\sqrt{\bar X/n}$。这些区间的近似步骤不同，有限样本端点也不同。教材给出的 $n=20$、$\bar x=2.6$、$s=1.635$ 下，式~\eqref{eq:study-poisson-score-ci} 约为 $[1.983,3.409]$，一般均值公式约为 $[1.883,3.317]$。单个算例中端点相近，并不能证明两个方法对所有参数都有相近的实际覆盖率。
\end{solve}

\begin{example}
\label{ex:study-exponential-mean-interval}
教材例 7.10：$X_1,\ldots,X_n$ 独立同分布于率参数为 $\lambda>0$ 的指数分布，其密度为 $f(x)=\lambda e^{-\lambda x}$（$x>0$），求均值 $\theta=1/\lambda$ 的置信区间，并与大样本方法比较。
\end{example}
\begin{solve}
先利用模型的精确分布。$2\lambda X_i$ 的密度为 $e^{-v/2}/2$（$v>0$），即 $\chi_2^2$ 密度；各项独立，卡方分布的可加性给出
\[
  Q=2\lambda\sum_{i=1}^nX_i
   =\frac{2n\bar X}\theta\sim\chi_{2n}^2.
\]
因此 $\theta$ 的精确等尾区间为
\begin{equation}
  \left[
    \frac{2n\bar X}{\chi_{2n}^2(\alpha/2)},\,
    \frac{2n\bar X}{\chi_{2n}^2(1-\alpha/2)}
  \right].
  \label{eq:study-exponential-mean-exact-ci}
\end{equation}
另一方面，指数分布满足 $\operatorname{Var}(X_i)=\theta^2$，故 $\sqrt n(\bar X-\theta)/\theta\xrightarrow{d}N(0,1)$。在 $\sqrt n>z=u_{\alpha/2}$ 时，对 $|\sqrt n(\bar X-\theta)/\theta|\le z$ 反解，得到近似区间
\begin{equation}
  \left[\frac{\bar X}{1+z/\sqrt n},\,
        \frac{\bar X}{1-z/\sqrt n}\right].
  \label{eq:study-exponential-mean-asymptotic-ci}
\end{equation}
若 $\sqrt n\le z$，第二个分母不再为正，不能继续照写该有限上端点，而应重新求解不等式。取 $n=20$、$\bar x=5$、$\alpha=0.05$，精确区间约为 $[3.370,8.186]$，近似区间约为 $[3.476,8.901]$；在指数模型成立时，前者具有有限样本精确覆盖保证。
\end{solve}

\section{自助法置信区间}
\label{sec:study-bootstrap-intervals}
% 教材第 253--257 页。

\emph{自助法（bootstrap）}用从现有样本重抽样所产生的波动近似原统计量的抽样波动。它的基本想法是：未知总体分布 $F$ 无法直接重复抽样，但可以先用经验分布
\[
  F_n(x)=\frac1n\sum_{i=1}^n\mathbf1_{\{X_i\le x\}}
\]
替代 $F$。给定观测值 $x_1,\ldots,x_n$ 后，$\mathbf1_{\{x_i\le x\}}$ 在条件满足时为 $1$、否则为 $0$；从 $F_n$ 抽样，就是把每个观测位置赋予概率 $1/n$，有放回地独立抽取 $n$ 次。即使若干观测值相同，它们对应的观测位置仍各有 $1/n$ 的抽取概率。

非参数自助法不必指定正态等参数分布，却仍有统计条件。这里默认原始数据来自独立同分布抽样，经验分布能合理代表总体，且所估计的统计特征适合这种近似。样本很小时，重抽样不能补回未观察到的尾部或少数群体；对于样本极值、参数边界、非光滑或退化统计量，普通自助法还可能不一致。对时间序列或其他相关数据，逐个观测独立重抽样会破坏相关结构，需要相应的重抽样设计。\emph[caution]{不要求某个参数分布，不等于没有适用条件，也不等于小样本下自动准确。}

\subsection{基本自助区间}
\label{subsec:study-basic-bootstrap-ci}

记参数 $\theta$ 的原始估计为 $\widehat\theta=t(x_1,\ldots,x_n)$。第 $b$ 次重抽样产生 $x_1^{*(b)},\ldots,x_n^{*(b)}$，并计算
\[
  \widehat\theta^{*(b)}
   =t(x_1^{*(b)},\ldots,x_n^{*(b)}),\qquad b=1,\ldots,B.
\]
各次重抽样条件独立。给定原样本时，一个有序重抽样序列有 $n^n$ 种位置选择；忽略顺序而只记录各观测位置被抽取的次数，则有 $\binom{2n-1}{n}$ 种计数向量，这些计数向量并非等概率。实际通常以有限的 $B$ 次模拟近似整个自助分布，$B$ 增大只减少模拟误差，并不增加原始数据量。

令 $\mathbb P^*$ 表示给定原样本后的重抽样概率，$q_\beta^*$ 表示 $\widehat\theta^*$ 的\emph{下尾} $\beta$ 分位数，例如定义为
\[
  q_\beta^*=\inf\{x:\mathbb P^*(\widehat\theta^*\le x)\ge\beta\}.
\]
实际以 $B$ 个自助估计的经验分位数近似它；这种下尾约定与式~\eqref{eq:study-upper-quantiles} 中的上尾分位数不同。自助法希望用给定数据后的 $\widehat\theta^*-\widehat\theta$ 分布近似原抽样中的 $\widehat\theta-\theta$ 分布。若这一近似成立且相关分位数足够稳定，就希望得到近似覆盖关系
\[
  \mathbb P_\theta\left\{
    q_{\alpha/2}^*-\widehat\theta
    \le\widehat\theta-\theta
    \le q_{1-\alpha/2}^*-\widehat\theta
  \right\}\approx1-\alpha.
\]
此处概率针对原始抽样，分位数本身也随原样本变化；它不是由重抽样次数足够大就自动成立的等式。对该事件反解 $\theta$，即可得到下面的区间。

\begin{definition}
\label{def:study-basic-bootstrap-ci}
由上述误差分布近似反解得到的
\begin{equation}
  [\,2\widehat\theta-q_{1-\alpha/2}^*,\,
      2\widehat\theta-q_{\alpha/2}^*\,]
  \label{eq:study-basic-bootstrap-ci}
\end{equation}
称为 $\theta$ 的名义置信水平为 $1-\alpha$ 的\emph{基本自助置信区间（basic bootstrap confidence interval）}。其覆盖性质取决于自助分布对抽样误差分布的近似质量。
\end{definition}

\begin{example}
\label{ex:study-basic-bootstrap-summary}
沿用教材例 7.11 所列的自助计算汇总值：样本均值取 $\widehat\theta=18.53$，$B=1000$ 个自助均值的 $2.5\%$ 和 $97.5\%$ 分位数取 $16.76$、$20.62$。计算基本自助区间。
\end{example}
\begin{solve}
由式~\eqref{eq:study-basic-bootstrap-ci}，区间为
\[
  [\,2\times18.53-20.62,\,2\times18.53-16.76\,]
  =[16.44,20.30].
\]
这里使用教材给出的已舍入汇总值，重点是围绕原估计反射分位数的构造；它不同于直接取 $[16.76,20.62]$ 的\emph{百分位数自助区间（percentile bootstrap interval）}。实际重新抽样时，端点会随随机种子、重复次数和经验分位数的约定略有变化，报告结果时应保留这些计算信息。
\end{solve}

\subsection{学生化自助区间与嵌套重抽样}
\label{subsec:study-bootstrap-t-ci}

若能估计 $\widehat\theta$ 的标准误，记其估计为 $\widehat{\operatorname{se}}$，则可先将估计误差标准化。对第 $b$ 个外层自助样本，计算该样本下的标准误估计 $\widehat{\operatorname{se}}^{*(b)}$，并构造
\begin{equation}
  T^{*(b)}=
  \frac{\widehat\theta^{*(b)}-\widehat\theta}
       {\widehat{\operatorname{se}}^{*(b)}}.
  \label{eq:study-bootstrap-studentized-statistic}
\end{equation}
设其自助分布的下尾分位数为 $c_{\alpha/2}^*,c_{1-\alpha/2}^*$，则\emph{自助 $t$ 置信区间（bootstrap-$t$ confidence interval）}为
\begin{equation}
  [\,\widehat\theta-c_{1-\alpha/2}^*
                         \widehat{\operatorname{se}},\,
      \widehat\theta-c_{\alpha/2}^*
                         \widehat{\operatorname{se}}\,].
  \label{eq:study-bootstrap-t-ci}
\end{equation}
“$t$”表示这种学生化构造，不意味着 $T^*$ 精确服从某个自由度的 $t$ 分布；分位数来自重抽样，而非直接查 $t$ 分布表。其依据是用 $T^*$ 的条件分布近似 $(\widehat\theta-\theta)/\widehat{\operatorname{se}}$ 的抽样分布，再反解参数。

若原始标准误没有便于使用的解析估计，可以用外层自助估计的样本标准差
\begin{equation}
  \widehat{\operatorname{se}}_B
  =\left\{\frac1{B-1}\sum_{b=1}^B
       (\widehat\theta^{*(b)}-\overline\theta^*)^2\right\}^{1/2},
  \qquad
  \overline\theta^*=\frac1B\sum_{b=1}^B\widehat\theta^{*(b)},
  \label{eq:study-bootstrap-standard-error}
\end{equation}
估计它。对每个外层样本的标准误，也可以从该外层样本再有放回抽取 $R$ 个内层样本，以这些内层估计的样本标准差估计 $\widehat{\operatorname{se}}^{*(b)}$。因此算法依次为：计算原估计；生成 $B$ 个外层样本及各自的估计、标准误和 $T^{*(b)}$；取 $T^*$ 的经验分位数；最后代入式~\eqref{eq:study-bootstrap-t-ci}。若使用嵌套抽样，内层统计量计算次数约为 $BR$，明显高于基本区间的 $B$ 次。

若统计量是样本均值，可以直接以 $S/\sqrt n$ 估计原始标准误，以每个外层样本的 $S^{*(b)}/\sqrt n$ 估计其标准误，无须再做内层抽样。若某次重抽样产生零标准误，式~\eqref{eq:study-bootstrap-studentized-statistic} 就无法直接使用，必须检查数据、统计量和算法的适用性，不能悄悄删掉这些重复后仍按原公式解释覆盖率。教材还提及标准正态自助区间、百分位数区间和偏差校正加速区间；它们使用不同的校准方式，不应仅凭同一次模拟中哪个区间最短来选择方法。

\section{置信限}
\label{sec:study-confidence-bounds}
% 教材第 257--258 页。

若只关心参数是否足够小或足够大，可以把覆盖概率集中在一侧，而不必同时控制两个有限端点。

\begin{definition}
\label{def:study-confidence-bounds}
若统计量 $U(X)$ 对所有允许参数值满足
\[
  \mathbb P_\theta\{\theta\le U(X)\}\ge1-\alpha,
\]
则称 $U(X)$ 为 $\theta$ 的置信水平至少为 $1-\alpha$ 的\emph{置信上限（upper confidence bound）}。若统计量 $L(X)$ 满足
\[
  \mathbb P_\theta\{L(X)\le\theta\}\ge1-\alpha,
\]
则称其为相应的\emph{置信下限（lower confidence bound）}。当等式对所有参数值成立时，得到精确的单侧置信系数。
\end{definition}

上限对应置信集合 $(-\infty,U]$，下限对应 $[L,\infty)$，并可与参数空间相交。二者分别具有各自的单侧覆盖保证；如果需要同时使用两个端点形成双侧区间，就必须重新分配两尾概率，不能把两个 $95\%$ 单侧界直接称为一个 $95\%$ 双侧区间。

正态总体均值在方差已知时，由
\[
  \mathbb P\left\{\frac{\sqrt n(\bar X-\mu)}\sigma\ge-u_\alpha\right\}
  =1-\alpha
\]
得 $\mu$ 的上限为 $\bar X+u_\alpha\sigma/\sqrt n$；同理，下限为 $\bar X-u_\alpha\sigma/\sqrt n$。方差未知且 $n\ge2$ 时，两个单侧界分别为
\begin{equation}
  L_\mu=\bar X-t_{n-1}(\alpha)\frac S{\sqrt n},\qquad
  U_\mu=\bar X+t_{n-1}(\alpha)\frac S{\sqrt n}.
  \label{eq:study-normal-mean-bounds}
\end{equation}
与双侧等尾区间比较，这里用 $\alpha$ 分位数，而不是 $\alpha/2$ 分位数。

对于正态总体方差，在均值未知时令 $V=(n-1)S^2$、$\nu=n-1$，则
\begin{equation}
  L_{\sigma^2}=\frac V{\chi_\nu^2(\alpha)},\qquad
  U_{\sigma^2}=\frac V{\chi_\nu^2(1-\alpha)}.
  \label{eq:study-normal-variance-bounds}
\end{equation}
其依据分别是事件 $V/\sigma^2\le\chi_\nu^2(\alpha)$ 和 $V/\sigma^2\ge\chi_\nu^2(1-\alpha)$，各有概率 $1-\alpha$。若均值已知为 $\mu_0$，则改取 $V=\sum(X_i-\mu_0)^2$、$\nu=n$；对端点开平方得到标准差的单侧界。卡方分布不对称，不能只记“把 $\alpha/2$ 换成 $\alpha$”，却忽略上尾分位数在上下端点中的位置。

\begin{example}
\label{ex:study-one-sided-normal-mean}
沿用教材例 7.13 的 $12$ 个观测值，只讨论其正态抽样模型下的总体均值：$\bar x=45.25$、$s\approx10.3934$，求均值的 $90\%$ 置信上限。
\end{example}
\begin{solve}
由式~\eqref{eq:study-normal-mean-bounds}，上限为
\[
  45.25+t_{11}(0.1)\frac{10.3934}{\sqrt{12}}
  \approx49.3407.
\]
这提供的是均值的单侧界；即使该上限低于某个给定阈值，也不能推出每个测量点都低于该阈值，更不能仅凭均值区间直接完成需要其他指标和规则的整体分类。
\end{solve}

大样本单侧均值界可以把式~\eqref{eq:study-general-mean-ci} 中的 $u_{\alpha/2}$ 改为 $u_\alpha$；比例的单侧得分界则取式~\eqref{eq:study-wilson-ci} 中 $z=u_\alpha$，分别使用所需的一个端点。自助法同样可用单侧分位数校准，例如基本自助上限为 $2\widehat\theta-q_\alpha^*$，下限为 $2\widehat\theta-q_{1-\alpha}^*$。这些方法各自继承其近似条件，不能因为只取一侧就省略分布或重抽样假设。

\section{扩展阅读：“足球赛会杀人”的真假}
\label{sec:study-football-claim}
% 教材第 258--260 页；以下把教材叙述作为统计推断案例讨论。

教材以一项足球比赛期间死亡计数的观察性研究说明，区间估计的解释取决于估计对象和概率模型。按教材列出的数据，某比赛日相关男性死亡计数为 $41$，前后各五天共十天的日均计数为 $27.2$。这些是教材案例提供的观察记录；单凭“比赛日的数值高于平日”或“高于某个区间端点”，不能直接推出观看比赛导致了死亡增加。

首先必须分清\emph{平均水平的置信区间}与\emph{一次新观测的预测区间（prediction interval）}。若假定各日计数 $Y_1,\ldots,Y_m$ 独立同分布于 $\operatorname{Pois}(\lambda)$，则 $\lambda$ 是每日平均计数，$\bar Y$ 的方差为 $\lambda/m$，而一个新日期计数 $Y_{\mathrm{new}}$ 的方差为 $\lambda$。当新日期独立于用于估计均值的日期时，
\begin{equation}
  \operatorname{Var}(Y_{\mathrm{new}}-\bar Y)
   =\lambda\left(1+\frac1m\right).
  \label{eq:study-poisson-prediction-variance}
\end{equation}
因此，一个单日计数超出平均水平的置信区间，并不等价于它超出单日计数应有的预测波动范围；后者要同时计入新观测的随机性和均值估计误差。

其次，泊松计算本身依赖建模假设。每天具有同一均值、相互独立，以及方差等于均值，都需要结合数据背景检查；季节、星期、其他共同因素或时间依赖可能改变结论。把估计得到的 $\widehat\lambda$ 当成真实已知参数来计算尾概率，只得到代入参数的概率计算，并未自动处理参数估计误差。若以“至少 $41$ 人”为事件，应计算 $\mathbb P_\lambda(Y\ge41)$，而“超过 $41$ 人”是 $\mathbb P_\lambda(Y\ge42)$；离散计数中两个事件不能混同。

最后，研究问题和异常模式是否事先指定，会影响推断。若先浏览许多比赛、日期、年龄组和病因，再挑选最突出的组合，单个组合的尾概率不能直接代表整个搜索过程的误报概率。教材把某个连续十一日的特定模式概率倒数解释为大约每隔若干天出现一次，这种直观说法还需要说明模式定义、扫描窗口及其重叠依赖，不能仅凭一次事件概率就当作现实中的独立等待时间结论。

本例的学习重点是明确可以支持到哪一层结论：观察记录可以描述变化；给定模型可以评价这种变化与模型的相容程度；因果判断则还需要暴露信息、可比对照和对混杂因素的处理。一个区间包含无效应值不能证明“没有效应”，一个区间排除无效应值也不能单独证明“存在因果效应”；统计不确定性与研究设计是相互联系的两部分。

\section*{本章总结}
\addcontentsline{toc}{section}{本章总结}

置信区间首先是一条重复使用的推断规则。它用随机端点覆盖固定参数，其可靠性由抽样模型下的覆盖率描述。观测后的区间宽度表达估计精度，但既不是个体观测的取值范围，也不是未来观测的预测范围；比较区间优劣必须同时检查覆盖要求、长度准则和计算代价。

构造方法有三条主线。枢轴变量法利用已知且不含未知参数的分布，正态均值、方差、独立两总体的均值差和方差比，以及指数均值，都提供了典型例子。大样本方法以极限分布近似有限样本分布，反解时保留未知方差函数可得到得分区间，先代入估计则常得到瓦尔德区间。自助法以经验分布重抽样近似抽样波动，基本区间反射误差分位数，自助 $t$ 区间进一步学生化。后两类方法的名义置信水平必须结合其近似条件解释。

实际计算前应明确参数、抽样单位和独立性，区分单侧与双侧、已知与未知方差、独立样本与配对样本；计算中核对上尾分位数、自由度和反解不等号的方向；计算后说明估计对象、单位、置信水平和精确或近似的性质。增加样本量能改善随机误差，却不能自动消除抽样偏差、模型失配或研究设计中的混杂。

\section*{习题}
\addcontentsline{toc}{section}{习题}

以下选取教材第七章中能突出不同构造方法及适用条件的习题，题号均指教材本章习题号。计算中所用标准差均以分母 $n-1$ 定义，结果在最后一步舍入。

\begin{example}
\label{ex:study-ci-exercise-glue}
教材习题 1。假设胶合木板的断开力服从正态分布，$10$ 次独立试验观测为
\[
  561,459,534,601,476,516,522,490,480,550
  \quad(\mathrm N).
\]
求总体平均断开力的 $90\%$ 置信区间。
\end{example}
\begin{solve}
由数据得 $\bar x=518.9\,\mathrm N$、$s\approx44.0995\,\mathrm N$。总体方差未知，使用自由度为 $9$ 的 $t$ 区间，得到
\[
  518.9\pm t_9(0.05)\frac{44.0995}{\sqrt{10}}
  \approx[493.34,544.46]\,\mathrm N.
\]
这里置信水平为 $90\%$，所以双侧各留 $5\%$ 的尾概率；不能误用上尾 $0.1$ 分位数。
\end{solve}

\begin{example}
\label{ex:study-ci-exercise-measurements}
教材习题 7。测量误差独立同分布于 $N(0,10^2)$。至少重复测量多少次，才能使平均误差绝对值小于 $2$ 的概率不低于 $99\%$？
\end{example}
\begin{solve}
平均误差的标准差为 $10/\sqrt n$，所求条件是
\[
  \mathbb P(|\bar X|<2)
   =\mathbb P\left(|Z|<\frac{2\sqrt n}{10}\right)\ge0.99.
\]
因此 $n\ge(10u_{0.005}/2)^2\approx165.87$，至少需要 $166$ 次。总体均值为零意味着没有系统误差；增加重复次数缩小的是这里假定的独立随机误差。
\end{solve}

\begin{example}
\label{ex:study-ci-exercise-paired-sales}
教材习题 11 第三问。同一组 $11$ 种商品分别记录经营策略改变前后的销售额，求“后减前”的总体平均差的 $95\%$ 置信区间。明确采用各商品的差值相互独立、同分布且服从正态分布的模型。
\end{example}
\begin{solve}
按商品配对，差值为
\[
  3.2,-4.5,0.7,6.7,-6.1,-3.7,8.9,3.2,25.6,-17.3,11.8
  \quad(\text{万元}).
\]
其样本均值为 $\bar d\approx2.5909$ 万元、样本标准差为 $s_d\approx11.1188$ 万元。用差值样本作单样本 $t$ 推断，得到
\[
  \bar d\pm t_{10}(0.025)\frac{s_d}{\sqrt{11}}
  \approx[-4.879,10.061]\text{ 万元}.
\]
区间包含 $0$，说明按该模型和方法尚不能把总体平均变化的方向确定为正或负，并不证明改变策略没有作用。同一商品的前后值通常相关，不能未经假设把两组数据当作独立样本；仅假设前后边缘分布分别正态，也不足以推出差值正态。
\end{solve}

\begin{example}
\label{ex:study-ci-exercise-variance-df}
教材习题 16。正态总体的 $10$ 个长度观测值为
\[
  49.5,50.4,49.7,51.1,49.4,49.7,50.8,49.9,50.3,50.0
  \quad(\mathrm{mm}).
\]
分别在总体均值已知为 $50\,\mathrm{mm}$ 和均值未知时，求总体方差的 $95\%$ 置信区间。
\end{example}
\begin{solve}
若均值已知为 $50$，则 $\sum(x_i-50)^2=2.900\,\mathrm{mm}^2$，由式~\eqref{eq:study-normal-known-mean-variance-pivot} 得
\[
  \left[\frac{2.900}{\chi_{10}^2(0.025)},\,
        \frac{2.900}{\chi_{10}^2(0.975)}\right]
  \approx[0.1416,0.8931]\,\mathrm{mm}^2.
\]
若均值未知，则 $\bar x=50.08\,\mathrm{mm}$，$\sum(x_i-\bar x)^2=2.836\,\mathrm{mm}^2$，由式~\eqref{eq:study-normal-variance-ci} 得
\[
  \left[\frac{2.836}{\chi_9^2(0.025)},\,
        \frac{2.836}{\chi_9^2(0.975)}\right]
  \approx[0.1491,1.0502]\,\mathrm{mm}^2.
\]
同一组数据下，是否知道总体均值会同时改变中心化的平方和与卡方自由度。
\end{solve}

\begin{example}
\label{ex:study-ci-exercise-uniform-maximum}
教材习题 19。设 $X_1,\ldots,X_n$ 独立同分布于均匀分布 $U(0,\theta)$，$\theta>0$，记 $M=\max_iX_i$。求 $c_n$，使 $[M,c_nM]$ 是 $\theta$ 的置信系数为 $1-\alpha$ 的区间。
\end{example}
\begin{solve}
对 $0\le v\le1$，独立性给出
\[
  \mathbb P_\theta(M/\theta\le v)
   =\prod_{i=1}^n\mathbb P_\theta(X_i\le v\theta)=v^n.
\]
因此 $M/\theta$ 是枢轴量，且 $M\le\theta$ 几乎必然。要求区间覆盖等价于 $M/\theta\ge1/c_n$，故需 $c_n\ge1$ 且
\[
  1-c_n^{-n}=1-\alpha,
  \qquad c_n=\alpha^{-1/n}.
\]
所以区间为 $[M,M\alpha^{-1/n}]$。这不是围绕点估计对称延伸的区间，说明区间形状应由抽样分布和参数空间决定。
\end{solve}

\begin{example}
\label{ex:study-ci-exercise-defects-bound}
教材习题 22。独立检验 $120$ 件产品，发现 $9$ 件次品。在二项模型下求次品率的点估计、近似 $95\%$ 双侧置信区间和近似 $95\%$ 置信上限。
\end{example}
\begin{solve}
点估计为 $\widehat p=9/120=0.075$。将 $z=u_{0.025}$ 代入式~\eqref{eq:study-wilson-ci}，得到双侧得分区间约为
\[
  [0.0400,0.1364].
\]
计算单侧上限时改用 $z=u_{0.05}$，并只取相应的上端点，得到约 $0.1246$。上限与双侧区间的上端点不同，是因为前者把全部 $\alpha=0.05$ 分配到同一侧，而后者只在该侧分配 $0.025$。次品数只有 $9$，这里仍需说明正态近似的局限；若要求有限样本覆盖保证，可按二项尾概率反演单侧上限。
\end{solve}

\begin{example}
\label{ex:study-ci-exercise-rare-defects}
教材习题 26。在 $1000$ 件独立检验的产品中发现 $5$ 件次品，求次品率 $p$ 的 $95\%$ 置信区间，并判断能否仅凭 $n=1000$ 就认定瓦尔德近似充分准确。
\end{example}
\begin{solve}
$\widehat p=0.005$，得分区间约为 $[0.00214,0.01165]$。虽然总样本量达到 $1000$，观测到的次品数却只有 $5$，且参数估计接近边界，因此不能仅凭 $n$ 很大就保证瓦尔德区间的覆盖率。这里应结合稀少计数检查近似，必要时使用式~\eqref{eq:study-binomial-exact-inversion} 的精确二项反演。总样本量和有效事件数提供的信息不同。
\end{solve}

\chapter{假设检验}
\label{chap:study-hypothesis-testing}

\section{问题的提法和基本概念}
\label{sec:study-testing-concepts}
% 待补充本节的学习思考与习题解答。

\section{正态总体参数检验}
\label{sec:study-normal-tests}
% 待补充本节的学习思考与习题解答。

\section{比例 $p$ 的检验}
\label{sec:study-proportion-test}
% 待补充本节的学习思考与习题解答。

\section{似然比检验}
\label{sec:study-likelihood-ratio}
% 待补充本节的学习思考与习题解答。

\section{$p$ 值}
\label{sec:study-p-value}
% 待补充本节的学习思考与习题解答。

\section{扩展阅读：多重假设检验}
\label{sec:study-multiple-testing}
% 待补充本节的学习思考与习题解答。

\section*{本章总结}
\addcontentsline{toc}{section}{本章总结}
% 待补充；整理习题时注明教材题号。

\section*{习题}
\addcontentsline{toc}{section}{习题}
% 待补充；整理习题时注明教材题号。

\chapter{非参数假设检验}
\label{chap:study-nonparametric-testing}

\section{拟合优度检验}
\label{sec:study-goodness-of-fit}
% 待补充本节的学习思考与习题解答。

\section{威尔科克森秩和检验}
\label{sec:study-wilcoxon}
% 待补充本节的学习思考与习题解答。

\section{符号检验}
\label{sec:study-sign-test}
% 待补充本节的学习思考与习题解答。

\section{其他非参数检验概述}
\label{sec:study-other-nonparametric}
% 待补充本节的学习思考与习题解答。

\section{扩展阅读：正态性检验}
\label{sec:study-normality-tests}
% 待补充本节的学习思考与习题解答。

\section*{本章总结}
\addcontentsline{toc}{section}{本章总结}
% 待补充；整理习题时注明教材题号。

\section*{习题}
\addcontentsline{toc}{section}{习题}
% 待补充；整理习题时注明教材题号。

\chapter{相关分析和回归分析}
\label{chap:study-correlation-regression}

\section{相关分析}
\label{sec:study-correlation}
% 待补充本节的学习思考与习题解答。

\section{回归分析}
\label{sec:study-regression}
% 待补充本节的学习思考与习题解答。

\section{多元回归中自变量的选择和模型诊断简述}
\label{sec:study-regression-diagnostics}
% 待补充本节的学习思考与习题解答。

\section{扩展阅读：相关与因果}
\label{sec:study-correlation-causation}
% 待补充本节的学习思考与习题解答。

\section{附录}
\label{sec:study-regression-appendix}
% 待补充本节的学习思考与习题解答。

\section*{本章总结}
\addcontentsline{toc}{section}{本章总结}
% 待补充；整理习题时注明教材题号。

\section*{习题}
\addcontentsline{toc}{section}{习题}
% 待补充；整理习题时注明教材题号。

\end{document}
